\chapter{Chiral Hochschild cohomology and Koszul duality}
\label{chap:deformation-theory}\label{ch:chiral-hochschild-koszul}

\index{deformation theory!chiral|textbf}
\index{Hochschild cohomology!chiral|textbf}
\index{three Hochschild theories!discipline}

\begin{remark}[Five Hochschild functors]
\label{rem:three-hochschild-discipline}

Each Hochschild construction begins with its source category.  This
chapter studies the chart chiral cochain object
\[
 C^\bullet_{\mathrm{ch}}(\cA,\cA)
 :=R\operatorname{End}_{\cA^e}(\cA)
\]
in the derived category of chiral \(\cA\)-bimodules.  Its shift by one
governs deformations of the chiral product.  Four companion functors
occur in the monograph:
\begin{enumerate}[label=\textup{(\alph*)}]
\item Bakalov--De Sole--Kac bounded vertex cochains
 \(C^\bullet_{\mathrm{ch},b}(V,V)\), formed from the bounded part of
 the translation-equivariant chiral operad \(P^{\mathrm{ch}}(V)\);
\item categorical Hochschild cochains
 \[
 C^\bullet_{\mathrm{cat}}(\mathcal C)
 :=R\operatorname{End}_{\operatorname{Fun}^{\mathrm{ex}}(\mathcal C,
 \mathcal C)}(\operatorname{Id}_{\mathcal C});
 \]
\item topological Hochschild homology
 \(\mathrm{THH}(R):=R\otimes^{L}_{R^e}R\), an
 \(S^1\)-equivariant spectrum associated with an \(E_1\)-ring
 spectrum~\(R\);
\item the coHochschild complex of a conilpotent chiral coalgebra~\(C\)
 with bicomodule coefficients~\(N\), whose ordered associative
 realization has coalgebraic variance and cohomology
 \(\operatorname{Cotor}_{C^e}(C_\Delta,N)\).
\end{enumerate}
The chart algebra, vertex algebra, dg category, ring spectrum, and
coalgebra determine five distinct domains.  A passage between two
domains is a specified natural transformation or quasi-isomorphism.
Throughout this chapter \(\ChirHoch^n(\cA)\) means
\(H^n(C^\bullet_{\mathrm{ch}}(\cA,\cA))\).
\end{remark}

\begin{definition}[Hochschild comparison datum;
\ClaimStatusDefinitional]
\label{def:theorem-h-comparison-datum}
Let \(V\) be a translation-equivariant vertex algebra and let
\(\cA_V\) be its chiral algebra on the curve.  A
\emph{bounded-to-chart comparison} is a cochain map
\[
 \chi_V^{\mathrm{bd}}
 \colon C^\bullet_{\mathrm{ch},b}(V,V)
 \longrightarrow C^\bullet_{\mathrm{ch}}(\cA_V,\cA_V).
\]
For a compactly generated factorization category \(\mathcal C\) with
compact generator~\(b\) and chart \(A_b=R\operatorname{End}(b)\), a
\emph{categorical-to-chart comparison} is a cochain map
\[
 \chi_b^{\mathrm{cat}}
 \colon C^\bullet_{\mathrm{cat}}(\mathcal C)
 \longrightarrow C^\bullet_{\mathrm{ch}}(A_b,A_b)
\]
induced once a factorization Eilenberg--Watts equivalence has been
specified.  A
bar--bicomodule equivalence similarly produces a map from the
coHochschild model to the chart complex.  Each transport theorem in
this chapter names the corresponding map and assumes or proves its
quasi-isomorphism.
\end{definition}

The Maurer--Cartan element
$\Theta_\cA\in\MC(\gAmod)$ determines the chiral product, and the
shifted complex
\[
 \mathbb T_\cA:=C^\bullet_{\mathrm{ch}}(\cA,\cA)[1]
\]
is its deformation complex.  Thus $H^1(\mathbb T_\cA)$ parametrises
first-order product deformations and $H^2(\mathbb T_\cA)$ receives
their primary extension obstructions.  Under the family support datum
$H_H(\cA;S)$, Theorem~H computes this complex by a family-specific
deformation retract.  The retract determines its completion and its
cohomological support.

Pairwise Koszul biduality concerns the represented
quadratic--Verdier reflection of the chart algebra.  A statement about the support of
$C^\bullet_{\mathrm{ch}}(\cA,\cA)$ requires an additional comparison
with a cochain model whose differential and contraction are explicit.
The support depends on the algebra and coefficient theory.
Bakalov--De Sole--Kac compute
\[
 \dim H^n_{\mathrm{ch},b}(B_{\mathfrak h},B_{\mathfrak h})
 =(2,1,0,0,\ldots)
\]
for a one-dimensional even space~\(\mathfrak h\), and compute the
bounded cohomology of \(\operatorname{Vir}_c\) as one-dimensional in
degrees \(0,2,3\) and supported on those degrees
\cite[Theorems~7.4 and~7.2]{BDSK21}.  A quasi-isomorphism
\(\chi_V^{\mathrm{bd}}\) transports either calculation to the chart
complex.  De Sole--Kac variational Poisson cohomology is nonvanishing
in unboundedly many degrees for Virasoro-type Poisson vertex algebras,
so every bounded-support statement of this chapter lives strictly on
the stated Koszul/bounded locus, with the comparison
\(\chi_V^{\mathrm{bd}}\) itself part of the hypothesis package
(cf.\ \cite{BDSK21}).

The cyclic deformation complex $\Defcyc(\cA)$
(Definition~\ref{def:cyclic-deformation-elementary}) is the cyclic
refinement of $C^\bullet_{\mathrm{ch}}(\cA,\cA)$.

The chiral Hochschild complexes
carry natural brace and $E_2$ structures. The convolution dg~Lie
algebra $\Convstr(\barB(\cA), \cA)$ governing deformations is the
strict model of the $L_\infty$ deformation object
$\Convinf(\barB(\cA), \cA)$
(Theorem~\ref{thm:operadic-homotopy-convolution}).
The modular cyclic deformation complex
$\Defcyc^{\mathrm{mod}}(\cA)$
(Definition~\ref{def:modular-cyclic-deformation-complex}),
strict model of $\Definfmod(\cA)$
(Remark~\ref{rem:modular-cyc-strictification}),
extends these structures to the full modular regime.

Two algebraic objects govern the open/closed dictionary:
\begin{enumerate}[label=\textup{(\alph*)}]
\item \emph{Bar/cobar} classifies \emph{twisting morphisms}:
 maps from a conilpotent coalgebra into a chiral algebra, equivalently
 maps out of the corresponding cobar algebra
 (Chapter~\ref{chap:bar-cobar}). For~$\cA$ itself this gives the
 inversion $\Omega\barB^{\mathrm{ch}}(\cA)\simeq \cA$.  A quadratic
 presentation defines the coalgebra $\cA^i$ and a twisting comparison
 \[
 q_\cA\colon\cA^i\longrightarrow\barB^{\mathrm{ch}}(\cA).
 \]
 Under the Koszul hypothesis $q_\cA$ is a quasi-isomorphism, and
 Verdier duality applied to $\cA^i$ defines the algebra $\cA^!$.
\item \emph{Chiral Hochschild cochains}
 $C^\bullet_{\mathrm{ch}}(\cA, \cA)$ form the cochain-level
 closed-sector endocentral object acting on~$\cA$ in the local
 chiral Swiss-cheese category
 (Theorem~\ref{thm:chiral-deligne-tamarkin}). The brace dg algebra
 structure gives the algebraic closed object over $\cA$ in the
 curve-level chiral Swiss-cheese comparison.  A physical
 factorisation algebra~$T$ enters through an open--closed comparison datum
 \textup{(}Definition~\ref{def:thqg-oca-datum}\textup{)}, in
 particular a typed brace dg comparison
 $\beta_T\colon \operatorname{Loc}_{\partial}
 \operatorname{Obs}^{\mathrm{bulk}}(T)
 \to C^\bullet_{\mathrm{ch}}(\cA,\cA)$.  When $\beta_T$ is a
 quasi-isomorphism, it identifies the physical bulk with the
 derived chiral centre
 (Proposition~\ref{prop:thqg-universal-action-not-reconstruction}).
 The complex itself is
 $Z^{\mathrm{der}}_{\mathrm{ch}}(\cA):=
 C^\bullet_{\mathrm{ch}}(\cA,\cA)$; its cohomology is
 $\ChirHoch^\bullet(\cA,\cA)$, and the ordinary chiral centre is
 $Z_{\mathrm{ch}}(\cA)=\ChirHoch^0(\cA,\cA)$
 (Definition~\ref{def:thqg-chiral-derived-center}).
\end{enumerate}
The open/closed MC element $\Theta^{\mathrm{oc}}_\cA$
(\S\ref{sec:thqg-open-closed-realization}) packages both:
its closed projection is the bar-intrinsic $\Theta_\cA$
(Theorem~\ref{thm:mc2-bar-intrinsic}), and its mixed projection
is the closed-sector action on the boundary through
$C^\bullet_{\mathrm{ch}}(\cA, \cA)$. The shadow obstruction tower
$(\kappa, \mathfrak{C}, \mathfrak{Q}, \ldots)$ lives in the closed
sector; the annulus trace
$\int_{S^1} \mathcal{M}
\simeq \mathrm{HH}_{\bullet,\mathrm{cat}}(\mathcal{M})$
(Theorem~\ref{thm:thqg-annulus-trace}) is the first modular shadow
of the open sector.

\section{Deformations of a chiral algebra}

Let $\mathcal{A}$ be a chiral algebra on a smooth projective curve~$X$,
with chiral product
$\mu\colon j_*j^*(\mathcal{A}\boxtimes\mathcal{A})\to\Delta_*\mathcal{A}$
(Chapter~\ref{ch:algebraic-foundations}).
\index{obstruction theory|textbf}
A \emph{first-order deformation} of~$\mathcal{A}$ is a
$\mathcal{D}_X$-module
$\mathcal{A}_\epsilon = \mathcal{A}\oplus\epsilon\,\mathcal{A}$
over $k[\epsilon]/(\epsilon^2)$, equipped with a chiral product
$\mu_\epsilon = \mu + \epsilon\,\phi$, where
$\phi\colon j_*j^*(\mathcal{A}\boxtimes\mathcal{A})\to\Delta_*\mathcal{A}$
is a $\mathcal{D}_X$-bimodule map.

The deformed product $\mu_\epsilon$ must satisfy the Borcherds
identity~\cite{Bor92,BD04}:
for sections $a,b,c$ of~$\mathcal{A}$ at pairwise distinct points
$z_1,z_2,z_3\in X$,
\[
\begin{aligned}
 &\operatorname{Res}_{z_1\!=z_2}\!\mu_\epsilon(a(z_1),\mu_\epsilon(b(z_2),c(z_3)))
 - \operatorname{Res}_{z_1\!=z_2}\!\mu_\epsilon(b(z_2),\mu_\epsilon(a(z_1),c(z_3)))\\
 &\qquad= \operatorname{Res}_{z_2\!=z_3}\!\mu_\epsilon(\mu_\epsilon(a(z_1),b(z_2)),c(z_3)).
\end{aligned}
\]
Expanding modulo~$\epsilon^2$ and using the Borcherds identity for~$\mu$,
the $\epsilon$-coefficient gives:
\begin{equation}\label{eq:deformation-cocycle}
\begin{split}
 \operatorname{Res}_{z_1\!=z_2}\!\bigl[\phi(a,\mu(b,c))
 &+ \mu(a,\phi(b,c))\\
 &- \phi(b,\mu(a,c))
 - \mu(b,\phi(a,c))\bigr]\\
 &= \operatorname{Res}_{z_2\!=z_3}\!\bigl[\phi(\mu(a,b),c)
 + \mu(\phi(a,b),c)\bigr].
\end{split}
\end{equation}
This is a cocycle condition: $\phi$ is closed in a complex whose
differential is built from~$\mu$ via the residue operations above.
That complex is the \emph{chiral Hochschild complex}
$C^{\bullet}_{\mathrm{chiral}}(\mathcal{A})$, constructed in
\S\ref{subsec:hochschild-duality}.

The cocycle condition~\eqref{eq:deformation-cocycle} involves three
marked points on~$X$ and residues as pairs collide. An $n$-cochain
in the chiral Hochschild complex involves $(n+2)$ points on~$X$;
the differential acts by residues along the collision divisors
$D_{ij}=\{z_i=z_j\}$ in~$X^{n+2}$. The complex lives on the
configuration space
$C_{n+2}(X)=\{(z_1,\ldots,z_{n+2})\in X^{n+2}: z_i\neq z_j\}$,
and the differential extends to the Fulton--MacPherson
compactification $\overline{C}_{n+2}(X)$ as the bar
differential does (Chapter~\ref{chap:bar-cobar}).

At genus $g\geq 1$, the fiberwise differential satisfies
$\dfib^{\,2}=\kappa(\cA)\cdot\omega_g$
(Chapter~\ref{ch:heisenberg-frame};
Convention~\ref{conv:higher-genus-differentials}), so the bar complex
is a curved $A_\infty$-algebra. The curvature scalar~$\kappa(\cA)$
obstructs lifting a genus-$0$ deformation to genus~$g$.
A degree-\(d\) perfect chain pairing for \((\cA,\cA^!)\) exchanges
the two cohomologies by
\[
 \ChirHoch^n(\cA)
 \cong\ChirHoch^{d-n}(\cA^!)^\vee\otimes\omega_X.
\]
Theorem~H supplies the family support through its deformation
retract; Definition~\ref{def:theorem-h-shift-normalization-package}
supplies this additional pairing statement.

\section{The chiral Hochschild cochain complex}

\subsection{The cochain spaces}

A degree-$n$ cochain involves $(n+2)$ marked points on~$X$: one output,
$n$ normalized input slots, and one evaluation point.  In the algebraic
Ext comparison the $n$ input factors are
$\overline{\mathcal A}^{\widehat\otimes n}$; the two remaining marked
points are the bimodule legs of the two-sided bar correspondence.

\begin{definition}[Fulton--MacPherson cochain model]
\label{def:chiral-hochschild-complex}
\index{Hochschild complex!chiral!geometric realization|textbf}
For a chiral algebra $\mathcal{A}$ on a smooth curve $X$, define
\[
C^n_{\mathrm{FM}}(\mathcal A)
=
R\Gamma\!\left(
\overline{C}_{n+2}(X),
R\mathcal{H}om_{\mathcal D_{X^{n+2}}}
\!\left(j_*j^*\mathcal A^{\boxtimes(n+2)},\Delta_*\mathcal A\right)
\otimes\Omega^n_{\overline C_{n+2}(X)}(\log D)
\right).
\]
where:
\begin{itemize}
\item $\overline{C}_{n+2}(X)$ is the Fulton--MacPherson compactification
\item $j: C_{n+2}(X) \to \overline{C}_{n+2}(X)$ is the open embedding
\item the internal Hom supplies the output in the diagonal module
  $\Delta_*\mathcal A$
\item $\Omega^n_{\overline{C}_{n+2}(X)}(\log D)$ are $n$-forms with logarithmic poles along the boundary divisor $D$
\end{itemize}
The total differential is the sum of the internal, chiral-product, and
logarithmic de~Rham differentials.  Under the geometric--algebraic
comparison of Theorem~\ref{thm:HH-config-space-formula}, this is a
model of the chiral Hochschild complex.  The algebraic Ext model
$\Ext^*_{\cA^e}(\cA,\cA)$
\textup{(}Definition~\textup{\ref{def:chirHoch-Ext}} in
Chapter~\textup{\ref{chap:koszul-pair-structure}}\textup{)} computes
the same chiral derived endomorphism centre under
Theorem~\textup{\ref{thm:geometric-chiral-hochschild}}; the bigraded
Verdier model is recorded in
Definition~\textup{\ref{def:bigraded-hochschild}} of
Chapter~\textup{\ref{chap:higher-genus}}. All three
models compute $Z^{\mathrm{der}}_{\mathrm{ch}}(\cA)$ under the
geometric--algebraic comparison hypotheses named in
Theorem~\textup{\ref{thm:hochschild-bar-cobar}}.  Its cohomological
support is determined by the differential of this total complex.
\end{definition}

\begin{definition}[Diagonal-bar chiral Hochschild and completed
\texorpdfstring{$RHH_{\mathrm{ch}}$}{RHHch};
\ClaimStatusDefinitional]
\label{def:diagonal-bar-rhh-chiral}
Let $\cA$ be a chiral algebra on the smooth curve~$X$.  The
\emph{diagonal bar resolution} is the two-sided chiral bar object
\[
\operatorname{Bar}_{\cA^e}^{\Delta}(\cA)
:=
\cA\otimes \bar B^{\mathrm{ch}}_X(\cA)\otimes \cA
\longrightarrow \cA
\]
in chiral $\cA$-$\cA$-bimodules. Its bar-degree-$p$ piece defines
\[
CH^{p,\bullet}_{\mathrm{ch}}(\cA)
:=
R\operatorname{Hom}_{\cA^e}
\!\left(
\cA\otimes \bar B^{\mathrm{ch},p}_X(\cA)\otimes\cA,\,
\cA
\right),
\]
realized geometrically by the Fulton--MacPherson model on
$\overline C_{p+2}(X)$ with the two external Hochschild marked points
kept as the diagonal bimodule legs. The uncompleted chiral
Hochschild cochain complex is the direct-sum totalization
\[
CH^\bullet_{\mathrm{ch}}(\cA)
:=
\operatorname{Tot}^{\oplus}_{p\ge0}
CH^{p,\bullet}_{\mathrm{ch}}(\cA)[-p].
\]
For a finite conformal-weight/PBW tower
$\{\cA_{\leq N}\}_{N\ge0}$ satisfying strict transition maps, the
completed chiral Hochschild complex is
\[
RHH_{\mathrm{ch}}(\cA)
:=
R\!\varprojlim_N\,
CH^\bullet_{\mathrm{ch}}(\cA_{\leq N}).
\]
In a finite-type situation this agrees with
$CH^\bullet_{\mathrm{ch}}(\cA)$. In a class-\(\mathsf M\) or
infinite-generator situation the symbol $RHH_{\mathrm{ch}}$ denotes
the completed finite-window object. Its cohomology is
\[
\ChirHoch^n(\cA):=H^n(RHH_{\mathrm{ch}}(\cA)).
\]
Accordingly, $C^\bullet_{\mathrm{ch}}(\cA,\cA)$ below denotes the
derived endomorphism complex in the ambient under discussion: the
direct-sum totalization in finite type and $RHH_{\mathrm{ch}}(\cA)$
in the completed finite-window ambient.
In the cohomological convention of this chapter,
$\ChirHoch^1(\cA)$ is the fixed-product derivation column
$\Der_{\mathrm{ch}}(\cA)/\Inn_{\mathrm{ch}}(\cA)$, while
$\ChirHoch^2(\cA)$ contains the binary product-deformation classes in
the unshifted convention.
The shifted formal-moduli convention reads binary product
deformations from the degree-$2$ cyclic deformation complex.  The
unshifted degree-$1$ column retains its fixed-product derivation
meaning.
\end{definition}

\begin{definition}[Family support datum for Theorem H;
\ClaimStatusDefinitional]
\label{def:theorem-h-collision-depth-package}
Fix a finite set \(S\subset\mathbb Z\).  A family support datum
\(H_H(\cA;S)\) consists of a complete filtered cochain complex
\(Q_\cA\), a quasi-isomorphism
\[
 \gamma_\cA\colon Q_\cA
 \xrightarrow{\ \sim\ }
 C^\bullet_{\mathrm{ch}}(\cA,\cA),
\]
and a complete complex \(K_{\cA,S}\) whose cochain terms occur in
degrees belonging to~\(S\), together
with filtration-continuous maps
\[
 K_{\cA,S}
 \mathop{\xrightarrow{\ \iota_\cA\ }}
 Q_\cA
 \mathop{\xrightarrow{\ p_\cA\ }}
 K_{\cA,S}
\]
and a degree \(-1\) map \(h_\cA\colon Q_\cA\to Q_\cA[-1]\)
satisfying
\begin{equation}\label{eq:theorem-h-global-retract}
 p_\cA\iota_\cA=\id_{K_{\cA,S}},
 \qquad
 d h_\cA+h_\cA d
 =\id_{Q_\cA}-\iota_\cA p_\cA.
\end{equation}
When \(Q_\cA\) is obtained from conformal-weight windows
\(Q_{\cA,N}\), the datum includes the comparison
\[
 Q_\cA\xrightarrow{\ \sim\ }R\!\varprojlim_N Q_{\cA,N}
\]
and the Milnor identities
\(\varprojlim_N^1H^{m-1}(Q_{\cA,N})=0\).
When the finite-window model is first constructed on ordered
configurations, the datum also includes compatible quasi-isomorphisms
\[
 \operatorname{av}_{\cA,N}\colon
 \bigl(Q^{\mathrm{ord}}_{\cA,N}\bigr)_{h\Sigma}
 \xrightarrow{\ \sim\ }Q^{\mathrm{sym}}_{\cA,N},
 \qquad
 R\!\varprojlim_NQ^{\mathrm{sym}}_{\cA,N}
 \xrightarrow{\ \sim\ }Q_\cA.
\]
Thus \(H_H(\cA;S)\) specifies the chain model, the completion, the
ordered-to-symmetric passage, and the deformation retract that
computes its support.
\end{definition}

\begin{definition}[Finite-window collision realization;
\ClaimStatusDefinitional]
\label{def:theorem-h-pbw-finite-window-lane}
Let \(\mathsf P_{p}\) be the incidence category of collision strata
of \(\overline C_{p+2}(X)\).  For a conformal-weight window~\(N\),
write \(T_{\cA,\sigma,N}\) for the residue-twisted coefficient
complex on \(\sigma\in\mathsf P_p\).  A
\emph{finite-window collision realization} consists of complexes
\(K_{\cA,\sigma,N}\) and maps
\[
 K_{\cA,\sigma,N}
 \mathop{\xrightarrow{\ \iota_{\sigma,N}\ }}
 T_{\cA,\sigma,N}
 \mathop{\xrightarrow{\ p_{\sigma,N}\ }}
 K_{\cA,\sigma,N},
 \qquad h_{\sigma,N}\colon
 T_{\cA,\sigma,N}\to T_{\cA,\sigma,N}[-1],
\]
satisfying the identities of
\eqref{eq:theorem-h-global-retract}.  The maps are natural for every
incidence morphism \(\sigma\to\tau\), compatible with the bar-face
maps between insertion degrees, compatible with the window maps
\(N+1\to N\), and equivariant for the ordered
\(\Sigma_{p+2}\)-action.  Their incidence--bar totalizations are
denoted
\[
 K_{\cA,N}:=\operatorname{Tot}^{\oplus}_{p\geq0}
 \operatorname{Tot}_{\sigma\in\mathsf P_p}
 K_{\cA,\sigma,N}[-p],
 \qquad
 Q^{\mathrm{ord}}_{\cA,N}:=
 \operatorname{Tot}^{\oplus}_{p\geq0}
 \operatorname{Tot}_{\sigma\in\mathsf P_p}
 T_{\cA,\sigma,N}[-p].
\]
PBW filtrations and quadratic Koszul complexes provide one method for
constructing these maps.  The deformation-retract identities and
their incidence compatibility are the hypotheses consumed by
Theorem~H.
\end{definition}

\begin{proposition}[Cohomology of a family support datum;
\ClaimStatusProvedHere]
\label{prop:theorem-h-kdh-criterion}
For every \(H_H(\cA;S)\), the maps \(\iota_\cA\) and \(p_\cA\)
induce inverse isomorphisms
\[
 H^\bullet(K_{\cA,S})
 \cong H^\bullet(Q_\cA)
 \cong \ChirHoch^\bullet(\cA).
\]
Consequently
\[
 \operatorname{Supp}\ChirHoch^\bullet(\cA)
 \subseteq S.
\]
The same conclusion holds for a completed finite-window realization
when its inverse-limit and averaging maps are the quasi-isomorphisms
specified in Definition~\ref{def:theorem-h-collision-depth-package}.
\end{proposition}

\begin{proof}
Equation~\eqref{eq:theorem-h-global-retract} exhibits
\(\iota_\cA p_\cA\) as chain-homotopic to
\(\id_{Q_\cA}\), while \(p_\cA\iota_\cA\) is the identity of
\(K_{\cA,S}\).  Hence \(\iota_\cA\) and \(p_\cA\) are homotopy
inverse quasi-isomorphisms.  Composition with \(\gamma_\cA\) gives
the asserted chart-cohomology isomorphism.  The Milnor exact sequence
and the averaging quasi-isomorphisms carry the same conclusion from
the finite windows to the completed symmetric complex.
\end{proof}

\begin{proposition}[coHochschild transport scope;
\ClaimStatusProvedHere]
\label{prop:cohochschild-transport-scope-theorem-h}
Let $A$ be a strongly admissible associative chiral algebra, let
$C=\Barch(A)$, and assume the bimodule--bicomodule equivalence
Theorem~\ref{thm:bimod-bicomod}. Then chiral Hochschild theory
transports to chiral coHochschild theory through that equivalence:
\[
\HH_{\mathrm{ch}}^\bullet(A,M)
\simeq
\coHoch_{\mathrm{ch}}^\bullet
\bigl(C,\KK_A^{\mathrm{bi}}(M)\bigr)
\]
for complete $A$-bimodules~$M$, and similarly in homology.  Thus a
coalgebra-side support theorem consists of a bar presentation
$C\simeq\Barch(A)$, the diagonal bimodule--bicomodule comparison, and
a family support datum for~$A$.
\end{proposition}

\begin{proof}
The displayed equivalence is exactly the cohomological transport
theorem
Theorem~\ref{thm:HH-coHH-cohomology}, whose proof applies
Theorem~\ref{thm:bimod-bicomod} to derived Hom over the diagonal.
The homological version is Theorem~\ref{thm:HH-coHH-homology}.  Strong
admissibility and the diagonal equivalence define the transport
functor $\KK_A^{\mathrm{bi}}$ and carry the deformation retract to the
coalgebraic complex.
\end{proof}

Explicitly, a degree $n$ cochain is a sum of expressions
\[
\phi = \sum_I a_0^{(I)}(z_0) \otimes a_1^{(I)}(z_1) \otimes \cdots \otimes a_n^{(I)}(z_n) \otimes a_{\infty}^{(I)}(z_{\infty}) \otimes \omega_I
\]
where $a_i^{(I)} \in \mathcal{A}$ and $\omega_I$ is an $n$-form on $\overline{C}_{n+2}(X)$ with logarithmic singularities.

\begin{lemma}[Three-point Arnold relation; \ClaimStatusProvedElsewhere]
\label{lem:arnold-three-point-hochschild}
For the logarithmic $1$-forms
$\eta_{ij}=d\log(z_i-z_j)$ on $C_3(X)$,
\begin{equation}\label{eq:arnold-three-point}
\eta_{12} \wedge \eta_{23}
+ \eta_{23} \wedge \eta_{31}
+ \eta_{31} \wedge \eta_{12} = 0.
\end{equation}
\end{lemma}

\begin{proof}[Attribution]
This is the three-point case of Arnold's relation
(Arnold~\cite{Arnold69}); it is also the $n=3$ specialization of
Theorem~\ref{thm:arnold-relations}.
\end{proof}

\subsection{The differential: three components}

The chiral Hochschild differential
$d_{\mathrm{Hoch}}: C^n_{\text{chiral}} \to C^{n+1}_{\text{chiral}}$
has three components reflecting the algebraic, geometric, and operadic
structures:

\begin{theorem}[The chiral Hochschild differential; \ClaimStatusProvedHere]
\label{thm:chiral-hochschild-differential}
The differential decomposes as
\[
d_{\mathrm{Hoch}} = d_{\text{int}} + d_{\text{fact}} + d_{\text{config}}
\]
where:
\begin{enumerate}
\item $d_{\text{int}}$: internal differential from the $\mathcal{D}_X$-module structure
\item $d_{\text{fact}}$: factorization using chiral multiplication
\item $d_{\text{config}}$: de Rham differential on configuration space
\end{enumerate}
\end{theorem}

\begin{proof}
We expand $d_{\mathrm{Hoch}}^{\,2}
= (d_{\text{int}} + d_{\text{fact}}
+ d_{\text{config}})^2$ into six terms (three squares and
three anticommutators) and show that each vanishes on a
generic cochain $\phi \in C^n_{\text{chiral}}$.

\emph{(I) Pure squares.}

\emph{(I.1)} $d_{\text{int}}^2(\phi) = 0$.
The internal differential $d_{\text{int}}$ acts on each
tensor factor $a_i$ by the $\mathcal{D}_X$-module differential
$d_\mathcal{A}$, alternating signs. Since $d_\mathcal{A}^2 = 0$
and the signs follow the standard Koszul rule, this is the
usual chain complex identity.

\emph{(I.2)} $d_{\text{config}}^2(\phi) = 0$.
$d_{\text{config}} = d_{\overline{C}_{n+2}(X)}$ is the
de~Rham differential on the compactified configuration space.
Since $d_{\mathrm{dR}}^2 = 0$, this is immediate.

\emph{(I.3)} $d_{\text{fact}}^2(\phi) = 0$.
Applied to $\phi \in C^n$, $d_{\text{fact}}^2$ composes two
chiral multiplications at distinct pairs of colliding points.
For each pair of collisions $(z_i = z_j)$ and $(z_k = z_l)$
with $\{i,j\} \cap \{k,l\} = \varnothing$, the two residue
extractions commute. For overlapping pairs (sharing a point),
the double residue
$\mathrm{Res}_{z_i = z_j}\mathrm{Res}_{z_j = z_k}$ is the
iterated OPE, which cancels with
$\mathrm{Res}_{z_j = z_k}\mathrm{Res}_{z_i = z_j}$ by
associativity of the chiral product: $\mu \circ (\mu \otimes
\mathrm{id}) = \mu \circ (\mathrm{id} \otimes \mu)$ on the
localizations $j_*j^*$ (the Borcherds identity).

\emph{(II) Mixed anticommutators.}

\emph{(II.1)} $\{d_{\text{int}}, d_{\text{fact}}\}(\phi) = 0$.
The chiral product $\mu\colon j_*j^*(\mathcal{A} \boxtimes
\mathcal{A}) \to \Delta_*\mathcal{A}$ is a morphism of
$\mathcal{D}_X$-modules, so $d_\mathcal{A} \circ \mu =
\mu \circ (d_\mathcal{A} \boxtimes \mathrm{id} +
\mathrm{id} \boxtimes d_\mathcal{A})$. Expanding the
anticommutator term by term, each
$d_\mathcal{A}(\mu(a_i, a_j))$ equals
$\mu(d_\mathcal{A} a_i, a_j) + (-1)^{|a_i|}
\mu(a_i, d_\mathcal{A} a_j)$,
and the signs from $d_{\text{int}}$ and $d_{\text{fact}}$
cancel pairwise.

\emph{(II.2)} $\{d_{\text{int}}, d_{\text{config}}\}(\phi) = 0$.
The logarithmic forms
$\omega \in \Omega^*_{\overline{C}_{n+2}}(\log D)$
are \emph{universal} functions of the configuration coordinates
$z_i$.  Thus $d_\mathcal{A}$ acts on the $a_i$ factors and
$d_{\mathrm{dR}}$ acts on $\omega$. Since the
two differentials act on different tensor factors:
$d_\mathcal{A}(a_i \otimes \omega) =
(d_\mathcal{A} a_i) \otimes \omega$ and
$d_{\mathrm{dR}}(a_i \otimes \omega) =
a_i \otimes d_{\mathrm{dR}}\omega$,
the anticommutator reduces to
$(d_\mathcal{A} \otimes d_{\mathrm{dR}} +
d_{\mathrm{dR}} \otimes d_\mathcal{A})(a_i \otimes \omega)
= 0$ by the Koszul sign rule for the tensor product of
chain complexes.

\emph{(II.3)} $\{d_{\text{fact}}, d_{\text{config}}\}(\phi) = 0$.
The cancellation requires the
Arnold--Orlik--Solomon relations.

On a degree-$1$ cochain
$\phi = a_1 \otimes a_2 \otimes a_3 \otimes
\eta_{12} \in C^1$ (the smallest nontrivial case,
with three points $z_1, z_2, z_3$), the two compositions act
as:
\begin{align*}
(d_{\text{config}} \circ d_{\text{fact}})(\phi)
&= d_{\text{config}}\Bigl(
 \mathrm{Res}_{z_1 = z_2}\bigl[a_1 a_2 \otimes a_3
 \otimes \eta_{12}\bigr] \Bigr), \\
(d_{\text{fact}} \circ d_{\text{config}})(\phi)
&= \sum_{i<j}
 \mathrm{Res}_{z_i = z_j}\bigl[a_1 \otimes a_2 \otimes a_3
 \otimes \partial^{\mathrm{Stokes}}\eta_{12}\bigr],
\end{align*}
where $\partial^{\mathrm{Stokes}}$ denotes the boundary
contribution of $d_{\text{config}}$ on the compactified
complex: the form $\eta_{12} = d\log(z_1 - z_2)$ is closed on
the open configuration space, $d_{\mathrm{dR}}(\eta_{12}) = 0$,
so $d_{\text{config}}$ acts on it purely through the Stokes
restriction to the boundary divisors of
$\overline{C}_3(X)$.

The key identity is the three-point Arnold relation
(Lemma~\ref{lem:arnold-three-point-hochschild}) in its
boundary-restriction form: on the collision divisor $D_{12}$
the logarithmic forms satisfy
\[
\eta_{13}\big|_{D_{12}}
\;=\;
\eta_{23}\big|_{D_{12}},
\]
both restricting to $d\log(z - z_3)$ for the collided point
$z = z_1 = z_2$.
In each residue term, $d_{\text{fact}}$ extracts the OPE
at a collision, and the Stokes term of $d_{\text{config}}$
produces the matching boundary form. The restriction identity
ensures these contributions
cancel pairwise, yielding
$\{d_{\text{fact}}, d_{\text{config}}\}(\phi) = 0$.

For general $n$, the argument extends by applying the
generalized Arnold relations on $\overline{C}_{n+2}(X)$
(Theorem~\ref{thm:arnold-relations}): every boundary
stratum of $\overline{C}_{n+2}$ where two points collide
contributes a residue (from $d_{\text{fact}}$) and a
boundary form (from $d_{\text{config}}$), and the
Arnold relations ensure these pair up with opposite signs.
\end{proof}

\subsection{Explicit formula for the differential}

For a cochain $\phi \in C^n_{\text{chiral}}$, the differential acts by:

\begin{align}
(d_{\text{int}}\phi)(&z_0,\ldots,z_{n+1}) = \sum_{i=0}^{n+1} (-1)^i d_{\mathcal{A}}(\phi(z_0,\ldots,\hat{z}_i,\ldots,z_{n+1})) \\
(d_{\text{fact}}\phi)(&z_0,\ldots,z_{n+1}) = \sum_{i=1}^n (-1)^i \text{Res}_{z_i = z_0} \phi(\mu(z_0,z_i),z_1,\ldots,\hat{z}_i,\ldots,z_{n+1}) \\
&+ \sum_{1 \leq i < j \leq n} (-1)^{i+j} \phi(z_0,\ldots,\mu(z_i,z_j),\ldots,\hat{z}_i,\ldots,\hat{z}_j,\ldots,z_{n+1}) \\
(d_{\text{config}}\phi)(&z_0,\ldots,z_{n+1}) = d_{\overline{C}_{n+2}}(\phi)
\end{align}

where $\hat{z}_i$ denotes omission and $\mu$ is the chiral multiplication.

\section{The diagonal bar resolution and the chiral Hochschild spectral sequence}

\subsection{The diagonal bimodule resolution}

\begin{theorem}[chiral Hochschild via bar-cobar; \ClaimStatusProvedHere]\label{thm:hochschild-bar-cobar}
\textup{[Regime: quadratic on the Koszul locus
\textup{(}Convention~\textup{\ref{conv:regime-tags})}.]}

Let $\mathcal{A}$ be a bounded PBW chiral Koszul algebra
\textup{(}Definition~\textup{\ref{def:chiral-koszul-pair})}. The
two-sided chiral bar object
\[
\operatorname{Bar}_{\mathcal{A}^{e}}(\mathcal{A})
:=
\mathcal{A}\otimes
\bar{B}^{\mathrm{ch}}(\mathcal{A})\otimes\mathcal{A}
\]
is a cofibrant replacement of the diagonal
$\mathcal{A}$-$\mathcal{A}$ bimodule. Hence
\[
C^{\bullet}_{\mathrm{ch}}(\mathcal{A},\mathcal{A})
\;\simeq\;
R\mathrm{Hom}_{\mathcal{A}^{e}}(\mathcal{A},\mathcal{A})
\;\simeq\;
\mathrm{Hom}_{\mathcal{A}^{e}}
\bigl(\operatorname{Bar}_{\mathcal{A}^{e}}(\mathcal{A}),\mathcal{A}\bigr).
\]
The algebra quasi-isomorphism
$\Omega^{\mathrm{ch}}\bar{B}^{\mathrm{ch}}(\mathcal{A})\simeq
\mathcal{A}$ is the bar--cobar inversion used to build this
resolution. The Koszul dual algebra remains the separate Verdier
object~$\mathcal{A}^{!}$.
\end{theorem}

\begin{proof}
The bar construction $\bar{B}^{\mathrm{ch}}(\mathcal{A})$ is a
cofree conilpotent chiral coalgebra. Tensoring it with the left and
right regular $\mathcal{A}$-modules gives the two-sided bar object
$\operatorname{Bar}_{\mathcal{A}^{e}}(\mathcal{A})$. Its differential
is the sum of the internal differential, the chiral residue
differential, and the two endpoint module actions. The standard
augmentation
\[
\operatorname{Bar}_{\mathcal{A}^{e}}(\mathcal{A})
\longrightarrow \mathcal{A}
\]
is the diagonal-bimodule version of the bar--cobar counit. On the
PBW Koszul locus the associated graded is the quadratic Koszul
resolution, so the augmentation is a quasi-isomorphism.

Applying $\mathrm{Hom}_{\mathcal{A}^{e}}(-,\mathcal{A})$ computes
$R\mathrm{Hom}_{\mathcal{A}^{e}}(\mathcal{A},\mathcal{A})$. The
Fulton--MacPherson realization of this Hom complex is precisely
$C^{\bullet}_{\mathrm{ch}}(\mathcal{A},\mathcal{A})$: an $n$-cochain
is a diagonal bimodule functional on the $(n+2)$-point two-sided bar
configuration, with the three differential pieces of
Theorem~\ref{thm:chiral-hochschild-differential}. Thus the displayed
quasi-isomorphisms compute the chiral derived centre and keep
$\barB^{\mathrm{ch}}(\mathcal{A})$, $\mathcal{A}^{i}$, and
$\mathcal{A}^{!}$ out of the Hochschild-cochain object.
\end{proof}

\subsection{The chiral Hochschild spectral sequence}

\begin{theorem}[chiral Hochschild spectral sequence; \ClaimStatusProvedHere]\label{thm:hochschild-spectral-sequence}
There exists a spectral sequence
\[
E_1^{n,q} = H^q(\overline{C}_{n+2}(X), \mathcal{A}^{\boxtimes (n+2)} \otimes \Omega^n_{\log}) \Rightarrow \ChirHoch^{n+q}(\mathcal{A})
\]
where $n$ is the bar degree (number of internal insertions) and $q$ is the sheaf cohomology degree, which are independent indices.
\end{theorem}

\begin{proof}
The chiral Hochschild complex $\ChirHoch^*(\mathcal{A})$ is a double complex with horizontal differential $d_{\mathrm{bar}}$ (the bar differential of Definition~\ref{def:bar-differential-complete}) and vertical differential $d_{\mathrm{sheaf}}$ (sheaf cohomology on configuration spaces). Filter by bar degree: $F^n = \bigoplus_{m \geq n} \ChirHoch^{m,*}$. The resulting spectral sequence has $E_0^{n,q} = \Gamma(\overline{C}_{n+2}(X), \mathcal{A}^{\boxtimes(n+2)} \otimes \Omega^n_{\log})$ with $d_0 = d_{\mathrm{sheaf}}$, so $E_1^{n,q} = H^q(\overline{C}_{n+2}(X), \mathcal{A}^{\boxtimes(n+2)} \otimes \Omega^n_{\log})$. Convergence follows from the bounded-below hypothesis on $\mathcal{A}$ and the finite cohomological dimension of configuration spaces.
\end{proof}

For the cyclic refinement relating chiral Hochschild to cyclic homology via Connes' operator, see Theorem~\ref{thm:HC-spectral-sequence}. For formal chiral algebras (quasi-isomorphic to their cohomology), the spectral sequence above degenerates at $E_1$, giving:
\[
\ChirHoch^k(\mathcal{A}) \cong \bigoplus_{n+q=k} H^q(\overline{C}_{n+2}(X), \mathcal{A}^{\boxtimes (n+2)} \otimes \Omega^n_{\log})
\]

\section{Koszul duality for chiral algebras}

\subsection{Quadratic chiral algebras and their duals}

\begin{definition}[Quadratic chiral algebra]
A chiral algebra $\mathcal{A}$ is \emph{quadratic} if it admits a presentation
\[
\mathcal{A} = T_{\text{chiral}}(\mathcal{V})/(R)
\]
where:
\begin{itemize}
\item $\mathcal{V}$ is a locally free $\mathcal{O}_X$-module of generators
\item $T_{\text{chiral}}(\mathcal{V})$ is the free chiral algebra on $\mathcal{V}$
\item $R \subset j_*j^*(\mathcal{V} \boxtimes \mathcal{V})$ consists of quadratic relations
\end{itemize}
\end{definition}

The free chiral algebra requires care. Following Beilinson--Drinfeld:

\begin{definition}[Free chiral algebra]
The free chiral algebra on $\mathcal{V}$ is
\[
T_{\text{chiral}}(\mathcal{V}) = \bigoplus_{n \geq 0} \pi_{n*}\left(j_*j^*\mathcal{V}^{\boxtimes n} \otimes \mathcal{D}_{C_n(X)/X}\right)^{\Sigma_n}
\]
where $\pi_n: C_n(X) \to X$ is the projection and $\mathcal{D}_{C_n(X)/X}$ denotes relative differential operators.
\end{definition}

\begin{definition}[Koszul dual]
The Koszul dual of a quadratic chiral algebra $\mathcal{A}$ is
\[
\mathcal{A}^! = T_{\text{chiral}}(\mathcal{V}^*)/(R^{\perp})
\]
where:
\begin{itemize}
\item $\mathcal{V}^* = \mathcal{H}om_{\mathcal{O}_X}(\mathcal{V}, \omega_X)$ is the dual shifted by the canonical bundle
\item $R^{\perp}$ consists of relations orthogonal to $R$ under the canonical pairing
\[
\langle \cdot, \cdot \rangle: j_*j^*(\mathcal{V}^* \boxtimes \mathcal{V}) \to j_*\omega_{X^2 \setminus \Delta}
\]
\end{itemize}
\end{definition}

\begin{remark}[Typed Koszul objects]\label{rem:koszul-dual-meaning}
Fix a quadratic presentation
$\cA=T_{\mathrm{ch}}(\mathcal V)/(R)$.  It determines the
presentation coalgebra and its canonical comparison
\[
 \cA^i:=C_X(s^{-1}\mathcal V,s^{-2}R),
 \qquad
 q_\cA\colon\cA^i\longrightarrow\bar B^{\mathrm{ch}}_X(\cA).
\]
The target is the full conilpotent dg bar coalgebra.  Its cohomology is
$H^\bullet(\bar B^{\mathrm{ch}}_X(\cA))$.  On the Koszul comparison
surface, $q_\cA$ is a quasi-isomorphism, so
$H^\bullet(q_\cA)$ identifies the cohomology of the presentation
coalgebra with bar cohomology.

Verdier duality forms the bar-dual algebra
\[
 K_X(\cA):=\mathbb D_{\Ran}\bar B^{\mathrm{ch}}_X(\cA),
 \qquad
 \mathbb D(q_\cA)\colon
 K_X(\cA)\longrightarrow\mathbb D_{\Ran}(\cA^i).
\]
A chosen chiral Koszul-pair structure supplies an algebra comparison
$\nu_\cA\colon K_X(\cA)\to\cA^!$; on a finite-type quadratic lane it
agrees with the composite through
$\mathbb D_{\Ran}(\cA^i)\simeq T_{\mathrm{ch}}(\mathcal V^*)/(R^\perp)$.
The full bar coalgebra also carries the independent inversion counit
$\Omega(\bar B^{\mathrm{ch}}_X(\cA))\xrightarrow{\sim}\cA$
of Theorem~\ref{thm:bar-cobar-inversion-qi}.  The typed comparison
diagram is therefore
\[
\cA^i\xrightarrow{\ q_\cA\ }\bar B^{\mathrm{ch}}_X(\cA),
\qquad
K_X(\cA)\xrightarrow{\ \mathbb D(q_\cA)\ }
\mathbb D_{\Ran}(\cA^i),
\qquad
K_X(\cA)\xrightarrow{\ \nu_\cA\ }\cA^!,
\]
with $H^\bullet(q_\cA)$ recording quadratic recognition.
\end{remark}

\begin{example}[Explicit correspondence for Heisenberg]
For the Heisenberg chiral algebra $\mathcal{H}_k$ at level~$k$ with generator $\alpha(z)$ and OPE:
\[\alpha(z)\alpha(w) \sim \frac{k}{(z-w)^2}\]

The Koszul dual $\mathcal{H}^!$ has:
\begin{itemize}
\item \emph{Dual generator}: $\alpha^*(z)$ with $\langle \alpha, \alpha^* \rangle = 1$ under residue pairing
\item \emph{Coproduct}:
 \[\Delta(\alpha^*) = \alpha^* \otimes 1 + 1 \otimes \alpha^* + \text{(higher order terms)}\]
 encoding the dual of the commutative algebra structure
\item \emph{Bar construction}: $\bar{B}^{\text{ch}}(\mathcal{H})$ consists of forms like:
 \[\alpha(z_1) \otimes \cdots \otimes \alpha(z_n) \otimes \eta_{12} \wedge \eta_{23} \wedge \cdots\]
 whose residues extract the coproduct coefficients
\end{itemize}

The cobar of the bar recovers $\mathcal{H}$ itself: $\Omega^{\text{ch}}(\bar{B}^{\text{ch}}(\mathcal{H})) \simeq \mathcal{H}$ (Theorem~\ref{thm:bar-cobar-inversion-qi}).
\end{example}

\begin{remark}[Koszul dual versus dual]
``Koszul dual'' emphasizes the derived character
(bar-cobar construction, quasi-isomorphisms, involutivity:
$(\mathcal{A}^!)^! \simeq \mathcal{A}$). For quadratic $\mathcal{A}$,
this recovers the classical quadratic dual.
\end{remark}

\subsection{The universal twisting morphism}

The universal twisting morphism
$\pi\colon \barB_X(\cA) \to \cA$
(Proposition~\ref{prop:universal-twisting-adjunction}) controls
all bar-cobar identifications in this chapter. Twisting morphisms
are MC elements in the convolution algebra $\mathrm{Hom}(C, A)$
(Definition~\ref{def:twisting-morphism}). The natural bijection
$\mathrm{Tw}(C, A) \cong \mathrm{Hom}(C, B(A)) \cong
\mathrm{Hom}(\Omega(C), A)$
mediates the bar-cobar adjunction
(Proposition~\ref{prop:universal-twisting-adjunction}(iii)); for the
geometric realization as a propagator integral, see
Remark~\ref{rem:twisting-as-propagator}.

\begin{example}[Twisting for fermion-boson duality]
\label{ex:fermion-boson-twisting}
For the Koszul pair (free fermions $\mathcal{F}$, $\beta\gamma$
system), the universal twisting morphism
$\pi\colon \barB^{\mathrm{ch}}(\mathcal{F}) \to \mathcal{F}$
takes the explicit form
$\pi(\psi^*)(z) = \operatorname{Res}_{w=z}\, \psi(w) \cdot \frac{dw}{z-w}$,
implementing the fermion-boson correspondence at the MC level
(Remark~\ref{rem:twisting-as-propagator}).
\end{example}

\subsection{Chain pairings between Koszul-dual charts}
\label{subsec:hochschild-duality}

\begin{definition}[Shifted chain pairing;
\ClaimStatusDefinitional]
\label{def:theorem-h-shift-normalization-package}
\index{Theorem H!chain pairing}
Let \(\mathscr D\) be the stable symmetric monoidal category in which
the completed chart cochains live, and let \(\omega_X\) be its
dualizing object.  A \emph{degree-\(d\) chain pairing} for a Koszul
pair \((\cA,\cA^!)\) is a morphism
\[
 \langle-,-\rangle_\cA\colon
 Q_\cA\otimes^{L}Q_{\cA^!}
 \longrightarrow \omega_X[d]
\]
in \(\mathscr D\).  Its adjoint is
\begin{equation}\label{eq:theorem-h-pairing-adjoint}
 \Phi_\cA\colon Q_\cA
 \longrightarrow
 R\operatorname{Hom}_{\mathscr D}
 \bigl(Q_{\cA^!},\omega_X[d]\bigr).
\end{equation}
The pairing is \emph{perfect} when \(\Phi_\cA\) is a
quasi-isomorphism.  The complex
\[
 \mathfrak D_\cA:=\operatorname{Cone}(\Phi_\cA)
\]
records the pairing cohomology.  A geometric construction of the
shift \(d\) specifies the Verdier convention, the collision
restriction maps, the fibre degree, and the totalization grading in
the same complex.  A shift \(d=2\) is therefore a theorem about this
pairing datum.
\end{definition}

\begin{lemma}[Perfect-pairing criterion; \ClaimStatusProvedHere]
\label{lem:hochschild-shift-computation}
\index{Hochschild cohomology!duality shift computation}
Let \((\cA,\cA^!)\) carry a degree-\(d\) chain pairing as in
Definition~\ref{def:theorem-h-shift-normalization-package}.  The
following conditions are equivalent:
\begin{enumerate}[label=\textup{(\roman*)}]
\item the pairing is perfect;
\item the adjoint \(\Phi_\cA\) of
\eqref{eq:theorem-h-pairing-adjoint} is a quasi-isomorphism;
\item the complex \(\mathfrak D_\cA=\operatorname{Cone}(\Phi_\cA)\)
is acyclic.
\end{enumerate}
Under these conditions there is an equivalence in \(\mathscr D\)
\[
 Q_\cA\simeq
 R\operatorname{Hom}_{\mathscr D}
 \bigl(Q_{\cA^!},\omega_X[d]\bigr).
\]
If the two complexes are perfect, its cohomology form is
\[
 H^n(Q_\cA)
 \cong
 H^{d-n}(Q_{\cA^!})^\vee\otimes\omega_X.
\]
\end{lemma}

\begin{proof}
The first two conditions agree by the definition of perfectness.  A
map of complexes is a quasi-isomorphism precisely when its cone is
acyclic, which gives the third.  The displayed equivalence is the
adjoint map itself.  Perfectness identifies derived duals with
ordinary duals on cohomology, and the shift~\([d]\) sends degree~\(n\)
to degree~\(d-n\).
\end{proof}

\begin{lemma}[Support under a convergent filtration;
\ClaimStatusProvedHere]
\label{lem:totalization-amplitude}
\index{Hochschild!totalization amplitude}
Let \((Q,F)\) be a complete, exhaustive filtered cochain complex whose
spectral sequence converges strongly to \(H^\bullet(Q)\).  Fix a set
\(S\subset\mathbb Z\), and assume
\[
 E_1^{p,q}(Q,F)=0
 \qquad\text{for }p+q\in\mathbb Z\setminus S.
\]
Then
\[
 \operatorname{Supp}H^\bullet(Q)\subseteq S.
\]
For a tower \(Q\simeq R\!\varprojlim_NQ_N\), the same conclusion
holds when every finite-window spectral sequence has this support and
\(\varprojlim_N^1H^{m-1}(Q_N)=0\) for every~\(m\).
\end{lemma}

\begin{proof}
Every page \(E_r^{p,q}\) is a subquotient of
\(E_1^{p,q}\).  Hence each page is supported on bidegrees whose total
degree lies in~\(S\).  Strong convergence identifies
\(E_\infty^{p,q}\) with the associated graded of
\(H^{p+q}(Q)\), proving the support assertion.  For an inverse-limit
tower, the Milnor exact sequence identifies the completed cohomology
with \(\varprojlim_NH^m(Q_N)\) under the displayed
\(\varprojlim^1\)-hypothesis, and the support is preserved.
\end{proof}

\begin{remark}[From local complexes to global support]
\label{rem:totalization-amplitude-non-triviality}
The local collision complexes assemble through their incidence maps.
A natural system of strong deformation retracts assembles to a global
retract by Proposition~\ref{prop:fm-tower-collapse}.  Alternatively,
a convergent collision filtration determines the global support from
its \(E_1\)-page by Lemma~\ref{lem:totalization-amplitude}.  PBW
Koszulness describes an associated-graded complex; the incidence
homotopies supply the chain-level passage to the chiral Hochschild
complex.
\end{remark}

\begin{lemma}[Quadratic-chiral Koszul transfer; \ClaimStatusConditional]
\label{lem:chiral-quadratic-koszul}
\index{Koszul duality!quadratic chiral transfer|textbf}
\index{Fresse!quadratic operadic Koszul}
Let $\cP$ be a quadratic chiral operad on the smooth curve~$X$
with generators in degree~$2$ and relations in degree~$3$, and
let $\cP^{!}$ denote its chiral operadic Koszul dual
\textup{(}constructed by $\cD_X$-linear duality on the space
of relations, with the chiral pairing induced by OPE residue
at the principal collision\textup{)}. Then there is a
degree-preserving chain map
\[
\sigma \;\colon\; \cP^{!}
\;\longrightarrow\;
T^{c}\!\bigl(s^{-1}\,\overline{\cP}\bigr)^{\vee},
\]
from $\cP^{!}$ into the linear dual of the cofree tensor
coalgebra on the desuspended augmentation ideal
$\overline{\cP} = \ker(\epsilon_\cP)$, intertwining the
operadic Koszul differential on the source with the bar
differential on the target. In equations,
\[
\sigma \circ d_{\mathrm{K}}
\;=\;
d_{\bar{B}}^{\vee} \circ \sigma.
\]
The map $\sigma$ is an isomorphism on the degree-$2$ generator
piece, and its induced map on cohomology is a quasi-isomorphism
whenever $\cP$ is chiral Koszul.
\end{lemma}

\begin{proof}
We transport the quadratic Koszul construction of
Loday--Vallette \cite[Ch.~7]{LV12} from the category of dg
symmetric operads over a field to the category of chiral
operads on~$X$ along the OPE-residue pairing. The transport
requires one non-tautological input, recorded in
Step~1 below: the OPE-residue pairing on the weight-$3$ free
operad piece is non-degenerate on the relation space~$R$.
All subsequent steps are formal consequences of Fresse's construction
once non-degeneracy is in hand.

\smallskip
\emph{Step 1 (non-degeneracy of the OPE-residue pairing on
$\cF(E)(3)$).}
We first prove that the chiral pairing
$\langle -,-\rangle_{\mathrm{OPE}}$, obtained by iterated OPE
residues at the principal collision on
$\overline{C}_{3}(X)$, is non-degenerate on the free-operad
degree-$3$ piece $\cF(E)(3)$. The argument is local: by
$\cD_{X^2}$-flatness of~$E$ and the smoothness of~$X$, the
free-operad piece $\cF(E)(3)$ is a locally free
$\cD_{X^3}$-module whose fibre at a generic point $p \in X$
is the classical free-operad piece on the fibre
$E_{p} \in \Vect$. Under this fibrewise identification the
OPE-residue pairing restricts to the classical composition
pairing of quadratic operadic Koszul duality
\cite[Ch.~7]{LV12}, which is non-degenerate on
$\cF(E_{p})(3)$. Since matrix rank is lower semicontinuous,
non-degeneracy at a generic point, together
with $\cD_{X^3}$-coherence, holds on a Zariski-open dense
locus.  Along the collision divisors the rank is determined by the
specialisation, which
factors through the local model $\FM_{3}(\C)$, on whose real
cohomology the induced pairing descends to the classical
composition pairing on the Arnold algebra
$H^*(\FM_3(\C);\R)$ via the cohomological
$\En$-formality quasi-isomorphism
\textup{(}Proposition~\textup{\ref{prop:en-formality}}\textup{;}
used here as a cohomology-level identification\textup{)}, and is
non-degenerate on that
Arnold algebra by the quadratic-duality pairing
\cite[Ch.~7]{LV12} applied to the degree-$3$ piece.
Thus $R^{\perp} \subset \cF(E)(3)$ is a well-defined
$\cD_{X^3}$-submodule of complementary rank, and the chiral
Koszul dual $\cP^{!}$ is well defined.

\smallskip
\emph{Step 2 (Fresse transport).}
The quadratic datum is a triple $(E,R,\gamma)$ with
$E \in \operatorname{Mod}(\cD_{X^2})$ the generating
space \textup{(}chiral degree~$2$\textup{)},
$R \subset \cF(E)(3)$ the relation space in free-operad
degree~$3$, and $\gamma$ the structural map. The chiral Koszul
dual has generator space $E^{\vee}[1]$ and relation space
$R^{\perp}$ as constructed in Step~1. The quadratic-duality construction
\cite[Ch.~7]{LV12} produces, for any
such quadratic datum over a base symmetric-monoidal category,
a canonical chain map
$\sigma_{\mathrm{Fresse}} \colon (P)^{!} \to
B(P)^{\vee}$, where $B$ is the operadic bar
construction. The transfer to the chiral setting replaces
$B$ by the chiral bar $\bar{B}^{\mathrm{ch}}$ and the
composition pairing by the OPE-residue pairing of Step~1.
The compatibility of $d_{\mathrm{K}}$ with $d_{\bar{B}}^{\vee}$
then holds because both sides are computed by the same
collision-residue formula
\textup{(}Theorem~\textup{\ref{thm:residue-formula}}\textup{)}.

On the degree-$2$ generator piece, $\sigma$ restricts to the
canonical identification $E^{\vee} \cong (s^{-1}\overline{\cP})^{\vee}_1$,
which is a degree-preserving isomorphism by construction.
When $\cP$ is chiral Koszul, the bar cohomology
$H^{*}(\bar{B}^{\mathrm{ch}}(\cP))$ is concentrated in
tensor degree~$1$ \textup{(}by the definition of chiral
Koszulness, equivalent
\textup{(}i\textup{)}\textup{)}, and the Fresse transfer
identifies this concentrated cohomology with $\cP^{!}$ via
$\sigma$. Hence $\sigma$ is a quasi-isomorphism whenever
$\cP$ is chiral Koszul, which is the statement of the lemma.
\end{proof}

\begin{proposition}[Incidence-compatible collision retracts;
\ClaimStatusProvedHere]
\label{prop:fm-tower-collapse}
\index{configuration space!collapse|textbf}
\index{Fulton--MacPherson!collision-depth spectral sequence}
Fix a finite weight window~\(N\).  Assume a
finite-window collision realization in the sense of
Definition~\ref{def:theorem-h-pbw-finite-window-lane}.  Then the
incidence totalizations carry maps
\[
 K_{\cA,N}
 \mathop{\xrightarrow{\ \iota_N\ }}
 Q^{\mathrm{ord}}_{\cA,N}
 \mathop{\xrightarrow{\ p_N\ }}
 K_{\cA,N},
 \qquad
 h_N\colon Q^{\mathrm{ord}}_{\cA,N}
 \longrightarrow Q^{\mathrm{ord}}_{\cA,N}[-1],
\]
where on total incidence--bar degree~\(r\) the homotopy is
\(h_N=(-1)^r h_{\sigma,N}\).  These maps satisfy
\begin{equation}\label{eq:fm-incidence-retract}
 p_N\iota_N=\id,
 \qquad
 d h_N+h_Nd=\id-\iota_Np_N.
\end{equation}
Compatibility with \(N+1\to N\) gives the same identities after
derived inverse limit.  If
\[
 K_{\cA,S}\simeq R\!\varprojlim_NK_{\cA,N}
\]
is supported in a finite set~\(S\), and if the averaging and chart
comparison maps of
Definition~\ref{def:theorem-h-collision-depth-package} are
quasi-isomorphisms, then these maps define \(H_H(\cA;S)\).

\end{proposition}

\begin{proof}
Naturality for the incidence category says that
\(\iota_{\sigma,N}\), \(p_{\sigma,N}\), and
\(h_{\sigma,N}\) are natural transformations of diagrams.  The
incidence--bar total differential is \(D=\delta+(-1)^r d\) in total
incidence--bar degree~\(r\).  The choice
\(h_N=(-1)^rh_{\sigma,N}\) makes the two
mixed terms \(\delta h_N+h_N\delta\) cancel by naturality, while the
vertical terms give \(dh_{\sigma,N}+h_{\sigma,N}d\).  This proves
\eqref{eq:fm-incidence-retract}.  Compatibility with the transition
maps makes these transformations maps of inverse systems, so derived
inverse limit preserves the identities.  The averaging and chart
comparison quasi-isomorphisms transport the resulting retract to
\(C^\bullet_{\mathrm{ch}}(\cA,\cA)\).  Its target is
\(K_{\cA,S}\), hence the resulting data are precisely
\(H_H(\cA;S)\).

\end{proof}

\begin{remark}[Quadratic data and collision homotopies]
\label{rem:pbw-parallel-consequences}
\index{PBW!serial consequences}
PBW Koszulness identifies the associated-graded residue complex, and
Orlik--Solomon Koszulness provides the contraction of its quadratic
Koszul factor.  A family calculation then lifts this contraction to
maps \((\iota_{\sigma,N},p_{\sigma,N},h_{\sigma,N})\) on the
residue-twisted coefficient complex and verifies their incidence and
completion compatibilities.  These maps form the exact input of
Proposition~\ref{prop:fm-tower-collapse}:
\[
 \text{quadratic contraction}
 \longrightarrow
 \text{residue-twisted stratum retracts}
 \longrightarrow
 \text{incidence totalization}.
\]
\end{remark}

\begin{proposition}[Support from finite-window collision retracts;
\ClaimStatusProvedHere]
\label{prop:theorem-h-finite-window-collision-retracts}
\index{Theorem H!finite-window collision retracts}
\textup{[Type signature: Open quadrant; chain presentation;
Beilinson level~$3$; hypothesis package: incidence-compatible
finite-window retracts, averaging, chart comparison, and
Mittag--Leffler completion.]}
Fix a finite set \(S\subset\mathbb Z\).  Assume that every bar degree
and every conformal-weight window carries a collision realization from
Definition~\ref{def:theorem-h-pbw-finite-window-lane}, and that the
incidence totalization \(K_{\cA,N}\) is supported in~\(S\).  Assume
also the quasi-isomorphisms
\[
 \bigl(Q^{\mathrm{ord}}_{\cA,N}\bigr)_{h\Sigma}
 \xrightarrow{\ \sim\ }Q^{\mathrm{sym}}_{\cA,N},
 \qquad
 R\!\varprojlim_NQ^{\mathrm{sym}}_{\cA,N}
 \xrightarrow{\ \sim\ }C^\bullet_{\mathrm{ch}}(\cA,\cA),
\]
and the completion identities
\[
 K_{\cA,S}\simeq R\!\varprojlim_NK_{\cA,N},
 \qquad
 \varprojlim_N^1H^{m-1}(Q^{\mathrm{sym}}_{\cA,N})=0.
\]
Then the collision maps define \(H_H(\cA;S)\), and
\[
 H^\bullet(K_{\cA,S})\cong\ChirHoch^\bullet(\cA),
 \qquad
 \operatorname{Supp}\ChirHoch^\bullet(\cA)\subseteq S.
\]
\end{proposition}

\begin{proof}
Proposition~\ref{prop:fm-tower-collapse} assembles the stratum
retracts and carries them through completion and averaging.  The
result is the family support datum \(H_H(\cA;S)\).
Proposition~\ref{prop:theorem-h-kdh-criterion} identifies its target
cohomology with \(\ChirHoch^\bullet(\cA)\) and gives the support
inclusion.
\end{proof}

\begin{proposition}[Hilbert series of a family support model;
\ClaimStatusProvedHere]
\label{prop:chirhoch-sharp-hilbert}
\index{chiral Hochschild!bigraded Hilbert series}
\index{Russian school!sharp Hilbert series}
Assume \(H_H(\cA;S)\) and finite-dimensional graded pieces of
\(H^n(K_{\cA,S})\).  Then the bigraded Hilbert series is
\begin{equation}\label{eq:chirhoch-sharp-hilbert}
 P_{\cA}(t,q)
 \;=\;
 \sum_{n\in S}\dim_qH^n(K_{\cA,S})\,t^n.
\end{equation}
For finite~\(S\), this is a Laurent polynomial in~\(t\).  If
\(S\subset\mathbb Z_{\geq0}\), it is a polynomial.  A perfect
degree-\(d\) pairing from
Definition~\ref{def:theorem-h-shift-normalization-package}, together
with a family support datum \(H_H(\cA^!;S^!)\) having
finite-dimensional graded pieces, adds the
palindromic relation
\[
 P_\cA(t,q)=t^dP_{\cA^!}(t^{-1},q)
\]
whenever the duality preserves the \(q\)-grading.
\end{proposition}

\begin{proof}
Proposition~\ref{prop:theorem-h-kdh-criterion} identifies
\(H^n(K_{\cA,S})\) with \(\ChirHoch^n(\cA)\).  Taking graded
dimensions gives the first formula.  A perfect degree-\(d\) pairing
identifies degree~\(n\) for \(\cA\) with the dual of degree~\(d-n\)
for \(\cA^!\), which gives the palindromic identity.
\end{proof}

\begin{lemma}[Curved-dual oscillator criterion;
\ClaimStatusConditional]
\label{lem:curved-dual-centre-heisenberg}
\index{curved Koszul dual!Heisenberg}
Let \(k\neq0\).  Assume the quadratic comparison for
\(\mathfrak H_k\) identifies its curved Koszul dual with
\[
 \mathfrak H_k^!
 \simeq
 \bigl(\mathrm{Sym}^{\mathrm{ch}}(V^*[1]),m_0=-k\omega\bigr).
\]
Assume further that, in every conformal-weight window \(N\), the
completed second-kind Hochschild complex admits a transition-compatible
quasi-isomorphism
\[
 \Xi_N\colon
 C_{\mathrm{Hoch},\leq N}^{\mathrm{II}}(\mathfrak H_k^!)
 \xrightarrow{\ \sim\ }
 \C\mathbf1\otimes
 \bigotimes_{n=1}^{N}
 \left[
  \C e_n\xrightarrow{\ -kn\ }\C f_n
 \right].
\]
Then each positive-weight summand is contractible, the inverse system
is strict Mittag--Leffler, and
\[
 H^0\!\left(
 C_{\mathrm{Hoch}}^{\mathrm{II}}(\mathfrak H_k^!)
 \right)\cong\C\mathbf1,
 \qquad
 R^1\!\varprojlim_N
 H^\bullet\!\left(
 C_{\mathrm{Hoch},\leq N}^{\mathrm{II}}(\mathfrak H_k^!)
 \right)=0.
\]
\end{lemma}

\begin{proof}
For \(1\leq n\leq N\), multiplication by \(-kn\) is invertible.  The
homotopy
\[
 h_n(f_n)=(-kn)^{-1}e_n,\qquad h_n(e_n)=0
\]
contracts the \(n\)-th two-term factor.  Tensoring these contractions
leaves the vacuum line.  The transition \(N+1\to N\) removes the last
contractible factor and is degreewise surjective.  The images
therefore stabilize in every cohomological degree, giving the
Mittag--Leffler assertion and the displayed inverse-limit cohomology.
\end{proof}

\begin{corollary}[Rank-one even superboson bounded cohomology;
\ClaimStatusProvedElsewhere]
\label{cor:chirhoch-heisenberg}
\index{Hochschild cohomology!Heisenberg}
Bakalov--De Sole--Kac compute the bounded vertex cohomology of the
rank-one even superboson \(B_{\mathfrak h}\) as
\[
 \dim H^n_{\mathrm{ch},b}(B_{\mathfrak h},B_{\mathfrak h})
 =(2,1,0,0,\ldots).
\]
If \(\chi_{B_{\mathfrak h}}^{\mathrm{bd}}\) is a quasi-isomorphism,
the same dimensions hold for
\(\ChirHoch^n(\cA_{B_{\mathfrak h}})\).
\end{corollary}

\begin{proof}
By \cite[Theorem~7.4]{BDSK21},
\[
 H^n_{\mathrm{ch},b}(B_{\mathfrak h})
 \cong
 \bigl(S^n(\Pi\mathfrak h)\bigr)^*
 \oplus
 \bigl(S^{n+1}(\Pi\mathfrak h)\bigr)^*.
\]
For a one-dimensional even \(\mathfrak h\), the parity shift
\(\Pi\mathfrak h\) is an odd line.  Its super-symmetric algebra is
the exterior algebra on that line, which gives dimensions~\(2\) in
degree zero and~\(1\) in degree one.  The comparison clause follows
from invariance of cohomology under quasi-isomorphism.
\end{proof}

\begin{remark}[Family parameter versus fixed fibre]
\label{rem:heisenberg-family-vs-fixed-fiber}
Corollary~\ref{cor:chirhoch-heisenberg} concerns the bounded cochain
complex of a fixed vertex algebra.  The parameter~\(k\) belongs to the
base of the Heisenberg family.  Over \(k\neq0\), the change of
generators \(\alpha\mapsto k^{-1/2}\alpha\) identifies the fibres
after choosing a square root.  The family parameter and the fixed-fibre
cohomology are recorded by separate complexes.
\end{remark}

\begin{proposition}[Heisenberg finite-window complexes and their
inverse limit; \ClaimStatusProvedHere]
\label{prop:heisenberg-theorem-h-window-limit}
\index{Theorem H!Heisenberg finite-window inverse limit}
\index{Heisenberg!Theorem H finite-window inverse limit}
\textup{[Type signature: Open quadrant; diagonal chiral Hochschild
chain presentation; Beilinson level~$3$; hypothesis package:
rank-one Heisenberg algebra at $k\neq0$ and finite-dimensional graded
realization.]}
Let \(k\neq0\), and let \(\mathfrak H_k\) be the rank-one
Heisenberg chiral algebra.  In the conformal-weight window
\(\leq N\), write
\[
F_{\leq N}\mathfrak H_k
=
\operatorname{span}\{
\alpha_{-m_1}\cdots\alpha_{-m_r}\mathbf 1
\mid
m_i\geq1,\ \sum_i m_i\leq N
\}.
\]
Let
\[
 C_N^\bullet(\mathfrak H_k)
 :=RHH_{\mathrm{ch}}(\mathfrak H_k)/
 F_{>N}^{\mathrm{wt}}RHH_{\mathrm{ch}}(\mathfrak H_k),
\]
and let
\[
 q_{N+1,N}\colon C_{N+1}^\bullet(\mathfrak H_k)
 \longrightarrow C_N^\bullet(\mathfrak H_k)
\]
be the chain map that projects away conformal weight~$N+1$.
For $M\geq N$, write
$q_{M,N}:=q_{N+1,N}\circ\cdots\circ q_{M,M-1}$.
Apply a finite-dimensional graded realization to the perfect
holonomic coefficient complexes.  Then:
\begin{enumerate}[label=\textup{(\roman*)}]
\item The vector space \(F_{\leq N}\mathfrak H_k\) is finite
 dimensional, with basis indexed by partitions of integers
 \(0,\ldots,N\).
\item Every normalized bar summand in the weight-$\leq N$ window has
 at most~$N$ non-vacuum inputs, because each normalized input has
 conformal weight at least~$1$.  Hence the diagonal finite-window
 Hochschild complex has bar length at most~$N$ and is a finite direct
 sum of perfect coefficient complexes.
\item Each $H^n(C_N^\bullet(\mathfrak H_k))$ is finite dimensional.
 For fixed $n$ and~$N$, the subspaces
 \[
  \operatorname{im}\!\left(
   H^n(C_M^\bullet(\mathfrak H_k))
   \xrightarrow{\,H^n(q_{M,N})\,}
   H^n(C_N^\bullet(\mathfrak H_k))
  \right),
  \qquad M\geq N,
 \]
 form a descending chain and therefore stabilize.  Consequently
 \[
  \varprojlim\nolimits_N^1
  H^{n-1}(C_N^\bullet(\mathfrak H_k))=0
  \qquad\text{for every }n.
 \]
\end{enumerate}
Thus the Heisenberg weight tower satisfies the degreewise
Mittag--Leffler hypothesis in the Milnor sequence, and the completed
definition gives
\[
 H^n\!\left(RHH_{\mathrm{ch}}(\mathfrak H_k)\right)
 \cong
 H^n\!\left(R\!\varprojlim_N C_N^\bullet(\mathfrak H_k)\right)
 \cong
 \varprojlim_N H^n(C_N^\bullet(\mathfrak H_k)).
\]
This establishes the finite-window and Mittag--Leffler parts of a
family support datum.  A chart support theorem additionally uses the
incidence-compatible collision retracts and the bounded-to-chart
comparison \(\chi_{B_{\mathfrak h}}^{\mathrm{bd}}\).
\end{proposition}

\begin{proof}
The Fock basis of the rank-one Heisenberg algebra is the usual
partition basis.  The number of monomials of total conformal weight at
most \(N\) is \(\sum_{m=0}^{N}p(m)\), hence finite.  The normalized
augmentation ideal begins in conformal weight~$1$, so a bar word with
\(p\) normalized entries has total weight at least \(p\); this proves
the bar-length bound.  Perfectness of the finite-window holonomic
coefficient complexes gives finite-dimensional realized cohomology.

The Hochschild differential preserves conformal weight, hence each
$q_{N+1,N}$ is a chain map.  For fixed $n$ and~$N$, the images of
$H^n(q_{M,N})$ lie in the finite-dimensional vector space
$H^n(C_N^\bullet)$ and decrease with~$M$.  Their stabilization is the
Mittag--Leffler condition, which gives the displayed vanishing of
$\varprojlim^1$.  This calculation supplies the completion component
of \(H_H(\mathfrak H_k;S)\).  The stratum homotopies and comparison
map supply its support component.
\end{proof}

\begin{corollary}[Affine $\mathfrak{sl}_{2}$ bounded-cohomology bound;
\ClaimStatusConjectured]
\label{cor:chirhoch-affine-sl2}
\index{Hochschild cohomology!affine sl_2}
Bakalov--De Sole--Kac conjecture, for
\(k\in\mathbb C\setminus\{-2\}\), the bounds
\[
 \dim H^n_{\mathrm{ch},b}
 \bigl(V_k(\mathfrak{sl}_2),V_k(\mathfrak{sl}_2)\bigr)
 \leq
 \dim\!\left(
 \bigwedge^n\mathfrak{sl}_2\oplus
 \bigwedge^{n+1}\mathfrak{sl}_2
 \right).
\]
A bounded-to-chart quasi-isomorphism transports these conjectural
bounds to \(\ChirHoch^n(\cA_{V_k(\mathfrak{sl}_2)})\).
\end{corollary}

\begin{proof}
This is the specialization \(\fg=\mathfrak{sl}_2\) of
\cite[Conjecture~7.5]{BDSK21}.  The final clause is transport along
\(\chi_{V_k(\mathfrak{sl}_2)}^{\mathrm{bd}}\).
\end{proof}

\begin{corollary}[Affine $\mathfrak{sl}_{3}$ bounded-cohomology bound;
\ClaimStatusConjectured]
\label{cor:chirhoch-affine-sl3}
\index{Hochschild cohomology!affine sl_3}
Bakalov--De Sole--Kac conjecture, for
\(k\in\mathbb C\setminus\{-3\}\), the bounds
\[
 \dim H^n_{\mathrm{ch},b}
 \bigl(V_k(\mathfrak{sl}_3),V_k(\mathfrak{sl}_3)\bigr)
 \leq
 \dim\!\left(
 \bigwedge^n\mathfrak{sl}_3\oplus
 \bigwedge^{n+1}\mathfrak{sl}_3
 \right).
\]
A bounded-to-chart quasi-isomorphism transports these conjectural
bounds to \(\ChirHoch^n(\cA_{V_k(\mathfrak{sl}_3)})\).
\end{corollary}

\begin{proof}
This is the specialization \(\fg=\mathfrak{sl}_3\) of
\cite[Conjecture~7.5]{BDSK21}.  The final clause is transport along
\(\chi_{V_k(\mathfrak{sl}_3)}^{\mathrm{bd}}\).
\end{proof}

\begin{corollary}[Virasoro bounded cohomology;
\ClaimStatusProvedElsewhere]
\label{cor:chirhoch-virasoro-hilbert}
\index{Hochschild cohomology!Virasoro Hilbert series}
Bakalov--De Sole--Kac compute, for every central charge~\(c\),
\[
 \dim H^n_{\mathrm{ch},b}(\operatorname{Vir}_c,
 \operatorname{Vir}_c)
 =\begin{cases}
 1,&n\in\{0,2,3\},\\
 0,&n\in\mathbb Z_{\geq0}\setminus\{0,2,3\}.
 \end{cases}
\]
Thus \(P^{\mathrm{bd}}_{\operatorname{Vir}_c}(t)=1+t^2+t^3\).
If \(\chi_{\operatorname{Vir}_c}^{\mathrm{bd}}\) is a
quasi-isomorphism, the same polynomial computes the chart complex.
\end{corollary}

\begin{proof}
The bounded cohomology calculation is
\cite[Theorem~7.2]{BDSK21}.  The chart statement follows from
invariance of cohomology under \(\chi_{\operatorname{Vir}_c}^{\mathrm{bd}}\).
\end{proof}

\begin{remark}[Computed and conjectural support calculations]
\label{rem:standard-family-theorem-h-computation-surfaces}
\index{Theorem H!standard-family support calculations}
The proven bounded vertex-cohomology calculations are the rank-one
even superboson of Corollary~\ref{cor:chirhoch-heisenberg} and the
Virasoro algebra of
Corollary~\ref{cor:chirhoch-virasoro-hilbert}.  The universal affine
bounds of Corollaries~\ref{cor:chirhoch-affine-sl2} and
\ref{cor:chirhoch-affine-sl3} retain the conjectural status of
\cite[Conjecture~7.5]{BDSK21}.  Each chart-algebra statement is
obtained from a bounded-to-chart quasi-isomorphism.  Principal
\(\mathcal W\)-algebras, admissible quotients, minimal models, and
logarithmic triplet algebras require their own complexes
\(K_{\cA,S}\) and their own comparison maps.
\end{remark}

\begin{lemma}[The pure-braid Koszul complex and the Arnold algebra;
\ClaimStatusProvedHere]
\label{lem:chiral-homotopy-transport}
\index{Shelton--Yuzvinsky!chiral transport}
\index{Koszul complex!Orlik--Solomon}
Let $m \geq 2$ and let $\operatorname{OS}(A_{m-1})$ be the
Orlik--Solomon algebra of the type-$A_{m-1}$ braid arrangement:
the graded-commutative algebra on the Arnold forms
$\omega_{ij} = d\log(z_i - z_j)$ modulo the three-term relation,
with $\operatorname{OS}(A_{m-1}) \cong H^{*}(\FM_m(\C);\Q)$
\cite{Arnold1969,OrlikSolomon1980}, zero differential, and
Poincar\'e polynomial $\prod_{j=1}^{m-1}(1+jt)$.
\begin{enumerate}[label=\textup{(\roman*)}]
\item The arrangement $A_{m-1}$ is supersolvable, so
$\operatorname{OS}(A_{m-1})$ is quadratic Koszul
\cite{SheltonYuzvinsky1997}. Equivalently
\textup{(}Priddy \cite{Priddy1970}\textup{)}, the Koszul complex
\[
 \operatorname{OS}(A_{m-1}) \otimes
 \bigl(\operatorname{OS}(A_{m-1})^{!}\bigr)^{\vee},
 \qquad d_{\mathrm K} = \textup{Priddy twist},
\]
is acyclic in positive degree: the supersolvable flag furnishes a
contracting homotopy of the \emph{Koszul complex} onto the ground
field.
\item The Arnold algebra carries the zero differential and therefore
 \[
 H^{\geq1}\bigl(\operatorname{OS}(A_{m-1}),0\bigr)
 =\operatorname{OS}(A_{m-1})^{\geq1}.
 \]
 Thus the Priddy homotopy belongs to the Koszul complex in
 \textup{(i)}, while the Arnold algebra retains its positive-degree
 cohomology.
\end{enumerate}
\end{lemma}

\begin{proof}
(i) Supersolvability of $A_{m-1}$ is the modular flag of
coordinate subarrangements; quadratic Koszulness is
Shelton--Yuzvinsky \cite{SheltonYuzvinsky1997}; the equivalence of
Koszulness with positive-degree acyclicity of the Koszul complex
is \cite{Priddy1970} (see also \cite[Ch.~3]{LV12}).
(ii) The identity follows from \(d=0\).  The Poincar\'e polynomial
\(\prod_{j=1}^{m-1}(1+jt)\) has degree \(m-1\), so the
positive-degree summand occurs for every \(m\geq2\).
\end{proof}

\begin{definition}[Nondegenerate collision symbol;
\ClaimStatusDefinitional]
\label{def:nondegenerate-collision-symbol}
Let (F) be the PBW--collision filtration on

\[
 T_{\cA,\sigma,N}
 =\left(\bigotimes_i\operatorname{OS}(A_{m_i-1})\right)
  \otimes M_{\cA,\sigma,N}.
\]
The leading residue operators along the internal edges of~\(\sigma\)
define a sequence \((R_e)_e\) on
\(\operatorname{gr}_F M_{\cA,\sigma,N}\).  The collision symbol is
\emph{nondegenerate} when this sequence is regular on every positive
Arnold-degree summand, so that its Koszul complex has cohomology in
fibre degree zero.  Equivalently, the symbol complex admits a
contraction
\[
 \left(\operatorname{gr}_F T_{\cA,\sigma,N},d_{\mathrm{symb}}\right)
 \;\rightleftarrows\;
 K^{(0)}_{\cA,\sigma,N}
\]
onto its fibre-degree-zero cohomology.  For a regular-OPE commutative
chiral algebra all (R_e) vanish, and the positive Arnold classes
identify the degenerate collision-symbol locus.
\end{definition}

\begin{conjecture}[Family-specific ordered collision retract;
\ClaimStatusConjectured]
\label{conj:ordered-twisted-tensor-acyclicity}
\index{twisted tensor complex!ordered chiral}
Fix a chiral algebra~\(\cA\), a finite weight window~\(N\), and a
collision stratum~\(\sigma\) of depth profile
\((m_1,\ldots,m_k)\) in
\(\overline C^{\mathrm{ord}}_{p+2}(X)\).  Its residue-twisted
coefficient complex is
\[
 T_{\cA,\sigma,N}
 =\left(\bigotimes_i\operatorname{OS}(A_{m_i-1})\right)
 \otimes M_{\cA,\sigma,N},
 \qquad d=d_{\mathrm{int}}+d_{\mathrm{res}}.
\]
\textup{[Type signature: Open quadrant; ordered logarithmic
Fulton--MacPherson presentation; Beilinson level~$3$; hypothesis
package: nondegenerate collision symbols, complete PBW--collision
filtration, finite weight windows, and incidence-compatible symbol
contractions.]}
Assume every collision symbol is nondegenerate in the sense of
Definition~\ref{def:nondegenerate-collision-symbol}.  The conjecture
asserts that the symbol contractions can be chosen compatibly with
edge contraction and lifted through the complete filtration to maps
\[
 K_{\cA,\sigma,N}
 \mathop{\xrightarrow{\ \iota_{\sigma,N}\ }}
 T_{\cA,\sigma,N}
 \mathop{\xrightarrow{\ p_{\sigma,N}\ }}
 K_{\cA,\sigma,N},
 \qquad
 dh_{\sigma,N}+h_{\sigma,N}d
 =\id-\iota_{\sigma,N}p_{\sigma,N},
\]
where (K_{\cA,\sigma,N}) is the homological-perturbation transfer of
(K^{(0)}_{\cA,\sigma,N}).  The maps are natural for collision
incidences and compatible with weight projections.  PBW Koszulness
constructs the associated-graded coefficient model; nondegeneracy of
the residue symbol contracts the positive Arnold fibre; completeness
permits the perturbation series.  Incidence compatibility is the
remaining family-specific assertion.
\end{conjecture}

\begin{proposition}[Two-point Heisenberg residue-twisted Arnold
contraction; \ClaimStatusProvedHere]
\label{prop:heisenberg-two-point-residue-twisted-acyclicity}
\index{Heisenberg!two-point residue-twisted contraction}
\index{Theorem H!Heisenberg two-point Arnold line}
\textup{[Type signature: Open quadrant; ordered logarithmic
Fulton--MacPherson presentation; Beilinson level~$3$ reached from the
level~$2$ bar; depth-one two-point collision stratum; hypothesis
package: rank-one Heisenberg OPE
\(\alpha_{(1)}\alpha=k\mathbf1\), \(k\neq0\), and the ordered
bar-residue sign convention of Theorem~\ref{thm:residue-formula}.]}
For the rank-one Heisenberg chiral algebra \(\mathfrak H_k\) with
\(k\neq0\), restrict
Conjecture~\ref{conj:ordered-twisted-tensor-acyclicity} to the
depth-one collision stratum with two generator inputs
\(\alpha,\alpha\).  The local positive-fibre summand is the two-term
complex
\[
 C^1_{2,\alpha}
 =\C\cdot\bigl([\alpha|\alpha]\otimes\eta_{12}\bigr)
 \xrightarrow{\ d_1\ }
 C^0_{2,\alpha}
 =\C\cdot\mathbf1,
 \qquad
 \eta_{12}=d\log(z_1-z_2),
\]
where \(d_1\) is the OPE-residue part of the collision-depth
differential.  With the ordered adjacent-face convention,
\[
 d_1\bigl([\alpha|\alpha]\otimes\eta_{12}\bigr)
 =
 \alpha_{(1)}\alpha
 =
 k\,\mathbf1.
\]
Consequently \(H^1(C^\bullet_{2,\alpha})=0\); the positive
\(\operatorname{OS}(A_1)\)-Arnold line in this central-current
two-point summand is contracted.

The proved scope is the displayed arity-\(2\) summand of the
residue-twisted mechanism.  Arbitrary Fock monomials, collision
clusters \(m\ge3\), multi-stratum depth profiles, ordered-to-symmetric
descent, and the completed curved second-kind complex of
Lemma~\ref{lem:curved-dual-centre-heisenberg} require their respective
complexes and homotopies.
\end{proposition}

\begin{proof}
Apply the arbitrary-mode residue formula
\eqref{eq:ordered-residue-arbitrary-mode} with \(N=2\),
\(I=\{1,2\}\), \(m=1\), \(a_1=a_2=\alpha\), and
\(\omega=\eta_{12}\).  The Heisenberg generator \(\alpha\) is even, so
the adjacent-face sign of Theorem~\ref{thm:residue-formula} is \(+1\).
The Poincar\'e residue removes the logarithmic normal factor
\(\eta_{12}\), while the Beilinson--Drinfeld mode projection supplies
\(\operatorname{pr}_1\mu_{\mathrm{BD}}(\alpha,\alpha)
=\alpha_{(1)}\alpha\).  The Heisenberg OPE gives
\(\alpha_{(1)}\alpha=k\mathbf1\).  Since \(k\neq0\), the displayed
differential is an isomorphism of one-dimensional vector spaces, so
the positive fibre cohomology \(H^1(C^\bullet_{2,\alpha})\) vanishes.
The calculation is confined to this collision stratum and this
coefficient summand.
\end{proof}

\begin{proposition}[Two-point Heisenberg weight-one polynomial
residue contraction; \ClaimStatusProvedHere]
\label{prop:heisenberg-two-point-weight-one-polynomial-residue}
\index{Heisenberg!two-point weight-one polynomial contraction}
\index{Theorem H!Heisenberg polynomial Arnold line}
\textup{[Type signature: Open quadrant; ordered logarithmic
Fulton--MacPherson presentation; Beilinson level~$3$ reached from the
level~$2$ bar; depth-one two-point collision stratum; hypothesis
package: rank-one Heisenberg mode relation
\([\alpha_m,\alpha_n]=mk\delta_{m+n,0}\mathbf1\), \(k\neq0\), and the
ordered bar-residue sign convention of
Theorem~\ref{thm:residue-formula}.]}
Let \(u_q=\alpha_{-1}^{q}\mathbf1=:\!\alpha^q\!:\) be the
unnormalized weight-one-current monomial in \(\mathfrak H_k\), with
\(q\geq1\).  In the depth-one two-point collision stratum, restrict the
residue-twisted tensor complex to the ordered summand
\[
 C^1_{2,q}
 =
 \C\cdot\bigl([\alpha|u_q]\otimes\eta_{12}\bigr)
 \xrightarrow{\ d_1\ }
 C^0_{2,q}
 =
 \C\cdot u_{q-1}.
\]
Then
\[
 d_1\bigl([\alpha|u_q]\otimes\eta_{12}\bigr)
 =
 \alpha_{(1)}u_q
 =
 qk\,u_{q-1}.
\]
For \(k\neq0\) this is an isomorphism of one-dimensional vector
spaces.  Hence \(H^1(C^\bullet_{2,q})=0\) for every \(q\geq1\).

This proves the positive Arnold fibre contraction on the two-point
weight-one polynomial Heisenberg string \(\C[\alpha_{-1}]\).  The
single-mode higher-oscillator polynomial extension is
Proposition~\ref{prop:heisenberg-two-point-single-mode-polynomial-residue}.
Mixed-mode oscillator monomials, clusters \(m\geq3\), multi-strata,
and ordered-to-symmetric descent require the full family collision
retract.
\end{proposition}

\begin{proof}
Apply \eqref{eq:ordered-residue-arbitrary-mode} with \(N=2\),
\(I=\{1,2\}\), \(m=1\), \(a_1=\alpha\), \(a_2=u_q\), and
\(\omega=\eta_{12}\).  The field \(\alpha\) is even, so the adjacent
ordered-face sign is \(+1\).  The residue removes the logarithmic
normal factor and leaves the first product \(\alpha_{(1)}u_q\).

In the rank-one Heisenberg mode algebra,
\[
 [\alpha_m,\alpha_n]=mk\,\delta_{m+n,0}\mathbf1,
 \qquad
 \alpha_r\mathbf1=0\quad(r\geq0).
\]
Therefore
\[
 \alpha_{(1)}u_q
 =
 \alpha_1\alpha_{-1}^{q}\mathbf1
 =
 [\alpha_1,\alpha_{-1}^{q}]\mathbf1
 =
 q[\alpha_1,\alpha_{-1}]\alpha_{-1}^{q-1}\mathbf1
 =
 qk\,u_{q-1}.
\]
The scalar \(qk\) is invertible over \(\C\) for \(q\geq1\) and
\(k\neq0\), so its positive-fibre cohomology is zero.
\end{proof}

\begin{proposition}[Two-point Heisenberg single-oscillator
arbitrary-mode residue contraction; \ClaimStatusProvedHere]
\label{prop:heisenberg-two-point-single-oscillator-residue}
\index{Heisenberg!two-point single-oscillator contraction}
\index{Theorem H!Heisenberg arbitrary-mode Arnold line}
\textup{[Type signature: Open quadrant; ordered logarithmic
Fulton--MacPherson presentation; Beilinson level~$3$ reached from the
level~$2$ bar; depth-one two-point collision stratum; hypothesis
package: rank-one Heisenberg mode relation
\([\alpha_m,\alpha_n]=mk\delta_{m+n,0}\mathbf1\), \(k\neq0\), and the
arbitrary-mode ordered residue formula
\eqref{eq:ordered-residue-arbitrary-mode}.]}
Let \(v_n=\alpha_{-n}\mathbf1\) with \(n\geq1\).  In the depth-one
two-point collision stratum, restrict the residue-twisted tensor
complex to the ordered \(m=n\) OPE-mode summand
\[
 C^1_{2,n}
 =
 \C\cdot\bigl([\alpha|v_n]\otimes\eta_{12}\bigr)
 \xrightarrow{\ d_1^{(n)}\ }
 C^0_{2,n}
 =
 \C\cdot\mathbf1 .
\]
Then
\[
 d_1^{(n)}\bigl([\alpha|v_n]\otimes\eta_{12}\bigr)
 =
 \alpha_{(n)}v_n
 =
 nk\,\mathbf1 .
\]
For \(k\neq0\) this is an isomorphism of one-dimensional vector
spaces.  Hence \(H^1(C^\bullet_{2,n})=0\) for every \(n\geq1\).
The case \(n=1\) recovers
Proposition~\ref{prop:heisenberg-two-point-residue-twisted-acyclicity};
the cases \(n>1\) are the single higher-oscillator
arbitrary-mode summands.

The proved scope is the two-point single-oscillator line
\(\{\alpha_{-n}\mathbf1:n\geq1\}\).  The single-mode powers
\(\alpha_{-n}^{q}\) are treated in
Proposition~\ref{prop:heisenberg-two-point-single-mode-polynomial-residue}.
Mixed-mode oscillator monomials, full Fock-window linear combinations,
clusters \(m\geq3\), multi-strata, and ordered-to-symmetric descent
require the full family collision retract.
\end{proposition}

\begin{proof}
Apply \eqref{eq:ordered-residue-arbitrary-mode} with \(N=2\),
\(I=\{1,2\}\), \(m=n\), \(a_1=\alpha\), \(a_2=v_n\), and
\(\omega=\eta_{12}\).  Since \(\alpha\) is even, the adjacent
ordered-face sign is \(+1\).  The Poincar\'e residue removes the
single logarithmic normal factor \(\eta_{12}\), and the
\(m=n\) Beilinson--Drinfeld mode projection leaves
\(\alpha_{(n)}v_n\).

The rank-one Heisenberg mode relation gives
\[
 [\alpha_m,\alpha_l]=mk\,\delta_{m+l,0}\mathbf1,
 \qquad
 \alpha_r\mathbf1=0\quad(r\geq0).
\]
Therefore
\[
 \alpha_{(n)}v_n
 =
 \alpha_n\alpha_{-n}\mathbf1
 =
 [\alpha_n,\alpha_{-n}]\mathbf1
 =
 nk\,\mathbf1 .
\]
The scalar \(nk\) is invertible over \(\C\) for \(n\geq1\) and
\(k\neq0\), so its positive-fibre cohomology is zero.
\end{proof}

\begin{proposition}[Two-point Heisenberg single-mode polynomial
arbitrary-mode residue contraction; \ClaimStatusProvedHere]
\label{prop:heisenberg-two-point-single-mode-polynomial-residue}
\index{Heisenberg!two-point single-mode polynomial contraction}
\index{Theorem H!Heisenberg single-mode polynomial Arnold line}
\textup{[Type signature: Open quadrant; ordered logarithmic
Fulton--MacPherson presentation; Beilinson level~$3$ reached from the
level~$2$ bar; depth-one two-point collision stratum; hypothesis
package: rank-one Heisenberg mode relation
\([\alpha_m,\alpha_n]=mk\delta_{m+n,0}\mathbf1\), \(k\neq0\), and the
arbitrary-mode ordered residue formula
\eqref{eq:ordered-residue-arbitrary-mode}.]}
Fix \(n\geq1\) and \(q\geq1\), and put
\[
 u_{n,q}=\alpha_{-n}^{q}\mathbf1,
 \qquad
 u_{n,q-1}=\alpha_{-n}^{q-1}\mathbf1 .
\]
In the depth-one two-point collision stratum, restrict the
residue-twisted tensor complex to the ordered \(m=n\) OPE-mode
summand
\[
 C^1_{2,n,q}
 =
 \C\cdot\bigl([\alpha|u_{n,q}]\otimes\eta_{12}\bigr)
 \xrightarrow{\ d_1^{(n)}\ }
 C^0_{2,n,q}
 =
 \C\cdot u_{n,q-1}.
\]
Then
\[
 d_1^{(n)}\bigl([\alpha|u_{n,q}]\otimes\eta_{12}\bigr)
 =
 \alpha_{(n)}u_{n,q}
 =
 q\,n\,k\,u_{n,q-1}.
\]
For \(k\neq0\) this is an isomorphism of one-dimensional vector
spaces.  Hence \(H^1(C^\bullet_{2,n,q})=0\) for every \(n,q\geq1\).
The case \(n=1\) is
Proposition~\ref{prop:heisenberg-two-point-weight-one-polynomial-residue},
and the case \(q=1\) is
Proposition~\ref{prop:heisenberg-two-point-single-oscillator-residue}.

This proves the two-point single-mode polynomial strings
\(\C[\alpha_{-n}]\), one mode \(n\) at a time.  Mixed-mode oscillator
monomials, full Fock-window linear combinations, clusters \(m\geq3\),
multi-strata, and ordered-to-symmetric descent require the complete
residue differential and the corresponding family collision retract.
\end{proposition}

\begin{proof}
Apply \eqref{eq:ordered-residue-arbitrary-mode} with \(N=2\),
\(I=\{1,2\}\), \(m=n\), \(a_1=\alpha\), \(a_2=u_{n,q}\), and
\(\omega=\eta_{12}\).  The adjacent ordered-face sign is \(+1\) since
\(\alpha\) is even.  The Poincar\'e residue removes the normal
logarithmic form, and the \(m=n\) mode projection leaves
\(\alpha_{(n)}u_{n,q}\).

Using
\[
 [\alpha_m,\alpha_l]=mk\,\delta_{m+l,0}\mathbf1,
 \qquad
 \alpha_r\mathbf1=0\quad(r\geq0),
\]
we compute
\[
 \alpha_{(n)}u_{n,q}
 =
 \alpha_n\alpha_{-n}^{q}\mathbf1
 =
 [\alpha_n,\alpha_{-n}^{q}]\mathbf1
 =
 q[\alpha_n,\alpha_{-n}]\alpha_{-n}^{q-1}\mathbf1
 =
 q\,n\,k\,u_{n,q-1}.
\]
The scalar \(qnk\) is invertible over \(\C\) for \(n,q\geq1\) and
\(k\neq0\), so its positive-fibre cohomology is zero.
\end{proof}

\begin{proposition}[Two-point Heisenberg mixed-mode residue formula;
\ClaimStatusProvedHere]
\label{prop:heisenberg-two-point-mixed-mode-residue-formula}
\index{Heisenberg!two-point mixed-mode residue formula}
\index{Theorem H!Heisenberg mixed-mode formula}
\textup{[Type signature: Open quadrant; ordered logarithmic
Fulton--MacPherson presentation; Beilinson level~$3$ reached from the
level~$2$ bar; depth-one two-point collision stratum; hypothesis
package: rank-one Heisenberg mode relation
\([\alpha_m,\alpha_n]=mk\delta_{m+n,0}\mathbf1\), finite Fock
support, \(k\neq0\) when acyclicity of a line is asserted, and the
arbitrary-mode ordered residue formula
\eqref{eq:ordered-residue-arbitrary-mode}.]}
Let \(\mathbf q=(q_1,q_2,\ldots)\) be a finite-support multi-index
with \(q_r\in\mathbb Z_{\ge0}\), and write
\[
 u_{\mathbf q}
 =
 \prod_{r\ge1}\alpha_{-r}^{q_r}\mathbf1,
 \qquad
 \mathbf e_r=(0,\ldots,0,1,0,\ldots).
\]
On the depth-one two-point collision stratum, the total
arbitrary-mode residue of the ordered positive Arnold generator is
\[
 d_1\bigl([\alpha|u_{\mathbf q}]\otimes\eta_{12}\bigr)
 =
 \sum_{r:q_r>0}
 d_1^{(r)}\bigl([\alpha|u_{\mathbf q}]\otimes\eta_{12}\bigr)
 =
 k\sum_{r:q_r>0} r q_r\,u_{\mathbf q-\mathbf e_r}.
\]
In particular, for a fixed non-vacuum monomial \(u_{\mathbf q}\) and
\(k\neq0\), the one-dimensional positive Arnold line maps
injectively to its image, hence the line-to-image two-term summand has
zero positive-fibre cohomology.

This exact mixed-mode formula proves the contraction on each displayed
line.  A full Fock-window contraction uses the complete graded
residue-twisted Koszul differential.  On a raw ungraded polynomial span
\(\C[x_1,x_2,\ldots]\), the formula is the vector field
\[
 L_k=k\sum_{r\ge1}r\,\partial_{x_r},
 \qquad
 x_r\leftrightarrow \alpha_{-r}\mathbf1 .
\]
Thus \(L_k(x_2-2x_1)=0\) for \(k\neq0\).  A full mixed-mode
Theorem-H contraction is therefore a contraction of the actual graded
residue-twisted Koszul complex, equipped with its finite-window
homotopy data.
\end{proposition}

\begin{proof}
For each \(r\) with \(q_r>0\), apply
\eqref{eq:ordered-residue-arbitrary-mode} with \(N=2\),
\(I=\{1,2\}\), \(m=r\), \(a_1=\alpha\), \(a_2=u_{\mathbf q}\), and
\(\omega=\eta_{12}\).  The adjacent sign is \(+1\).  The form residue
removes \(\eta_{12}\), and the \(m=r\) mode projection leaves
\(\alpha_{(r)}u_{\mathbf q}\).

The Heisenberg relation gives
\[
 \alpha_{(r)}u_{\mathbf q}
 =
 \alpha_r
 \Bigl(\prod_{s\ge1}\alpha_{-s}^{q_s}\Bigr)\mathbf1
 =
 [\alpha_r,\alpha_{-r}^{q_r}]
 \prod_{s\neq r}\alpha_{-s}^{q_s}\mathbf1
 =
 r k q_r\,u_{\mathbf q-\mathbf e_r},
\]
because \([\alpha_r,\alpha_{-s}]=0\) for \(s\neq r\) and
\(\alpha_r\mathbf1=0\).  Summing over the finitely many
\(r\) in the support of \(\mathbf q\) gives the displayed formula.
If \(\mathbf q\neq0\) and \(k\neq0\), at least one coefficient
\(rkq_r\) is nonzero, and the target monomials
\(\{u_{\mathbf q-\mathbf e_r}:q_r>0\}\) are distinct.  Hence the
image vector is nonzero and the map from the one-dimensional source
line has zero kernel.

The raw-kernel statement is the direct translation of the same formula
to \(L_k\): \(L_k(x_2)=2k\) and \(L_k(x_1)=k\), so
\(L_k(x_2-2x_1)=0\).  This proves the stated obstruction to promoting
the formula alone into a full mixed-mode acyclicity theorem.
\end{proof}

\begin{remark}[Chain-level versus cohomology-level scope]
\label{rem:sigma-inverse-chain-level}
At each stratum Theorem~H consumes a deformation retract of the full
residue-twisted coefficient complex.  Shelton--Yuzvinsky Koszulness
contracts the quadratic Koszul complex
\(\operatorname{OS}\otimes(\operatorname{OS}^{!})^{\vee}\).
The family calculation lifts this contraction across the OPE-residue
twist and verifies naturality for collision incidences.  The lifted
maps are the data of
Definition~\ref{def:theorem-h-pbw-finite-window-lane}.
\end{remark}

\begin{theorem}[Ordered chiral Hochschild support;
\ClaimStatusProvedHere]
\label{thm:hochschild-concentration-E1}
\index{Hochschild!ordered-bar concentration}
\index{Orlik--Solomon!pure braid Koszulness}
\textup{[Type signature: Open quadrant; ordered chain presentation;
Beilinson level~$3$; hypothesis package: incidence-compatible
finite-window retracts and Mittag--Leffler completion.]}
Fix \(S\subset\mathbb Z\).  Suppose the ordered finite-window
collision complexes admit the realization of
Definition~\ref{def:theorem-h-pbw-finite-window-lane}, with
\(K^{\mathrm{ord}}_{\cA,S}\simeq
R\!\varprojlim_NK_{\cA,N}\) supported in~\(S\).  Then
\[
 H^\bullet\!\left(Q^{\mathrm{ord}}_\cA\right)
 \cong H^\bullet\!\left(K^{\mathrm{ord}}_{\cA,S}\right),
 \qquad
 \operatorname{Supp}H^\bullet(Q^{\mathrm{ord}}_\cA)
 \subseteq S.
\]
An averaging quasi-isomorphism
\((Q^{\mathrm{ord}}_\cA)_{h\Sigma}\xrightarrow{\sim}Q_\cA\)
carries this retract to the symmetric chart complex.
\end{theorem}

\begin{proof}
Proposition~\ref{prop:fm-tower-collapse} totalizes the natural
stratum retracts.  Passing to the compatible inverse limit gives a
strong deformation retract of \(Q^{\mathrm{ord}}_\cA\) onto
\(K^{\mathrm{ord}}_{\cA,S}\).  The homotopy identities give the
cohomology isomorphism and its support.  Homotopy coinvariants and the
averaging quasi-isomorphism transport the maps to the symmetric
complex.
\end{proof}

\begin{remark}[Degree-one bounded vertex cohomology]
\label{rem:chirhoch-derivation-types}
\index{Hochschild!derivation-type breakdown}
Theorem~\ref{thm:main-koszul-hoch} reads degree one from the selected
cochain complex.  Bakalov--De Sole--Kac compute
\[
 \dim H^1_{\mathrm{ch},b}(B_{\mathfrak h})=1
 \quad(\dim\mathfrak h_{\bar0}=1),
 \qquad
 H^1_{\mathrm{ch},b}(\operatorname{Vir}_c)=0
\]
\cite[Theorems~7.4 and~7.2]{BDSK21}.  For universal affine vertex
algebras their degree-one statement belongs to
\cite[Conjecture~7.5]{BDSK21}.  The comparison
\(\chi_V^{\mathrm{bd}}\) determines the corresponding chart result.

\end{remark}

\begin{proposition}[The first two deformation degrees;
\ClaimStatusConditional]
\label{prop:smooth-formal-moduli-standard}
\index{deformation theory!unobstructedness}
\index{formal moduli!smoothness}
Let \(\cA\) be a chiral algebra and assume the binary component of its
cyclic deformation complex is identified with the normalized chiral
Hochschild complex.  Then
\[
\begin{array}{rcl}
\textup{infinitesimal fixed-product automorphisms}
&=&
\ChirHoch^1(\cA)
=\Der_{\mathrm{ch}}(\cA)/\Inn_{\mathrm{ch}}(\cA),\\[2mm]
\textup{binary chiral-product deformations}
&\subseteq&
\ChirHoch^2(\cA)
\cong H^2\!\bigl(\pi_{2,0}(\Defcyc^{\mathrm{mod}}(\cA))\bigr).
\end{array}
\]
The first line is the stabilizer complex of the fixed product.  The
second line is the linearized Maurer--Cartan complex for binary product
deformations.  A degree-two cocycle defines a formal product tangent
when it lies in the cyclic summand, and its integration is governed by
the higher brackets of the deformation complex.
\end{proposition}

\begin{proof}
The normalized degree-one differential identifies derivations modulo
inner derivations.  Binary cyclic coderivations form the degree-two
component of the unshifted Hochschild complex and the degree-one
component after the standard Maurer--Cartan shift.  Gauge
one-cochains give its coboundaries.  The higher \(L_\infty\)-brackets
are the extension equations for a first-order class.
\end{proof}

\begin{remark}[Meaning of the first-order lanes across archetypes;
\ClaimStatusConditional]
\label{rem:smooth-formal-moduli-archetype}
Proposition~\ref{prop:smooth-formal-moduli-standard} separates the
fixed-product stabilizer from the binary product tangent for every
family.  Numerical dimensions enter after constructing the relevant
cochain model and its support datum.  The bounded superboson and
Virasoro calculations supply two such examples; affine,
Drinfeld--Sokolov, free-field, and logarithmic families carry their own
comparison maps and deformation retracts.  Conformal-vector and scalar
shadow parameters belong to their enriched parameter spaces.
\end{remark}

\begin{corollary}[Averaging to symmetric ChirHoch;
\ClaimStatusConditional]
\label{cor:hochschild-averaging-symmetric}
\index{Hochschild!averaging step}
Fix a finite set \(S\subset\mathbb Z\), and assume the ordered
support theorem
Theorem~\ref{thm:hochschild-concentration-E1} with target
\(K^{\mathrm{ord}}_{\cA,S}\).  Suppose the completed averaging map
\[
 \operatorname{av}_\cA\colon
 \bigl(Q^{\mathrm{ord}}_\cA\bigr)_{h\Sigma}
 \xrightarrow{\ \sim\ }C^\bullet_{\mathrm{ch}}(\cA,\cA)
\]
is a quasi-isomorphism and the ordered retract is
\(\Sigma\)-equivariant.  Then
\[
 H^\bullet\!\left(
 (K^{\mathrm{ord}}_{\cA,S})_{h\Sigma}
 \right)
 \cong\ChirHoch^\bullet(\cA),
 \qquad
 \operatorname{Supp}\ChirHoch^\bullet(\cA)\subseteq S.
\]
\end{corollary}

\begin{proof}
Equivariance transports the ordered homotopy identities to homotopy
coinvariants.  Composition with \(\operatorname{av}_\cA\) gives a
quasi-isomorphism from the averaged support model to the symmetric
chart complex.  Its cohomology is supported in~\(S\).
\end{proof}

\begin{remark}[Scope of the averaging corollary]
\label{rem:hochschild-averaging-scope}
Corollary~\textup{\ref{cor:hochschild-averaging-symmetric}} consumes
an actual quasi-isomorphism of completed cochain complexes.  A PBW
filtration can construct its associated-graded form; braid monodromy
and the \(R\)-matrix determine its lift to singular OPEs.  The chiral
Yangian therefore asks for its own completed averaging map and its own
family support model.
\end{remark}

\begin{conjecture}[Completed Yangian support comparison;
\ClaimStatusConjectured]
\label{conj:hochschild-concentration-E1-only}
\index{Hochschild!conjectural symmetric concentration}
For the chiral Yangian \(Y_\hbar(\fg)^{\mathrm{ch}}\), there is a
finite set \(S_Y\), an ordered family support model
\(K^{\mathrm{ord}}_{Y,S_Y}\), and a completed \(R\)-matrix averaging
quasi-isomorphism
\[
 \bigl(Q_Y^{\mathrm{ord}}\bigr)_{h\Sigma}
 \xrightarrow{\ \sim\ }
 C^\bullet_{\mathrm{ch}}(Y_\hbar(\fg)^{\mathrm{ch}},
 Y_\hbar(\fg)^{\mathrm{ch}}).
\]
These data determine the symmetric support by
Corollary~\ref{cor:hochschild-averaging-symmetric}.
\end{conjecture}

\begin{conjecture}[Completed support at admissible and minimal levels;
\ClaimStatusConjectured]
\label{conj:hochschild-concentration-class-M-non-generic}
\index{Hochschild!non-generic level conjecture}
For the Virasoro algebra $\mathrm{Vir}_{c}$ at minimal-model
central charges $c = c_{p,q} = 1 - 6(p-q)^{2}/(pq)$ and for
the principal W-algebra $\Walg^{k}(\fg)$ at admissible
levels, the coderived diagonal complex
\(RHH_{\mathrm{ch}}^{\mathrm{cod},\Delta}(\cA)\) admits a complete
family support model \(K^{\mathrm{cod}}_{\cA,S}\).  The model is
compatible with the Shapovalov filtration, the simple-quotient map,
and the inverse system of conformal-weight windows.  Its support set
\(S\) is determined by the resulting Kac-determinant calculation.
\end{conjecture}

\begin{definition}[Admissible Hochschild obstruction map;
\ClaimStatusDefinitional]
\label{def:admissible-hochschild-obstruction-map}
Let $\cA_{\mathrm{univ}}$ be a universal affine or principal
$\cW$-algebra in the generic PBW finite-window ambient, and let
\[
q_{\mathrm{adm}}\colon \cA_{\mathrm{univ}}\twoheadrightarrow
\cA_{\mathrm{adm}}
\]
be an admissible simple quotient at a Kac--Kazhdan/Shapovalov
singular parameter. The quotient induces a filtered map of completed
diagonal Hochschild complexes
\[
\Phi^{\mathrm{adm}}_{\mathrm H}(\cA)\colon
RHH_{\mathrm{ch}}(\cA_{\mathrm{univ}})
\longrightarrow
RHH_{\mathrm{ch}}^{\mathrm{cod},\Delta}(\cA_{\mathrm{adm}}),
\]
where the target is the coderived/completed diagonal complex of the
simple quotient. Its cone
\[
\mathcal O^{\mathrm{adm}}_{\mathrm H}(\cA)
:=
\operatorname{Cone}\!\bigl(
\Phi^{\mathrm{adm}}_{\mathrm H}(\cA)
\bigr)
\]
is the admissible Hochschild obstruction complex. A finite
Theorem~H statement for $\cA_{\mathrm{adm}}$ requires the high-degree
cohomology of this cone, together with the positive-depth
$\mathrm{KD}_{\mathrm H}$ obstruction, to vanish.  The quasi-isomorphism
status of $\Phi^{\mathrm{adm}}_{\mathrm H}$ is determined by this
admissible finite-window calculation.
\end{definition}

\begin{definition}[Cartan bar obstruction for admissible
\texorpdfstring{$\mathfrak{sl}_3$}{sl3}; \ClaimStatusDefinitional]
\label{def:cartan-bar-obstruction-sl3-admissible}
For an admissible quotient of $V_k(\mathfrak{sl}_3)$, let
$\mathfrak h\subset\mathfrak{sl}_3$ be the Cartan subalgebra and let
$RHH_{\mathrm{ch}}^{\mathfrak h,\Delta}$ denote the completed
diagonal Hochschild subcomplex obtained by restricting the diagonal
bar resolution to Cartan-current insertions. The \emph{Cartan bar
obstruction} is
\[
\mathfrak c_{\mathrm H,\mathfrak{sl}_3}^{\mathrm{Cartan}}
:=
H^{\ge3}\!\left(
\operatorname{Cone}\!\left(
RHH_{\mathrm{ch}}^{\mathfrak h,\Delta}
\longrightarrow
RHH_{\mathrm{ch}}^{\mathrm{cod},\Delta}
(V_k(\mathfrak{sl}_3)_{\mathrm{adm}})
\right)
\right).
\]
It measures the part of the admissible high-degree tail already
visible in the Cartan bar subcomplex.
\end{definition}

\begin{proposition}[Cartan obstruction injectivity criterion;
\ClaimStatusConditional]
\label{prop:cartan-obstruction-injectivity-criterion}
The Cartan obstruction of
Definition~\ref{def:cartan-bar-obstruction-sl3-admissible} injects
into the full admissible Hochschild obstruction precisely when the
relative cone
\[
\operatorname{Cone}\!\left(
RHH_{\mathrm{ch}}^{\mathfrak h,\Delta}
\longrightarrow
RHH_{\mathrm{ch}}^{\mathrm{cod},\Delta}
(V_k(\mathfrak{sl}_3)_{\mathrm{adm}})
\right)
\]
has connecting image zero on the Cartan degree-$3$ class. Equivalently,
the boundary map
\[
H^2\!\left(
RHH_{\mathrm{ch}}^{\mathrm{cod},\Delta}/
RHH_{\mathrm{ch}}^{\mathfrak h,\Delta}
\right)
\longrightarrow
H^3\!\left(RHH_{\mathrm{ch}}^{\mathfrak h,\Delta}\right)
\]
vanishes on the Cartan obstruction class.  Its value is a
finite-window calculation in the displayed relative complex.
\end{proposition}

\begin{proof}
Apply the long exact cohomology sequence to the short exact sequence
of filtered complexes
\[
RHH_{\mathrm{ch}}^{\mathfrak h,\Delta}
\longrightarrow
RHH_{\mathrm{ch}}^{\mathrm{cod},\Delta}
\longrightarrow
RHH_{\mathrm{ch}}^{\mathrm{cod},\Delta}/
RHH_{\mathrm{ch}}^{\mathfrak h,\Delta}.
\]
The kernel of the map from Cartan degree-three cohomology to the full
complex is the image of the displayed connecting homomorphism.  The
named class therefore injects precisely when its image under this
boundary criterion is zero.
\end{proof}

\begin{remark}[The ordered \(E_1\)-chiral presentation]
\label{rem:E1-scope-hochschild-concentration}
The label in Theorem~\ref{thm:hochschild-concentration-E1} records the
ordered \(E_1\)-chiral bar presentation.  Its chain model retains braid
monodromy on ordered Fulton--MacPherson configurations.  The averaging
map of Corollary~\ref{cor:hochschild-averaging-symmetric} carries an
equivariant ordered retract to symmetric chiral cochains.  A chiral
Yangian \(Y_\hbar(\fg)^{\mathrm{ch}}\) therefore asks for an ordered
family support complex and a completed \(R\)-matrix averaging
quasi-isomorphism.
\end{remark}

\begin{remark}[Collision diagrams and their retracts]
\label{rem:fm-collapse-nontrivial}
The spaces \(\overline C_{p+2}(X)\) form an unbounded tower, and each
collision stratum contributes a residue-twisted coefficient complex.
Orlik--Solomon Koszulness supplies a contraction for the quadratic
Koszul factor.  The family calculation lifts it to the full
OPE-residue differential and verifies compatibility with every
incidence map.  Proposition~\ref{prop:fm-tower-collapse} totalizes
these homotopies.  The resulting target \(K_{\cA,S}\) determines the
support set~\(S\), which may contain degree three as in the bounded
Virasoro calculation.
\end{remark}

\begin{lemma}[Chiral Hochschild descent; \ClaimStatusConditional]
\label{lem:chirhoch-descent}
\index{chiral Hochschild!descent from bar}
Let \(\cA\) be an augmented chiral algebra for which the two-sided
bar augmentation
\[
 \cA\otimes\bar B_X^{\mathrm{ch}}(\cA)\otimes\cA
 \longrightarrow\cA
\]
is a semifree resolution in chiral \(\cA\)-bimodules.  Then the
augmentation induces a quasi-isomorphism
\begin{equation}
\label{eq:chirhoch-descent}
 Q_\cA^{\mathrm{bar}}
 :=\operatorname{Tot}^{\oplus}_{p\geq0}
 R\operatorname{Hom}_{\cA^e}
 \bigl(\cA\otimes\bar B_X^{\mathrm{ch},p}(\cA)
 \otimes\cA,\cA\bigr)[-p]
 \xrightarrow{\ \sim\ }
 C^\bullet_{\mathrm{ch}}(\cA,\cA).
\end{equation}
Suppose each bar piece is Verdier-dualizable and carries a perfect
evaluation pairing compatible with the bar differential, collision
incidences, and symmetric descent.  These pairings define a second
comparison
\[
 \nu_\cA\colon Q_\cA^{\mathrm{bar}}
 \longrightarrow
 \mathsf{DHH}_\Sigma(\cA)
 :=\bigl(\mathbb D_{\Ran}
 \bar B_X^{\mathrm{ch}}(\cA)\bigr)_{\Sigma,\Delta}.
\]
When \(\nu_\cA\) is a quasi-isomorphism it realizes the
Verdier-diagonal model.  For a quadratic Koszul chart, the independent
maps defining the Koszul dual are
\[
 \cA^i\xrightarrow{\ q_\cA\ }
 \bar B_X^{\mathrm{ch}}(\cA),
 \qquad
 \cA^!:=\mathbb D_X(\cA^i).
\]
\end{lemma}

\begin{proof}
Applying derived Hom from a semifree resolution computes
\(R\operatorname{End}_{\cA^e}(\cA)\), which proves
\eqref{eq:chirhoch-descent}.  A compatible perfect evaluation pairing
on every bar piece adjoints to \(\nu_\cA\); its assumed
quasi-isomorphism gives the Verdier-diagonal realization.  The map
\(q_\cA\) belongs to the quadratic Koszul comparison, and
\(\mathbb D_X\) then forms the dual algebra \(\cA^!\).
\end{proof}

\begin{remark}[The four maps in Theorem H]
\label{rem:theorem-H-filter-exactness}
Theorem~H is carried by four explicit maps:
\[
 Q_\cA\xrightarrow{\ \gamma_\cA\ }
 C^\bullet_{\mathrm{ch}}(\cA,\cA),
 \qquad
 K_{\cA,S}\mathop{\xrightarrow{\ \iota_\cA\ }}Q_\cA
 \mathop{\xrightarrow{\ p_\cA\ }}K_{\cA,S},
\]
together with the homotopy \(h_\cA\) satisfying
\(dh_\cA+h_\cA d=\id-\iota_\cA p_\cA\).  A finite-window
construction also carries the inverse-limit comparison and the
ordered-to-symmetric averaging map.  These maps determine the
cohomological support.  A Koszul-dual pairing is an additional map
\(\Phi_\cA\) from
Definition~\ref{def:theorem-h-shift-normalization-package}; its
perfectness gives the separate duality statement of
Lemma~\ref{lem:hochschild-shift-computation}.  For the affine,
Virasoro, and \(\cW\)-algebra families the four maps are at present
pure hypothesis: no finite-window collision realization of
Definition~\ref{def:theorem-h-pbw-finite-window-lane} has been
constructed for any of them, and the corresponding lanes of Theorem~H
carry no proved content beyond the bounded calculations of
\cite{BDSK21} together with the hypothesis \(\chi_V^{\mathrm{bd}}\).
\end{remark}

\begin{theorem}[Family-indexed chiral Hochschild support
\textup{(}Theorem~H\textup{)}; \ClaimStatusConditional]
\label{thm:main-koszul-hoch}
\index{Hochschild cohomology!family support|textbf}
\index{Theorem H!strong deformation retract}

\textup{[Type signature: Open quadrant; chain presentation;
Beilinson level~$3$; hypothesis package \(H_H(\cA;S)\).]}

Let \(S\subset\mathbb Z\) be finite and let \(\cA\) carry the family
support datum \(H_H(\cA;S)\) of
Definition~\ref{def:theorem-h-collision-depth-package}.  Then
\[
 \gamma_\cA\iota_\cA\colon K_{\cA,S}
 \xrightarrow{\ \sim\ }C^\bullet_{\mathrm{ch}}(\cA,\cA)
\]
is a quasi-isomorphism, and
\begin{equation}\label{eq:theorem-h-family-support}
 \ChirHoch^n(\cA)\cong H^n(K_{\cA,S}),
 \qquad
 \operatorname{Supp}\ChirHoch^\bullet(\cA)\subseteq S.
\end{equation}

Let \(V\) be a translation-equivariant vertex algebra and assume the
bounded-to-chart map
\(\chi_V^{\mathrm{bd}}\) of
Definition~\ref{def:theorem-h-comparison-datum} is a
quasi-isomorphism.  Bakalov--De Sole--Kac then give the following
chart supports:
\begin{align*}
 V=B_{\mathfrak h},\quad \dim\mathfrak h_{\bar0}=1,
 \ \mathfrak h_{\bar1}=0
 &:\quad
 \dim\ChirHoch^n(\cA_V)=
 \begin{cases}
 2,&n=0,\\
 1,&n=1,\\
 0,&n\geq2,
 \end{cases}\\
 V=\operatorname{Vir}_c
 &:\quad
 \dim\ChirHoch^n(\cA_V)=
 \begin{cases}
 1,&n\in\{0,2,3\},\\
 0,&n\in\mathbb Z_{\geq0}\setminus\{0,2,3\}.
 \end{cases}
\end{align*}

Suppose in addition that \(\cA^!\) carries
\(H_H(\cA^!;S^!)\) and
that \((\cA,\cA^!)\) carries a perfect degree-\(d\) chain pairing.
Then the support theorem is accompanied by the separate equivalence
\begin{equation}\label{eq:hochschild-H-derived}
 Q_\cA
 \simeq
 R\operatorname{Hom}_{\mathscr D}
 \bigl(Q_{\cA^!},\omega_X[d]\bigr),
\end{equation}
and hence
\[
 \ChirHoch^n(\cA)
 \cong
 \ChirHoch^{d-n}(\cA^!)^\vee\otimes\omega_X.
\]
The geometric pairing determines~\(d\); a degree-two pairing gives
the reflection \(n\leftrightarrow2-n\).
\end{theorem}

\begin{proof}
Equation~\eqref{eq:theorem-h-global-retract} makes
\(\iota_\cA\) and \(p_\cA\) homotopy inverse.  Composition with
\(\gamma_\cA\) proves \eqref{eq:theorem-h-family-support}.

Bakalov--De Sole--Kac prove
\[
 H^n_{\mathrm{ch},b}(B_{\mathfrak h})
 \cong
 \bigl(S^n(\Pi\mathfrak h)\bigr)^*
 \oplus
 \bigl(S^{n+1}(\Pi\mathfrak h)\bigr)^*
\]
and compute \(H^n_{\mathrm{ch},b}(\operatorname{Vir}_c)\) as a
one-dimensional space in degrees \(0,2,3\), with its remaining
degrees equal to the zero space
\cite[Theorems~7.4 and~7.2]{BDSK21}.  For a one-dimensional even
\(\mathfrak h\), the space \(\Pi\mathfrak h\) is odd, so its
super-symmetric powers occur in degrees~\(0\) and~\(1\); this gives
the vector \((2,1,0,\ldots)\).  The quasi-isomorphism
\(\chi_V^{\mathrm{bd}}\) carries these calculations to the chart
complex.

The final clause is Lemma~\ref{lem:hochschild-shift-computation}
applied to the assumed perfect chain pairing.  Taking cohomology of
\eqref{eq:hochschild-H-derived} gives the displayed degree reflection.
\end{proof}

\begin{remark}[Diagonal resolution, Koszul reflection, pairing]
Under the semifree hypothesis of
Lemma~\ref{lem:chirhoch-descent}, the diagonal two-sided bar resolution
computes \(C^\bullet_{\mathrm{ch}}(\cA,\cA)\).  The quadratic comparison
\(q_\cA\colon\cA^i\to\bar B^{\mathrm{ch}}(\cA)\) and Verdier duality
construct \(\cA^!\).  A chain pairing
\(Q_\cA\otimes Q_{\cA^!}\to\omega_X[d]\) compares the two
Hochschild complexes.  These three maps carry the inversion, the
Koszul reflection, and the Hochschild duality in their respective
categories.
\end{remark}

\begin{theorem}[Finite support and the Hochschild--Hilbert polynomial;
\ClaimStatusConditional]
\label{thm:hochschild-polynomial-growth}
\index{Hochschild cohomology!polynomial growth|textbf}
Assume \(H_H(\cA;S)\), where
\(S\subset\mathbb Z_{\geq0}\) is finite, and assume
\(H^n(K_{\cA,S})\) is finite-dimensional.  Then
\begin{equation}\label{eq:hoch-hilbert-polynomial}
 P_\cA(t)
 :=\sum_{n\geq0}\dim\ChirHoch^n(\cA)t^n
 =\sum_{n\in S}\dim H^n(K_{\cA,S})t^n.
\end{equation}
Thus \(P_\cA(t)\) is a polynomial of degree at most \(\max S\).
If \(\cA^!\) carries a family support datum
\(H_H(\cA^!;S^!)\) with finite-dimensional cohomology and
\((\cA,\cA^!)\) carries a perfect degree-\(d\) chain pairing, then
\[
 P_\cA(t)=t^dP_{\cA^!}(t^{-1}).
\]
\end{theorem}

\begin{proof}
Theorem~\ref{thm:main-koszul-hoch} identifies
\(\ChirHoch^n(\cA)\) with \(H^n(K_{\cA,S})\), giving
\eqref{eq:hoch-hilbert-polynomial}.  A perfect degree-\(d\) pairing
identifies the degree-\(n\) term with the dual degree-\(d-n\) term;
taking dimensions gives the palindromic identity.
\end{proof}

\begin{remark}[Deformation-theoretic reading of the amplitude]
The cyclic deformation object gives the first three columns
\[
Z(\cA)=\ChirHoch^0(\cA),\qquad
\operatorname{aut}^{\inf}_{\cA}
=\ChirHoch^1(\cA)
=\Der_{\mathrm{ch}}(\cA)/\Inn_{\mathrm{ch}}(\cA),\qquad
T^{\mathrm{prod}}_\cA
\subseteq
\ChirHoch^2(\cA).
\]
Here \(\ChirHoch^1\) is the fixed-product automorphism space, and
the shifted degree-two cyclic column contains the binary product
tangent.  Higher Hochschild groups contain the higher extension data
of the deformation problem.  The family support complex
\(K_{\cA,S}\) determines which of these columns occur.  In
particular, the bounded Virasoro calculation has a degree-three class,
while the rank-one even superboson is supported in degrees zero and
one.  Formula~\eqref{eq:hoch-hilbert-polynomial} records the support
of the chosen family model.
\end{remark}

\begin{remark}[Geometric substrate (Volume~II)]
\label{rem:theorem-h-lagrangian}
\index{Lagrangian self-intersection!Theorem H}
Given an HKR comparison for a Lagrangian embedding
\(\mathcal L\hookrightarrow\mathcal M\), the normal complex of the
derived self-intersection
\(\mathcal L\times_{\mathcal M}\mathcal L\) provides a candidate
for \(K_{\cA,S}\).  The HKR quasi-isomorphism then transports its
support to the chiral Hochschild complex.
\end{remark}

\begin{remark}[Chiral vs classical Hochschild: structural vs dimensional]
\label{rem:critical-level-dimensional-divergence}
Three cohomology theories already display different support:
\begin{enumerate}[label=\textup{(\roman*)},nosep]
\item Bakalov--De Sole--Kac bounded vertex cohomology of
 \(\operatorname{Vir}_c\) has one-dimensional groups in degrees
 \(0,2,3\) \cite[Theorem~7.2]{BDSK21}.
\item Classical Hochschild cohomology of the first Weyl algebra has
 \(\HH^0(A_1)=\bC\) and positive-degree groups equal to the zero
 space \cite{Sridharan61}.
\item Goncharova's calculation of the continuous cohomology of the
 positive Virasoro algebra \(L_+\) gives two classes in every positive
 degree, in weights \((3q^2\mp q)/2\); see Fuks,
 \emph{Cohomology of Infinite-Dimensional Lie Algebras}.
\end{enumerate}
Each theory retains the operations and variance of its source.  A
comparison with \(C^\bullet_{\mathrm{ch}}(\cA,\cA)\) is carried by a
map from Definition~\ref{def:theorem-h-comparison-datum}.  At affine
critical level the Feigin--Frenkel centre
\(\operatorname{Fun}(\operatorname{Op}_{\fg^\vee}(D))\) belongs to
degree zero of a completed chart model.  A critical-level Theorem~H
therefore asks for a completed complex \(Q_\cA\), a support model
\(K_{\cA,S}\), and the homotopy of
\eqref{eq:theorem-h-global-retract}.
\end{remark}

\begin{corollary}[Deformation-obstruction exchange at genus \texorpdfstring{$0$}{0}; \ClaimStatusConditional]\label{cor:def-obs-exchange-genus0}
\index{deformation-obstruction!genus-0 exchange}
Let $(\mathcal{A}, \mathcal{A}^!)$ be a PBW chiral Koszul pair on a
smooth projective curve~$X$, carrying family support data on both
sides and a perfect degree-two chain pairing from
Definition~\textup{\ref{def:theorem-h-shift-normalization-package}}.
Then degree~$2$ of the unshifted curve-level chiral Hochschild complex of
$\mathcal{A}$ is dual to the ordinary chiral centre of
$\mathcal{A}^!$:
\begin{equation}\label{eq:def-obs-genus0}
\ChirHoch^2(\mathcal{A}) \;\cong\; \ChirHoch^0(\mathcal{A}^!)^{\vee} \otimes \omega_X.
\end{equation}
In Gerstenhaber product-deformation grading this is the
product-deformation degree; in the shifted formal-moduli convention
used for the MC problem, the product-moduli tangent is the shifted
degree-$2$ cyclic class, while unshifted
$\ChirHoch^1(\mathcal A)$ remains the fixed-product derivation degree.
More generally, the genus-$0$ curve-level chiral deformation complex of
$\mathcal{A}$ is dual to the corresponding curve-level obstruction
complex of $\mathcal{A}^!$ with a shift by~$2$.
\end{corollary}

\begin{proof}
Set $n = 2$ in Theorem~\ref{thm:main-koszul-hoch}:
\[
\ChirHoch^2(\mathcal{A})
\cong
\ChirHoch^{0}(\mathcal{A}^!)^\vee \otimes \omega_X .
\]
Here $\ChirHoch^2(\mathcal{A})$ is the curve-level chiral
product-deformation degree, equivalently the primary obstruction
target after shifting to the formal-moduli convention, and
$\ChirHoch^0(\mathcal{A}^!)=Z(\mathcal{A}^!)$ is the ordinary
chiral centre of the Koszul dual.  Its domain is the chart cochain
complex; topological, categorical, and product-ambient CY Hochschild
complexes enter through their respective comparison maps.
\end{proof}

\begin{remark}\label{rem:def-obs-higher-genus}
At genus $g \geq 1$, the exchange is Theorem~\ref{thm:quantum-complementarity-main} (Theorem~C): $Q_g(\mathcal{A}) \oplus Q_g(\mathcal{A}^!) \cong H^*(\overline{\mathcal{M}}_g, Z(\mathcal{A}))$. The genus-$0$ corollary is the seed from which higher-genus complementarity grows via the Faber--Pandharipande mechanism (Theorem~\ref{thm:genus-universality}).
\end{remark}

\begin{example}[Heisenberg specialization]\label{ex:heisenberg-curved-specialization}
For the Heisenberg algebra $\mathcal{H}_k$, the curved
$A_\infty$ structure on the Koszul dual
$\mathrm{Sym}^{\mathrm{ch}}(V^*)$ has
$m_0=-k\omega$ and $m_n=0$ for $n\geq3$
(\S\ref{sec:frame-genus1}).  This describes the curved Koszul-dual
chart.  The Hochschild support is supplied separately by a family
support datum.  For the rank-one even superboson, the
Bakalov--De Sole--Kac calculation and a bounded-to-chart
quasi-isomorphism give the polynomial $2+t$ of
Example~\ref{ex:HH-heisenberg-complete}.
\end{example}


%================================================================
% FORMALITY, MODULI, AND OPERADIC PERSPECTIVES ON KOSZULNESS
%================================================================

\subsection{Formality, moduli, and operadic perspectives on Koszulness}
\label{subsec:formality-moduli-operadic}

The modular
convolution algebra $\gAmod$ carries the universal Maurer--Cartan
element $\Theta_{\cA}$
(Theorem~\ref{thm:mc2-bar-intrinsic}); every structural property
of~$\cA$ is a property of~$\Theta_{\cA}$. Four
formality-type conditions on~$\gAmod$ and their consequences for
Koszulness follow.

\medskip

\begin{proposition}[Conditional $\Etwo$-formality criterion for chiral Hochschild cohomology;
\ClaimStatusConditional]
\label{prop:e2-formality-hochschild}
\index{formality!$\Etwo$-chiral Hochschild}
\index{Hochschild cohomology!$\Etwo$-formality}
Let $\cA$ be a chirally Koszul algebra on a smooth projective
curve~$X$. Assume a brace-compatible contraction, equivalently an
$\Etwo$ quasi-isomorphism from $\ChirHoch^\bullet(\cA)$ to its
cohomology Gerstenhaber algebra. Then:
\begin{enumerate}[label=\textup{(\roman*)}]
\item $\ChirHoch^*(\cA)$ is formal as an $\Etwo$-algebra:
 the $\Etwo$ quasi-isomorphism type of $\ChirHoch^*(\cA)$ is
 determined by the graded algebra $H^*(\ChirHoch(\cA))$ with its
 cup product and degree-$1$ bracket.

\item All higher $\Etwo$-operations \textup{(}braces of
 degree~$\geq 3$, Massey products, $A_\infty$ corrections\textup{)}
 are cohomologically trivial.
\end{enumerate}
\end{proposition}

\begin{proof}
Choose the brace-compatible contraction in the hypothesis. Homological
transfer then identifies the $\Etwo$-algebra
$\ChirHoch^\bullet(\cA)$ with the cohomology Gerstenhaber algebra
equipped with the transferred $\Etwo$-structure. The additional
hypothesis says precisely that this transferred structure is
quasi-isomorphic, as an $\Etwo$-algebra, to the strict Gerstenhaber
model determined by the cup product and bracket. This gives~(i).

For~(ii), the same $\Etwo$ quasi-isomorphism sends each higher brace
or Massey operation of arity at least~$3$ to zero in the strict
Gerstenhaber model; equivalently, those operations are coboundaries in
the original brace dg algebra.  When a family support datum
$H_H(\cA;S)$ is also supplied,
Theorem~\ref{thm:hochschild-polynomial-growth} locates the cohomology
in~$S$.  The support datum and the brace-compatible contraction are
independent structures: the first locates the cohomology, and the
second contracts the transferred higher operations.
\end{proof}

\section{Beilinson--Drinfeld diagonal factorization cochains and
chiral Hochschild}
\label{sec:bd-fact-vs-chirhoch}

\index{factorization cohomology!vs.\ chiral Hochschild|textbf}
\index{Beilinson--Drinfeld!factorization cohomology comparison|textbf}
\index{Ran space!diagonal bimodule}

There are two a priori distinct Ext-theoretic cohomologies of a
chiral algebra~$\cA$ on a smooth projective
curve~$X$:
\begin{enumerate}[label=\textup{(\roman*)}]
\item the \emph{chiral Hochschild cohomology}
$\ChirHoch^\bullet(\cA) =
\Ext^\bullet_{\cA^e}(\cA,\cA)$
\textup{(}Definition~\textup{\ref{def:chiral-hochschild-complex}},
geometric FM model;
Definition~\textup{\ref{def:chirHoch-Ext}} in
Chapter~\textup{\ref{chap:koszul-pair-structure}},
algebraic $\Ext$-model\textup{)},
which is chain-level and native to
$\mathcal{D}_X$-module theory;
\item the \emph{Beilinson--Drinfeld factorization cohomology}
of the diagonal bimodule,
\[
\mathcal{C}^\bullet_{\mathrm{fact}}(\cA,\cA)
\;:=\;
R\mathrm{Hom}_{\operatorname{Fact}(X)}
\bigl(\cA,\, \cA^{\mathrm{diag}}\bigr),
\]
the derived $\operatorname{Fact}(X)$-internal Hom from $\cA$ (as
factorization algebra on $\operatorname{Ran}(X)$) to itself as
a factorization bimodule
\textup{(}Beilinson--Drinfeld~\textup{\cite[\S3.4]{BD04}};
Francis--Gaitsgory
$(\infty,1)$-enhancement~\textup{\cite[\S5.5.2]{HA}},
Ayala--Francis~\textup{\cite[\S3-\S4]{AF15}}\textup{)}, which
is $(\infty,1)$-categorical and native to $\operatorname{Ran}(X)$.
\end{enumerate}
The question of when these coincide is structurally parallel to
the bar--cobar identification
$\Bbar_X(\cA) \simeq \int_X\cA$
\textup{(}Proposition~\textup{\ref{prop:bar-fh}} in
Chapter~\textup{\ref{chap:bar-cobar}}\textup{)} and its dual
$\Omega_X(\cC) \simeq \int^X\cC$
\textup{(}Proposition~\textup{\ref{prop:cobar-fh}} in
Chapter~\textup{\ref{chap:bar-cobar}}\textup{)}: bar/cobar are
factorization homology/cohomology of the unit bimodule when the
BD comparison hypotheses are satisfied, whereas
$\ChirHoch$ and $\mathcal{C}^\bullet_{\mathrm{fact}}(\cA,\cA)$
are the cohomological Ext of the diagonal bimodule.

\begin{theorem}[Diagonal descent for chart Hochschild cochains;
\ClaimStatusConditional]
\label{thm:bd-fact-vs-chirhoch}
\index{Theorem H!BD factorization cohomology refinement}
\textup{[Type signature: Open quadrant; Ran and
Fulton--MacPherson chain presentations; Beilinson level~$3$;
hypothesis package $D_H(\cA;S)$.]}

Let $X$ be a smooth projective curve, let $S\subset\mathbb Z$ be
finite, and let $\cA$ carry the family support datum $H_H(\cA;S)$.
The diagonal-descent datum $D_H(\cA;S)$ begins with a filtered cochain
map
\begin{equation}\label{eq:bd-chirhoch-comparison}
c_{\cA}\colon
\mathcal{C}^\bullet_{\mathrm{fact}}(\cA,\cA)
\;\longrightarrow\;
 C^\bullet_{\mathrm{ch}}(\cA,\cA).
\end{equation}
It consists further of:
\begin{enumerate}[label=\textup{(D\arabic*)}]
\item complete, exhaustive filtrations on the source and target of
$c_\cA$ whose spectral sequences converge strongly;
\item finite-window, incidence-compatible stratum retracts as in
Definition~\ref{def:theorem-h-pbw-finite-window-lane}, intertwined by
the associated-graded comparison;
\item a quasi-isomorphism between the corresponding retract targets,
compatible with the inverse systems in the conformal-weight window.
\end{enumerate}
Then:
\begin{enumerate}[label=\textup{(\alph*)}]
\item The morphism $c_\cA$ is a filtered quasi-isomorphism.  It
identifies the Ran and Fulton--MacPherson presentations in the derived
category while retaining their respective chain models.
\item The induced map
\[
 H^n(c_{\cA})\colon
 H^n\bigl(\mathcal C^\bullet_{\mathrm{fact}}(\cA,\cA)\bigr)
 \xrightarrow{\ \cong\ }
 \ChirHoch^n(\cA)
\]
is an isomorphism for every $n$, and
\[
 \operatorname{Supp}
 H^\bullet\bigl(\mathcal C^\bullet_{\mathrm{fact}}(\cA,\cA)\bigr)
 \subseteq S.
\]
\item In $\mathrm D(\mathrm{Fact}(X))$ the comparison represents an
equivalence
\[
 \mathcal C^\bullet_{\mathrm{fact}}(\cA,\cA)
 \simeq RHH_{\mathrm{ch}}(\cA).
\]
\item Suppose $\cA^!$ carries $H_H(\cA^!;S^!)$ and
$D_H(\cA^!;S^!)$, and suppose the two chart complexes carry a perfect
degree-$d$ pairing whose adjoint square commutes with $c_\cA$ and
$c_{\cA^!}$ in the derived category.  Then
\[
 H^n\bigl(\mathcal C^\bullet_{\mathrm{fact}}(\cA,\cA)\bigr)
 \cong
 H^{d-n}\bigl(\mathcal C^\bullet_{\mathrm{fact}}(\cA^!,\cA^!)\bigr)^\vee
 \otimes\omega_X.
\]
\end{enumerate}
\end{theorem}

\begin{proof}
The comparison map and its compatibility with the two differentials
belong to $D_H(\cA;S)$.  Its naturality is taken with respect to
morphisms equipped with compatible diagonal-bimodule maps.

\emph{Part~(a).}\;
The stratum retracts in~(D2) totalize by
Proposition~\ref{prop:fm-tower-collapse}.  Condition~(D3) makes the
associated-graded morphism of the two filtered complexes a
quasi-isomorphism.  Strong convergence in~(D1) and the comparison
theorem for filtered complexes show that $c_\cA$ is a filtered
quasi-isomorphism.

\emph{Part~(b).}\;
Part~(a) identifies the two cohomologies.  The family support datum
$H_H(\cA;S)$ and Theorem~\ref{thm:main-koszul-hoch} identify the chart
cohomology with $H^\bullet(K_{\cA,S})$, giving the asserted support.

\emph{Part~(c).}\;
Localization of $\mathrm{Fact}(X)$ at quasi-isomorphisms carries the
filtered quasi-isomorphism of part~(a) to the displayed equivalence in
$\mathrm D(\mathrm{Fact}(X))$.

\emph{Part~(d).}\;
The perfect degree-$d$ pairing gives the chart duality of
Theorem~\ref{thm:main-koszul-hoch}:
\[
\ChirHoch^n(\cA)
\cong
\ChirHoch^{d-n}(\cA^!)^\vee\otimes\omega_X.
\]
Commutativity of the adjoint square transports this isomorphism across
$c_\cA$ and $c_{\cA^!}$.
\end{proof}

\begin{remark}[Ran and Fulton--MacPherson presentations]
\label{rem:bd-fact-chirhoch-two-lanes}
\index{factorization cohomology!two lanes}
The Ran internal-Hom and the Fulton--MacPherson chart retain their
respective operations and filtrations.  Under $D_H(\cA;S)$,
Theorem~\ref{thm:bd-fact-vs-chirhoch}(a) relates them by a filtered
quasi-isomorphism, part~(b) identifies their cohomology, and part~(c)
identifies their images in the derived factorization category.
Chapter~\ref{chap:higher-genus} uses the cohomological comparison;
the Ran-derived enhancements there and the topological factorization
comparisons in Volume~II use the derived comparison.
\end{remark}

\begin{corollary}[Theorem~H transported along diagonal factorization
descent;
\ClaimStatusConditional]
\label{cor:thm-H-bd-extension}
\index{Theorem H!BD extension}
Let $\cA$ carry $H_H(\cA;S)$ and $D_H(\cA;S)$.  Then
\begin{enumerate}[label=\textup{(\roman*)}]
\item
\[
 \operatorname{Supp}
 H^\bullet\bigl(\mathcal C^\bullet_{\mathrm{fact}}(\cA,\cA)\bigr)
 \subseteq S.
\]
\item If $\cA^!$ carries $H_H(\cA^!;S^!)$ and
$D_H(\cA^!;S^!)$, and the chart complexes carry a perfect degree-$d$
pairing compatible with the diagonal-descent maps, then
\[
 H^n\bigl(\mathcal C^\bullet_{\mathrm{fact}}(\cA,\cA)\bigr)
 \cong
 H^{d-n}\bigl(\mathcal C^\bullet_{\mathrm{fact}}(\cA^!,\cA^!)\bigr)^\vee
 \otimes\omega_X.
\]
\item Under the finite-dimensionality hypotheses of
Theorem~\ref{thm:hochschild-polynomial-growth}, the factorization
Hilbert series is a polynomial and the perfect pairing in~\textup{(ii)}
gives
\[
 P^{\mathrm{fact}}_{\cA}(t)
 =t^dP^{\mathrm{fact}}_{\cA^!}(t^{-1}).
\]
\end{enumerate}
\end{corollary}

\begin{proof}
Immediate from
Theorem~\textup{\ref{thm:bd-fact-vs-chirhoch}}\textup{(b, d)}
and Theorem~\textup{\ref{thm:hochschild-polynomial-growth}}.
\end{proof}

\begin{remark}[Cross-volume: Vol II topological factorization]
\label{rem:vol2-topological-factorization-comparison}
\index{cross-volume!Vol II SC$^{\mathrm{ch,top}}$ coproducts}
Volume~II defines the topological factorization algebra
$\mathrm{SC}^{\mathrm{ch,top}}$ and identifies its bar differential
with holomorphic factorization and its coproduct with topological
factorization (Vol~II Parts~I--VII, with the cross-volume scope
discipline recorded there). When the diagonal comparison of
Theorem~\textup{\ref{thm:bd-fact-vs-chirhoch}}(c) is available, the
$(\infty,1)$-enhancement $\mathcal{C}^\bullet_{\mathrm{fact}}(\cA,\cA)\simeq
RHH_{\mathrm{ch}}(\cA)$ intertwines the Vol~II
$\mathrm{SC}^{\mathrm{ch,top}}$-coproduct on the BD side with the
brace coproduct on the chiral Hochschild side
\textup{(}Theorem~\textup{\ref{thm:chiral-deligne-tamarkin}}\textup{)}.
The Vol~I $\leftrightarrow$ Vol~II compatibility therefore rests on
this diagonal comparison: it is the typed map carrying the Vol~I
chart brace coproduct to the Vol~II topological coproduct.
\end{remark}

\begin{remark}[Support and $E_2$-formality]
\label{rem:e2-formality-vs-thmH}
The family datum $H_H(\cA;S)$ locates the cohomology in~$S$.
Proposition~\ref{prop:e2-formality-hochschild} supplies the additional
brace-compatible contraction that identifies the transferred
$\Etwo$-structure with its strict Gerstenhaber model.  Under that
contraction, the secondary Borcherds operations $F_n$ for $n\geq3$
from Theorem~\ref{thm:cech-hca} become exact.  A converse direction is
therefore a construction of this brace contraction from Koszul or
scalar-orbit data; Theorem~\ref{thm:convolution-formality-one-channel}
formulates one conditional realization.
\end{remark}

\begin{theorem}[Scalar universal class implies convolution
formality along its distinguished orbit;
\ClaimStatusConditional]
\label{thm:convolution-formality-one-channel}
\index{formality!convolution algebra}
\index{one-channel!scalar-orbit formality problem}
\index{scalar saturation|see{one-channel}}
Let $\cA$ be a modular Koszul chiral algebra.
The modular convolution algebra
\[
\gAmod
\;=\;
\prod_{g,n}\,
\operatorname{Hom}_{\Sigma_n}\!\bigl(
 C_*(\overline{\cM}_{g,n}),\;
 \operatorname{End}_{\cA}(n)
\bigr)
\]
carries the five-component differential
$D = d_{\mathrm{int}} + [\tau,-] + d_{\mathrm{sew}}
 + d_{\mathrm{pf}} + \hbar\Delta$
and the universal Maurer--Cartan element
$\Theta_{\cA} := D_{\cA} - d_0
\in \MC(\gAmod)$
\textup{(}Theorem~\textup{\ref{thm:mc2-bar-intrinsic})}.
Assume the full universal Maurer--Cartan element is scalar:
\begin{equation}\label{eq:one-channel-formality}
\Theta_{\cA}
\;=\;
\kappa(\cA) \cdot \eta \otimes \Lambda.
\end{equation}
Then the quantum $L_\infty$-structure is formal along the
distinguished scalar Maurer--Cartan orbit of\/~$\Theta_{\cA}$:
the higher transferred brackets vanish on the one-dimensional
scalar line $\mathbb{C}\!\cdot\!\eta \otimes \Lambda$, and the MC
equation restricts there to the scalar genus identities.
A converse requires a derivation of the scalar identity
\eqref{eq:one-channel-formality} from the one-channel minimal-model
formula $\Theta_{\cA}^{\min}=\eta\otimes\Gamma_\cA$; this is the
remaining tautological-purity obligation.
\end{theorem}

\begin{proof}
Throughout this proof we write $\kappa$ for $\kappa(\cA)$.
Suppose $\Theta_{\cA} = \kappa \cdot \eta \otimes \Lambda$.
By Theorem~\ref{thm:operadic-complexity-detailed},
one-channel forces shadow depth $r_{\max}(\cA) = 2$,
hence $A_\infty$-depth $d_\infty(\cA) = 2$ and
$L_\infty$-formality level $f_\infty(\cA) = 2$.
The genus-$0$ transferred $L_\infty$ brackets therefore
satisfy $\ell_k^{(0),\mathrm{tr}} = 0$ for $k \geq 3$:
the genus-$0$ convolution algebra is formal.

For the full scalar orbit, the shadow
obstruction tower terminates at degree~$2$:
$\Theta_{\cA}^{(\le r)} = \kappa \cdot \eta \otimes \Lambda$
for all~$r$. The obstruction classes
$o^{(r+1)}(\cA) \in H^2(F^{r+1}\gAmod / F^{r+2}\gAmod, d_2)$
vanish for all $r \geq 2$
(Theorem~\ref{thm:nms-all-degree-master-equation}).
Since the genus-$g$ quantum brackets $\ell_k^{(g)}$ for
$g \geq 1$ are induced by sewing on $\overline{\cM}_{g,n}$,
and sewing applied to the scalar element
$\kappa \cdot \eta \otimes \Lambda$ produces the scalar
genus-$g$ curvature $\kappa^g \cdot \eta_g \otimes \Lambda_g$
(which lies in the image of the differential), the transferred
higher-genus brackets
$\ell_k^{(g),\mathrm{tr}}$ also vanish for $k \geq 3$.
This gives formality of the distinguished scalar orbit of
$\Theta_{\cA}$ inside~$\gAmod$.
\end{proof}

\begin{remark}[Scalar-orbit formality and scalar saturation]
\label{rem:convolution-formality-cg}
\index{Chriss--Ginzburg principle!convolution formality}
\emph{Forward direction} (writing $\kappa$ for $\kappa(\cA)$).\;
If $\Theta_{\cA} = \kappa \cdot \eta \otimes \Lambda$, the
MC equation $D \Theta + \tfrac{1}{2}[\Theta, \Theta] = 0$
reduces to a scalar identity at each $(g,n)$:
$[\eta \otimes \Lambda, \eta \otimes \Lambda]$ factors through the
metric contraction, so the higher $L_\infty$-brackets $\ell_k$
($k \geq 3$) act trivially on
$\bC \cdot \eta \otimes \Lambda$. Hence
$\gAmod$ is formal along the orbit of~$\Theta_{\cA}$.

\emph{Converse / open direction.}\;
The converse problem is to derive scalar saturation from formality of
the full quantum $L_\infty$-algebra and to construct formality from
one-channel data. The strongest proved input is the one-channel
minimal-model statement
$\Theta_{\cA}^{\min}=\eta\otimes\Gamma_{\cA}$
\textup{(}Corollary~\textup{\ref{cor:scalar-saturation}}\textup{)};
identifying $\Gamma_{\cA}$ with
$\kappa(\cA)\Lambda$ is the open tautological-purity step outside
the proved scalar locus.

\emph{Relation to the shadow obstruction tower.}\;
The non-scalar universal classes
of~\S\ref{subsec:non-scalar-theta} measure the deviation of
$\Theta_{\cA}$ from the scalar line. The cubic shadow~$C$, quartic resonance class~$Q$,
and all higher shadows $\mathrm{Sh}_r$ are the
\emph{obstructions to formality}: the higher-order
corrections preventing $\Theta_{\cA}$ from lying on the
scalar line.

\emph{Pillar~C interaction} (Mok~\cite{Mok25}):
the planted-forest coefficient algebra $G_{\mathrm{pf}}$
(Definition~\ref{def:planted-forest-coefficient-algebra}) controls
the non-formal corrections. Vanishing planted-forest
contributions removes one source of deviation from the scalar
orbit.  A converse criterion for full convolution formality requires
control of the remaining planted-forest contributions
(Remark~\ref{rem:conditional-homotopy-why}).
\end{remark}

\begin{remark}[Hierarchy: Koszulness, modular Koszulness, formality]
\label{rem:formality-vs-koszul}
Three properties of increasing strength:
\begin{enumerate}[label=\textup{(\alph*)}]
\item \emph{Koszulness} (genus~$0$): the bar spectral sequence
 collapses at~$E_2$
 (Theorem~\ref{thm:bar-concentration}).

\item \emph{Modular Koszulness} (all genera): MK3 extends PBW
 collapse; the scalar trace $\kappa(\cA)$ is the sole survivor
 in the modular bar cohomology at $g \geq 1$.

\item \emph{Scalar-orbit convolution formality} (structural, on the
 proved scalar locus): when the full universal class is scalar,
 the quantum $L_\infty$-structure is formal along its
 distinguished scalar orbit, by
 Theorem~\ref{thm:convolution-formality-one-channel}.
\end{enumerate}
On the proved scalar locus, the implication
$\text{(c)} \Rightarrow \text{(b)} \Rightarrow \text{(a)}$
holds. On its complement, the first converse-type problem is
exactly the non-scalar shadow obstruction tower
(\S\ref{subsec:non-scalar-theta}); the second is the
higher-genus curvature.
\end{remark}

\begin{proposition}[MC tangent support from the chart;
\ClaimStatusConditional]
\label{prop:mc-tangent-amplitude}
\index{Maurer--Cartan!tangent amplitude}
Let \(\cA\) be a chiral algebra on a smooth projective curve with
universal Maurer--Cartan element
\(\Theta_\cA\in\operatorname{MC}(\gAmod)\).  Assume
\(H_H(\cA;S)\) and the chart-to-shadow quasi-isomorphism
\(\rho_\cA\) of Proposition~\ref{prop:hochschild-shadow-projection}.
Then the curve-level tangent summand satisfies
\[
 H^n\bigl(\mathbb T^{\mathrm{curve}}_{\Theta_\cA}\bigr)
 \cong H^n(K_{\cA,S}),
 \qquad
 \operatorname{Supp}
 H^\bullet\bigl(\mathbb T^{\mathrm{curve}}_{\Theta_\cA}\bigr)
 \subseteq S.
\]
After passage to a minimal \(L_\infty\)-model, smoothness in the curve
direction is equivalent to the vanishing of the Kuranishi maps
\begin{equation}\label{eq:mc-smoothness}
 \kappa_m\colon
 \operatorname{Sym}^m
 H^1\bigl(\mathbb T^{\mathrm{curve}}_{\Theta_\cA}\bigr)
 \longrightarrow
 H^2\bigl(\mathbb T^{\mathrm{curve}}_{\Theta_\cA}\bigr),
 \qquad m\geq2.
\end{equation}
In particular, \(2\notin S\) makes every displayed target the zero
space.
\end{proposition}

\begin{proof}
The chart-to-shadow map identifies the curve-level tangent summand
with \(C^\bullet_{\mathrm{ch}}(\cA,\cA)\), and Theorem~H identifies the
latter with \(K_{\cA,S}\).  Homotopy transfer places the higher
brackets on cohomology.  The Maurer--Cartan equation in the minimal
model is precisely the family of Kuranishi maps
\eqref{eq:mc-smoothness}.
\end{proof}

\begin{remark}[The obstruction target and the Kuranishi maps]
\label{rem:mc-smoothness-pbw}
At genus zero, the chart comparison identifies the obstruction target
with
\[
 H^2(K_{\cA,S})\cong\ChirHoch^2(\cA).
\]
If \((\cA,\cA^!)\) carries a perfect degree-two chain pairing, this
becomes
\[
 \ChirHoch^2(\cA)
 \cong Z(\cA^!)^\vee\otimes\omega_X.
\]
At positive genus, the corresponding targets are cohomology groups of
the modular bar complex and are governed by its own filtration and
support data.  Smoothness is the vanishing of the transferred
Kuranishi maps into these targets.  A converse construction of
Koszulness from smoothness would associate a nonzero Kuranishi map to
each surviving class of the bar spectral sequence.
\end{remark}

\begin{theorem}[Bifunctor obstruction decomposition;
\ClaimStatusConditional]
\label{thm:bifunctor-obstruction-decomposition}%
\index{bifunctor obstruction!decomposition}
\index{Robert-Nicoud--Wierstra!no-bifunctor obstruction}
The convolution $sL_\infty$-algebra
$\operatorname{hom}_\alpha(\barB^{\mathrm{ch}}(\cA),\, \cA)$
is functorial in each slot separately but fails as a bifunctor
\textup{(\cite{RNW19}, Theorem~6.6;
Remark~\textup{\ref{rem:two-level-convention})}}.
For a chirally Koszul algebra~$\cA$:
\begin{enumerate}[label=\textup{(\roman*)}]
\item The degree-$3$ obstruction class
\[
\mathrm{obs}_3
\;\in\;
H^2\!\bigl(
 \operatorname{hom}_\alpha^{(3)}
 (\barB^{\mathrm{ch}}(\cA),\, \cA)
\bigr)
\]
is nonzero in general, but is immaterial for the MC3 categorical
lift: Koszulness makes one-slot factorization sufficient.

\item Every two-sided $\infty$-morphism that fails to exist as a
 bifunctor morphism decomposes as a one-sided coalgebra morphism
 followed by a one-sided algebra morphism, with the composition
 independent of ordering up to homotopy.
\end{enumerate}
\end{theorem}

\begin{proof}
Part~(i): the obstruction class $\mathrm{obs}_3$ is the failure of
$\operatorname{hom}_\alpha$ to extend to $\infty$-morphisms in both
slots simultaneously (\cite{RNW19}, Theorem~6.6). For chirally
Koszul~$\cA$, the bar-cobar counit
$\Omega(\barB(\cA)) \xrightarrow{\sim} \cA$ is a quasi-isomorphism;
the one-slot convention
(Remark~\ref{rem:two-level-convention}) fixes the cooperad slot and
varies the algebra by $\infty$-morphisms. The bar-cobar inversion
gives a canonical strict model, so both-slot variation is
unnecessary: the MC3 categorical lift proceeds one slot at a time
(Theorem~\ref{thm:mc3-type-a-resolution}).

Part~(ii): a two-sided $\infty$-morphism $(f, g)$ with $f$ on the
coalgebra side and $g$ on the algebra side satisfies
$\partial f + f \star g = 0$; by the one-slot convention, this
factors as $f$~followed by~$g$ with the composition determined up
to homotopy by the contractibility of the space of
$\infty$-morphisms between quasi-isomorphic objects.
\end{proof}

\begin{remark}[Heisenberg test case and the MC3 frontier]
\label{rem:bifunctor-analysis}
\emph{Testability.}\;
For the Heisenberg algebra $\cH_k$, the obstruction
$\mathrm{obs}_3$ lives in the two-dimensional space
$H^2(\operatorname{hom}_\alpha^{(3)}
 (\barB^{\mathrm{ch}}(\cH_k), \cH_k))$.
An explicit computation: $\barB^{\mathrm{ch}}(\cH_k)$ is
concentrated in bar degree~$1$
(Theorem~\ref{thm:bar-concentration}),
so $\operatorname{hom}_\alpha^{(3)}$ reduces to maps from
the triple self-product of the cogenerators. The obstruction is
the failure of the Leibniz rule for the simultaneous
coalgebra-algebra functoriality, which is nonzero even for this
simplest example because the Koszul duality functor itself is
not strict-monoidal.

\emph{Structural origin.}\;
The no-bifunctor theorem of~\cite{RNW19} (\S6) is universal:
$\mathrm{obs}_3 \neq 0$ for \emph{any} cooperad-operad pair.
The content of part~(ii) is that Koszulness provides a
\emph{one-slot workaround}: the Quillen equivalence
$\Omega_\kappa \dashv \barB_\kappa$
(Theorem~\ref{thm:quillen-equivalence-chiral})
transports $\infty$-morphisms through one slot at a time.
This is why the MC3 resolution
(Theorem~\ref{thm:categorical-cg-all-types}, Corollary~\ref{cor:mc3-all-types}) succeeds: each step
on the proved all-types domain
\textup{(}categorical prefundamental Clebsch--Gordan via
multiplicity-free $\ell$-weights together with the generated-core DK
comparison\textup{)} is a one-slot operation. The later chromatic,
completion, and compact-extension steps belong to the separate
post-CG hypotheses, reduced in type~$A$ and open in the remaining
types.
The all-types generalization therefore proceeds one slot at a time on
the evaluation-generated core; the decomposition property of
part~(ii) explains why this is sufficient on that proved domain.
\end{remark}

\begin{remark}[Operadic inheritance of Koszulness]
\label{rem:operadic-inheritance}
\index{Koszulness!operadic inheritance}
The descent of chiral Koszulness from Koszulness of the chiral
operad $P^{\mathrm{ch}}$ is a separate theorem, not an implicit
consequence of Theorem~\ref{thm:e1-chiral-koszul-duality}. Its proof
requires formality of $\overline{\mathrm{FM}}_n(X)$ as operadic
objects on curves and the transfer statement recorded in
Theorem~\ref{thm:operadic-complexity}; the associative case is the
one-slot specialization implemented by
Lemma~\ref{lem:operadic-koszul-transfer}.
\end{remark}

\section{Example: boson-fermion duality}

\subsection{The free boson chiral algebra}

The free boson $\mathcal{B} = \mathcal{H}_k$ at level~$k$ on a curve $X$ is defined as follows:

\emph{As a chiral algebra.}
The free boson $\mathcal{B}$ is the chiral envelope of the Heisenberg Lie$^*$ algebra
$\omega_X$ (with Lie$^*$ bracket given by the residue pairing). Its vacuum module is
$\mathrm{Sym}(\omega_X \otimes t^{-1}\mathbb{C}[t^{-1}])$, a polynomial algebra on
infinitely many generators $\alpha_{-n}$ ($n \geq 1$).

\emph{Generator.} The field $\alpha(z)$ generates $\mathcal{B}$ with conformal weight $h = 1$.

\emph{Chiral multiplication.} Determined by the OPE
\[
\alpha(z_1) \alpha(z_2) = \frac{k}{(z_1-z_2)^2} + \text{regular}
\]

In terms of modes $\alpha(z) = \sum_{n \in \mathbb{Z}} \alpha_n z^{-n-1}$:
\[
[\alpha_m, \alpha_n] = k \cdot m \delta_{m+n,0}
\]
The central charge is $c = 1$ (at level $k = 1$; in general $c = 1$ for any~$k$, since the Heisenberg algebra has $c = 1$ independent of level).

\emph{Vacuum representation.} The Fock space
\[
\mathcal{F}_{\mathcal{B}} = \mathbb{C}[\alpha_{-1}, \alpha_{-2}, \ldots] |0\rangle
\]
with $\alpha_n |0\rangle = 0$ for $n \geq 0$.

\subsection{The free fermion chiral algebra}

The free fermion $\mathcal{F}$ has:

\emph{Generators.} Two fermionic fields $\psi(z), \psi^*(z)$ with $h = 1/2$.

\emph{Relations.} The OPEs
\begin{align}
\psi(z_1)\psi^*(z_2) &= \frac{1}{z_1-z_2} + \text{regular} \\
\psi(z_1)\psi(z_2) &= 0 + \text{regular} \\
\psi^*(z_1)\psi^*(z_2) &= 0 + \text{regular}
\end{align}

In modes (half-integer for Neveu--Schwarz sector):
\begin{align}
\{\psi_r, \psi^*_s\} &= \delta_{r+s,0} \\
\{\psi_r, \psi_s\} &= 0 \\
\{\psi^*_r, \psi^*_s\} &= 0
\end{align}

\emph{Fock space.}
\[
\mathcal{F}_{\mathcal{F}} = \Lambda^{\bullet}(\psi_{-1/2}, \psi_{-3/2}, \ldots, \psi^*_{-1/2}, \psi^*_{-3/2}, \ldots) |0\rangle
\]

\subsection{Boson--fermion correspondence as a lattice VOA extension}

\begin{theorem}[Boson-fermion correspondence via lattice VOA; \ClaimStatusProvedElsewhere{} \cite{FK80}]\label{thm:boson-fermion-lattice}
The free boson $\mathcal{B}$ and free fermion $\mathcal{F}$ are related by bosonization through the rank-1 lattice vertex algebra $V_{\mathbb{Z}}$:
\[
V_{\mathbb{Z}} \supset \mathcal{B}, \qquad \mathcal{F} \cong V_{\mathbb{Z}} \quad \text{(as super vertex algebras)}
\]
\end{theorem}

\begin{remark}
Bosonization is the extension morphism
$\mathcal B\hookrightarrow V_{\mathbb Z}\cong\mathcal F$, obtained by
adjoining the charge sectors.  Koszul duality is the bar--Verdier
construction.  For $\kappa\neq0$ it sends the Heisenberg chart to the
curved second-kind branch
$\bigl(\mathrm{Sym}^{\mathrm{ch}}(V^*[1]),m_0=-\kappa\omega\bigr)$
under the quadratic comparison of
Lemma~\ref{lem:curved-dual-centre-heisenberg}.  Thus the lattice
extension and the Koszul-dual construction have different source and
target categories.  The vertex operators
$\psi={:}e^{i\phi}{:}$ and $\psi^*={:}e^{-i\phi}{:}$ belong to the
charge-$\pm1$ summands of $V_{\mathbb Z}$.
\end{remark}

\begin{proof}
The argument has three levels:

\emph{Level 1: The Lattice Extension}

The free boson $\mathcal{B}$ is generated by $\alpha(z) = i\partial\phi(z)$ with OPE $\alpha(z)\alpha(w) \sim k/(z-w)^2$. The rank-1 lattice VOA $V_{\mathbb{Z}}$ extends $\mathcal{B}$ by adjoining vertex operators $e^{in\phi}$ for $n \in \mathbb{Z}$:
\[
V_{\mathbb{Z}} = \bigoplus_{n \in \mathbb{Z}} \mathcal{F}_n, \qquad \mathcal{F}_n = \mathcal{B} \cdot {:}e^{in\phi}{:}
\]
The charge-$(\pm 1)$ sectors produce the fermionic fields $\psi = {:}e^{i\phi}{:}$ and $\psi^* = {:}e^{-i\phi}{:}$.

\emph{Level 2: Bosonization}

The explicit isomorphism is given by bosonization:
\begin{align}
\psi(z) &= :e^{i\phi(z)}: \\
\psi^*(z) &= :e^{-i\phi(z)}: \\
\alpha(z) &= i\partial\phi(z)
\end{align}

where $\phi$ is the bosonic field with $\phi(z)\phi(w) \sim -\log(z-w)$.

This realizes the isomorphism at the level of vertex operators:
\[
Y_{\mathcal{F}}(\psi, z) = :e^{i\int^z \alpha}: \quad \text{(fermion as exponential of boson)}
\]

\emph{Level 3: Bar Complex of the Heisenberg}

The quadratic presentation gives a comparison
$q_{\mathcal H_\kappa}\colon\mathcal H_\kappa^i\to
\bar B^{\mathrm{ch}}(\mathcal H_\kappa)$.  When this comparison is a
quasi-isomorphism, Verdier duality forms the curved second-kind algebra
\[
\mathcal{H}_\kappa^!
\simeq(\mathrm{Sym}^{\mathrm{ch}}(V^*[1]),m_0=-\kappa\omega)
\]
(Part~\ref{part:bar-cobar-engine}).  The cobar equivalence
$\Omega(\bar B(\mathcal H_\kappa))\simeq\mathcal H_\kappa$ is the
separate inversion map.  These maps give the bar--Verdier type
signature of the Heisenberg chart.
\end{proof}

\subsection{Boson and fermion chart complexes and the lattice comparison}

\begin{computation}[Rank-one even superboson benchmark;
\ClaimStatusProvedElsewhere]
\label{comp:boson-hochschild}
For the rank-one even superboson \(B_{\mathfrak h}\),
Bakalov--De Sole--Kac compute
\[
 H^n_{\mathrm{ch},b}(B_{\mathfrak h},B_{\mathfrak h})
 \cong
 \bigl(S^n(\Pi\mathfrak h)\bigr)^*
 \oplus
 \bigl(S^{n+1}(\Pi\mathfrak h)\bigr)^*
 \qquad\text{\cite[Theorem~7.4]{BDSK21}}.
\]
Thus its bounded cohomology has dimensions
\[
 (2,1,0,0,\ldots).
\]
A bounded-to-chart quasi-isomorphism
\(\chi_{B_{\mathfrak h}}^{\mathrm{bd}}\) transports these dimensions
to the curve chart.  The curved second-kind dual of
Lemma~\ref{lem:curved-dual-centre-heisenberg}, under its stated
oscillator comparison, is a separate Koszul-dual object; its curvature
enters a Hochschild duality statement through
the explicit chain pairing of
Definition~\ref{def:theorem-h-shift-normalization-package}.
\end{computation}

\begin{computation}[Free-fermion chart support;
\ClaimStatusConditional]
\label{comp:fermion-hochschild}
For the free fermion \(\mathcal F\), the simple pole
\[
 \psi(z)\psi^*(w)\sim\frac{1}{z-w}
\]
contributes a zero-mode residue to the collision differential.  Its
curve-level groups are computed by an incidence-compatible family
support datum \(H_H(\cA_{\mathcal F};S_{\mathcal F})\).  A bounded
vertex-cohomology computation gives the same chart groups after
\(\chi_{\mathcal F}^{\mathrm{bd}}\) has been established as a
quasi-isomorphism.  These are precisely the two routes recorded in
Example~\ref{ex:HH-fermion-complete}.
\end{computation}

\begin{remark}[Pole order and the chart differential]
\label{rem:boson-fermion-hochschild-comparison}
A simple pole contributes a zero-mode collision face, whereas a double
pole contributes the corresponding first-product face.  This
difference changes the local chart differential.  The global
cohomology is obtained after totalizing every collision stratum and
passing through the chosen completion.

The boson--fermion correspondence is realized by a lattice vertex
algebra extension.  The Koszul dual of the nonzero-level Heisenberg
chart is instead the curved second-kind branch
\[
 \bigl(\mathrm{Sym}^{\mathrm{ch}}(V^*[1]),m_0=-k\omega\bigr).
\]
Thus lattice extension and Koszul duality are represented by two
different functors and two different comparison maps.
\end{remark}

\subsection{The cyclic deformation complex}
\label{subsec:cyclic-deformation}

\begin{definition}[Cyclic deformation complex: elementary models]
\label{def:cyclic-deformation-elementary}
\index{cyclic deformation complex|textbf}
\index{Def cyc@$\operatorname{Def}_{\mathrm{cyc}}$|textbf}
Let $\cA$ be a chiral algebra on~$X$. The homotopy-invariant
\emph{cyclic deformation object} of~$\cA$ is a complete filtered cyclic
$L_\infty$-algebra controlling cyclic deformations of~$\cA$.
A strict presentation is the
\emph{cyclic deformation complex} $\Defcyc(\cA)$
used for calculations below.

\smallskip\noindent
\textbf{Abstract definition (H-level).}
$\operatorname{Def}_{\mathrm{cyc}}(\cA)$ is the Lie algebra object in
the $\infty$-category of cochain complexes whose Maurer--Cartan
elements parametrize flat cyclic deformations of~$\cA$ over
Artinian bases. For the full construction, which requires a cyclic
$L_\infty$ algebra over the modular operad, see
Theorem~\ref{thm:universal-MC}.

\smallskip\noindent
\textbf{Concrete models (M-level).}
For the examples treated in this monograph, explicit models are:

\begin{enumerate}[label=\textup{(\alph*)}]
\item \emph{Bar coderivations.}
 $\operatorname{Def}_{\mathrm{cyc}}(\cA)_{\mathrm{bar}}
 := \operatorname{CoDer}^{\mathrm{cyc}}(\bar{B}_X(\cA))$,
 the complex of cyclic coderivations of the bar construction.
 Its Maurer--Cartan elements are exactly the curved deformations
 of the bar differential.
\item \emph{Cyclic chiral Hochschild complex.}
 $\operatorname{Def}_{\mathrm{cyc}}(\cA)_{\mathrm{Hoch}}
 := CC^*_{\mathrm{ch}}(\cA)^{\mathrm{cyc}}$,
 the cyclic-invariant part of the chiral Hochschild cochain complex
 (the complex controlling $S^1$-equivariant deformations).
\item \emph{Stable-graph operations.}
 For $\cA$ modular Koszul
 (Definition~\ref{def:modular-koszul-chiral}), the genus-$g$
 deformations are generated by stable-graph operations
 $\gamma \in \Gamma_{g,n}$, and the Maurer--Cartan equation
 splits into genus-graded components whose leading term is the
 curvature $\mcurv{g} = \kappa(\cA) \cdot \lambda_g$
 \textup{(UNIFORM-WEIGHT)}.
\end{enumerate}

The Heisenberg specialization (Chapter~\ref{ch:heisenberg-frame})
is computed via model~(a): the cyclic coderivation complex is
one-dimensional at genus~$1$, generated by the Eisenstein element
$\kappa \cdot E_2(\tau)$.
\end{definition}

\begin{remark}[Quantum \texorpdfstring{$L_\infty$}{L-infinity} versus cyclic \texorpdfstring{$L_\infty$}{L-infinity}]
\label{rem:quantum-vs-cyclic-linf}
\index{quantum L-infinity algebra@quantum $L_\infty$ algebra|textbf}
\index{loop homotopy algebra|see{quantum $L_\infty$ algebra}}
Three levels arise as the genus filtration opens:
(i)~\emph{cyclic $L_\infty$} ($g = 0$, tree-level operations);
(ii)~\emph{quantum $L_\infty$} (all $g$; Braun--Lazarev~\cite{BL15},
Markl~\cite{Markl06}: operations $m^{(g)}_n$ on stable graphs
satisfying the QME $\sum_g \hbar^g m^{(g)} = 0$);
(iii)~\emph{modular homotopy type} (coupled to factorization on
$\operatorname{Ran}(X)$,
Definition~\ref{def:modular-homotopy-theory-intro}).
The lift (i)$\to$(ii) is obstructed by $\Delta(m_0)$ (unimodularity);
Koszul chiral algebras carry full modular homotopy types via
$\Theta_\cA$ (Theorem~\ref{thm:universal-MC}).
\end{remark}

\begin{proposition}[Genus-\texorpdfstring{$0$}{0} cyclic coderivation complex;
\ClaimStatusProvedHere]\label{prop:genus0-cyclic-coderivation}
\index{cyclic deformation complex!genus-$0$ construction}
Let $\cA$ be a Koszul chiral algebra with non-degenerate
invariant bilinear form
$\langle{-},{-}\rangle_{\cA}\colon \cA\otimes\cA\to\omega_X$.
Let $\bar{B}_X(\cA) = T^c(s^{-1}\bar{\cA})$ be the reduced bar
coalgebra
\textup{(}Convention~\textup{\ref{conv:bar-coalgebra-identity}}\textup{)}.
Then:
\begin{enumerate}[label=\textup{(\roman*)}]
\item \emph{Cyclic coalgebra structure.}
 The form $\langle{-},{-}\rangle_{\cA}$ induces a non-degenerate
 cyclic pairing $\eta$ on $T^c(s^{-1}\bar{\cA})$ via
 \begin{equation}\label{eq:cyclic-coalgebra-pairing}
 \eta\bigl(s^{-1}a_1\,|\cdots|\,s^{-1}a_n,\;
 s^{-1}b_1\,|\cdots|\,s^{-1}b_n\bigr)
 \;=\;
 \sum_{j=0}^{n-1}\varepsilon_j\,
 \prod_{i=1}^{n}
 \langle a_{j+i},\,b_i\rangle_{\cA},
 \end{equation}
 where indices are cyclic modulo~$n$ and $\varepsilon_j$ is
 the Koszul sign of the cyclic permutation by~$j$ steps.

\item \emph{Cyclic coderivation sub-Lie-algebra.}
 The space
 $\operatorname{CoDer}^{\mathrm{cyc}}(\bar{B}_X(\cA))
 := \{f\in\operatorname{CoDer}(\bar{B}_X(\cA))\;:\;
 \eta(f(x),y)+(-1)^{|f||x|}\eta(x,f(y))=0\;\forall\,x,y\}$
 is a sub-Lie-algebra of
 $\operatorname{CoDer}(\bar{B}_X(\cA))$.

\item \emph{Bar differential is cyclic.}
 The bar differential $d_{\barB}\in
 \operatorname{CoDer}(\bar{B}_X(\cA))$ lies in
 $\operatorname{CoDer}^{\mathrm{cyc}}(\bar{B}_X(\cA))$.

\item \emph{Cyclic dg Lie algebra.}
 The degree-shifted complex
 $\operatorname{CoDer}^{\mathrm{cyc}}(\bar{B}_X(\cA))[1]$
 is a cyclic dg Lie algebra: the bracket $l_2 = [-,-]$
 is $\eta$-invariant and $l_1 = [d_{\barB},-]$ is
 $\eta$-skew. The Maurer--Cartan locus inherits a
 $(-1)$-shifted symplectic structure.
\end{enumerate}
\end{proposition}

\begin{proof}
(i)\enspace
The cyclic pairing~\eqref{eq:cyclic-coalgebra-pairing} is the
standard construction for cyclic coalgebras:
each bar-length-$n$ component of~$\eta$ is obtained by cyclically
summing the tensor product of~$\langle{-},{-}\rangle_{\cA}$ over
the $n$~factors. Non-degeneracy of~$\eta$ on each bar-length
component follows from non-degeneracy
of~$\langle{-},{-}\rangle_{\cA}$.

(ii)\enspace
For $f,g\in\operatorname{CoDer}^{\mathrm{cyc}}$, the
commutator $[f,g] = f\circ g - (-1)^{|f||g|}g\circ f$ satisfies
the cyclic identity
\begin{align*}
\eta([f,g](x),y) &= \eta(f(g(x)),y) -
(-1)^{|f||g|}\eta(g(f(x)),y)\\
&= -(-1)^{|f|(|g|+|x|)}\eta(g(x),f(y))
 + (-1)^{|f||g|+|g||x|}\eta(f(x),g(y))\\
&= -(-1)^{(|f|+|g|)|x|}\eta(x,[f,g](y)),
\end{align*}
using the cyclic property of~$f$ and~$g$ in the second step and
regrouping signs. Hence $[f,g]$ is cyclic.

(iii)\enspace
The bar differential $d_{\barB}$ is a coderivation determined by
the chiral multiplication and higher operations of~$\cA$.
The cyclicity of $d_{\barB}$ is equivalent to the invariance of
$\langle{-},{-}\rangle_{\cA}$ under the chiral operations:
for the leading term,
$\langle m_2(a,b),c\rangle_{\cA}
= -(-1)^{|a||b|}\langle m_2(b,a),c\rangle_{\cA}
= \langle a, m_2(b,c)\rangle_{\cA}$,
and the higher-order terms satisfy the same identity by the
$A_\infty$ compatibility of the invariant form.

(iv)\enspace Parts~(ii) and~(iii) together make
$\operatorname{CoDer}^{\mathrm{cyc}}[1]$ a cyclic dg Lie algebra
in the sense of Kontsevich--Soibelman~\cite{KontsevichSoibelman}:
$l_1 = [d_{\barB},-]$ is skew because $d_{\barB}$ is cyclic and
the bracket is $\eta$-invariant, while $l_2 = [-,-]$ is invariant
by~(ii). The $(-1)$-shifted symplectic structure on the
Maurer--Cartan formal moduli problem is then the standard
consequence (cf.\ Loday--Vallette~\cite{LV12}, \S10.4).
\end{proof}

\begin{proposition}[Killing cocycle \texorpdfstring{$L_\infty$}{L-infinity} extension;
\ClaimStatusProvedHere]\label{prop:killing-linf-extension}
\index{Killing cocycle!$L_\infty$ extension}
\index{$L_\infty$-algebra!Killing cocycle extension}
Let $\mathfrak{g}$ be a finite-dimensional Lie algebra with
non-degenerate invariant symmetric bilinear form~$\kappa_{\mathfrak{g}}$.
Define the Killing $3$-cocycle
$\phi(a,b,c) := \kappa_{\mathfrak{g}}([a,b],c) \in \Lambda^3(\mathfrak{g}^*)$.
Then:
\begin{enumerate}[label=\textup{(\roman*)}]
\item $\phi$ is a Chevalley--Eilenberg $3$-cocycle:
 $\delta\phi = 0$.
\item The data $l_1 = 0$, $l_2 = [-,-]$,
 $l_3 = \phi$ \textup{(}valued in a one-dimensional marker
 direction~$\eta$\textup{)},
 and $l_n = 0$ for $n \geq 4$ defines an $L_\infty$-algebra
 structure on $\mathfrak{g} \oplus \mathbb{C}\cdot\eta$,
 where $\deg(\mathfrak{g}) = 0$ and $\deg(\eta) = 2$.
\item The pairing $\kappa_{\mathfrak{g}}$ on~$\mathfrak{g}$ extended by
 $\langle\eta,\eta\rangle = 1$ and
 $\langle\mathfrak{g},\eta\rangle = 0$ makes this a
 \emph{cyclic} $L_\infty$-algebra: $l_2$ is
 $\kappa_{\mathfrak{g}}$-invariant, the $4$-form
 $\langle l_3(a,b,c), d\rangle$ vanishes on
 $\mathfrak{g}^{\otimes 4}$, and the restriction
 $\langle l_3(a,b,c),\eta\rangle = \phi(a,b,c)$ is
 totally antisymmetric.
\end{enumerate}
\end{proposition}

\begin{proof}
(i)\enspace
The CE~differential of~$\phi$ evaluates as
\begin{align*}
(\delta\phi)(a,b,c,d)
&= \textstyle\sum_{i<j}(-1)^{i+j}\,
 \phi([x_i,x_j],x_k,x_l) \\
&= \textstyle\sum_{i<j}(-1)^{i+j}\,
 \kappa_{\mathfrak{g}}([[x_i,x_j],x_k],x_l).
\end{align*}
Using ad-invariance $\kappa_{\mathfrak{g}}([u,v],w) = \kappa_{\mathfrak{g}}(u,[v,w])$
repeatedly, this equals
$2\bigl(\kappa_{\mathfrak{g}}([a,c],[b,d]) - \kappa_{\mathfrak{g}}([a,b],[c,d])
 - \kappa_{\mathfrak{g}}([a,d],[b,c])\bigr)$.
Applying ad-invariance once more:
$\kappa_{\mathfrak{g}}([a,c],[b,d]) - \kappa_{\mathfrak{g}}([a,b],[c,d])
 - \kappa_{\mathfrak{g}}([a,d],[b,c])
= \kappa_{\mathfrak{g}}\bigl(a,\,
 [c,[b,d]] - [b,[c,d]] - [d,[b,c]]\bigr) = 0$
by the Jacobi identity $[b,[c,d]] + [c,[d,b]] + [d,[b,c]] = 0$.

(ii)\enspace
The $L_\infty$ homotopy Jacobi identities hold at
each degree~$n$.

\emph{Degrees $1$--$3$.}
With $l_1 = 0$, degrees~$1$ and~$2$ hold trivially.
At degree~$3$, the identity reduces to the Jacobi identity
for~$l_2$, since the only $l_3$-involving term is
$l_1\circ l_3 = 0$.

\emph{Degree~$4$.}
The identity (with $l_1 = l_4 = 0$) reads
\begin{equation}\label{eq:degree4-linf}
\begin{split}
\textstyle\sum_{\sigma \in \mathrm{Sh}(1,3)}
 &\,\varepsilon(\sigma)\,l_2\bigl(x_{\sigma(1)},\,
 l_3(x_{\sigma(2)},x_{\sigma(3)},x_{\sigma(4)})\bigr)\\
&{}+ \textstyle\sum_{\sigma \in \mathrm{Sh}(2,2)}
 \varepsilon(\sigma)\,l_3\bigl(l_2(x_{\sigma(1)},
 x_{\sigma(2)}),\,x_{\sigma(3)},x_{\sigma(4)}\bigr) = 0.
\end{split}
\end{equation}
The first sum vanishes: $l_3$ outputs in the
$\eta$-direction and $l_2(x,\eta) = 0$ for all
$x \in \mathfrak{g}$. The second sum equals
$(\delta\phi)(a,b,c,d)\cdot\eta = 0$ by~(i).

\emph{Degrees $\geq 5$.}
With $l_n = 0$ for $n \geq 4$, every surviving term in the
degree-$n$ identity involves a composition
$l_3(\dots, l_3(\dots), \dots)$, which requires $l_3$ to
accept $\eta$ as input. Since $l_3$ is defined only on
$\mathfrak{g}^{\otimes 3}$, all such terms vanish.

(iii)\enspace
The bracket $l_2 = [-,-]$ is $\kappa_{\mathfrak{g}}$-invariant by the
ad-invariance of~$\kappa_{\mathfrak{g}}$
(Proposition~\ref{prop:genus0-cyclic-coderivation}(ii)).
For~$l_3$: on $\mathfrak{g}^{\otimes 4}$, the $4$-form
$\langle l_3(a,b,c),d\rangle = \phi(a,b,c)\langle\eta,d\rangle
= 0$ since $\langle\eta,\mathfrak{g}\rangle = 0$.
The restriction
$\langle l_3(a,b,c),\eta\rangle = \phi(a,b,c)$ inherits total
antisymmetry from the ad-invariance and antisymmetry
of~$\kappa_{\mathfrak{g}}\circ\mathrm{ad}$.
\end{proof}

\begin{corollary}[Kac--Moody cyclic deformation complex;
\ClaimStatusConditional]\label{cor:km-cyclic-deformation}
\index{Kac--Moody algebra!cyclic deformation complex}
For $\widehat{\mathfrak{g}}_k$ with $k\neq -h^\vee$ and
$\mathfrak{g}$ simple, the genus-$0$ cyclic deformation complex
$L_0 := \operatorname{CoDer}^{\mathrm{cyc}}(\bar{B}
(\widehat{\mathfrak{g}}_k))[1]$
of Proposition~\textup{\ref{prop:genus0-cyclic-coderivation}}
has the following explicit data:
\begin{enumerate}[label=\textup{(\roman*)}]
\item On generators, $l_2 = \mathrm{ad}$ \textup{(}adjoint
 bracket of~$\mathfrak{g}$\textup{)} and the cyclic pairing
 restricts to the level-$k$ Killing form
 $\kappa_k = k\,\kappa_{\mathfrak{g}} + \kappa_{\mathrm{Killing}}$.
\item The adjoint Casimir
 $C_2 = \sum_i \mathrm{ad}(e_i)\circ\mathrm{ad}(e^i)$
 acts as $2h^\vee\cdot\mathrm{id}$ on
 $s^{-1}\bar{\mathfrak{g}} \subset L_0$.
\item The unique cyclic infinitesimal deformation is
 the Killing $3$-cocycle
 $l_3(a,b,c) = \kappa_{\mathfrak{g}}([a,b],c)$, corresponding to the
 generator of $H^2_{\mathrm{cyc}}(\mathfrak{g},\mathfrak{g})
 \cong\mathbb{C}$
 \textup{(}Proposition~\textup{\ref{prop:cyclic-ce-identification}}\textup{)}.
\item The scalar shadow
 $\kappa(\widehat{\mathfrak{g}}_k) =
 \dim(\mathfrak{g})\,(k+h^\vee)/(2h^\vee)$
 is recovered from the $l_2$ bracket and the cyclic pairing via
 the two-channel decomposition
 \textup{(}Remark~\textup{\ref{rem:mc2-status}},
 items~\textup{(vi)--(ix)}\textup{)}.
\end{enumerate}
\end{corollary}

\begin{proof}
(i)\enspace
At genus~$0$, the PBW filtration identifies the generator-level
structure of~$L_0$ with the adjoint representation
of~$\mathfrak{g}$. The bar differential has leading term
$d_2(sa\,|\,sb) = s[a,b]$, so $l_2$ on generators is the adjoint
bracket. The invariant form on $\widehat{\mathfrak{g}}_k$ is
$\kappa_k$, which restricts to the level-$k$ normalized Killing
form on~$\mathfrak{g}$.

(ii)\enspace
On the adjoint representation, $C_2 = \sum_i \mathrm{ad}(e_i)
\circ\mathrm{ad}(e^i)$ has eigenvalue~$2h^\vee$ (standard: the
Casimir eigenvalue on the adjoint representation equals twice the
dual Coxeter number for simple~$\mathfrak{g}$).

(iii)\enspace
By Proposition~\ref{prop:cyclic-ce-identification},
$H^2_{\mathrm{cyc}}(\mathfrak{g},\mathfrak{g})\cong
H^3(\mathfrak{g})\cong\mathbb{C}$, generated by the Killing
$3$-cocycle $\phi(a,b,c)=\kappa_{\mathfrak{g}}([a,b],c)$. This is the unique
infinitesimal cyclic deformation direction: it defines
an~$l_3$ bracket on~$L_0$ satisfying all homotopy Jacobi
identities and cyclic symmetry by
Proposition~\ref{prop:killing-linf-extension}.
The computations in
Remark~\ref{rem:mc2-status}, items~(v)--(vi),
(viii)--(ix) confirm this independently for
$\mathfrak{sl}_2$, $\mathfrak{sl}_3$,
$\mathfrak{sp}_4$, and~$\mathfrak{g}_2$.

(iv)\enspace
The two-channel decomposition extracts
$\kappa_{\mathrm{dp}}(\widehat{\mathfrak{g}}_k) = \dim(\mathfrak{g})\cdot k/(2h^\vee)$ from the
double-pole channel and $\dim(\mathfrak{g})/2$ from the
simple-pole channel. Their sum gives
$\kappa(\widehat{\mathfrak{g}}_k) = \kappa_{\mathrm{dp}}(\widehat{\mathfrak{g}}_k) + \dim(\mathfrak{g})/2
= \dim(\mathfrak{g})(k+h^\vee)/(2h^\vee)$, matching the
formula of Theorem~\ref{thm:modular-characteristic}.
\end{proof}

\begin{definition}[Cyclic deformation complex: bar-intrinsic
 form; \ClaimStatusDefinitional]
\label{def:cyclic-deformation-bar}
\index{cyclic deformation complex!bar-intrinsic form}
Let $\cA$ be a modular Koszul chiral algebra
(Definition~\ref{def:modular-koszul-chiral}) carrying the trace/duality
data needed to define the chain-level Verdier/BV pairing. The
\emph{bar-intrinsic cyclic deformation complex} is
\begin{equation}\label{eq:def-cyc-bar}
\Defcyc(\cA)
\;:=\;
\operatorname{CoDer}^{\mathrm{cyc}}\!\bigl(
 \widehat{\barB}_X(\cA)
\bigr)[1],
\end{equation}
the degree-shifted \emph{cyclic $L_\infty$-algebra} of continuous
cyclic coderivations of the completed reduced bar coalgebra
$\widehat{\barB}_X(\cA)$; ``cyclic'' means skew-adjoint for the
canonical pairing induced by Verdier duality/BV trace.
The $L_\infty$ structure has brackets $\ell_1 = [d_{\barB},-]$,
$\ell_2 = [-,-]$ (the coderivation bracket), and higher brackets
$\ell_n$ for $n \geq 3$ induced by the transferred $A_\infty$
structure on $\widehat{\barB}_X(\cA)$
\cite[Theorem~10.3.8]{LV12}.
For Koszul~$\cA$, the convolution $L_\infty$ is \emph{formal}
(the higher brackets $\ell_n = 0$ for $n \geq 3$ on the
cohomological model), so the strict dg~Lie structure suffices for
computations
(Theorem~\ref{thm:cyclic-linf-graph}).

Four properties:
\begin{enumerate}[label=\textup{(\roman*)}]
\item \emph{Filtered-completeness:} the bar-length filtration makes
 $\Defcyc(\cA)$ complete and pronilpotent.
\item \emph{Functoriality:} the assignment $\cA \mapsto \Defcyc(\cA)$
 is functorial in morphisms of augmented chiral algebras.
\item \emph{Cyclic pairing:} the Verdier/BV trace induces a
 $(-1)$-shifted symplectic structure on the formal moduli problem
 $\operatorname{MC}(\Defcyc(\cA))$.
\item \emph{Curved MC control:} Maurer--Cartan elements of
 $\Defcyc(\cA)$ parametrize curved deformations of the bar
 differential, so the genus tower is controlled by a single MC
 deformation.
\end{enumerate}
The degree shift~$[1]$ makes the Lie bracket of degree~$0$, matching
the convention for tangent Lie algebras of formal moduli problems
(Lurie~\cite{HA}). Equivalently,
$\Defcyc(\cA)$ is the tangent dg Lie algebra of cyclic deformations
of~$\cA$ as a modular factorization algebra; it presents the abstract
H-level object of Definition~\ref{def:cyclic-deformation-elementary}
via bar-coalgebra coderivations.

At genus~$0$, the cyclic coderivation complex is fully constructed
(Proposition~\ref{prop:genus0-cyclic-coderivation}), and its
Kac--Moody specialization is computed explicitly
(Corollary~\ref{cor:km-cyclic-deformation}).
\textbf{For affine Kac--Moody algebras}, $\Defcyc(\widehat{\mathfrak{g}}_k)$
is a well-defined cyclic $L_\infty$-algebra with
$H^2 \cong \mathbb{C}$: Theorem~\ref{thm:mc2-1-km} identifies the
coderivation differential with the cyclic CE~differential and proves
Hypothesis~\ref{mc2-hyp:cyclic} for this family.
The full modular extension for general modular Koszul chiral
algebras is now established by
Theorem~\ref{thm:cyclic-linf-graph}: the chiral graph complex
construction equips $\Defcyc(\cA)$ with a cyclic $L_\infty$
structure for all Koszul chiral algebras with non-degenerate
invariant form.
\end{definition}

\begin{definition}[Modular cyclic deformation complex;
\ClaimStatusProvedHere]
\label{def:modular-cyclic-deformation-complex}
\index{modular cyclic deformation complex|textbf}
\index{cyclic deformation complex!modular extension}
For a chiral algebra~$\cA$ with bar
complex~$\barB(\cA)$, the \emph{modular cyclic deformation complex} is
\begin{equation}\label{eq:modular-cyc-def-complex}
\Defcyc^{\mathrm{mod}}(\cA)
\;:=\;
\prod_{\substack{g,n \\ 2g-2+n>0}}
\operatorname{CoDer}^{\mathrm{cyc}}\!\bigl(
 \barB^{(g,n)}(\cA)
\bigr)[1],
\end{equation}
where $\barB^{(g,n)}(\cA)$ is the genus-$g$, $n$-marked bar complex
and the product ranges over stable pairs~$(g,n)$.
The differential is the sum $d_{\mathrm{int}} + d_{\mathrm{sew}}$
where:
\begin{itemize}
\item $d_{\mathrm{int}}$ is the internal cyclic coderivation
 differential on each component;
\item $d_{\mathrm{sew}}$ is the clutching/sewing operation induced by
 the boundary maps
 $\partial\colon \overline{\mathcal{M}}_{g,n} \supset
 \Delta_{g_1,S_1;\,g_2,S_2} \to
 \overline{\mathcal{M}}_{g_1,|S_1|+1} \times
 \overline{\mathcal{M}}_{g_2,|S_2|+1}$.
\end{itemize}

The Lie bracket is \emph{graph composition}: for
$\xi \in \operatorname{CoDer}^{\mathrm{cyc}}(\barB^{(g_1,n_1)}(\cA))$
and
$\eta \in \operatorname{CoDer}^{\mathrm{cyc}}(\barB^{(g_2,n_2)}(\cA))$,
\begin{equation}\label{eq:modular-cyc-graph-bracket}
[\xi,\eta]
\;:=\;
\sum_{\Gamma}\xi \circ_\Gamma \eta,
\end{equation}
summed over stable graphs~$\Gamma$ connecting an output of~$\xi$ to an
input of~$\eta$ with cross-polarization on internal edges.
\end{definition}

\begin{remark}[Strictification of the modular deformation object]
\label{rem:modular-cyc-strictification}
The modular cyclic deformation complex $\Defcyc^{\mathrm{mod}}(\cA)$
is the strict model of the homotopy-invariant modular deformation
object $\Definfmod(\cA)$
\textup{(}Theorem~\textup{\ref{thm:modular-homotopy-convolution})}.
The cyclic coderivation description is one model, obtained from the
cofree resolution of~$\barB(\cA)$; a different choice of
contracting homotopy produces an $L_\infty$-quasi-isomorphic
deformation complex. For affine Kac--Moody algebras, the
strict model is already formal: the universal class
$\Theta^{\mathrm{str}}_{\widehat{\mathfrak{g}}_k}$ satisfies the strict
MC equation with all higher $L_\infty$-brackets vanishing
\textup{(}Theorem~\textup{\ref{thm:km-strictification})}.
\end{remark}

\begin{proposition}[Genus truncation;
\ClaimStatusProvedHere]
\label{prop:modular-deformation-truncation}
\index{modular cyclic deformation complex!genus truncation}
The genus-$0$ truncation
$\Defcyc^{\mathrm{mod}}(\cA)\big|_{g=0} = \Defcyc(\cA)$
recovers the elementary cyclic deformation complex of
Definition~\textup{\ref{def:cyclic-deformation-bar}}.
The genus grading induces a complete descending filtration
$F^{\bullet}\Defcyc^{\mathrm{mod}}(\cA)$ with
$\gr_g$ computing deformations at fixed genus.
\end{proposition}

\begin{proof}
At $g = 0$ the only stable pairs are $(0,n)$ with $n \geq 3$.
The internal differential~$d_{\mathrm{int}}$ on each
$\operatorname{CoDer}^{\mathrm{cyc}}(\barB^{(0,n)}(\cA))[1]$
coincides with the coderivation differential of
Definition~\ref{def:cyclic-deformation-bar}, and the sewing
differential~$d_{\mathrm{sew}}$ reduces to the graph composition
already present in the genus-$0$ bracket.
The genus filtration $F^g := \prod_{h \geq g}$ is complete because
the product topology is, and the associated graded
$\gr_g = \operatorname{CoDer}^{\mathrm{cyc}}(\barB^{(g,\bullet)}(\cA))[1]$
computes deformations at fixed genus by construction.
\end{proof}

\begin{remark}[Modular cyclic deformation complex as MC ambient]
\label{rem:chriss-ginzburg-modular-cyc}
\index{Chriss--Ginzburg principle}
\index{Maurer--Cartan!modular cyclic}
The modular cyclic deformation complex is the ambient home for
$\Theta_\cA$. At the dg~level,
the Lie bracket and differential give a cyclic dg~Lie algebra;
Theorem~\ref{thm:cyclic-linf-graph} extends this to a full
\emph{cyclic $L_\infty$-algebra}, and
Theorem~\ref{thm:modular-quantum-linfty} promotes it to a
\emph{quantum $L_\infty$-algebra} with operations
$\{\ell_n^{(g)}\}_{n \geq 1,\, g \geq 0}$: here
$\ell_1^{(0)} = D$, $\ell_2^{(0)} = [-,-]$,
$\ell_n^{(0)}$ for $n \geq 3$ are the higher operadic brackets from
homotopy transfer through $C_*(\overline{\mathcal{M}}_{0,n+1})$
(Loday--Vallette~\cite[Thm~10.3.8]{LV12}), and
$\ell_n^{(g)}$ for $g \geq 1$ are the higher-genus Feynman
amplitudes from stable-graph sums of type~$(g,n)$
(Getzler--Kapranov~\cite[\S5]{GeK98},
Chuang--Lazarev~\cite{CL10}).

The modular MC problem: find
$\Theta_\cA \in \MC(\Defcyc^{\mathrm{mod}}(\cA))$ satisfying the
quantum $L_\infty$ Maurer--Cartan equation
\[
\sum_{n \geq 1} \frac{1}{n!}\,\ell_n^{(0)}
 (\Theta_\cA,\dots,\Theta_\cA)
\;+\;
\sum_{g \geq 1} \hbar^{g}\,
 \ell_1^{(g)}(\Theta_\cA)
\;=\; 0,
\]
whose leading terms $\ell_1^{(0)}(\Theta) + \tfrac{1}{2}\ell_2^{(0)}(\Theta,\Theta) = 0$ give
the classical dg~MC equation
$d\Theta_\cA + \tfrac{1}{2}[\Theta_\cA, \Theta_\cA] = 0$,
the strict shadow of the completed $L_\infty$-MC equation
(Theorem~\ref{thm:modular-homotopy-convolution}).
Theorem~\ref{thm:universal-MC} establishes existence of~$\Theta_\cA$
for Kac--Moody at non-critical level,
Virasoro, and principal~$\mathcal{W}_N$ at generic level.
All modular invariants ($\kappa$, $\Delta$, shadows
$\operatorname{Sh}_r$) are projections of~$\Theta_\cA$ onto graded
components; the higher brackets $\ell_n^{(g)}$ encode
the Feynman-amplitude content beyond the leading dg~terms.
\end{remark}

\begin{proposition}[Chart-to-shadow comparison;
\ClaimStatusConditional]
\label{prop:hochschild-shadow-projection}
\index{chiral Hochschild cohomology!shadow projection}
\index{shadow algebra!degree-2 projection}

Let \(\cA\) carry \(H_H(\cA;S)\), and suppose the cyclic modular
deformation complex is equipped with a quasi-isomorphism
\[
 \rho_{\cA}\colon
 C^\bullet_{\mathrm{ch}}(\cA,\cA)
 \xrightarrow{\ \sim\ }
 \pi_{2,0}\bigl(\Defcyc^{\mathrm{mod}}(\cA)\bigr).
\]
Then
\[
 \ChirHoch^n(\cA)
 \cong
 H^n\!\left(\pi_{2,0}
 \bigl(\Defcyc^{\mathrm{mod}}(\cA)\bigr)\right),
 \qquad
 \operatorname{Supp}\ChirHoch^\bullet(\cA)\subseteq S.
\]
The higher genus-zero projections
\(\pi_{r,0}(\Defcyc^{\mathrm{mod}}(\cA))\), \(r\geq3\), govern the
transferred higher operations.  A brace-compatible contraction of
these projections gives the \(E_2\)-formality statement of
Proposition~\ref{prop:e2-formality-hochschild}.
\end{proposition}

\begin{proof}
Theorem~\ref{thm:main-koszul-hoch} identifies the chart complex with
\(K_{\cA,S}\).  Composition with \(\rho_\cA\) gives the displayed
cohomological identification and its support.  Homological transfer
along a brace-compatible contraction gives the final assertion.
\end{proof}

\begin{theorem}[MC2-1 for Kac--Moody algebras;
\ClaimStatusConditional]\label{thm:mc2-1-km}
\index{MC2!Kac--Moody resolution}
For $\widehat{\mathfrak{g}}_k$ at non-critical level ($k \neq -h^\vee$)
with $\mathfrak{g}$ simple, Hypothesis~\textup{\ref{mc2-hyp:cyclic}}
is satisfied:
\begin{enumerate}[label=\textup{(\roman*)}]
\item \emph{Cyclic $L_\infty$ structure.}
 The bar-intrinsic cyclic deformation complex
 \[
 \Defcyc(\widehat{\mathfrak{g}}_k)
 := \operatorname{CoDer}^{\mathrm{cyc}}(\widehat{\barB}_X
 (\widehat{\mathfrak{g}}_k))[1]
 \]
 is a well-defined complete cyclic $L_\infty$-algebra
 \textup{(}a fortiori: a cyclic dg~Lie algebra\textup{)}.

\item \emph{CE identification.}
 The coderivation differential $l_1 = [d_{\barB},-]$ is canonically
 identified with the cyclic Chevalley--Eilenberg differential:
 \begin{equation}\label{eq:coder-ce-ident}
 \bigl(\Defcyc(\widehat{\mathfrak{g}}_k),\, l_1\bigr)
 \;\cong\;
 \bigl(C^{*-1}_{\mathrm{cyc}}(\mathfrak{g},\mathfrak{g}),\,
 \delta_{\mathrm{CE}}\bigr).
 \end{equation}
 In particular,
 $H^n(\Defcyc(\widehat{\mathfrak{g}}_k), l_1) \cong
 H^n_{\mathrm{cyc}}(\mathfrak{g},\mathfrak{g})$
 for all $n \geq 0$.

\item \emph{One-dimensional obstruction.}
 $H^2(\Defcyc(\widehat{\mathfrak{g}}_k), l_1) \cong \mathbb{C}$,
 generated by the Killing $3$-cocycle
 $\phi(a,b,c) = \kappa([a,b],c)$.
\end{enumerate}
\end{theorem}

\begin{proof}
(i)\enspace
By Proposition~\ref{prop:genus0-cyclic-coderivation},
$\operatorname{CoDer}^{\mathrm{cyc}}(\barB(\widehat{\mathfrak{g}}_k))[1]$
is a cyclic dg Lie algebra.
The bar coalgebra $\barB = T^c(s^{-1}\bar{\mathfrak{g}})$ is
conilpotent, and its bar-length filtration is complete and separated.
Continuous coderivations of the completed bar coalgebra
$\widehat{\barB}$ are determined by their bar-length components
$f_p\colon (s^{-1}\bar{\mathfrak{g}})^{\otimes p}
\to s^{-1}\bar{\mathfrak{g}}$, forming the product
$\prod_{p \geq 1}
\operatorname{Hom}^{\mathrm{cyc}}
((s^{-1}\bar{\mathfrak{g}})^{\otimes p},\,
s^{-1}\bar{\mathfrak{g}})$.
The bracket of a bar-length-$p$ and a bar-length-$q$ coderivation
has bar-length $p + q - 1$ (by the pre-Lie composition formula),
and the cyclic pairing~$\eta$ pairs bar-length~$p$ with
bar-length~$p$; both converge in the completed product.
A cyclic dg~Lie algebra is in particular a cyclic
$L_\infty$-algebra (with $l_n = 0$ for $n \geq 3$).

(ii)\enspace
The reduced bar differential of $\widehat{\mathfrak{g}}_k$ on
$\barB = T^c(s^{-1}\bar{\mathfrak{g}})$ has cogeneration components
$d_n\colon (s^{-1}\bar{\mathfrak{g}})^{\otimes n}
\to s^{-1}\bar{\mathfrak{g}}$.
Only $d_2$ is nonzero: $d_2(a,b) = [a,b]$ encodes the simple-pole
OPE $J^a(z)J^b(w) \sim f^{ab}_c J^c/(z{-}w)$, while the double-pole
data $k\kappa^{ab}/(z{-}w)^2$ maps to the vacuum (augmentation
sector) and does not contribute to the \emph{reduced} bar
differential; higher poles are absent for affine Kac--Moody.

For a coderivation $f = (f_p)_{p \geq 1}$ and the bar differential
$d_{\barB}$ with sole component~$d_2$, the pre-Lie composition
formula gives:
\[
[d_2, f_p]
\;=\; d_2 \circ f_p \;-\; (-1)^{|d_2||f_p|}\, f_p \circ d_2,
\]

which has bar-length $p + 2 - 1 = p + 1$. Hence $l_1 = [d_{\barB},-]$
maps bar-length-$p$ to bar-length-$(p{+}1)$, with no bar-length-lowering
terms. Explicitly, on a bar-length-$p$ cyclic map
$f_p\colon (s^{-1}\bar{\mathfrak{g}})^{\otimes p}
\to s^{-1}\bar{\mathfrak{g}}$:
\begin{align*}
(d_2 \circ f_p)(a_1, \ldots, a_{p+1})
&= \sum_{i=1}^{p} \varepsilon_i\,
 [f_p(a_1, \ldots, \widehat{a}_i, \ldots, a_{p+1}),\, a_i], \\
(f_p \circ d_2)(a_1, \ldots, a_{p+1})
&= \sum_{i < j} \varepsilon_{ij}\,
 f_p(a_1, \ldots, [a_i, a_j], \ldots, \widehat{a}_j,
 \ldots, a_{p+1}),
\end{align*}

where $\varepsilon_i, \varepsilon_{ij}$ are the standard Koszul signs.
These are exactly the two terms of the Chevalley--Eilenberg differential
$\delta_{\mathrm{CE}}$ on
$C^p_{\mathrm{cyc}}(\mathfrak{g},\mathfrak{g})
= \operatorname{Hom}^{\mathrm{cyc}}
((s^{-1}\bar{\mathfrak{g}})^{\otimes p},\,
s^{-1}\bar{\mathfrak{g}})$,
with the cyclic-invariance property preserved because $d_2 = [-,-]$
is $\kappa$-invariant (Proposition~\ref{prop:genus0-cyclic-coderivation}(iii)).
The identification~\eqref{eq:coder-ce-ident} follows.

(iii)\enspace
By Proposition~\ref{prop:cyclic-ce-identification},
$H^n_{\mathrm{cyc}}(\mathfrak{g},\mathfrak{g})
\cong H^{n+1}(\mathfrak{g})$.
For simple~$\mathfrak{g}$, $H^3(\mathfrak{g}) = \mathbb{C}$,
generated by the Killing $3$-form
$\phi(a,b,c) = \kappa([a,b],c) \in \Lambda^3(\mathfrak{g}^*)$.
Under the isomorphism of~(ii), this corresponds to the
unique cyclic infinitesimal deformation direction
(Corollary~\ref{cor:km-cyclic-deformation}(iii)).
\end{proof}

\begin{corollary}[Minimal cyclic \texorpdfstring{$L_\infty$}{L-infinity} model for Kac--Moody;
\ClaimStatusConditional]\label{cor:km-minimal-linf}
\index{minimal model!Kac--Moody cyclic $L_\infty$}
Theorem~\textup{\ref{thm:mc2-1-km}} and the homotopy transfer
theorem for cyclic $L_\infty$-algebras
\textup{(}Kontsevich--Soibelman~\textup{\cite{KontsevichSoibelman}},
Loday--Vallette~\textup{\cite{LV12}} \S10.3\textup{)}
give a minimal cyclic $L_\infty$ model for
$\Defcyc(\widehat{\mathfrak{g}}_k)$ on the cohomology
$H^*(\Defcyc, l_1) \cong \mathfrak{g} \oplus \mathbb{C}\cdot\eta$
\textup{(}degrees $0$ and~$2$\textup{)}.
The transferred structure is exactly
the Killing cocycle extension of
Proposition~\textup{\ref{prop:killing-linf-extension}}:
\begin{equation}\label{eq:km-minimal-linf-structure}
l_2^{\mathrm{tr}} = [-,-]_{\mathfrak{g}},\quad
l_3^{\mathrm{tr}} = \phi,\quad
l_n^{\mathrm{tr}} = 0 \;\text{for}\; n \geq 4.
\end{equation}
\end{corollary}

\begin{proof}
The homotopy transfer theorem produces a minimal $L_\infty$-algebra
$(H^*(\Defcyc, l_1),\, l_n^{\mathrm{tr}})$ quasi-isomorphic to the
original cyclic dg~Lie algebra. By
Theorem~\ref{thm:mc2-1-km}(ii), $H^0 \cong \mathfrak{g}$ (the
adjoint representation) and $H^2 \cong \mathbb{C}$.

\emph{$l_2^{\mathrm{tr}} = \mathrm{ad}$.}
The transferred binary bracket $l_2^{\mathrm{tr}}(a,b) =
p\,l_2(i(a), i(b))$, where $i\colon \mathfrak{g} \hookrightarrow
\Defcyc$ is the inclusion into bar-length~$1$ and $p$ is the
projection. Since bar-length-$1$ coderivations are linear
endomorphisms and $l_2$ is the commutator, $l_2^{\mathrm{tr}}(a,b)
= [a,b]$.

\emph{$l_3^{\mathrm{tr}} = \phi$.}
The transferred ternary bracket on $\mathfrak{g}^{\otimes 3}$
lands in $H^2 \cong \mathbb{C}\cdot\eta$, so it defines an element of
$\operatorname{Hom}(\Lambda^3(\mathfrak{g}),\, \mathbb{C})
\cong H^3(\mathfrak{g}) = \mathbb{C}$.
The homotopy transfer formula
$l_3^{\mathrm{tr}}(a,b,c) = p\,l_2(h\,l_2(i(a),i(b)),\, i(c))
+ \text{(cyclic)}$
is an explicit $3$-cocycle on~$\mathfrak{g}$, hence proportional
to~$\phi$. It is nonzero because the contracting homotopy~$h$
maps the bar-length-$2$ image of $l_2(i(a),i(b))$, which carries
the Killing-form data via the cyclic pairing, nontrivially
into the obstruction direction.
Normalizing the generator of~$H^2$ to agree with
Proposition~\ref{prop:killing-linf-extension}
gives $l_3^{\mathrm{tr}} = \phi$.

\emph{$l_n^{\mathrm{tr}} = 0$ for $n \geq 4$.}
For $n = 4$: the transfer formula
(Loday--Vallette \cite{LV12} Theorem~10.3.1) gives
$l_4^{\mathrm{tr}}$ as a composite of $l_2^{\mathrm{tr}}, l_3^{\mathrm{tr}}$
and the homotopy~$h$. Every such composite lands in
$H^2 \cong \mathbb{C}\cdot\eta$, and composing further with
$l_2^{\mathrm{tr}}$ on this $1$-dimensional space vanishes
(since $\eta \notin \mathfrak{g}$). Equivalently, the degree-$4$
homotopy Jacobi identity $\partial l_3^{\mathrm{tr}} +
\{l_2^{\mathrm{tr}}, l_4^{\mathrm{tr}}\} = 0$ combined with
$\partial l_3^{\mathrm{tr}} = \delta\phi = 0$
(Proposition~\ref{prop:killing-linf-extension}(i)) forces
$l_4^{\mathrm{tr}} = 0$.
For $n \geq 5$: every tree in the transfer formula with $n$ leaves
must route at least one input through $H^2 \cong \mathbb{C}\cdot\eta$;
since $l_3^{\mathrm{tr}}$ requires all three inputs in~$\mathfrak{g}$
and $\eta \notin \mathfrak{g}$, the composite vanishes.
\end{proof}

\begin{remark}[Scope of Theorem~\ref{thm:mc2-1-km}]
\label{rem:mc2-1-km-scope}
This resolves Hypothesis~\ref{mc2-hyp:cyclic} for
$\widehat{\mathfrak{g}}_k$ at non-critical levels; all three MC2
hypotheses are now resolved by
Theorem~\ref{thm:mc2-full-resolution}. The minimal model
(Corollary~\ref{cor:km-minimal-linf}) gives the Killing extension,
computed independently in Remark~\ref{rem:mc2-status}(v)--(ix).
\end{remark}

\begin{remark}[Why introduce \texorpdfstring{$\operatorname{Def}_{\mathrm{cyc}}$}{Def_cyc}
now]\label{rem:def-cyc-role}
The genus tower, curvature scalars $\mcurv{g}$, and obstruction
classes are all shadows of the shadow obstruction tower
$\Theta_\cA^{\leq r}$
(Definition~\ref{def:shadow-postnikov-tower})
in $\operatorname{Def}_{\mathrm{cyc}}(\cA) \widehat{\otimes}
R\Gamma(\overline{\mathcal{M}}_{g,\bullet}, \mathbb{Q})$;
at the scalar level (Theorem~\ref{thm:explicit-theta}):
$\Theta_\cA^{\leq 2} = \kappa(\cA) \cdot \eta \otimes \Lambda$
(here $\Theta_\cA^{\leq 2}$ is \emph{proved} for all Koszul families;
the full tower $\Theta_\cA = \varprojlim_r \Theta_\cA^{\leq r}$
exists by the bar-intrinsic construction;
see Theorem~\ref{thm:mc2-bar-intrinsic}).
The MC2 extension problem is packaged by
Proposition~\ref{prop:mc2-reduction-principle} into three tasks;
the clutching package is resolved by
Proposition~\ref{prop:geometric-modular-operadic-mc}, and the
surviving one-channel obstruction is further reduced through a
sequence of propositions
(\ref{prop:tautological-line-support-criterion}--\ref{prop:one-channel-normalization-criterion})
to a normalized scalar comparison.
\end{remark}

\subsection{The chiral graph complex and cyclic
\texorpdfstring{$L_\infty$}{L-infinity} structure}
\label{subsec:chiral-graph-complex}
\index{graph complex!chiral}
\index{cyclic $L_\infty$-algebra!graph complex construction}

The dg Lie algebra $\Defcyc(\cA)$ of
Proposition~\ref{prop:genus0-cyclic-coderivation} is the genus-$0$ shadow
of a cyclic $L_\infty$-algebra whose higher
brackets $l_n$ ($n \geq 3$) arise from sums over chiral
graphs. This adapts the Kontsevich--Soibelman
construction~\cite{KontsevichSoibelman} to chiral algebras on algebraic
curves via configuration space integrals on
Fulton--MacPherson compactifications.

\begin{definition}[Chiral graphs; \ClaimStatusDefinitional]\label{def:chiral-graph}
\index{chiral graph}
A \emph{chiral graph} of type $(n, L)$ is a datum
$\Gamma = (V_{\mathrm{ext}}, V_{\mathrm{int}}, E, \sigma)$ where:
\begin{enumerate}[label=\textup{(\roman*)}]
\item $V_{\mathrm{ext}} = \{v_1, \ldots, v_n\}$ is an ordered set of
 $n$~\emph{external vertices} (legs);
\item $V_{\mathrm{int}}$ is a finite set of \emph{internal vertices},
 each of valence $\geq 3$;
\item $E$ is a set of edges connecting vertices in
 $V = V_{\mathrm{ext}} \sqcup V_{\mathrm{int}}$, with the graph
 $(V, E)$ connected;
\item $\sigma$ is a cyclic ordering of the half-edges at each internal
 vertex (ribbon structure);
\item the loop number is $L = |E| - |V| + 1$.
\end{enumerate}
The automorphism group $\operatorname{Aut}(\Gamma)$ consists of
graph automorphisms preserving the external vertex labeling and the
ribbon structure. Let $\mathfrak{G}_{n,L}$ denote the set of
isomorphism classes of chiral graphs of type $(n, L)$.
\end{definition}

\begin{definition}[Graph amplitude; \ClaimStatusDefinitional]\label{def:graph-amplitude}
\index{graph amplitude!chiral}
Let $\cA$ be a Koszul chiral algebra on a smooth projective curve~$X$
with non-degenerate invariant form
$\langle{-},{-}\rangle_{\cA}\colon \cA\otimes\cA\to\omega_X$,
and let $S(z,w)$ denote the Szeg\H{o} kernel on~$X$: the unique
section of
$\omega_X^{1/2} \boxtimes \omega_X^{1/2}(\Delta)$
with a simple pole along the diagonal and no other singularities.

For a chiral graph $\Gamma \in \mathfrak{G}_{n,L}$ and cyclic
coderivations $f_1, \ldots, f_n \in
\operatorname{CoDer}^{\mathrm{cyc}}(\barB_X(\cA))$, the
\emph{graph amplitude} is
\begin{equation}\label{eq:graph-amplitude}
w_\Gamma(f_1, \ldots, f_n)
\;:=\;
\int_{\ConfigSpace{|V_{\mathrm{int}}|}}
\biggl(\prod_{e = (v_a, v_b) \in E_{\mathrm{int}}}
 S(z_a, z_b)\biggr)
\,\biggl(\prod_{i=1}^{n} \iota_{v_i}(f_i)\biggr)
\,\biggl(\prod_{v \in V_{\mathrm{int}}} m_{\mathrm{val}(v)}\biggr),
\end{equation}

where:

\begin{itemize}
\item $\ConfigSpace{|V_{\mathrm{int}}|}$ is the Fulton--MacPherson
 compactification of the configuration of internal vertex positions
 on~$X$;
\item $E_{\mathrm{int}} \subset E$ is the set of internal edges
 (both endpoints in $V_{\mathrm{int}}$);
\item $\iota_{v_i}(f_i)$ inserts the coderivation $f_i$ at the
 external vertex~$v_i$ via the bar-length projection
 $\operatorname{pr}_1\colon T^c(s^{-1}\bar{\cA}) \to
 s^{-1}\bar{\cA}$;
\item $m_{\mathrm{val}(v)}$ is the $A_\infty$ operation of
 degree $\mathrm{val}(v)$ at each internal vertex~$v$, contracted
 with $\langle{-},{-}\rangle_{\cA}$ along each internal edge;
\item signs follow the Koszul convention determined by the ribbon
 structure of~$\Gamma$
 (Appendix~\ref{app:signs}, \S\ref{app:sign-conventions}).
\end{itemize}
When $V_{\mathrm{int}} = \emptyset$ (all vertices external), the
integral is absent and the amplitude is a direct algebraic contraction.
\end{definition}

\begin{construction}[$L_\infty$ brackets from graph
sums]\label{constr:linf-from-graphs}
\index{$L_\infty$-algebra!graph sum construction}
Define the $n$-ary bracket on
$\Defcyc(\cA) = \operatorname{CoDer}^{\mathrm{cyc}}(\barB_X(\cA))[1]$
by
\begin{equation}\label{eq:linf-graph-sum}
l_n(f_1, \ldots, f_n)
\;:=\;
\sum_{\Gamma \in \mathfrak{G}_{n,0}}
\frac{(-1)^{\epsilon(\Gamma)}}
{\lvert\operatorname{Aut}(\Gamma)\rvert}
\, w_\Gamma(f_1, \ldots, f_n),
\end{equation}
where the sum runs over \emph{tree-level} chiral graphs ($L = 0$) with
$n$~external legs, and $\epsilon(\Gamma)$ is the sign from the
canonical orientation of~$\Gamma$ determined by the ribbon structure
and edge ordering.
\end{construction}

The graph sum~\eqref{eq:linf-graph-sum} is well-defined for Koszul
chiral algebras:
(a)~$\mathfrak{G}_{n,0}$ is finite at each degree~$n$ (a tree on $n + k$ vertices with $n$~external legs has
$k$~internal vertices, and valence~$\geq 3$ forces
$k \leq n - 2$);
(b)~the integrals over~$\ConfigSpace{k}$ converge by
regularity of forms on the FM~compactification;
(c)~the bar-length filtration makes the sum locally finite in each
filtered degree.

\begin{proposition}[Fay trisecant identity with prime-form normalisation {\cite[Corollary~2.5]{Fay73}};
\ClaimStatusProvedElsewhere]\label{prop:fay-trisecant}
\index{Fay trisecant identity}
\index{Szeg\H{o} kernel!Fay identity}
Let $X$ be a smooth projective curve of genus~$g \geq 0$ and
$S(z,w)$ a Szeg\H{o} kernel on~$X$: the section of
$\omega_X^{1/2} \boxtimes \omega_X^{1/2}(\Delta)$ with simple pole
along the diagonal and the chosen theta-characteristic
normalisation. Fay's trisecant identity is a section-level
prime-form/theta addition identity on $C_3(X)$; after pullback to
the three-point FM boundary and contraction with the
Arnold--Kohno/IHX tensor relation it cancels the three local
Stokes contributions in the chiral graph complex. It is not the
bare scalar identity
$S_{12}S_{23}+S_{23}S_{31}+S_{31}S_{12}=0$ for an unnormalised
kernel.

In the elliptic normalisation with periods $(1,\tau)$ one concrete
specialisation is
\begin{equation}\label{eq:algebraic-fay}
 \wp(a,\tau)-\wp(b,\tau)
 =
 -\,\frac{
 \vartheta_1(a+b,\tau)\,\vartheta_1(a-b,\tau)}
 {\vartheta_1(a,\tau)^2\,\vartheta_1(b,\tau)^2}
 \,\vartheta_1'(0,\tau)^2 .
\end{equation}
At genus~$0$, the same Fay identity degenerates to the
partial-fraction identity for $1/(z_1-z_2)(z_2-z_3)$.
\end{proposition}

\begin{proof}
This is the trisecant identity of Fay~\cite[Corollary~2.5]{Fay73},
proved via the addition formula for theta functions on~$X$ and
the prime form~$E(z,w)$. The displayed torus formula is the
theta-prime-normalised genus-$1$ specialisation; the
$\vartheta_1'(0,\tau)^2$ factor is part of the identity.
\end{proof}

\begin{proposition}[Stokes regularity for graph amplitudes on FM
compactifications;
\ClaimStatusProvedHere]\label{prop:stokes-regularity-FM}
\index{FM compactification!Stokes regularity}
\index{graph amplitude!boundary regularity}
Let $X$ be a smooth projective curve, $\cA$ a Koszul chiral algebra
on~$X$ with non-degenerate invariant form, and
$\omega_\Gamma \in \Omega^*(\ConfigSpace{k})$
a graph amplitude form~\eqref{eq:graph-amplitude} for a chiral
graph~$\Gamma$ with $k$~internal vertices. Then:
\begin{enumerate}[label=\textup{(\alph*)}]
\item \emph{Logarithmic regularity.}
 The form $\omega_\Gamma$ extends to a form with at most logarithmic
 singularities along the boundary divisors $D_S \subset
 \ConfigSpace{k}$
 \textup{(}$S \subset \{1,\ldots,k\}$, $|S|\geq 2$\textup{)}.
 In particular, $\omega_\Gamma$ is locally integrable on
 $\ConfigSpace{k}$.
\item \emph{Stokes identity.}
 \begin{equation}\label{eq:stokes-FM}
 \int_{\ConfigSpace{k}} d\omega_\Gamma
 \;=\;
 \sum_{S} \int_{D_S} \omega_\Gamma,
 \end{equation}
 where the sum runs over boundary divisors.
\item \emph{Boundary correspondence.}
 The restriction of $\omega_\Gamma$ to the boundary divisor
 $D_S \cong \ConfigSpace{|S|} \times \ConfigSpace{k-|S|+1}$
 computes the product of graph amplitudes obtained by collapsing
 the subset~$S$ to a single vertex.
\end{enumerate}
\end{proposition}

\begin{proof}
\emph{(a) Logarithmic regularity.}
The FM compactification $\ConfigSpace{k}$ is a smooth algebraic
variety with simple normal crossing boundary
divisor~\cite{FM94}. The graph amplitude form $\omega_\Gamma$ is a
product of Szeg\H{o} propagators $S(z_a, z_b)$, one per internal
edge, contracted with $A_\infty$ operations and the invariant form
of~$\cA$. Each factor $S(z_a, z_b)$ is a section of
$\omega_X^{1/2} \boxtimes \omega_X^{1/2}(\Delta_{ab})$
with a simple pole along the diagonal $\Delta_{ab}$.

In the FM compactification, the diagonal $\Delta_{ab}$ is replaced
by the exceptional divisor $D_{\{a,b\}}$. In the local coordinates
$(r_S, \theta_S)$ of the FM blowup at $D_S$
\cite[§2]{FM94}, the propagator has the expansion
$S(z_a, z_b) = r_S^{-1} \cdot \phi_S + O(1)$ where $\phi_S$ is a
smooth section on $D_S$. This is a log pole:
$S(z_a,z_b) \, dz_a$ contributes $dr_S/r_S$ (a log
$1$-form) plus smooth terms. Hence the product over edges gives a
form with logarithmic poles along each $D_S$ meeting the graph's
edges, precisely a section of $\Omega^*(\log D)$ on
$\ConfigSpace{k}$, where $D = \bigcup_S D_S$.

\emph{(b) Stokes identity.}
For log forms on a smooth variety with simple normal crossing
boundary divisor, the residue form of Stokes' theorem holds:
$\int d\omega = \sum_S \int_{D_S} \operatorname{Res}_{D_S} \omega$.
This is the content of the Leray residue calculus for log de~Rham
complexes~\cite{Deligne70}. Alternatively, the real-analytic
version follows from Stokes' theorem on manifolds with corners
applied to the real-oriented blowup of
$\ConfigSpace{k}$~\cite{AS94}, as used in Chern--Simons
perturbation theory. The application to configuration-space
integrals with propagator insertions is carried out
in~\cite{Kon03} and~\cite{costello-renormalization}.

\emph{(c) Boundary correspondence.}
When the subset $S = \{z_{a_1}, \ldots, z_{a_s}\}$ of internal
vertex positions collide, the FM blowup produces the exceptional
divisor $D_S \cong \ConfigSpace{|S|} \times \ConfigSpace{k-|S|+1}$,
where the first factor parametrizes the relative positions of the
colliding points and the second parametrizes the remaining
configuration with $S$ collapsed to a single point~\cite[§3]{FM94}.
The residue of $\omega_\Gamma$ along $D_S$ localizes the
propagators within $S$: the edges internal to $S$ contribute the
graph amplitude of the subgraph $\Gamma|_S$, while the edges
connecting $S$ to the complement contribute the amplitude of the
quotient graph $\Gamma/S$ with $S$ collapsed to a single vertex
of valence $\sum_{a \in S} \mathrm{val}(a) - 2(|E_S|)$, where
$|E_S|$ counts edges internal to~$S$. The product of these two
amplitudes is exactly the graph amplitude of the collapsed graph,
as required.
\end{proof}

\begin{theorem}[Cyclic \texorpdfstring{$L_\infty$}{L-infinity} structure via chiral graph
complex; \ClaimStatusProvedHere]\label{thm:cyclic-linf-graph}
\index{cyclic $L_\infty$-algebra!from chiral graph complex}
Let $\cA$ be a Koszul chiral algebra on a smooth projective
curve~$X$ with non-degenerate invariant form
$\langle{-},{-}\rangle_{\cA}$.
The two analytic inputs are now unconditional:
\begin{enumerate}[label=\textup{(H\arabic*)}]
\item\label{hyp:stokes-FM}
 \emph{Stokes regularity on FM compactifications}
 \textup{(}Proposition~\textup{\ref{prop:stokes-regularity-FM}}\textup{)}.
\item\label{hyp:fay-propagator}
 \emph{Fay trisecant identity for the Szeg\H{o} kernel}
 \textup{(}Proposition~\textup{\ref{prop:fay-trisecant}}\textup{)}.
\end{enumerate}
The brackets $\{l_n\}_{n \geq 1}$ of
Construction~\textup{\ref{constr:linf-from-graphs}} equip
$\Defcyc(\cA)$ with the structure of a cyclic $L_\infty$-algebra:
\begin{enumerate}[label=\textup{(\alph*)}]
\item \emph{Homotopy Jacobi.}
 For each degree~$N \geq 1$:
 \begin{equation}\label{eq:homotopy-jacobi-graph}
 \sum_{\substack{i + j = N + 1 \\
 \sigma \in \operatorname{Sh}(i,N-i)}}
 \varepsilon(\sigma)\,
 l_j\bigl(l_i(f_{\sigma(1)}, \ldots, f_{\sigma(i)}),\,
 f_{\sigma(i+1)}, \ldots, f_{\sigma(N)}\bigr) \;=\; 0.
 \end{equation}
\item \emph{Recovery.}
 $l_1 = [d_{\barB}, -]$ and
 $l_2 = [-,-]$ coincide with the dg Lie brackets of
 Proposition~\textup{\ref{prop:genus0-cyclic-coderivation}}.
\item \emph{Cyclic symmetry.}
 The pairing~$\eta$ of
 Proposition~\textup{\ref{prop:genus0-cyclic-coderivation}(i)}
 satisfies, for all~$n$:
 \begin{equation}\label{eq:cyclic-symmetry-graph}
 \langle l_n(f_1, \ldots, f_n), f_{n+1}\rangle_\eta
 \;=\;
 (-1)^{n + |f_1|(|f_2|+\cdots+|f_{n+1}|)}
 \langle l_n(f_2, \ldots, f_{n+1}), f_1\rangle_\eta.
 \end{equation}
 In particular, the formal moduli problem
 $\operatorname{MC}(\Defcyc(\cA))$ inherits a $(-1)$-shifted
 symplectic structure.
\end{enumerate}
\end{theorem}

\begin{proof}
(a)\enspace
The homotopy Jacobi identity at degree~$N$ translates into the
vanishing of the sum
$\mathrm{Jac}_N = \sum l_j \circ l_i$
(with appropriate signs and unshuffles).
Express each $l_i$ and $l_j$ as a graph sum. The composition
$l_j \circ l_i$ corresponds to gluing: attach the output of a
tree~$\Gamma_i \in \mathfrak{G}_{i,0}$ to one input of a
tree~$\Gamma_j \in \mathfrak{G}_{j,0}$ to produce a graph
$\Gamma_i \circ_v \Gamma_j$ with exactly one internal edge
joining the two subtrees.

The integral over $\ConfigSpace{k}$ of the
glued graph computes a boundary contribution: the internal edge
carries a propagator $S(z_a, z_b)$ whose pole as
$z_a \to z_b$ localises on a boundary divisor~$D_S$.
By Proposition~\ref{prop:stokes-regularity-FM}(b), the sum of
all boundary contributions equals
$\int_{\ConfigSpace{k}} d\,\omega_N$ for an explicit
form~$\omega_N$. Stokes' theorem yields
$\int d\omega_N = \int_\partial \omega_N$; the boundary strata
are exactly the gluing terms appearing in the Jacobi sum.
By Proposition~\ref{prop:stokes-regularity-FM}(c), each boundary
integral computes the product of graph amplitudes from collapsing
a subset to a vertex. These products correspond bijectively to the
$(l_j \circ l_i)$-terms in $\mathrm{Jac}_N$, and cancel pairwise
by the antisymmetry of~$S$ and the skew-symmetry of the coderivation
bracket.

For the terms where three propagators meet at a common vertex, the
Fay trisecant identity
(Proposition~\ref{prop:fay-trisecant}) ensures the correct
cancellation after the prime-form normalisation and the
Arnold--Kohno/IHX tensor contraction. This replaces the IHX
relation of classical graph homology~\cite{LV12} with its
chiral-algebraic analogue on~$X$: the three graphs related by an
IHX move produce amplitudes whose section-level sum vanishes by
Fay's identity, whose elliptic scalar normalisation is recorded in
\eqref{eq:algebraic-fay}. The proof does not use the false
unnormalised scalar product
$S_{12}S_{23}+S_{23}S_{31}+S_{31}S_{12}=0$.

(b)\enspace
At degree $n = 1$: the unique element of $\mathfrak{G}_{1,0}$ is the
single-vertex graph with one external leg and no internal vertices.
Its amplitude is the identity insertion $f_1 \mapsto f_1$, acted on
by the bar differential: $l_1(f_1) = [d_{\barB}, f_1]$.

At degree $n = 2$: $\mathfrak{G}_{2,0}$ contains a single tree, the
edge connecting $v_1$ to~$v_2$, with no internal vertices. Its
amplitude is the commutator $[f_1, f_2]$ of coderivations.

(c)\enspace
Cyclic symmetry. Rotating the inputs $(f_1, \ldots, f_{n+1})$
by one step produces two signs:
\begin{enumerate}
\item \emph{Koszul sign from reordering.} Moving
 $f_1$ (degree $|f_1|$) past
 $f_2 \otimes \cdots \otimes f_{n+1}$
 (total degree $|f_2| + \cdots + |f_{n+1}|$) in the
 graded tensor product gives $(-1)^{|f_1|(|f_2|+\cdots+|f_{n+1}|)}$.
\item \emph{Degree of $l_n$.} The operation $l_n$ has degree
 $2 - n$ (cohomological convention). Moving it past the inner
 product $\langle{-},{-}\rangle_\eta$ (degree~$0$) and
 accounting for the $n$-fold desuspension gives a factor
 $(-1)^n$.
\end{enumerate}
The product of these two factors is exactly the sign
in~\eqref{eq:cyclic-symmetry-graph}. The ribbon structure
on each graph $\Gamma \in \mathfrak{G}_{n,0}$ induces a cyclic
ordering of external legs, and $\mathrm{ad}$-invariance of
$\langle{-},{-}\rangle_\cA$ ensures that the amplitude
$w_\Gamma$ respects the cyclic rotation with this sign.
\end{proof}

\begin{corollary}[Recovery of the Killing cocycle extension;
\ClaimStatusProvedHere]\label{cor:killing-recovery-graph}
\index{Killing cocycle!recovery from graph complex}
By Theorem~\textup{\ref{thm:cyclic-linf-graph}},
for $\widehat{\mathfrak{g}}_k$ with $\mathfrak{g}$ simple and
$k \neq -h^\vee$:
\begin{enumerate}[label=\textup{(\roman*)}]
\item On the generators
 $s^{-1}\bar{\mathfrak{g}} \subset \Defcyc(\widehat{\mathfrak{g}}_k)$:
 $l_3(a,b,c) = \kappa([a,b],c)$, recovering the Killing $3$-cocycle of
 Proposition~\textup{\ref{prop:killing-linf-extension}}.
\item $l_n = 0$ on generators for all $n \geq 4$.
\end{enumerate}
\end{corollary}

\begin{proof}
(i)\enspace
The tree graphs in $\mathfrak{G}_{3,0}$ contributing to $l_3$ on
generators are the tripod graphs: three external legs
meeting at one trivalent internal vertex. (No other tree with
$3$~leaves has all internal valences~$\geq 3$.)
The internal vertex carries the chiral bracket~$m_2$ and the
edges contract via the invariant form, giving
$w_{\text{tripod}}(a,b,c) = \langle [a,b], c \rangle_\kappa
= \kappa([a,b],c)$.
No configuration-space integration is needed (the single internal
vertex contributes a $0$-dimensional moduli space).

(ii)\enspace
Any tree in $\mathfrak{G}_{n,0}$ with $n \geq 4$~leaves
and all internal valences~$\geq 3$ has at least $n - 2 \geq 2$
internal vertices and at least one internal vertex of
valence~$\geq 3$. For
$\cA = \widehat{\mathfrak{g}}_k$, the $A_\infty$ operations
satisfy $m_k = 0$ for $k \geq 3$ on the generators (the OPE is at
most quadratic: $m_1 = 0$, $m_2 = \mu$, $m_k = 0$ for
$k \geq 3$). Every tree with $n \geq 4$ external legs requires
internal vertices with total valence exceeding what $m_2$ alone can
supply, forcing at least one $m_k$ with $k \geq 3$ to appear, which
vanishes.

Alternatively, this follows from
Proposition~\ref{prop:killing-linf-extension}: the degree-$\geq 5$
homotopy Jacobi identities are automatically satisfied by the
truncated system $l_1 = 0$, $l_2 = [-,-]$, $l_3 = \phi$,
$l_n = 0$ for $n \geq 4$, so the graph-complex construction
agrees with this truncation on generators.
\end{proof}

\begin{remark}[Relation to known constructions]
\label{rem:graph-complex-context}
\index{Kontsevich--Soibelman!chiral adaptation}
(i)~On $X = \mathbb{A}^1$, the graph sum reduces to the
Kontsevich--Soibelman $L_\infty$ structure on
$\operatorname{CoDer}(BA)[1]$~\cite{KontsevichSoibelman}.
(ii)~In the Costello--Gwilliam formalism~\cite{CG17,CG-vol2}, the
graph sum is perturbative quantization of a holomorphic CS-type
theory on~$X$.
(iii)~The full graph complex
$\bigoplus_L \mathfrak{G}_{n,L}$ carries
Getzler--Kapranov~\cite{GeK98} modular-operadic structure; the
$L > 0$ part encodes the genus tower.
(iv)~Both analytic hypotheses are unconditional:
Proposition~\ref{prop:stokes-regularity-FM} and
Proposition~\ref{prop:fay-trisecant}.
\end{remark}

\begin{remark}[Physical interpretation]\label{rem:graph-complex-physics}
Each chiral graph $\Gamma$ is a tree-level Feynman diagram: external
legs carry deformation inputs, internal edges carry the Szeg\H{o}
propagator $S(z,w)$, and internal vertices carry $A_\infty$
operations. The homotopy Jacobi identity is the Ward identity,
mediated by Stokes' theorem on $\ConfigSpace{k}$. At
$L > 0$ (loop level), $L = g$
loop diagrams produce the genus-$g$ obstruction; the cyclic pairing ensures the BV QME order by
order~\cite{costello-renormalization}.
\end{remark}

\subsection{The modular deformation complex}
\label{subsec:modular-deformation-complex}
\index{modular deformation complex|textbf}

The genus-$0$ cyclic deformation complex
$\Defcyc(\cA)$ (Definition~\ref{def:cyclic-deformation-elementary})
is the tree-level shadow of a modular quantum $L_\infty$-algebra.
The modular version assembles all genera:

\begin{definition}[Modular deformation complex]
\label{def:modular-deformation-complex}
\index{modular deformation complex|textbf}
\index{quantum $L_\infty$ algebra!modular deformation complex}
For a chiral algebra~$\cA$ on a smooth projective curve~$X$, the
\emph{modular deformation complex} is the \emph{quantum
$L_\infty$-algebra} (modular convolution $L_\infty$-algebra)
\begin{equation}\label{eq:modular-deformation-complex}
\operatorname{Def}^{\mathrm{mod}}(\cA)
\;:=\;
\prod_{g,n}
\operatorname{Hom}_{\Sigma_n}\!\bigl(
 C_*(\overline{\mathcal{M}}_{g,n}),\;
 \operatorname{End}_\cA(g,n)
\bigr),
\end{equation}
where:
\begin{itemize}
\item $C_*(\overline{\mathcal{M}}_{g,n})$ denotes chains on the
 Deligne--Mumford compactification, and the product runs over
 stable pairs $(g,n)$ with $2g - 2 + n > 0$;
\item $\operatorname{End}_\cA(g,n)$ is the endomorphism space of
 $\cA^{\otimes n}$ with genus-$g$ curvature;
\item the $L_\infty$ brackets arise from \emph{graph gluing}:
 the binary bracket $\ell_2$ composes endomorphisms along separating
 clutching maps; the BV operator $\ell_1^{(g \geq 1)} = \hbar^g\Delta$
 contracts legs along non-separating clutching maps; the higher
 brackets $\ell_n^{(g)}$ are genus-$g$ stable-graph Feynman amplitudes
 (Theorem~\ref{thm:modular-quantum-linfty}).
\end{itemize}
A Maurer--Cartan element
$\Theta \in
\operatorname{MC}(\operatorname{Def}^{\mathrm{mod}}(\cA))$ is
equivalent to a modular homotopy type on the underlying graded
object of~$\cA$
(Definition~\ref{def:modular-homotopy-abstract}).
The MC equation
$d\Theta + \frac{1}{2}[\Theta,\Theta] = 0$
encodes clutching compatibility at every genus.
\end{definition}

\begin{remark}[Homotopy-invariant deformation complex]
\label{rem:deformation-complex-linfty}
The modular deformation complex of Definition~\ref{def:modular-deformation-complex}
is the strict dg~Lie model of a complete filtered cyclic modular $L_\infty$
algebra $\mathrm{Def}_\infty^{\mathrm{mod}}(\cA)$. The higher brackets
encode homotopy-coherent deformation data: $\ell_3$ governs
obstructions to extending first-order deformations, $\ell_4$ governs
secondary obstructions. The coderivation model computes the
same formal moduli problem \cite{KontsevichSoibelman}.

By Tamarkin's theorem \cite{Tamarkin00}, the deformation complex of a
chiral algebra (viewed as an algebra over the chiral operad, which is a
2-algebra in the sense of $E_n$-operads) carries a canonical $E_2$
structure, a homotopy Gerstenhaber algebra. The bar functor is an
$E_2$-coalgebra functor, not a dg~coalgebra functor; this is the algebraic
source of the bar-cobar adjunction's naturality.

Theorem~H identifies the curve-level tangent complex with a
family-indexed support model $K_{\cA,S}$.  When $3\notin S$, the
degree-three extension group is zero.  The transferred
$L_\infty$-brackets are controlled by the additional brace-compatible
contraction of Proposition~\ref{prop:e2-formality-hochschild}.
\end{remark}

\begin{remark}[Modular versus cyclic]
\label{rem:modular-vs-cyclic-complex}
The restriction of the strict model
$\operatorname{Def}^{\mathrm{mod}}(\cA)$ to genus~$0$
recovers $\Defcyc(\cA)$, and the same is true for the transferred
homotopy object $\Definfmod(\cA)$ after forgetting every
positive-genus Taylor component. At $g \geq 1$, the coefficients are
chains on $\overline{\mathcal{M}}_{g,n}$, so stable-graph geometry
enters the deformation object itself rather than merely its differential
(Remark~\ref{rem:convolution-hierarchy}).
\end{remark}

\begin{construction}[Modular homotopy deformation object]
\label{const:modular-homotopy-deformation-object}
\index{modular deformation!homotopy object|textbf}
Filter $\operatorname{Def}^{\mathrm{mod}}(\cA)$ by genus and
bar length. Choose any complete filtered contraction of this
quantum $L_\infty$-algebra onto a smaller filtered complex.
By the operadic homotopy transfer theorem
(Theorem~\ref{thm:htt-operadic}), the transferred higher
brackets $\{\ell_n\}_{n\ge 1}$ endow the target with a cyclic
$L_\infty$-structure. Its filtered $L_\infty$ quasi-isomorphism class
is denoted
\[
\Definfmod(\cA).
\]
The formula in
Definition~\ref{def:modular-deformation-complex} is one explicit
strict chart for this invariant object.
\end{construction}

\begin{proposition}[Strictification principle for modular deformation theory; \ClaimStatusProvedHere]
\label{prop:modular-strictification-principle}
\index{strictification!modular deformation|textbf}
The strict modular deformation complex
$\operatorname{Def}^{\mathrm{mod}}(\cA)$ is a strict model
of $\Definfmod(\cA)$. Consequently:
\begin{enumerate}[label=\textup{(\roman*)}]
\item complete filtered $L_\infty$ quasi-isomorphic choices of strict
 model determine equivalent Maurer--Cartan deformation functors;
\item the genus-zero truncation of
 $\Definfmod(\cA)$ recovers the cyclic
 deformation object $\Defcyc(\cA)$;
\item the canonical class $\Theta_\cA$ is a Maurer--Cartan element whose
 Taylor components simultaneously govern the modular bar differential,
 clutching, and the higher brackets on every transferred model.
\end{enumerate}
\end{proposition}

\begin{proof}
The first statement is exactly the output of operadic homotopy transfer
for complete filtered dg~Lie algebras
(Theorem~\ref{thm:htt-operadic}, applied to the dg~Lie operad in the
complete filtered category).
Maurer--Cartan deformation functors are invariant under filtered
$L_\infty$ quasi-isomorphism
\cite{KontsevichSoibelman,Get09}. Restricting the stable-graph
formulas to trees gives genus~$0$. The formulas defining
$D_{\mathrm{tot}}$ and the clutching operators are written as
coderivation actions of $\Theta_\cA$, so the transferred higher
brackets are the same data expressed on a different model.
\end{proof}

\subsection{Concrete modular constructions: from stable graphs to quantum homotopy theory}
\label{subsec:concrete-modular-constructions}
\index{modular homotopy theory!concrete constructions|textbf}
\index{stable graph!bar differential components}

The passage from tree-level to modular homotopy theory
(\S\ref{subsec:tree-to-modular}) replaces rooted trees by stable
graphs. At genus $g \geq 1$, curvature appears:
$\dfib^{\,2} \neq 0$, because propagators circulating around
non-contractible cycles of~$\Sigma_g$ produce period integrals.
The organizing principle: every codimension-$1$ boundary stratum of
$\overline{\mathcal{M}}_{g,n}$ contributes one component of the
modular bar differential; the Maurer--Cartan equation is $\partial^2
= 0$ read through the Feynman transform; and the categorical
logarithm $\barB^{\mathrm{tot}}_X$ of
Remark~\ref{rem:nilpotence-periodicity} acquires its strict
nilpotence $D_{\mathrm{tot}}^2 = 0$ by this mechanism.

\subsubsection{The three-component modular bar differential}
\label{subsubsec:three-component-differential}
\index{modular bar differential!three components|textbf}

Recall from Definition~\ref{def:feynman-transform} that the modular bar
complex is the Feynman transform
$\barB^{\mathrm{mod}}(\cA) = \mathrm{FT}_{\mathcal{M}od}(\cA)$,
with underlying graded space the stable-graph sum
of~\eqref{eq:bar-graph-sum}. Its differential decomposes into three
geometrically distinct components.

\begin{definition}[Three-component modular bar differential]
\label{def:three-component-modular-bar}
\index{modular bar differential!d_int@$d_{\mathrm{int}}$}
\index{modular bar differential!d_sep@$d_{\mathrm{sep}}$}
\index{modular bar differential!d_nonsep@$d_{\mathrm{ns}}$}
Let $\Gamma \in \mathcal{G}^{\mathrm{st}}_{g,n}$ be a stable graph
of genus~$g$ with $n$~legs. On the summand
\[
\barB^{(\Gamma)}(\cA)
\;:=\;
\bigotimes_{v \in V(\Gamma)} \cA(g_v, \operatorname{val}(v))
\;\otimes\; \operatorname{or}(E(\Gamma)),
\]
the modular bar differential is
$d_{\mathrm{mod}} = d_{\mathrm{int}} + d_{\mathrm{sep}} + d_{\mathrm{ns}}$,
where:
\begin{enumerate}[label=\textup{(\roman*)}]
\item \emph{Internal differential} $d_{\mathrm{int}}$.
 At each vertex~$v$ of~$\Gamma$ carrying genus label~$g_v$ and
 valence~$\operatorname{val}(v)$, the factor
 $\cA(g_v, \operatorname{val}(v))$ carries its own differential
 $d_v$ built from the $A_\infty$ operations $\{m_k\}$ and the
 genus-$g_v$ propagator on the fiber curve~$\Sigma_{g_v}$.
 The internal differential acts vertex-by-vertex:
 \begin{equation}\label{eq:d-int}
 d_{\mathrm{int}}
 \;=\;
 \sum_{v \in V(\Gamma)}
 \operatorname{id} \otimes \cdots \otimes d_v
 \otimes \cdots \otimes \operatorname{id}.
 \end{equation}
 It preserves the underlying graph~$\Gamma$.

\item \emph{Separating boundary differential} $d_{\mathrm{sep}}$.
 For each edge $e \in E(\Gamma)$ connecting distinct vertices
 $v_1, v_2$ (or each partition of the legs of a single vertex
 creating a separating node), the clutching map
 $\xi_{e}^{\mathrm{sep}}\colon
 \overline{\mathcal{M}}_{g_1, n_1+1} \times
 \overline{\mathcal{M}}_{g_2, n_2+1} \to
 \overline{\mathcal{M}}_{g, n}$
 produces a new graph~$\Gamma/e$ by contracting edge~$e$,
 merging the two endpoint vertices:
 \begin{equation}\label{eq:d-sep}
 d_{\mathrm{sep}} = \sum_{e \in E_{\mathrm{sep}}(\Gamma)}
 \pm\, \iota_e^* \circ \mu_e,
 \end{equation}
 where $\mu_e$ contracts the propagator along edge~$e$ using the
 invariant pairing
 $\langle -, - \rangle_\cA\colon \cA \otimes \cA \to \omega_X$,
 $\iota_e^*$ is the boundary pullback, and the sign is determined
 by the orientation $\operatorname{or}(E(\Gamma))$.
 Geometrically: $d_{\mathrm{sep}}$ implements the separating
 degenerations of the underlying curve.

\item \emph{Non-separating boundary differential} $d_{\mathrm{ns}}$.
 For each loop edge $e$ at vertex~$v$ in~$\Gamma$
 (an edge with both endpoints at~$v$),
 contracting~$e$ produces a graph $\Gamma / e$ in which
 the vertex genus increases by~$1$ ($g_v \mapsto g_v + 1$)
 and the loop disappears ($b_1 \mapsto b_1 - 1$):
 \begin{equation}\label{eq:d-nonsep}
 d_{\mathrm{ns}} = \sum_{\substack{e \in E(\Gamma) \\
 e \text{ loop at } v}}
 \pm\, \mu_e,
 \end{equation}
 where $\mu_e$ contracts the loop edge using the invariant
 pairing $\langle -, - \rangle_\cA$, absorbing the handle
 into the vertex genus. The sign is determined by
 $\operatorname{or}(E(\Gamma))$.
 Geometrically: $d_{\mathrm{ns}}$ contracts non-separating
 nodes, replacing a loop in the dual graph with a handle on
 the curve.
\end{enumerate}
\end{definition}

\begin{proposition}[$d_{\mathrm{mod}}^2 = 0$;
\ClaimStatusConditional]\label{prop:d-mod-squared-zero}
\index{modular bar differential!nilpotence}
The modular bar differential satisfies $d_{\mathrm{mod}}^2 = 0$.
\end{proposition}

\begin{proof}
Expand $d_{\mathrm{mod}}^2 = (d_{\mathrm{int}} + d_{\mathrm{sep}}
+ d_{\mathrm{ns}})^2$ into nine terms grouped in three classes.

By Theorem~\ref{thm:prism-higher-genus}(iii),
$d_{\mathrm{mod}}^2 = 0$ follows from the modular operad axioms
(Definition~\ref{def:modular-operad}). The three components give the
chain-level proof.

\emph{Diagonal terms.}
$d_{\mathrm{int}}^2 = 0$: here $d_{\mathrm{int}}$ is the
Feynman-transform internal differential, which acts on
$\barB^{(\Gamma)}(\cA) = \bigotimes_v \cA(g_v,
\operatorname{val}(v)) \otimes \operatorname{or}(E(\Gamma))$
as the total differential on vertex factors: the
\emph{strict} differential $\Dg{g_v}$ at each vertex
(Convention~\ref{conv:higher-genus-differentials}(ii)), which
satisfies $(\Dg{g_v})^2 = 0$ by
Theorem~\ref{thm:quantum-diff-squares-zero}. (The
\emph{fiberwise} curved differential $\dfib$ with
$\dfib^{\,2} = \kappa \cdot \omega_g$ is a separate object,
obtained by restricting to a single fiber.)
$d_{\mathrm{sep}}^2 = 0$ and $d_{\mathrm{ns}}^2 = 0$:
the boundary of a boundary stratum decomposes into
codimension-$2$ strata that cancel in pairs
($\partial^2 = 0$ on $\overline{\mathcal{M}}_{g,n}$).

\emph{Cross terms.}
$d_{\mathrm{int}} d_{\mathrm{sep}} + d_{\mathrm{sep}}
d_{\mathrm{int}} = 0$: edge contraction commutes with the vertex
differential because $d_v$ and the contraction $\mu_e$ act on
disjoint tensor factors (the contracting edge connects two
\emph{different} vertices, and $d_v$ acts on one).
$d_{\mathrm{int}} d_{\mathrm{ns}} + d_{\mathrm{ns}}
d_{\mathrm{int}} = 0$: the self-sewing $\Delta_v$ commutes with
$d_v$ because $\Delta_v$ pairs two marked points via the
$\operatorname{ad}$-invariant form $\langle -,- \rangle_\cA$,
which is a chain map.
$d_{\mathrm{sep}} d_{\mathrm{ns}} + d_{\mathrm{ns}}
d_{\mathrm{sep}} = 0$: codimension-$2$ corner cancellation
in the boundary stratification of $\overline{\mathcal{M}}_{g,n}$.
\end{proof}

\begin{remark}[Genus decomposition and exhaustiveness]
\label{rem:genus-decomposition}
\index{modular bar differential!genus accounting}
These three components exhaust the codimension-$1$ boundary
contributions to the modular bar differential.
A boundary divisor $D_\Gamma$ of
$\overline{\mathcal{M}}_{g,n}$ is indexed by a stable graph~$\Gamma$ with one edge.
Removing the edge either disconnects~$\Gamma$ (separating,
$d_{\mathrm{sep}}$) or does not (non-separating,
$d_{\mathrm{ns}}$). The third component $d_{\mathrm{int}}$ is
the vertex differential on each
component.

The three components act on the genus decomposition
$g = \sum_v g_v + b_1(\Gamma)$ as follows:
$d_{\mathrm{int}}$ preserves the graph
(hence preserves $b_1$ and all $g_v$);
$d_{\mathrm{sep}}$ contracts an edge between distinct components
(preserving $b_1$ and total genus);
$d_{\mathrm{ns}}$ absorbs a loop edge into a vertex
($g_v \mapsto g_v + 1$, $b_1 \mapsto b_1 - 1$, total
genus preserved).
The total genus $g = \sum_v g_v + b_1(\Gamma)$ is an invariant
of all three.
\end{remark}

\subsubsection{Quantum \texorpdfstring{$L_\infty$}{L-infinity} operations from loop-level graphs}
\label{subsubsec:quantum-linf-from-graphs}
\index{quantum L-infinity algebra@quantum $L_\infty$ algebra!from loop graphs|textbf}

The cyclic $L_\infty$ structure of
Theorem~\ref{thm:cyclic-linf-graph} is built from tree graphs
($L = 0$ loops). The next step extends it to loop-level graphs.
The operations below are written in $A_\infty$ form
(non-symmetric, with ordered inputs); the cyclic pairing
$\langle -, - \rangle_\cA$ symmetrizes them into quantum
$L_\infty$ operations \cite{BL15, Markl06}.

\begin{definition}[Loop-level graph operations]
\label{def:loop-level-graph-ops}
\index{loop-level graph operations}
For $n \geq 0$ and $L \geq 0$, let
$\mathfrak{G}_{n,L}$ be the set of connected graphs with
$n$~external legs, $L$~loops ($b_1 = L$), and vertex genus
labels $\{g_v\}_{v \in V}$ satisfying the stability condition
$2g_v - 2 + \operatorname{val}(v) > 0$ at each vertex.
For each $\Gamma \in \mathfrak{G}_{n,L}$, define the
\emph{loop-level graph amplitude}
\begin{equation}\label{eq:loop-amplitude}
w_\Gamma^{(L)}(f_1, \ldots, f_n)
\;=\;
\int_{\ConfigSpace{|V_{\mathrm{int}}|}(\Sigma)}
\prod_{e \in E_{\mathrm{int}}(\Gamma)} S(z_{s(e)}, z_{t(e)})
\;\cdot\;
\prod_{v \in V(\Gamma)} m_{\mathrm{val}(v)}
\bigl(\text{inputs at } v\bigr),
\end{equation}
where $\Sigma$ is a smooth projective curve,
$V_{\mathrm{int}}$ denotes internal vertices, $S(z,w)$ is
the Szeg\H{o} propagator on~$\Sigma$, and
$m_{\mathrm{val}(v)}$ is the $A_\infty$ vertex operation.

The \emph{genus-$g$ quantum $L_\infty$ operation} is
\begin{equation}\label{eq:quantum-linf-operation}
m_n^{(g)}(f_1, \ldots, f_n)
\;=\;
\sum_{L \geq 0}\;
\sum_{\substack{\Gamma \in \mathfrak{G}_{n,L} \\
 \sum_v g_v + L = g}}
\frac{1}{|\mathrm{Aut}(\Gamma)|}
\, w_\Gamma^{(L)}(f_1, \ldots, f_n).
\end{equation}
The sum runs over all stable graphs whose total genus
$\sum_v g_v + L$ equals~$g$.

For $L \geq 1$, the loop integrals require regularization: the
Szeg\H{o} propagator $S(z,w)$ has a pole along the diagonal, and
loop edges produce $\int S(z,z')\, dz$ with both variables
integrated. Convergence is ensured by the FM compactification
$\overline{C}_n(\Sigma)$, which resolves diagonal singularities
into normal-crossing divisors with logarithmic poles
(Proposition~\ref{prop:stokes-regularity-FM}); the
weight grading on~$\cA$ provides an additional finiteness
bound at each weight.
\end{definition}

\begin{theorem}[Quantum master equation;
\ClaimStatusConditional]\label{thm:quantum-master-equation}
\index{quantum master equation!from stable-graph operations}
\index{QME|see{quantum master equation}}
The genus-graded operations $\{m_n^{(g)}\}_{g,n}$ satisfy the
\emph{quantum master equation} (QME): for each $N \geq 0$ and
$G \geq 0$,
\begin{equation}\label{eq:QME}
\sum_{\substack{r + s + t = N \\ g_1 + g_2 = G}}
(-1)^{rs+t}\,
m_{r+1+t}^{(g_1)}\!\bigl(\mathrm{id}^{\otimes r}
\otimes m_s^{(g_2)} \otimes \mathrm{id}^{\otimes t}\bigr)
\;+\;
\sum_{i < j}
(-1)^{\epsilon_{ij}}
m_{N-1}^{(G-1)}\bigl(\Delta_{ij}\bigr)
\;=\; 0,
\end{equation}
where $\Delta_{ij}$ denotes the contraction of inputs $f_i, f_j$
using $\langle -, - \rangle_\cA$ (the self-sewing/non-separating
term), $\epsilon_{ij}$ is the Koszul sign, and the first sum is
the genus-graded $A_\infty$ composition relation (separating term).
\end{theorem}

\begin{proof}
The proof extends the Stokes argument of
Theorem~\ref{thm:cyclic-linf-graph} from trees to graphs with loops.

\emph{Step 1: Two sources of the QME.}
The QME has two structurally distinct contributions,
corresponding to the two types of codimension-$1$ boundary
strata of $\overline{\mathcal{M}}_{g,n}$
(Definition~\ref{def:modular-operad}):

\emph{Separating terms} (first sum in~\eqref{eq:QME}).
For each graph $\Gamma$ contributing at genus~$G$ and degree~$N$,
the graph amplitude integrates a form over
$\ConfigSpace{|V_{\mathrm{int}}|}(\Sigma)$.
Collision of internal vertices produces codimension-$1$ boundary
strata of the configuration space
(Proposition~\ref{prop:stokes-regularity-FM}).
By Stokes' theorem, these boundary contributions give
$m^{(g_1)} \circ m^{(g_2)}$ compositions with $g_1 + g_2 = G$,
exactly as in the tree-level proof of
Theorem~\ref{thm:cyclic-linf-graph}.

\emph{Non-separating terms} (second sum in~\eqref{eq:QME}).
These terms come from the \emph{graph structure}. Contracting a loop edge in
$\Gamma$ (the $d_{\mathrm{ns}}$ component of
Definition~\ref{def:three-component-modular-bar})
pairs two half-edges at a vertex via the invariant form
$\langle -, - \rangle_\cA$, producing the self-sewing contraction
$\Delta_{ij}$ at genus $G - 1$.
Analytically, this corresponds to the period integral
$\oint_\gamma S(z, w)\, dz$ along a homology cycle of $\Sigma$,
which evaluates to a scalar (the period) times the contraction.

\emph{Step 2: $\partial^2 = 0$ on moduli.}
Both contributions are components of the single identity
$d_{\mathrm{mod}}^2 = 0$
(Proposition~\ref{prop:d-mod-squared-zero}),
which is the modular operad boundary relation
$\partial^2 = 0$ on $\overline{\mathcal{M}}_{g,n}$
transported through the Feynman transform.
The separating terms correspond to
$d_{\mathrm{sep}}^2 + d_{\mathrm{int}} d_{\mathrm{sep}}
+ d_{\mathrm{sep}} d_{\mathrm{int}} = 0$;
the non-separating terms to
$d_{\mathrm{ns}}^2 + d_{\mathrm{int}} d_{\mathrm{ns}}
+ d_{\mathrm{ns}} d_{\mathrm{int}} = 0$;
the cross terms to
$d_{\mathrm{sep}} d_{\mathrm{ns}}
+ d_{\mathrm{ns}} d_{\mathrm{sep}} = 0$
(codimension-$2$ corner cancellation).
\end{proof}

\begin{remark}[The unimodularity obstruction]
\label{rem:unimodularity-obstruction}
\index{unimodularity obstruction|textbf}
\index{quantum L-infinity algebra@quantum $L_\infty$ algebra!unimodularity}
The quantum $L_\infty$ structure does not exist for all cyclic
$L_\infty$ algebras. The obstruction to promoting a cyclic
$L_\infty$ algebra to a quantum $L_\infty$ algebra (i.e., adding
the non-separating terms $m_n^{(g)}$ for $g \geq 1$) is the
\emph{unimodularity condition}
\begin{equation}\label{eq:unimodularity}
\Delta(m_0) = 0,
\end{equation}
where $\Delta\colon \operatorname{End}(\cA) \to \mathbb{C}$ is the
BV operator (supertrace of the second-order part).
Concretely: the non-separating self-sewing $\Delta_{ij}$ applied to
the curvature $m_0$ must vanish, or else the genus-$1$ QME has no
solution.

For Koszul chiral algebras, this obstruction is automatically
bypassed:
\begin{enumerate}[label=\textup{(\alph*)}]
\item Central curvature
 (Remark~\ref{rem:voa-central-curvature}):
 $m_0^{(g)} = \kappa(\cA) \cdot \omega_g \cdot \operatorname{id}$
 is a scalar multiple of the identity, hence
 $\Delta(m_0^{(g)}) = \kappa(\cA) \cdot \omega_g
 \cdot \operatorname{str}(\operatorname{id})$.
\item The supertrace is computed weight-by-weight.
 At each conformal weight~$h$, the bar complex
 $\barB_{[h]}$ is finite-dimensional and the Koszul
 reciprocity $h_{\cA}(t) \cdot h_{\cA^!}(-t) = 1$
 \cite[Theorem~3.5.1]{LV12} gives
 $\sum_n (-1)^n \dim \barB^n_{[h]} = 0$ for $h \geq 1$
 (the alternating sum vanishes at each positive weight by
 the functional equation). At weight~$0$,
 $\barB_{[0]} = \mathbb{k}$ contributes
 $\operatorname{str}(\operatorname{id}_{[0]}) = 1$, but
 the curvature $m_0^{(g)}$ has positive weight, so
 $\Delta(m_0^{(g)})_{[h]} = \kappa \cdot \omega_g \cdot
 \operatorname{str}(\operatorname{id}_{[h]}) = 0$
 for all $h$ in the support of~$m_0^{(g)}$.
 Hence $\Delta(m_0^{(g)}) = 0$.
 \item Therefore the named BV trace obstruction vanishes on the
 Koszul locus. A completed quantum $L_\infty$ structure, and hence a
 completed modular homotopy type
 (Definition~\ref{def:modular-homotopy-abstract}), requires the
 completed QME and the higher-loop obstruction package.
\end{enumerate}
By this mechanism, the universal MC
class~$\Theta_\cA$
(Theorem~\ref{thm:universal-MC}) exists for the Koszul locus.
\end{remark}

\begin{remark}[The QME as $\partial^2 = 0$ on moduli]
\label{rem:qme-boundary-squared}
\index{quantum master equation!geometric origin}
The QME~\eqref{eq:QME} has a single geometric origin:
$\partial^2 = 0$ on chains of $\overline{\mathcal{M}}_{g,n}$.
The two sums correspond to the two types of codimension-$1$
boundary strata:

\smallskip
\begin{center}
\renewcommand{\arraystretch}{1.3}
\begin{tabular}{lll}
\toprule
\emph{Boundary type} & \emph{Differential component}
 & \emph{QME term} \\
\midrule
Separating node & $d_{\mathrm{sep}}$
 & $\sum m^{(g_1)} \circ m^{(g_2)}$ \\
Non-separating node & $d_{\mathrm{ns}}$
 & $\sum m^{(G-1)}(\Delta_{ij})$ \\
\bottomrule
\end{tabular}
\end{center}
\smallskip

\noindent
The first sum splits the curve into two components of genera
$g_1 + g_2 = G$; the second sum creates a handle, lowering genus by~$1$.
The QME is $\partial^2 = 0$ transported through the Feynman transform.
In physics: the first term is the tree-level factorization of
amplitudes; the second is the BV operator $\Delta$ (the supertrace
over the propagator loop). The entire quantum $L_\infty$ structure
is the shadow of the modular operad axioms on
$\{\overline{\mathcal{M}}_{g,n}\}$.
\end{remark}

\subsubsection{The total differential from
\texorpdfstring{$\Theta_\cA$}{Theta-A}}
\label{subsubsec:total-differential-from-theta}
\index{total differential!from Maurer--Cartan class|textbf}
\index{completed bar complex!total differential}

The universal MC class $\Theta_\cA$ assembles the genus-graded bar
differentials into a single nilpotent on the completed bar
complex.

\begin{theorem}[Total differential from the MC class;
\ClaimStatusConditional]\label{thm:total-differential-from-mc}
Let $\cA$ be a modular Koszul chiral algebra for which
$\Theta_\cA \in
\operatorname{MC}(\operatorname{Def}^{\mathrm{mod}}(\cA))$
is constructed
\textup{(}Theorem~\textup{\ref{thm:universal-MC})}.
Define the \emph{total completed bar complex}
\begin{equation}\label{eq:completed-bar-total}
\barB^{\mathrm{tot}}_X(\cA)
\;:=\;
\prod_{g \geq 0} \hbar^{g}\,
\barB_X^{(g)}(\cA)
\end{equation}
with total differential
\begin{equation}\label{eq:D-total-from-theta}
D_{\mathrm{tot}}
\;:=\;
\dzero \;+\; [\Theta_\cA,\, -],
\end{equation}
where $\dzero$ is the genus-$0$ collision differential and
$[\Theta_\cA, -]$ is the coderivation action of $\Theta_\cA$ on
the bar coalgebra, expanded genus-by-genus as
\begin{equation}\label{eq:D-total-expanded}
D_{\mathrm{tot}}\big|_{\barB^{(g)}}
\;=\;
\dzero \;+\;
\sum_{h=1}^{g} \hbar^{h}\,
\bigl[\Theta_\cA^{(h)},\, -\bigr],
\end{equation}
where $\Theta_\cA^{(h)}$ is the genus-$h$ component of the MC
class.
Then:
\begin{enumerate}[label=\textup{(\roman*)}]
\item $D_{\mathrm{tot}}^2 = 0$ on
 $\barB^{\mathrm{tot}}_X(\cA)$.
\item At fixed genus~$g$, the restriction recovers the
 genus-$g$ total corrected differential:
 $D_{\mathrm{tot}}|_{\barB^{(g)}} = \Dg{g}$
 \textup{(}Convention~\textup{\ref{conv:higher-genus-differentials}(ii))}.
\item The fiberwise curved differential at genus~$g$ is
 $\dfib^{(g)} := D_{\mathrm{tot}}|_{\barB^{(g)}}
 \big|_{[\Sigma_g]}$
 (restriction to the fiber over $[\Sigma_g] \in
 \overline{\mathcal{M}}_g$, with the moduli-variation components
 $d_{\mathrm{sep}}$, $d_{\mathrm{ns}}$ of
 Definition~\textup{\ref{def:three-component-modular-bar}}
 dropped). It satisfies the curvature identity
 $(\dfib^{(g)})^2 = \kappa(\cA) \cdot \omega_g \cdot
 \operatorname{id}$.
\end{enumerate}
\end{theorem}

\begin{proof}
\emph{Part~(i).}
The MC equation for $\Theta_\cA$ in
$\operatorname{Def}^{\mathrm{mod}}(\cA)$
(Definition~\ref{def:modular-deformation-complex}) states
\begin{equation}\label{eq:mc-for-dtot}
\dzero(\Theta_\cA) + \tfrac{1}{2}[\Theta_\cA, \Theta_\cA] = 0.
\end{equation}
Compute:
\begin{align*}
D_{\mathrm{tot}}^2
&= (\dzero + [\Theta_\cA, -])^2 \\
&= \underbrace{\dzero^2}_{=\,0}
 + \dzero \circ [\Theta_\cA, -]
 + [\Theta_\cA, -] \circ \dzero
 + [\Theta_\cA, [\Theta_\cA, -]] \\
&= [\dzero(\Theta_\cA), -]
 + \tfrac{1}{2}[[\Theta_\cA, \Theta_\cA], -] \\
&= [\dzero(\Theta_\cA) + \tfrac{1}{2}[\Theta_\cA, \Theta_\cA],\, -]
\;=\; 0,
\end{align*}
where the second line uses the graded Jacobi identity for
coderivation brackets, the third line uses the coderivation identity
$\dzero \circ [\Theta, -] + [\Theta, -] \circ \dzero
= [\dzero(\Theta), -]$, and the final step is the
MC equation~\eqref{eq:mc-for-dtot}.

\emph{Part~(ii).}
By construction,
$D_{\mathrm{tot}}|_{\barB^{(g)}} = \dzero +
\sum_{h=1}^g \hbar^h [\Theta^{(h)}, -]$
(equation~\eqref{eq:D-total-expanded}).
The coderivation $[\Theta^{(h)}, -]$ acts on the bar coalgebra
as the genus-$h$ quantum correction to the differential,
which is precisely the period-integral correction $\hbar^h d_h$
of Convention~\ref{conv:higher-genus-differentials}(ii):
the MC class $\Theta^{(h)}$ encodes the same data as the
genus-$h$ propagator correction (by
Theorem~\ref{thm:universal-theta}).
Hence $D_{\mathrm{tot}}|_{\barB^{(g)}} = \Dg{g}$.

\emph{Part~(iii).}
On the fiber over $[\Sigma_g]$, the separating and non-separating
components $d_{\mathrm{sep}}$, $d_{\mathrm{ns}}$ vanish (they
require varying the curve), leaving
$\dfib^{(g)} = d_{\mathrm{int}}|_{[\Sigma_g]}$. Compute:
$(\dfib^{(g)})^2 = (d_{\mathrm{int}})^2|_{[\Sigma_g]}
= \sum_v d_v^2 = \sum_v \kappa(\cA) \cdot \omega_{g_v}$.
For a single-vertex graph ($g_v = g$, the smooth locus of
$\overline{\mathcal{M}}_g$), this gives $\kappa(\cA) \cdot
\omega_g$ directly.
On a multi-vertex graph, $\sum_v \omega_{g_v} = \omega_g$ by
additivity of the Arnold defect across irreducible components
of a nodal curve~\eqref{eq:arnold-defect-hg}.
\end{proof}

\begin{remark}[$D_{\mathrm{tot}}$ as twisted connection; $\Theta_\cA$ as curved twisting cochain]
\label{rem:D-tot-as-connection}%
\label{rem:theta-as-curved-twisting}%
\index{total differential!as twisted connection}
\index{twisting morphism!curved|textbf}
\index{Maurer--Cartan element!as curved twisting cochain}
Formula~\eqref{eq:D-total-from-theta} is a twisted connection:
$\dzero$ is the flat connection on the bar coalgebra,
$\Theta_\cA$ is a gauge field valued in the modular deformation
complex, and $D_{\mathrm{tot}} = \dzero + [\Theta_\cA, -]$ is
the covariant derivative.
The MC equation
$\dzero(\Theta) + \tfrac{1}{2}[\Theta, \Theta] = 0$
is the flatness condition $D_{\mathrm{tot}}^2 = 0$.

This has a precise formulation in the language of twisting cochains
(Remark~\ref{rem:theta-modular-twisting}):
the genus-$0$ canonical twisting morphism $\tau_0$ extends to a
\emph{modular twisting cochain}
$\tau_{\mathrm{mod}} = \tau_0 + \Theta_\cA$
in the modular convolution algebra, whose genus-$g$ curvature
$\kappa(\cA)\cdot\omega_g$ is the Arnold defect.
The cyclic $L_\infty$ operations $\{l_n\}$ on
$H^*_{\mathrm{cyc}}(\cA,\cA)$ are the Taylor coefficients of the
coderivation $[\Theta_\cA, -]$: the $L_\infty$ structure \emph{is}
the twisting cochain, read as an algebra structure on the
deformation complex.
\end{remark}

\begin{remark}[Clutching compatibility of $D_{\mathrm{tot}}$]
\label{rem:clutching-D-tot}
\index{clutching compatibility!total differential}
Since $\Theta_\cA$ is an MC element in the modular convolution
algebra $\widehat{L}_{\Gmod}$
(Proposition~\ref{prop:geometric-modular-operadic-mc}),
it satisfies clutching
compatibility~\eqref{eq:clutching-factorization}:
\[
\xi_{\mathrm{sep}}^*(\Theta^{(g)})
= \sum_{g_1+g_2=g} \Theta^{(g_1)} \star \Theta^{(g_2)},
\qquad
\xi_{\mathrm{ns}}^*(\Theta^{(g+1)})
= \Delta_{\mathrm{ns}}(\Theta^{(g)}).
\]
This means the total differential $D_{\mathrm{tot}}$ is
compatible with the modular operad structure: restricting to a
boundary stratum (separating or non-separating degeneration) gives
the tensor product (resp.\ self-sewing) of the differentials on the
two (resp.\ one) irreducible component(s). Equivalently,
$D_{\mathrm{tot}}$ is a \emph{modular bar differential} in the
sense of Definition~\ref{def:three-component-modular-bar}:
$D_{\mathrm{tot}} = d_{\mathrm{int}} + d_{\mathrm{sep}} +
d_{\mathrm{ns}}$ with the three components induced by~$\Theta_\cA$.
\end{remark}

\subsubsection{Virtual bar family and characteristic hierarchy}
\label{subsubsec:virtual-bar-k-class}
\index{virtual bar family|textbf}
\index{K-theory@$K$-theory!virtual bar class}

\begin{definition}[Virtual bar $K$-class]
\label{def:virtual-bar-k-class}
For a modular Koszul chiral algebra $\cA$, the
\emph{virtual bar family at genus~$g$} is the derived pushforward
\begin{equation}\label{eq:virtual-bar-k-class}
\mathbf{V}_{\cA}^{(g)}
\;:=\;
R\pi_{g*}\, \barB^{(g)}_X(\cA)
\;\in\; D^b(\overline{\mathcal{M}}_g),
\end{equation}
where $\pi_g\colon \overline{\mathcal{C}}_g \to
\overline{\mathcal{M}}_g$ is the universal curve. Its class in
$K$-theory is the \emph{virtual bar $K$-class}
\begin{equation}\label{eq:V-A-K-class}
V_{\cA}
\;:=\;
[R\pi_{g*}\, \barB^{(g)}_X(\cA)]
\;\in\; K^0(\overline{\mathcal{M}}_g).
\end{equation}
\end{definition}

\begin{proposition}[Characteristic hierarchy;
\ClaimStatusConditional]\label{prop:characteristic-hierarchy}
\index{characteristic hierarchy|textbf}
\index{modular characteristic package!K-theoretic}
The virtual bar $K$-class $V_\cA$ determines a three-level
hierarchy of invariants:
\begin{enumerate}[label=\textup{(L\arabic*)}]
\item \emph{Scalar level.}
 The first Chern class of the determinant line bundle recovers the
 scalar modular characteristic:
 \[
 c_1(\det V_\cA)
 \;=\;
 \kappa(\cA) \cdot \lambda
 \;\in\; \operatorname{Pic}(\overline{\mathcal{M}}_g),
 \]
 where $\lambda = c_1(\det \mathbb{E}^\vee)$ is the Hodge class.
 This is Theorem~\textup{\ref{thm:modular-characteristic}}:
 $\kappa(\cA)$ is the scalar modular characteristic.

\item \emph{Spectral level.}
 The Chern character of $V_\cA$ determines the spectral
 discriminant:
 \[
 \Delta_\cA(x)
 \;=\;
 \det(1 - x \cdot T_{\mathrm{rec}})
 \;=\;
 \prod_i (1 - \lambda_i x),
 \]
 where the $\lambda_i$ are the Chern roots of $V_\cA$
 and $T_{\mathrm{rec}}$ is the recursion operator
 \textup{(}Theorem~\textup{\ref{thm:discriminant-spectral})}.
 This encodes the full bar cohomology dimensions at all weights.

\item \emph{Holonomy level.}
 The flat connection on $V_\cA$ (from the Gauss--Manin connection
 on $R\pi_{g*}$) determines the periodicity profile
 \[
 \Pi_\cA \;=\; (M_\cA,\, Q_\cA,\, G_\cA)
 \;\;:\;\;
 \text{monodromy, quasiperiod, Galois data}
 \]
 \textup{(}Remark~\textup{\ref{rem:periodicity-triple})}.
 For rational central charge, the monodromy group
 $M_\cA \leq \mathrm{GL}(V_\cA)$ has finite image
 (the $T$-matrix has finite order on the character space).
\end{enumerate}
The three levels correspond to $\det$, $\mathrm{ch}$, and
$\mathrm{Hol}$ respectively: the standard filtration of
information in a $K$-class with connection.
\end{proposition}

\begin{proof}
\emph{(L1)}\enspace
By Grothendieck--Riemann--Roch on the universal curve
$\pi_g\colon \overline{\mathcal{C}}_g \to \overline{\mathcal{M}}_g$
(\S\ref{subsec:why-ahat}):
\[
c_1(\det R\pi_{g*} \barB^{(g)})
\;=\;
\pi_{g*}\!\bigl(\mathrm{ch}_1(\barB|_{\mathrm{fiber}})
\cdot \mathrm{Td}(T_\pi)\bigr).
\]
The fiber Chern character $\mathrm{ch}_1(\barB|_{\Sigma_g})
= \kappa(\cA)$ (Theorem~\ref{thm:modular-characteristic}),
and $\pi_{g*}(\mathrm{Td}_1(T_\pi)) = \lambda$
(Mumford's formula \cite{Mumford83}).
Hence $c_1(\det V_\cA) = \kappa(\cA) \cdot \lambda$.

\emph{(L2)}\enspace
The full Chern character $\mathrm{ch}(V_\cA)$
encodes the dimensions $\dim H^n(\barB^{(g)}_{[h]})$ at each
weight~$h$ and bar degree~$n$, which are the coefficients of the
spectral discriminant generating function
(Theorem~\ref{thm:ds-bar-gf-discriminant}).

\emph{(L3)}\enspace
The Gauss--Manin connection on
$\barB^{(g)}_X(\cA) \to \overline{\mathcal{M}}_g$
is flat (its curvature is the Kodaira--Spencer class, which acts
nilpotently on cohomology by
Theorem~\ref{thm:kodaira-spencer-chiral-complete}).
The monodromy representation records the action
of $\pi_1(\mathcal{M}_g) = \mathrm{MCG}(\Sigma_g)$ on the bar
cohomology fibers.
\end{proof}

\begin{remark}[Three readings of the virtual bar class]
\label{rem:three-readings-V}
\index{virtual bar family!three readings}
The hierarchy (L1)--(L3) reflects the three standard ways of
extracting invariants from a $K$-class: Euler characteristic,
Chern character, holonomy. In the categorical logarithm paradigm:
\begin{itemize}
\item \emph{(L1)} $\det V_\cA = \kappa \cdot \lambda$
 is the \emph{leading coefficient} of the logarithm
 (Theorem~D$_{\mathrm{scal}}$).
\item \emph{(L2)} $\mathrm{ch}(V_\cA) \leftrightarrow \Delta_\cA(x)$
 is the \emph{Taylor expansion} of the logarithm
 (Theorem~D$_\Delta$).
\item \emph{(L3)} $\mathrm{Hol}(V_\cA) = \Pi_\cA$
 is the \emph{monodromy} of the logarithm, the ambiguity in
 choosing a branch (the periodicity profile).
\end{itemize}
This realizes the passage
$\kappa \to \Delta \to \Pi$
(\emph{scalar $\to$ spectral $\to$ holonomy}) as the natural
filtration of information in the $K$-class $V_\cA$, just as the
passage $c_1 \to \mathrm{ch} \to \nabla$ is the natural filtration
of information in a vector bundle with connection. The Heisenberg
specialization
(Chapter~\ref{ch:heisenberg-frame},
 Definition~\ref{def:frame-modular-package})
illustrates all three levels: $\kappa = k$ (the level),
$\Delta = 1/(1-x)$ (free boson),
$\Pi = (\mathbb{Z}, k\mathbb{Z}, \{1\})$ (abelian periodicity).
The precise minimal periodicity of the bar cohomology dimensions is
still open; the proved hierarchy records only the finite-image
monodromy consequence at rational central charge.
\end{remark}

\subsubsection{Chiral homology from the modular bar complex}
\label{subsubsec:chiral-homology-recovery}
\index{chiral homology!from modular bar complex|textbf}

\begin{theorem}[Chiral homology recovery;
\ClaimStatusConditional]\label{thm:chiral-homology-recovery}
\index{chiral homology!bar complex identification}
Let $\cA$ be a modular Koszul chiral algebra on a smooth projective
curve $X$ of genus~$g$. Then:
\begin{enumerate}[label=\textup{(\roman*)}]
\item \emph{Fiberwise identification.}
 The chiral homology of~$\cA$ on the fixed curve $\Sigma_g$ is
 computed by the bar complex:
 \begin{equation}\label{eq:chiral-homology-bar}
 H^{\mathrm{ch}}_n(\Sigma_g, \cA)
 \;\cong\;
 H^n\bigl(\barB^{(g)}_X(\cA)\big|_{[\Sigma_g]},\;
 \Dg{g}\big|_{[\Sigma_g]}\bigr)
 \end{equation}
 for all $n \geq 0$, where $\Dg{g}\big|_{[\Sigma_g]}$ is the
 square-zero flat comparison differential at
 $[\Sigma_g] \in \overline{\mathcal{M}}_g$. When
 $\kappa(\cA) \neq 0$, the curved fiber operator $\dfib^{(\Sigma_g)}$
 is interpreted only coderivedly.

\item \emph{Center sheaf.}
 The local system on $\overline{\mathcal{M}}_g$ whose fiber
 over $[\Sigma_g]$ is $H^{\mathrm{ch}}_*(\Sigma_g, \cA)$ is the
 center sheaf:
 \begin{equation}\label{eq:center-sheaf-bar}
 \mathcal{Z}_\cA
 \;=\;
 R^*\!\pi_{g*}\,
 \bigl(\barB^{(g)}_X(\cA),\, \Dg{g}\bigr)
 \;\in\;
 \operatorname{Loc}(\overline{\mathcal{M}}_g).
 \end{equation}

\item \emph{Global sections.}
 The ambient complementarity complex
 \textup{(}Theorem~\textup{\ref{thm:quantum-complementarity-main})}
 is recovered as
 \begin{equation}\label{eq:ambient-from-bar}
 C_g(\cA) \;=\; R\Gamma(\overline{\mathcal{M}}_g,\,
 \mathcal{Z}_\cA)
 \;\cong\;
 H^*\bigl(\barB^{(g)}_X(\cA),\, D_{\mathrm{tot}}\bigr),
 \end{equation}
 and the Verdier involution $\sigma$ on $C_g(\cA)$ arises from
 Verdier duality $\mathbb{D}_{\operatorname{Ran}}$ on the bar
 coalgebra
 \textup{(}Theorem~\textup{\ref{thm:higher-genus-inversion})}.

\item \emph{Genus tower.}
 The full genus tower of chiral homology is assembled by the
 completed bar complex:
 \begin{equation}\label{eq:chiral-homology-tower}
 \bigoplus_{g \geq 0} \hbar^{g}\,
 H^{\mathrm{ch}}_*(\Sigma_g, \cA)
 \;\cong\;
 H^*\bigl(\barB^{\mathrm{tot}}_X(\cA),\, D_{\mathrm{tot}}\bigr),
 \end{equation}
 with the genus filtration corresponding to the $\hbar$-adic
 filtration on the completed bar.
\end{enumerate}
\end{theorem}

\begin{proof}
\emph{Part~(i).}
On a fixed curve $\Sigma_g$, the chiral homology
$H^{\mathrm{ch}}_*(\Sigma_g, \cA)$ is by definition
\cite[3.4.4]{BD04} the derived pushforward of $\cA^{\boxtimes n}$
along the Ran space $\operatorname{Ran}(\Sigma_g)$. The bar
construction computes this pushforward
(Theorem~\ref{thm:geometric-equals-operadic-bar} for genus~$0$;
Theorem~\ref{thm:prism-higher-genus} for the extension to all
genera via the Feynman transform): the bar complex on
$\overline{C}_n(\Sigma_g)$ with the flat comparison differential
$\Dg{g}\big|_{[\Sigma_g]}$ computes the same ordinary cohomology as
the Ran-space pushforward. If $\kappa(\cA) \neq 0$, the raw fiberwise
operator $\dfib^{(\Sigma_g)}$ is curved and therefore belongs on the
coderived side rather than in ordinary cohomology.

\emph{Part~(ii).}
The assignment $[\Sigma_g] \mapsto
H^{\mathrm{ch}}_*(\Sigma_g, \cA)$ varies in a local system because
the flat comparison bar complex is a square-zero family over
$\overline{\mathcal{M}}_g$
(the bar family of~\eqref{eq:bar-family}), and the formation of
cohomology of a flat family produces a local system. The resulting
local system is precisely the center sheaf $\mathcal{Z}_\cA$ of
Beilinson--Drinfeld~\cite[3.4.10]{BD04}.

\emph{Part~(iii).}
Global sections of $\mathcal{Z}_\cA$ over
$\overline{\mathcal{M}}_g$ compute
$R\Gamma(\overline{\mathcal{M}}_g, \mathcal{Z}_\cA)$
by definition.
By the Leray spectral sequence for
$\pi_g\colon \overline{\mathcal{C}}_g \to \overline{\mathcal{M}}_g$,
this agrees with
$H^*(\barB^{(g)}_X(\cA), D_{\mathrm{tot}})$ when the total
differential is used (incorporating both the fiberwise part and
the moduli variation). The Verdier involution on $C_g(\cA)$
arises from the Verdier duality
$\mathbb{D}_{\operatorname{Ran}} \barB_X(\cA) \simeq
\barB_X(\cA^!)$
(Theorem~\ref{thm:higher-genus-inversion}), which exchanges
$\cA$ and~$\cA^!$ on the bar level and induces $\sigma$ on
global sections.

\emph{Part~(iv).}
The genus-completed product~\eqref{eq:completed-bar-total} is
an inverse limit $\barB^{\mathrm{tot}} = \varprojlim_g
\bigoplus_{h \leq g} \hbar^h \barB^{(h)}$. Each finite
truncation has well-defined cohomology by parts~(i)--(iii).
The transition maps are surjective on cohomology (the genus
filtration is separated and complete by
Proposition~\ref{prop:geometric-modular-operadic-mc}(a)),
so the Mittag-Leffler condition is satisfied and
$H^*(\barB^{\mathrm{tot}}, D_{\mathrm{tot}})
\cong \varprojlim_g H^*(\bigoplus_{h \leq g}
\barB^{(h)}, D_{\mathrm{tot}})$
by \cite[Proposition~3.5.7]{Weibel94}.
Theorem~\ref{thm:total-differential-from-mc}(i) gives
$D_{\mathrm{tot}}^2 = 0$.
\end{proof}

\begin{remark}[The full recovery diagram]
\label{rem:full-recovery-diagram}
\index{modular homotopy theory!recovery diagram}
The constructions of this subsection fit into a single commutative
diagram:
\[
\begin{tikzcd}[column sep=large, row sep=large]
\text{Quantum } L_\infty
 \arrow[r, "\text{MC}"]
 \arrow[d, "g{=}0"']
& \Theta_\cA \in
 \operatorname{MC}(\operatorname{Def}^{\mathrm{mod}})
 \arrow[r, "\text{twist}"]
 \arrow[d, "\text{clutch}"]
& (\barB^{\mathrm{tot}},\, D_{\mathrm{tot}})
 \arrow[d, "R\pi_{g*}"]
\\
\text{Cyclic } L_\infty
 \arrow[r, "\text{Kill.}"]
& \xi^*(\Theta^{(g)})
& \mathcal{Z}_\cA
 \arrow[r, "R\Gamma"]
& C_g = Q_g \oplus Q_g^!
\end{tikzcd}
\]

\noindent
Three paths through this diagram recover three aspects of the
modular homotopy type:

\smallskip
\emph{Top row} (``algebraic path''): the quantum $L_\infty$
operations, packaged as the MC element $\Theta_\cA$, twist the
genus-$0$ bar differential to produce $D_{\mathrm{tot}}$. This is
the bar-coalgebra route: the modular homotopy type is the
\emph{totalization} of the genus tower.

\smallskip
\emph{Right column} (``geometric path''): the total bar complex
pushes forward along the universal curve to produce the center
sheaf $\mathcal{Z}_\cA$; its global sections are the
complementarity complex $C_g(\cA) = Q_g \oplus Q_g^!$. This
is the sheaf-on-moduli route: the modular homotopy type is the
\emph{variation} of the bar complex over
$\overline{\mathcal{M}}_g$.

\smallskip
\emph{Bottom row} (``clutching path''): the genus-$0$ cyclic
$L_\infty$ produces the tree-level structure; the clutching
maps $\xi_{\mathrm{sep}}^*$, $\xi_{\mathrm{ns}}^*$ reconstruct
the full center sheaf from its boundary behavior: the
Getzler--Kapranov route, where the modular homotopy type is the
\emph{modular operad structure} on the collection
$\{H^{\mathrm{ch}}_*(\Sigma_g, \cA)\}$.

\smallskip\noindent
All three paths commute. The single geometric fact underlying all of
them is $\partial^2 = 0$ on $\overline{\mathcal{M}}_{g,n}$, read
through the Feynman transform.
\end{remark}

\subsection{Non-scalar universal classes}
\label{subsec:non-scalar-theta}
\index{Theta_A@$\Theta_\cA$!non-scalar|textbf}

Corollary~\ref{cor:scalar-saturation} gives only one-channel
line concentration in general:
\[
\Theta_\cA^{\min} = \eta \otimes \Gamma_{\cA}
\qquad\text{whenever } \dim H^2_{\mathrm{cyc}}(\cA,\cA) = 1.
\]
On the proved scalar locus this further specializes to
\[
\Theta_\cA^{\min} = \kappa(\cA) \cdot \eta \otimes \Lambda
\qquad\text{with } \Gamma_{\cA}=\kappa(\cA)\Lambda.
\]
We construct the universal non-scalar class
for algebras with
$\dim H^2_{\mathrm{cyc}}(\cA,\cA) \geq 2$.

\begin{construction}[HPT non-scalar MC recursion; \ClaimStatusDefinitional]
\label{constr:non-scalar-hpt}
\index{homotopy perturbation theory!non-scalar MC}
Let $L = \Defcyc(\cA)$ be the cyclic deformation complex
with homotopy transfer data $(i,p,h,\ell_n)$ from a
minimal model. Choose a basis
$\eta_1, \ldots, \eta_r \in H^2_{\mathrm{cyc}}(\cA,\cA)$
and modular coefficients $U_1, \ldots, U_r \in \mathfrak{m}_{R}$
(the maximal ideal of the modular deformation ring).
The universal non-scalar MC element is built recursively:
\begin{equation}\label{eq:hpt-nonscalar-recursion}
\Theta^{(1)} = i\!\Bigl(\sum_{i=1}^r \eta_i \otimes U_i\Bigr),
\qquad
\Theta^{(N)} = -h\!\Biggl(
 \sum_{n \geq 2} \frac{1}{n!}
 \sum_{\substack{N_1 + \cdots + N_n = N \\ N_j \geq 1}}
 \ell_n\bigl(\Theta^{(N_1)}, \ldots, \Theta^{(N_n)}\bigr)
\Biggr).
\end{equation}
The obstruction equations cut out the modular deformation locus:
\[
p\!\Biggl(
 \sum_{n \geq 2} \frac{1}{n!}\,
 \ell_n(\Theta, \ldots, \Theta)
\Biggr) = 0.
\]
\end{construction}

\begin{example}[Two-channel explicit formula]
\label{ex:two-channel-nonscalar}
\index{Theta_A@$\Theta_\cA$!two-channel formula}
For $r = 2$, with basis $\eta_1, \eta_2$, the first
terms of the non-scalar universal class are:
\begin{equation}\label{eq:two-channel-explicit}
\begin{split}
\Theta_\cA^{\mathrm{ns}}
&= \eta_1 \otimes U_1
 + \eta_2 \otimes U_2
 - \tfrac{1}{2}\,h\ell_2(\eta_1,\eta_2)
 \otimes U_1 U_2 \\
&\quad - \tfrac{1}{6}\,h\ell_3(\eta_1,\eta_1,\eta_2)
 \otimes U_1^2 U_2
 - \tfrac{1}{6}\,h\ell_3(\eta_1,\eta_2,\eta_2)
 \otimes U_1 U_2^2
 + O(4).
\end{split}
\end{equation}
The class is genuinely non-scalar if and only if some mixed
coefficient $U_i U_j$ survives in gauge-normal form.
\end{example}

\begin{proposition}[Non-scalar criterion;
\ClaimStatusProvedHere]
\label{prop:non-scalar-criterion}
\index{Theta_A@$\Theta_\cA$!non-scalar criterion}
The universal class $\Theta_\cA$ is genuinely non-scalar
if and only if, in minimal cyclic model form, some mixed
transferred bracket
$\ell_n(\eta_{i_1}, \ldots, \eta_{i_n})$ with at least two
distinct basis indices produces a homotopy correction not
removable by gauge transformation.
Equivalently: after reduction to gauge-normal form,
some mixed monomial
$U_{i_1} \cdots U_{i_k}$ with $|\{i_1, \ldots, i_k\}| \geq 2$
has nonzero coefficient.
\end{proposition}

\begin{proof}
The gauge group acts by
$\Theta \mapsto e^{g} \cdot \Theta :=
\sum_{n=0}^\infty \frac{1}{n!}\, \ell_n(g, \ldots, g, \Theta)$
for $g \in L^1$. In gauge-normal form
$\operatorname{Im}(h) \cap \operatorname{Im}(\Theta) = 0$,
the recursion~\eqref{eq:hpt-nonscalar-recursion} has no
further gauge freedom. A non-scalar monomial $U_i U_j$
in~$\Theta^{(2)}$ arises from
$h\ell_2(\eta_i, \eta_j) \neq 0$ in the homotopy-transfer
data. Since $h$ is a homotopy (not a chain map), this
correction is algebraically independent of the scalar terms.
\end{proof}

\begin{remark}[Non-scalar landscape]
\label{rem:non-scalar-landscape}
\index{Theta_A@$\Theta_\cA$!non-scalar landscape}
All accessible families considered here (Kac--Moody at generic
level, Virasoro, principal $\mathcal{W}_N$) have
$\dim H^2_{\mathrm{cyc}}(\cA,\cA) = 1$, so their minimal
Maurer--Cartan class lies on a single cyclic line. On the proved
uniform-weight scalar locus this line is
$\kappa(\cA)\Lambda$; for principal $\mathcal{W}_N$ with $N\geq 3$
that stronger identification remains open. Non-scalar
$\Theta_\cA$ requires $\dim H^2_{\mathrm{cyc}} \geq 2$; test loci
include admissible-level Kac--Moody, products, and extended
algebras.
\end{remark}

\section{Stratified periodicity}

\subsection{Overview: three sources of periodicity}

The Hochschild cohomology of chiral algebras can exhibit three types of periodicity: \emph{Type I (Modular)}, from rational central charge and modular transformations; \emph{Type II (Quantum)}, from quantum groups at roots of unity; and \emph{Type III (Geometric)}, from the topology of the underlying curve.

\begin{remark}\label{rem:periodicity-scope}
The unconditional core is the root-of-unity calculation inside the
quantum-group realization and the structural theorem that independent
periodicity sources combine by lcm once their transport hypotheses are
met.
Explicit modular periods, sharp geometric periods, and higher-rank
closed formulas are conjectural outside the rank-$1$/minimal-model test
cases.
\end{remark}

\subsection{Type I: modular periodicity from rational conformal weights}

\subsubsection{The mechanism}

\begin{remark}[Periodicity as exponential of the scalar characteristic]
\label{rem:periodicity-as-exponential}
\index{nilpotence-periodicity correspondence!Type I application}
The scalar characteristic $\kappa(\cA)$ is the infinitesimal datum.
For $c=p'/q'$, the integer
$N_c=24q'/\gcd(p',24)$ kills the vacuum phase
$\exp(-2\pi i c/24)$. The full modular $T$-order is instead
\[
N_T(\cA)=\operatorname{lcm}_{i}\operatorname{den}
\bigl(h_i-c/24\bigr),
\]
where $h_i$ runs over the conformal weights of the simple modules and
the denominator is taken in lowest terms. Thus rational central charge
does not determine the full $T$-matrix period; the conformal weights are
part of the datum.
\end{remark}

\begin{conjecture}[Modular periodicity for minimal models; \ClaimStatusConjectured]
\label{conj:modular-periodicity-minimal}
\index{periodicity!modular!minimal models|textbf}
Let $M(p,q)$ be a Virasoro minimal model with $2 \leq p < q$, $\gcd(p,q) = 1$,
and central charge $c = 1 - 6(p-q)^2/(pq)$. Set
\[
N_T(M(p,q))
= \operatorname{lcm}_{1\le r\le p-1,\;1\le s\le q-1}
\operatorname{den}\bigl(h_{r,s}-c/24\bigr),
\qquad
h_{r,s}=\frac{(qr-ps)^2-(q-p)^2}{4pq}.
\]
Then for all bar degrees $n$ and all sufficiently large conformal weight $h$:
\[
\dim H^n(\bar{B}_{[h+N_T]}^{\mathrm{ch}}(M(p,q)))
\;=\;
\dim H^n(\bar{B}_{[h]}^{\mathrm{ch}}(M(p,q))).
\]
\end{conjecture}

\begin{evidence}
The following calculation separates the unconditional modular phase
statement from the unproved bar-periodicity assertion.  The convolution
with $1/\eta$ does not preserve eventual periodicity of Fourier
coefficients; Computation~\ref{comp:m58-weight-space-periodicity}
exhibits this obstruction for $M(5,8)$. The $T$-matrix periodicity is
unconditional.

\smallskip\noindent
The remaining implication would require periodic OPE matrices after
passing to the chiral bar differential.

\emph{Step~1: $T$-matrix periodicity.}
The minimal model $M(p,q)$ has exactly $(p-1)(q-1)/2$ irreducible modules
$V_{r,s}$ with $1 \leq r \leq p-1$, $1 \leq s \leq q-1$, subject to the
identification $V_{r,s} = V_{p-r,q-s}$. Their conformal weights are
\[
h_{r,s} \;=\; \frac{(qr - ps)^2 - (q-p)^2}{4pq}.
\]
The $T$-matrix is diagonal: $T_{(r,s),(r',s')} = \delta_{(r,s),(r',s')}
\, e^{2\pi i(h_{r,s} - c/24)}$. Write
\[
\lambda_{r,s} := e^{2\pi i(h_{r,s} - c/24)}.
\]
Then $\lambda_{r,s}^N = 1$ is equivalent to
$N(h_{r,s} - c/24) \in \mathbb{Z}$. Since $h_{r,s} - c/24$ is rational
with denominator dividing $24pq$, the integer $N_T(M(p,q))$ defined
above gives $T^{N_T}=\mathrm{Id}$ on the finite-dimensional character
space.

$\tau \mapsto \tau + 1$ multiplies the character by the phase
$e^{-2\pi i c/24}$. The smaller integer
$N_c=24q'/\gcd(p',24)$, for $c=p'/q'$ in lowest terms, kills this
vacuum phase:
\[
\frac{cN_c}{24} = \frac{p'N_c}{24q'} = \frac{p'}{\gcd(p',24)} \in \mathbb{Z},
\]
but $N_c$ need not be the full $T$-matrix order.

\emph{Step~2: Periodicity of weight-space dimensions.}
The character of each irreducible module $V_{r,s}$ is given by the
Rocha-Caridi formula:
\[
\mathrm{ch}_{V_{r,s}}(\tau) \;=\; \frac{\Theta_{r,s}(\tau)}{\eta(\tau)}
\]
where $\Theta_{r,s} = \theta_{qr-ps,\,2pq} - \theta_{qr+ps,\,2pq}$ is a
difference of level-$2pq$ theta functions and $\eta$ is the Dedekind eta
function. The Fourier coefficients of the level-$\ell$ theta function
$\theta_{m,\ell}(\tau) = \sum_{n \in \mathbb{Z}} q^{(m + 2\ell n)^2/(4\ell)}$
are supported on the arithmetic progression $\{(m + 2\ell n)^2/(4\ell) : n \in \mathbb{Z}\}$
and are periodic only on the corresponding finite discriminant group.

Now $\dim V_{r,s,\,[k]}$ is the coefficient of $q^{h_{r,s} + k - c/24}$ in
$\mathrm{ch}_{V_{r,s}}$. This coefficient equals the convolution of the
theta-function coefficients with the partition function $p(\cdot)$
(from $1/\eta$). The theta-function numerator $\Theta_{r,s}$ has exact
period dividing $2pq$ on the discriminant group. This numerator
period is separate from both the vacuum phase period $N_c$ and the
full $T$-matrix order~$N_T$.

This is the end of the unconditional argument. Multiplication by
$1/\eta$ convolves the theta numerator with the partition function and
does not preserve eventual periodicity of coefficients. Thus
$T^{N_T}=\mathrm{id}$ on characters is a statement about phases in a finite
modular representation, not a periodicity theorem for conformal-weight
spaces.

\begin{remark}[Denominator obstruction to stabilization]
\label{rem:step2-stabilization-threshold}
The convolution $\dim V_{r,s,[k]} = \sum_j a_\Theta(k-j)\, p(j)$ does not preserve eventual periodicity because $p(j)$ grows sub-exponentially.

\begin{computation}[M(5,8) weight-space dimensions; \ClaimStatusProvedHere]
\label{comp:m58-weight-space-periodicity}
For $M(5,8)$ ($c = -7/20$, $14$ modules, $N_T = 480$), the Rocha-Caridi formula
$\dim V_{r,s,[k]} = \sum_{n \in \mathbb{Z}} [p(k - \delta^+_n) - p(k - \delta^-_n)]$
(where $\delta^\pm_n$ are the standard null-vector exponents)
gives monotonically increasing sequences for all $14$ modules through $k = 80$. No eventual periodicity is detected at any period dividing $N_T = 480$ or $pq = 40$.
\end{computation}

Bar \emph{cohomology} periodicity is a separate bar-differential
question and requires computing OPE structure constants of $M(5,8)$.
\end{remark}

\emph{Step~3: Conditional lift to bar cohomology.}
The bar complex $\bar{B}^{(n)}_{[h]}(M(p,q))$ in degree~$n$ at conformal
weight~$h$ has dimension given by a finite sum over module labels weighted
by fusion multiplicities:
\[
\dim \bar{B}^{(n)}_{[h]} \;=\; \sum_{\substack{(r_1,s_1),\ldots,(r_n,s_n) \\
k_1 + \cdots + k_n = h - \sum_j h_{r_j,s_j}}}
\prod_{j=1}^n \dim V_{r_j,s_j,\,[k_j]}
\cdot F(r_1 s_1, \ldots, r_n s_n)
\]
where $F$ encodes the fusion multiplicities from the Verlinde formula
(depending only on module labels, not on~$h$), and the sum is over
non-negative integers~$k_j$. If one is given, for all sufficiently
large~$k$, linear isomorphisms
\[
\iota_{r,s,k}\colon V_{r,s,[k]}\xrightarrow{\sim}V_{r,s,[k+N_T]}
\]
which identify the OPE matrices entering the chiral bar differential,
then tensoring the maps~$\iota_{r,s,k}$ gives a chain isomorphism
\[
\bar B^{(n)}_{[h]}(M(p,q))\xrightarrow{\sim}
\bar B^{(n)}_{[h+N_T]}(M(p,q))
\qquad(h\gg0).
\]
Consequently
\[
\dim H^n(\bar{B}_{[h+N_T]}^{\mathrm{ch}})
= \dim H^n(\bar{B}_{[h]}^{\mathrm{ch}})
\]
in that range. The conjecture is precisely the assertion that such
eventual bar-differential identifications exist; they are not supplied
by the modular $T$-matrix alone.
\end{evidence}

\begin{conjecture}[Modular periodicity for general rational chiral algebras; \ClaimStatusConjectured]
\label{conj:modular-periodicity}
\index{periodicity!modular|textbf}
Let $\cA$ be a rational chiral algebra with central charge $c$ and
simple-module conformal weights $\{h_i\}$. Set
\[
N_T(\cA)=\operatorname{lcm}_i\operatorname{den}(h_i-c/24).
\]
Then for all bar degrees $n$ and all sufficiently large conformal weight $h$:
\[
\dim H^n(\bar{B}_{[h+N_T]}^{\mathrm{ch}}(\cA))
\;=\;
\dim H^n(\bar{B}_{[h]}^{\mathrm{ch}}(\cA)).
\]
\end{conjecture}

\begin{remark}[Scope and examples]\label{rem:modular-periodicity-scope}
For minimal models and WZW, the finite-order $T$-matrix is proved.
The obstruction is the passage from modular phases to eventual
bar-differential periodicity after the $1/\eta$ denominator is inserted.
For example, Ising has $N_T=48$, while
$\widehat{\mathfrak{sl}}_2$ at $k=1$ has $N_T=24$.
\end{remark}

\begin{conjecture}[Modular periodicity for WZW models; \ClaimStatusConjectured]
\label{conj:modular-periodicity-wzw}
\index{periodicity!modular!WZW models|textbf}
\index{WZW model!modular periodicity|textbf}
Let $\widehat{\mathfrak{g}}_k$ be an affine Kac--Moody algebra at
positive integer level $k \geq 1$, with $\mathfrak{g}$ a finite-dimensional
simple Lie algebra of rank~$r$. For an integrable highest weight
$\lambda\in P_k^+$ put
$h_\lambda=\langle\lambda,\lambda+2\rho\rangle/2(k+h^\vee)$ and set
\[
N_T(\widehat{\mathfrak g}_k)
=\operatorname{lcm}_{\lambda\in P_k^+}
\operatorname{den}(h_\lambda-c/24),
\qquad
c=\frac{k\dim(\mathfrak g)}{k+h^\vee}.
\]
Then for all bar degrees~$n$ and all sufficiently large conformal
weight~$h$:
\[
\dim H^n(\bar{B}_{[h+N_T]}^{\mathrm{ch}}(\widehat{\mathfrak{g}}_k))
\;=\;
\dim H^n(\bar{B}_{[h]}^{\mathrm{ch}}(\widehat{\mathfrak{g}}_k)).
\]
\end{conjecture}

\begin{evidence}
The Weyl--Kac character formula (Kac~\cite[\S12.7]{Kac}) gives the
finite-order $T$-matrix: the conformal weights
$h_\lambda=\langle\lambda,\lambda+2\rho\rangle/2(k+h^\vee)$ are rational,
so a common multiple of their denominators kills the $T$-phases. The
Weyl--Kac numerator is theta-periodic on the affine discriminant group,
but the affine denominator again inserts a partition-function
convolution. A root-of-unity tensor-category comparison can control the
fusion coefficients; it does not by itself identify the chiral bar
differential in successive conformal weights. Thus the WZW conjecture
has the same unresolved bar-differential condition as the minimal-model
case.
\end{evidence}

\begin{remark}[Scope]\label{rem:wzw-periodicity-scope}
$T$-matrix periodicity (Step~1) is unconditional for both minimal
models and WZW. The denominator convolution with $1/\eta$ does not
preserve eventual periodicity. General rational chiral algebras remain
open.
\end{remark}

\subsubsection{Koszul dual behavior}

\begin{conjecture}[Reflected vacuum-phase periodicity; \ClaimStatusConjectured]\label{conj:reflected-modular-periodicity}
\index{reflected modular periodicity}
If $\mathcal{A}$ has rational central charge $c = p/q$ and its
Koszul dual $\mathcal{A}^!$ has central charge $c' = p'/q'$, define
the vacuum phase periods
\[
N_c=\frac{24q}{\gcd(p,24)},\qquad
N_{c'}=\frac{24q'}{\gcd(p',24)}.
\]
Then
\begin{equation}\label{eq:period-sum}
\frac{1}{N_c} + \frac{1}{N_{c'}}
= \frac{\gcd(p,24)}{24q}
 + \frac{\gcd(p',24)}{24q'}.
\end{equation}
In the Virasoro same-family case $c+c'=26$, one has
$c'=(26q-p)/q$ and $q'=q$ after reduction. Then
\[
\frac{1}{N_c}+\frac{1}{N_{c'}}
=\frac{\gcd(p,24)+\gcd(26q-p,24)}{24q}.
\]
Thus $\frac{1}{N_c} + \frac{1}{N_{c'}} = \frac{1}{12}$ in this case if and
only if $\gcd(p,24)+\gcd(26q-p,24)=2q$.
The full $T$-matrix orders $N_T(\cA)$ and $N_T(\cA^!)$ also depend on
the conformal weights; no reciprocal formula for them is asserted here.
\end{conjecture}

\begin{remark}[Precise hypotheses]
\label{rem:precise-reflected-hypotheses}
The conjecture requires the following hypotheses to be well-posed:
\begin{enumerate}
\item $\cA$ is a rational chiral algebra (finitely many irreducible modules) with $c = p/q \in \mathbb{Q}$, $\gcd(p,q) = 1$.
\item $\cA$ admits a Koszul dual $\cA^!$ in the sense of Definition~\ref{def:chiral-koszul-pair}, with $\cA^!$ also rational.
\item The vacuum phase periods $N_c$, $N_{c'}$ kill
 $\exp(-2\pi i c/24)$ and $\exp(-2\pi i c'/24)$ respectively.
 The full $T$-matrix orders are the lcm of the denominators
 $\operatorname{den}(h_i-c/24)$ and
 $\operatorname{den}(h'_j-c'/24)$.
\item When $\cA$ arises from an affine Lie algebra at level $k$, the dual central charge $c' = c(\cA^!)$ is determined from the dual level $k' = -k - 2h^\vee$ via the Sugawara formula $c' = k'\dim(\mathfrak{g})/(k' + h^\vee)$.
\end{enumerate}
The relation~\eqref{eq:period-sum} holds for the Heisenberg algebra
($c=c'=1$, $N_c=N_{c'}=24$, sum $=1/12$). For
$\widehat{\mathfrak{sl}}_2$ at level $k=1$, affine level duality gives
$k'=-5$ and hence $c'=5$; both vacuum phase periods are~$24$.
For a Virasoro same-family pair $c=5$, $c'=21$, one gets
$N_c=24$, $N_{c'}=8$, and
$1/N_c+1/N_{c'}=1/6\ne1/12$.
\end{remark}

\begin{remark}[Scope]
The formula~\eqref{eq:period-sum} is a vacuum-phase identity once
$c'$ is rational. It does not imply Conjecture~\ref{conj:modular-periodicity},
because full bar periodicity also sees conformal weights and the chiral
bar differential.
\end{remark}

\subsection{Type II: quantum group periodicity}

\subsubsection{The quantum group structure}

\begin{proposition}[Quantum periodicity under admissible KL/DS transport;
\ClaimStatusConditional]
\label{prop:periodicity-quantum-input}
\index{periodicity!quantum|textbf}
Let $\mathcal{W}^k(\mathfrak{g})$ be the W-algebra at admissible level $k = -h^{\vee} + p/q$ where $h^{\vee}$ is the dual Coxeter number, $p \geq h^\vee$, $q \geq h^\vee$, $\gcd(p,q) = 1$. Then the quantum integers, quantum dimensions, and fusion coefficients (transported through the KL/DS interface) recur with period dividing
\[
M = 2h^{\vee}pq/\gcd(p,q,h^{\vee}).
\]
This records only the Type~II quantum input to the periodicity profile $(N_{\mathrm{mod}}, N_{\mathrm{quant}})$; upgrading to full bar complex periodicity requires a separate OPE transport theorem.
\end{proposition}

\begin{proof}
Assume a braided comparison from the admissible category to the
semisimplified tilting quotient
$\mathcal{C}(U_{q_{\mathrm{KL}}}(\mathfrak g))$, and assume that the
Drinfeld--Sokolov functor transports the relevant fusion coefficients
to the W-algebra block with the chosen OPE normalization. These are
the hypotheses in the phrase ``KL/DS transport.''

At admissible level $k = -h^\vee + p/q$, we write $q_{\mathrm{KL}} = \exp(\pi i/(k + h^\vee)) = \exp(\pi i \cdot q/p)$ for the Kazhdan--Lusztig quantum parameter. Since $\gcd(p,q) = 1$, this $q_{\mathrm{KL}}$ is a root of unity of order dividing $2p$.

The quantum integers $[n]_{q_{\mathrm{KL}}} = (q_{\mathrm{KL}}^n - q_{\mathrm{KL}}^{-n})/(q_{\mathrm{KL}} - q_{\mathrm{KL}}^{-1})$ satisfy $[n + 2\ell]_{q_{\mathrm{KL}}} = [n]_{q_{\mathrm{KL}}}$ where $\ell$ is the order of $q_{\mathrm{KL}}^2 = \exp(2\pi i q/p)$ as a root of unity; since $\gcd(p,q) = 1$, this order is $\ell = p$.

The quantum dimension of the simple module $L(\lambda)$ is the Weyl character formula:
\[
\dim_q L(\lambda) = \prod_{\alpha \in \Delta^+}
\frac{[\langle \lambda + \rho, \alpha^\vee \rangle]_{q_{\mathrm{KL}}}}
 {[\langle \rho, \alpha^\vee \rangle]_{q_{\mathrm{KL}}}}.
\]
This is a product of $|\Delta^+|$ quantum integers with arguments
determined by the weight~$\lambda$. The periodicity of each factor
gives $\dim_q L(\lambda + 2\ell\rho) = \dim_q L(\lambda)$ whenever both
weights lie in the transported quantum-group realization. Fusion
coefficients are functions of the same finite root-of-unity data
through the Verlinde formula.

The admissible weights are indexed by pairs $(r,s)$ with $0 \leq r \leq p - h^\vee$ and $0 \leq s \leq q - h^\vee$ (for generic simple~$\mathfrak{g}$; the precise range depends on $\Delta^+$). Combining the three sources
\begin{enumerate}[label=(\alph*)]
\item quantum integer period: $2\ell = 2p$,
\item affine Weyl denominator shifts, bounded by the Coxeter period $h^\vee$,
\item $q$-restricted weight lattice: period~$q$,
\end{enumerate}
we obtain a common period dividing
$M = 2h^\vee pq / \gcd(p,q,h^\vee)$ for the quantum integers, quantum
dimensions, and fusion data transported through the admissible-level
KL/DS interface.
\end{proof}

\begin{example}[Quantum periodicity for
\texorpdfstring{$\widehat{\mathfrak{sl}}_2$}{sl-hat_2} at \texorpdfstring{$k = -1/2$}{k = -1/2}]
\label{ex:hh-sl2-admissible}
\index{Hochschild cohomology!admissible level}
\index{sl2-hat@$\widehat{\mathfrak{sl}}_2$!Hochschild cohomology}
At the admissible level $k = -1/2$ ($p = 3$, $q = 2$,
$h^\vee = 2$), the quantum parameter is
$q_{\mathrm{KL}} = \exp(\pi i \cdot 2/3) = e^{2\pi i/3}$
(a primitive third root of unity). In the transported semisimplified
$U_{q_{\mathrm{KL}}}(\mathfrak{sl}_2)$ category, the trivial and
fundamental weights have quantum dimensions
\[
\dim_q(\mathcal{L}_0) = 1, \qquad
\dim_q(\mathcal{L}_1) = [2]_q = q + q^{-1} = -1
\]
(using $q = e^{2\pi i/3}$, so $[2]_q = 2\cos(2\pi/3) = -1$).

The quantum integer periodicity has $\ell = 3$
($[n+6]_q = [n]_q$, since $q^2 = e^{4\pi i/3}$ has order~$3$).
The period is $M = 2 \cdot 2 \cdot 3 \cdot 2 / \gcd(3,2,2) = 24$.

The root-of-unity calculation takes place in the transported
semisimplified quantum-group category.  The curve-level chiral
Hochschild complex of the admissible simple quotient is the coderived
diagonal complex of
Conjecture~\ref{conj:hochschild-concentration-class-M-non-generic};
its support is determined by the corresponding Shapovalov and
finite-window calculation.  The proposed period $M=24$ concerns the
weight-graded bar complex:
$\dim H^n(\bar{B}^{\mathrm{ch}}_{[h+24]}) = \dim H^n(\bar{B}^{\mathrm{ch}}_{[h]})$
for $h \gg 0$ (Conjecture~\ref{conj:modular-periodicity}).
\end{example}

\subsection{Type III: geometric tautological depth}

\subsubsection{Genus dependence and nilpotent depth}

The geometric contribution to bar cohomology from higher genus provides \emph{bounded tautological depth} rather than periodicity: the Kodaira--Spencer operator is nilpotent and the bar spectral sequence terminates after finitely many steps.

\begin{theorem}[Geometric tautological depth bound; \ClaimStatusConditional]
\label{thm:geometric-periodicity-weak}
\index{periodicity!geometric!tautological depth|textbf}
\index{tautological depth!geometric|textbf}
For a chiral algebra $\mathcal{A}$ on a genus $g \geq 2$ curve $X$,
the Kodaira--Spencer image of the Hodge class is nilpotent:
\[
\kappa(c_1(\mathbb{E}))^{3g-2} = 0
\quad \text{in } \ChirHoch^*(\mathcal{A}).
\]
In particular, the geometric direction of the bar spectral sequence
has bounded tautological depth: at most $3g - 3$ nontrivial geometric
filtration layers can contribute.

At $g = 0$, no Hodge bundle contributes; at $g = 1$, the relation
$c_1(\mathbb{E}) = \frac{1}{12}\delta$ absorbs the Hodge
contribution into the modular data.
\end{theorem}

\begin{proof}
The Kodaira--Spencer map (Theorem~\ref{thm:kodaira-spencer-chiral-complete}) provides a ring homomorphism
\[
\kappa \colon H^*(\overline{\mathcal{M}}_g, \mathbb{Q}) \longrightarrow \ChirHoch^*(\mathcal{A})
\]
whose image controls the geometric contributions to bar cohomology.
The geometric operator $\sigma_{\mathrm{geom}}$ acts via
multiplication by $\kappa(\lambda)$, where
$\lambda = c_1(\mathbb{E})$ is the first Chern class of the Hodge
bundle $\mathbb{E} \to \overline{\mathcal{M}}_g$.

Since $\dim_{\mathbb{C}} \overline{\mathcal{M}}_g = 3g - 3$, every cohomology class of degree exceeding $2(3g-3) = 6g - 6$ vanishes. In particular,
\[
\lambda^{3g-2} = c_1(\mathbb{E})^{3g-2} = 0 \quad \text{in } H^{2(3g-2)}(\overline{\mathcal{M}}_g, \mathbb{Q}),
\]
since $2(3g-2) = 6g - 4 > 6g - 6$ for all $g \geq 2$. Applying $\kappa$:
\[
\kappa(\lambda)^{3g-2} = \kappa(\lambda^{3g-2}) = \kappa(0) = 0 \quad \text{in } \ChirHoch^*(\mathcal{A}).
\]
Therefore $\sigma_{\mathrm{geom}}$ is nilpotent of index at most $3g - 2$, giving at most $3g - 3$ nontrivial filtration steps. At $g = 0$ there is no Hodge bundle; at $g = 1$, $c_1(\mathbb{E}) = \frac{1}{12}\delta$ absorbs into $N_{\mathrm{mod}}$.
\end{proof}

\begin{remark}[Nilpotent depth vs.\ periodicity]
\label{rem:nilpotent-vs-periodic}
\index{tautological depth!vs.\ periodicity}
$\sigma_{\mathrm{geom}}$ is nilpotent: $\sigma^N = 0$ for some~$N$.
No periodic identity $\sigma^N=\mathrm{id}$ occurs. The geometric
contribution provides a stabilization threshold
(cf.\ Remark~\ref{rem:stratified-periodicity}).
\end{remark}

\begin{theorem}[Sharp geometric depth on smooth moduli; \ClaimStatusProvedHere]
\label{thm:geometric-depth-smooth}
\index{tautological depth!sharp bound on $\mathcal{M}_g$|textbf}
\index{Graber--Vakil vanishing}
For a chiral algebra $\mathcal{A}$ on a smooth genus $g \geq 2$
curve, suppose the Kodaira--Spencer map factors through the
cohomology of the smooth moduli space:
\[
\kappa \colon H^*(\mathcal{M}_g, \mathbb{Q}) \longrightarrow
\ChirHoch^*(\mathcal{A}).
\]
Then the nilpotency index of $\kappa(\lambda)$ is at most~$g - 1$:
\begin{equation}\label{eq:sharp-geometric-depth}
\kappa(c_1(\mathbb{E}))^{g-1} = 0
\quad \text{in } \ChirHoch^*(\mathcal{A}).
\end{equation}
In particular:
\begin{enumerate}[label=\textup{(\alph*)}]
\item At genus~$2$, the Hodge class contribution vanishes entirely:
 $\kappa(\lambda) = 0$.
\item At genus~$3$, the geometric depth is at most~$2$
 (compared to~$7$ from the compactified bound).
\item In general, the sharp bound $g - 1$ improves the
 compactified bound $3g - 2$ by a factor of approximately~$3$.
\end{enumerate}
\end{theorem}

\begin{proof}
The tautological ring $R^*(\mathcal{M}_g)$ of the smooth
moduli space satisfies the Graber--Vakil vanishing
theorem~\cite{GraVak05} (extending Looijenga~\cite{Looij95}):
\begin{equation}\label{eq:graber-vakil-vanishing}
R^d(\mathcal{M}_g) = 0 \quad \text{for } d > g - 2.
\end{equation}
The Hodge class $\lambda = c_1(\mathbb{E})$ is a
tautological class of degree~$1$.
On the smooth moduli space, the Mumford relation reduces to
$\lambda|_{\mathcal{M}_g} = \kappa_1/12$
(since the boundary class $\delta$ vanishes on
$\mathcal{M}_g$). Therefore
\[
\lambda^{g-1} \in R^{g-1}(\mathcal{M}_g) = 0
\]
since $g - 1 > g - 2$. Applying $\kappa$:
\[
\kappa(\lambda)^{g-1} = \kappa(\lambda^{g-1}) = 0.
\]

For genus~$2$: $g - 1 = 1$, so $\lambda^1 = \lambda \in
R^1(\mathcal{M}_2)$. Since $R^d(\mathcal{M}_2) = 0$ for
$d > 0$, we get $\lambda|_{\mathcal{M}_2} = 0$, hence
$\kappa(\lambda) = 0$.
This means the Hodge class produces no geometric correction
to bar cohomology at genus~$2$: the genus-$2$ bar complex, at
the level of the Hodge class contribution, is determined
entirely by the genus-$0$ data.
\end{proof}

\begin{remark}[Factorization hypothesis]
\label{rem:factorization-hypothesis-geom}
The factorization through $H^*(\mathcal{M}_g)$ holds on smooth
curves. On $\overline{\mathcal{M}}_g$, boundary divisor classes
$\delta_i$ can increase nilpotent depth beyond $g-1$ (but never
beyond $3g-2$, Theorem~\ref{thm:geometric-periodicity-weak}).
\end{remark}

\begin{conjecture}[Geometric amplitude on compactified moduli; \ClaimStatusConjectured]
\label{conj:geometric-periodicity}
\index{tautological depth!sharp bound on $\overline{\mathcal{M}}_g$|textbf}
For a Koszul chiral algebra~$\mathcal{A}$
on $\overline{\mathcal{M}}_g$
with factorization extensions to the boundary, the nilpotency
index of the geometric operator $\kappa(\lambda)$ on
$\overline{\mathcal{M}}_g$ is controlled by the Mumford
relation $12\lambda = \kappa_1 + \delta$:
the boundary contribution $\kappa(\delta)$ increases the
nilpotency index beyond the interior bound $g - 1$ by at most
the depth of the boundary tautological ring. Precisely, there
exists a function $\mathrm{nil}_{\partial}(g)$, depending only
on the genus, such that:
\[
\mathrm{nil}(\kappa(\lambda)) \;\leq\; (g-1) + \mathrm{nil}_{\partial}(g),
\qquad \mathrm{nil}_{\partial}(g) \leq 2g - 1.
\]
The additive structure reflects the Mumford decomposition
$12\lambda = \kappa_1 + \delta$, with the interior contributing
at most $g - 1$ (proved) and the boundary contributing
$\mathrm{nil}_{\partial}(g)$ additional layers.
\end{conjecture}

\begin{remark}[Scope and provenance]\label{rem:geometric-periodicity-scope}
The interior bound $g-1$ is proved
(Theorem~\ref{thm:geometric-depth-smooth}, Graber--Vakil); the
compactified bound $3g-2$ is proved
(Theorem~\ref{thm:geometric-periodicity-weak}). The conjecture
interpolates via the Mumford relation. The geometric contribution
is nilpotent ($\sigma^N = 0$), not periodic ($\sigma^N = \mathrm{id}$);
true periodicity arises from Types~I and~II only.
\end{remark}

\subsubsection{Examples at different genera}

\begin{example}[Genus \texorpdfstring{$0$}{0}: the sphere]
No geometric tautological depth (simply connected, no moduli).
\end{example}

\begin{example}[Genus \texorpdfstring{$1$}{1}: the torus]
For elliptic curve $E_{\tau}$:
\begin{itemize}
\item Period lattice $\Lambda = \mathbb{Z} + \tau\mathbb{Z}$
\item Four spin structures (fermions have period 8)
\item Modular parameter $\tau$ gives $SL_2(\mathbb{Z})$ action
\end{itemize}

Free fermion on $E_{\tau}$: the weight-graded bar complex cohomology satisfies
\[
\dim H^n(\bar{B}^{\mathrm{ch}}_{[h+8]}(\mathcal{F},\, E_{\tau}))
\;=\;
\dim H^n(\bar{B}^{\mathrm{ch}}_{[h]}(\mathcal{F},\, E_{\tau}))
\quad\text{for } h \gg 0.
\]
The period~$8$ comes from: $4$ spin structures $\times$ $2$ (fermion parity).
This periodicity is indexed by conformal weight.  The Hochschild-degree
support is determined separately by a family support datum
$H_H(\mathcal F;S_{\mathcal F})$.
The Hodge class is absorbed into the modular data via
$c_1(\mathbb{E}) = \frac{1}{12}\delta$.
\end{example}

\begin{example}[Genus 2: Hodge class vanishing]
\label{ex:genus2-hodge-vanishing}
By Theorem~\ref{thm:geometric-depth-smooth}, the Hodge class
contribution vanishes entirely at genus~$2$:
$\kappa(\lambda) = 0$ in $\ChirHoch^*(\mathcal{A})$. This follows from
the Graber--Vakil vanishing $R^d(\mathcal{M}_2) = 0$ for $d > 0$.
The genus-$2$ bar complex receives no geometric corrections from
$c_1(\mathbb{E})$; any higher-genus modifications at $g = 2$ must
arise from other tautological classes or from boundary
contributions on $\overline{\mathcal{M}}_2$ (which reduce to
$\mathcal{M}_{1,1} \times \mathcal{M}_{1,1}$ and
$\mathcal{M}_{1,2}$ data).
\end{example}

\subsection{Stratified periodicity}

This subsection separates the unconditional structural part of the
periodicity story from the conditional numerical values.

\begin{conjecture}[Stratified periodicity package (partial); \ClaimStatusConjectured]
\label{conj:complete-periodicity-classification}
\index{periodicity!classification (partial)|textbf}
\index{periodicity!stratified}
For a rational chiral algebra $\mathcal{A}$ on a genus~$g$ curve with
central charge $c$, rational simple-module conformal weights
$\{h_i\}$, and quantum group level inducing period~$M$, the bar
cohomology dimensions are eventually periodic in conformal weight:
\[
\text{Period}(\mathcal{A}) \mid
\text{lcm}(N_{\text{modular}}, N_{\text{quantum}})
\]
where:
\begin{align}
N_{\text{modular}} &= N_T(\cA)
=\operatorname{lcm}_i\operatorname{den}(h_i-c/24) \quad
 \text{(Conjectures~\ref{conj:modular-periodicity-minimal},
 \ref{conj:modular-periodicity-wzw} for minimal/WZW;
 Conjecture~\ref{conj:modular-periodicity} in general)} \\
N_{\text{quantum}} &= M
 \text{ (from the quantum-group realization under admissible transport)}
\end{align}
and the conformal weight threshold beyond which this periodicity
holds is controlled by the geometric tautological depth
$\mathrm{nil}(\kappa(\lambda)) \leq g - 1$
(Theorem~\ref{thm:geometric-depth-smooth}).
\end{conjecture}

\begin{evidence}
The unconditional content is structural rather than numerical.
What is proved is the tautological depth bound and the root-of-unity
calculation inside the Type~II quantum-group realization. What remains
conjectural is the sharp Type~I bar-periodicity input and the
admissible transport from the quantum-group realization to the chiral
bar differential.

The bar spectral sequence
\[
E_2^{p,q} = H^q(\overline{C}_{p+2}(X))
 \otimes H^0(\mathcal{A}^{\otimes(p+2)})
\]
admits a tridegree decomposition
$\bigl(d_{\mathrm{mod}}, d_{\mathrm{quant}}, d_{\mathrm{geom}}\bigr)$,
where the three indices record respectively the representation-theoretic
weight, the fusion-theoretic degree, and the configuration-space
cohomological degree.

\emph{Step~1: Inputs.}
\begin{enumerate}
\item \emph{Modular (Type~I).} For the rational families treated
 above, Step~1 of
 Conjectures~\ref{conj:modular-periodicity-minimal}
 and~\ref{conj:modular-periodicity-wzw} proves only the
 $T$-matrix phase bound on the finite-dimensional character space.
 Upgrading that phase bound to eventual periodicity of the bar
 cohomology with order~$N_{\mathrm{mod}}$ is precisely the
 conjectural input of the Type~I modular-periodicity package.
\item \emph{Quantum (Type~II).} The quantum-integer periodicity
 $[n + 2\ell]_q = [n]_q$ is proved in the admissible-level
 KL/DS regime of Proposition~\ref{prop:periodicity-quantum-input}, so
 the Type~II profile contributes a proved
 $N_{\mathrm{quant}}$ factor.
\item \emph{Geometric (Type~III).}
 Theorem~\ref{thm:geometric-depth-smooth} proves
 $\kappa(\lambda)^{g-1} = 0$ on smooth moduli, giving a
 sharp tautological depth bound. This is a \emph{nilpotent
 stabilization}, not a periodicity source: it bounds the
 transient before modular/quantum periodicity takes hold,
 rather than contributing a periodic factor to the lcm
 (Remark~\ref{rem:nilpotent-vs-periodic}).
\end{enumerate}

\emph{Step~2: Structural mechanism.}
The genuine periodicity operators
$\sigma_{\mathrm{mod}}$ and $\sigma_{\mathrm{quant}}$
act on the first two tridegree directions. Their combined order
divides
$\mathrm{lcm}(N_{\mathrm{mod}}, N_{\mathrm{quant}})$.
The geometric direction is \emph{bounded}, not periodic: the
nilpotent operator $\sigma_{\mathrm{geom}} = \kappa(\lambda)$
satisfies $\sigma_{\mathrm{geom}}^{g-1} = 0$
(Theorem~\ref{thm:geometric-depth-smooth}), so the
geometric spectral sequence has at most $g - 2$ nontrivial
filtration steps.

The structural separation is:
\begin{itemize}
\item the modular input acts on the representation-theoretic component
 via $T$-eigenvalues / conformal weights (periodic, conjectural);
\item the quantum input acts on the fusion-theoretic component via
 quantum dimensions and Clebsch--Gordan data (periodic after the
 admissible transport hypothesis);
\item the geometric input acts on the configuration-space component via
 tautological classes on $\mathcal{M}_g$ (\emph{nilpotent}, proved).
\end{itemize}

\emph{Step~3: Combined structure.}
The genuine period of $\mathcal{A}$ is governed by the periodic
operators only:
$\mathrm{Period}(\mathcal{A}) \mid \mathrm{lcm}(N_{\mathrm{mod}},
N_{\mathrm{quant}})$. The geometric contribution controls
the \emph{stabilization threshold}: bar cohomology dimensions
become $\mathrm{Period}(\mathcal{A})$-periodic in conformal weight
once the weight exceeds a threshold determined by the tautological
depth. On $\overline{\mathcal{M}}_g$ the depth is at most $3g-2$
(Theorem~\ref{thm:geometric-periodicity-weak}); on $\mathcal{M}_g$
the sharp bound is $g-1$
(Theorem~\ref{thm:geometric-depth-smooth}).
\end{evidence}

\begin{remark}[Scope]\label{rem:periodicity-classification-scope}
Conjectured because the modular bound
(Conjecture~\ref{conj:modular-periodicity}) is open. The
root-of-unity quantum calculation is explicit under the admissible
transport hypothesis; the geometric contribution is nilpotent depth
(proved), not a periodic source. The lcm formula
$\mathrm{Period}(\cA) \mid
\mathrm{lcm}(N_{\mathrm{mod}}, N_{\mathrm{quant}})$ is the correct
statement.
\end{remark}

\begin{remark}[Periodicity profile refinement]\label{rem:periodicity-triple-deformation}
\index{periodicity profile!refinement}
The scalar period is the shadow of the richer profile
$\Pi(\cA) = (M_\cA, Q_\cA, G_\cA)$
(Remark~\ref{rem:periodicity-triple}):
$M_\cA$ carries the modular $T$-eigenvalue action (conjectural),
$Q_\cA$ encodes quantum CG periodicity under the admissible transport
hypothesis, and $G_\cA$
encodes the nilpotent tautological depth
(Theorem~\ref{thm:geometric-depth-smooth}).
\end{remark}

\begin{remark}[Stratified periodicity]\label{rem:stratified-periodicity}
\index{periodicity!stratified reformulation}
The mature formulation separates periodic from nilpotent contributions:
(a)~the lcm mechanism for the two genuine periodic sources is proved
(Proposition~\ref{prop:periodicity-exchange-koszul}); the quantum
root-of-unity calculation is explicit under admissible transport;
geometric depth $\leq g-1$
(Theorem~\ref{thm:geometric-depth-smooth}) is proved.
(b)~Minimal models and WZW have explicit $T$-matrix periodicity,
but bar-periodicity remains conjectural
(Conjectures~\ref{conj:modular-periodicity-minimal},
\ref{conj:modular-periodicity-wzw}: denominator convolution and
bar-differential compatibility).
(c)~For general Koszul chiral algebras, the period divides
$\mathrm{lcm}(N_{\mathrm{mod}}, N_{\mathrm{quant}})$; the spectral
radius of $T_{\mathrm{rec}}$ bounds $N_{\mathrm{mod}}$
(Theorem~\ref{thm:dominant-branch-point}).
No periodic factor is assigned to the nilpotent geometric direction.
\end{remark}

\subsection{Koszul duality and periodicity interaction}

\begin{proposition}[Periodicity-profile transport under Koszul duality;
 \ClaimStatusConditional]
\index{periodicity!exchange under duality}
\label{prop:periodicity-exchange-koszul}
Let $(\cA, \cA^!)$ be a Koszul dual pair of rational chiral algebras on a genus~$g$ curve~$X$. Then:
\begin{enumerate}
\item The geometric tautological bound is common to the pair:
 on $\mathcal M_g$ both nilpotence indices are at most $g-1$, and on
 $\overline{\mathcal M}_g$ both are at most $3g-2$.
\item The dual vacuum phase period is determined by the dual central
 charge:
 $N_{c'} = 24q'/\gcd(p', 24)$ where
 $c' = c(\cA^!) = p'/q'$ in lowest terms. The full dual
 $T$-order is
 $N_T(\cA^!)=\operatorname{lcm}_j\operatorname{den}(h'_j-c'/24)$.
\item If both algebras lie in the admissible-level KL/DS regime of
 Type~II above, then the quantum periodicity input is preserved:
 $N_{\text{quant}}(\cA) = N_{\text{quant}}(\cA^!)$.
\end{enumerate}
In particular, Koszul duality transports the vacuum phase bound and,
under the Type~II hypotheses, the quantum input; the geometric term is
the common ambient tautological nilpotence bound. It does not assert
equality of scalar periods, since the modular component changes with
the dual central charge. Whenever both periodic components are defined,
the total period of $\cA^!$ divides
$\mathrm{lcm}(N_T(\cA^!), N_{\text{quant}})$.
\end{proposition}

\begin{proof}
\emph{(1) Common geometric tautological bound.}
Both $\cA$ and $\cA^!$ live over the same curve $X$. The Hodge class
$\lambda=c_1(\mathbb E)$ is a class on the same moduli space, and
Theorems~\ref{thm:geometric-periodicity-weak} and
\ref{thm:geometric-depth-smooth} bound every image of this class in
chiral Hochschild cohomology by the same ambient tautological
nilpotence. The actual image may vanish earlier for one algebra than
for the other; only the bound is intrinsic to the curve.

\emph{(2) Vacuum phase and full $T$-order of the dual.}
The dual central charge $c'$ is determined by the level duality
$k \mapsto k' = -k - 2h^\vee$ via the Sugawara formula, giving
$c'=p'/q'$ in lowest terms. The integer
$N_{c'}=24q'/\gcd(p',24)$ kills the vacuum phase
$e^{-2\pi i c'/24}$. The full $T$-matrix of~$\cA^!$ has eigenvalues
$e^{2\pi i(h'_j-c'/24)}$ and therefore has order
$N_T(\cA^!)=\operatorname{lcm}_j\operatorname{den}(h'_j-c'/24)$.
This finite $T$-matrix order depends on rationality of the conformal
weights, not only on rationality of~$c'$.

\emph{(3) Quantum input in the admissible-level regime.}
In the Type~II admissible-level KL/DS regime, the quantum component is
encoded by quantum integers, quantum dimensions, and fusion data.
Koszul duality sends level $k$ to $k' = -k - 2h^\vee$, hence
$q \mapsto q^{-1}$. Since $[n]_{q^{-1}} = [n]_q$ for all $n$, the
Type~II quantum input is preserved under duality whenever that regime
applies.
\end{proof}

\begin{remark}[Vacuum phase exchange]\label{rem:periodicity-exchange-harmonic}
% backward-compatible label removed (was conj:periodicity-exchange-harmonic)
Proposition~\ref{prop:periodicity-exchange-koszul} establishes
unconditionally that the geometric tautological bound is common and
that the dual vacuum phase bound is determined by the dual central
charge. In the admissible-level KL/DS regime it also preserves the
Type~II quantum input. The sum of reciprocal vacuum phase periods is given by
Conjecture~\ref{conj:reflected-modular-periodicity}: writing $c = p/q$
and $c' = p'/q'$,
\[
\frac{1}{N_c} + \frac{1}{N_{c'}}
= \frac{\gcd(p,24)}{24q}
 + \frac{\gcd(p',24)}{24q'}.
\]
In the Virasoro same-family case $c+c'=26$ this specializes to the
criterion
$\gcd(p,24)+\gcd(26q-p,24)=2q$. For $q=1$ it requires both $p$ and
$26-p$ to be coprime to~$24$, e.g.\ $c=1$, $c'=25$ or $c=7$,
$c'=19$, but not $c=5$, $c'=21$ since $\gcd(21,24)=3$.
\end{remark}



\section{Physical applications}

\subsection{Marginal deformations in CFT}

In the curve-level chiral chart, $\ChirHoch^2(\mathcal{A})$ contains
first-order product-deformation classes of the OPE algebra in the
unshifted Gerstenhaber convention; the formal MC tangent convention is
shifted, so the product tangent lies in degree one of the shifted cyclic
complex, whereas unshifted $\ChirHoch^1$ is the fixed-product
automorphism degree. A CFT deformation
$S \mapsto S+\lambda\int_\Sigma\phi(z,\bar z)\,d^2z$ is exactly
marginal only after three independent conditions hold: $\phi$ has
weight~$(1,1)$, its chiral component defines a
$\ChirHoch^2$-cocycle, and the Kuranishi obstruction maps from
Proposition~\ref{prop:mc-tangent-amplitude} vanish. Theorem~H removes
the ambiguity in the ambient curve-level support by identifying it
with $H^\bullet(K_{\cA,S})$.  All-order integrability is then the
vanishing of the corresponding Kuranishi maps.

\begin{example}[Exactly marginal deformations]
For the rank-one even superboson, Bakalov--De Sole--Kac give bounded
cohomology polynomial $2+t$; a bounded-to-chart quasi-isomorphism
transports this calculation to the curve.  The physical radius belongs
to the enriched CFT parameter atlas, and fixed nonzero Heisenberg
levels are related by rescaling after a choice of square root.
For Virasoro, their bounded calculation is supported in degrees
$0,2,3$; its chart interpretation likewise uses the comparison
$\chi_{\operatorname{Vir}_c}^{\mathrm{bd}}$.
Four-dimensional
gauge-coupling moduli belong to a different field-theoretic
Hochschild/BV complex, connected to the Vol~I chart by a separate
comparison map.
\end{example}

\subsection{String field theory}

The cyclic Hochschild/BV complex of open string field theory gives
string vertices.  A chiral-shadow comparison map connects it to the
curve-level chart complex of Theorem~H.

\begin{theorem}[String field theory from Hochschild {\cite{Zwi93}}; \ClaimStatusProvedElsewhere]\label{thm:string-field-theory-hochschild}
The operations
$m_n: \mathcal{A}^{\otimes n} \to \mathcal{A}[2-n]$
extracted from the cyclic open-string Hochschild/BV complex satisfy:
\[
\sum_{\substack{r+s+t=n \\ r,t \geq 0,\; s \geq 1}} (-1)^{rs+t}\, m_{r+1+t}(\mathrm{id}^{\otimes r} \otimes m_s \otimes \mathrm{id}^{\otimes t}) = 0
\]
These give:
\begin{itemize}
\item $m_1$: BRST operator $Q$
\item $m_2$: String multiplication
\item $m_3$: Four-string vertex
\item Higher $m_n$: Contact terms
\end{itemize}
\end{theorem}

The action:
\[
S[\Psi] = \frac{1}{2}\langle \Psi, Q\Psi \rangle + \sum_{n \geq 3} \frac{1}{n!}\langle \Psi, m_n(\Psi,\ldots,\Psi)\rangle
\]

\subsection{Holographic duality}

\begin{conjecture}[Holographic Koszul duality; \ClaimStatusConjectured]
\label{conj:holographic-koszul-deformation}
A holographic Koszul duality datum for AdS$_3$/CFT$_2$ consists of:
\begin{itemize}
\item a holomorphic factorisation algebra
 $\mathcal{A}_{\mathrm{bulk}}$ on a three-manifold with boundary;
\item a boundary chiral algebra
 $\mathcal{A}_{\mathrm{bdry}}$ obtained by factorisation restriction
 to the boundary curve;
\item a finite-type bar coalgebra
 $B^{\mathrm{ord}}(\mathcal{A}_{\mathrm{bulk}})$ whose Verdier-dual
 algebra satisfies
 \[
   \bigl(B^{\mathrm{ord}}(\mathcal{A}_{\mathrm{bulk}})\bigr)^\vee
   \simeq \mathcal{A}_{\mathrm{bdry}}
 \]
 in the completed boundary factorisation category.
\end{itemize}
\end{conjecture}

\begin{remark}[Scope]
This conjecture lies outside the present monograph: it is a physical
interpretation, not a theorem. A theorem requires a chosen
Costello--Li factorisation algebra on the bulk, a boundary restriction
functor, finite-type Verdier duality, and a proof that the boundary
central-charge constraint is an anomaly-cancellation identity in this
factorisation category. The relation
$c_{\text{bulk}}+c_{\text{boundary}}=26$ requires the open/closed
comparison in addition to Koszul duality.
The named open/closed comparison is the brace dg map
$\beta_T$ of
Proposition~\ref{prop:thqg-universal-action-not-reconstruction}.  A
quasi-isomorphism $\beta_T$ in the chosen completion identifies the
physical theory with the cochain derived centre.
\end{remark}

The conjectural holographic dictionary has three inputs: the central
charge relation $c_{\mathrm{bulk}}+c_{\mathrm{boundary}}=26$, a perfect
chain pairing between the curve-level chart complexes, and the
open/closed twisting morphism.

\section{Computing chiral Hochschild cohomology via bar-cobar resolution}
\label{sec:hochschild-via-bar-cobar-complete}

\subsection{The diagonal bar resolution}

Chiral Hochschild cohomology is computed from the diagonal
$\mathcal A$-$\mathcal A$ bimodule resolution
\[
\operatorname{Bar}_{\mathcal A^e}(\mathcal A)
=
\mathcal A\otimes\bar B(\mathcal A)\otimes\mathcal A.
\]
The bar--cobar inversion
$\Omega(\bar B(\mathcal A))\simeq\mathcal A$ supplies the
cofibrancy of this resolution on the Koszul locus. The Hochschild
cochain complex is obtained by applying
$\mathrm{Hom}_{\mathcal A^e}(-,\mathcal A)$ to this diagonal
resolution; the Koszul dual algebra is the Verdier object
$\mathcal A^!$.

\subsection{The fundamental quasi-isomorphism}

\begin{theorem}[Bar-cobar resolution; \ClaimStatusConditional]\label{thm:bar-cobar-resolution}
\textup{[Regime: quadratic at genus~$0$; curved-central at genus~$g \geq 1$
on the Koszul locus
\textup{(}Convention~\textup{\ref{conv:regime-tags})}.]}

For any chiral Koszul algebra $\mathcal{A}$ on a curve $X$
\textup{(}Definition~\textup{\ref{def:chiral-koszul-pair})},
the counit is a quasi-isomorphism:
\[\Omega^{\text{geom}}(\bar{B}^{\text{geom}}(\mathcal{A})) \xrightarrow{\sim} \mathcal{A}\]
\end{theorem}

\begin{proof}
This is a consequence of
Theorem~\ref{thm:bar-cobar-isomorphism-main}~(III), which establishes
$\Omega^{\mathrm{ch}}(\bar{B}^{\mathrm{ch}}(\mathcal{A})) \simeq
\mathcal{A}$ for Koszul chiral algebras; the filtered proof is in
Chapter~\ref{chap:higher-genus}.

The bar complex $\bar{B}^{\mathrm{geom}}(\mathcal{A})$ is a conilpotent
chiral coalgebra with $\dzero^{\,2} = 0$
\textup{(}Theorem~\textup{\ref{thm:bar-nilpotency-complete}}\textup{)}.
The cobar construction $\Omega^{\mathrm{ch}}(\bar{B}^{\mathrm{ch}}(\mathcal{A}))$
is a free chiral algebra
\textup{(}Theorem~\textup{\ref{thm:cobar-free}}\textup{)}, hence the
composite is a free resolution.

The augmentation
$\epsilon\colon \Omega^{\mathrm{ch}}(\bar{B}^{\mathrm{ch}}(\mathcal{A}))
\to \mathcal{A}$ is a quasi-isomorphism by the following spectral
sequence argument: the bar filtration has associated graded governed
by the quadratic coalgebra $\mathcal A^i$ through the
quasi-isomorphism
$q_\mathcal A\colon\mathcal A^i\to\bar B^{\mathrm{ch}}(\mathcal A)$.
Verdier duality produces
$K_X(\mathcal A)=\mathbb D_{\Ran}\bar B^{\mathrm{ch}}(\mathcal A)$
and
$\mathbb D(q_\mathcal A)\colon K_X(\mathcal A)\to
\mathbb D_{\Ran}(\mathcal A^i)$; the chosen pair map
$\nu_\mathcal A\colon K_X(\mathcal A)\to\mathcal A^!$ supplies the
Koszul dual algebra. The associated Koszul
spectral sequence collapses at $E_2$ for Koszul algebras
(Theorem~\ref{thm:spectral-sequence-collapse}), and the abutment gives
$H^*(\Omega(\bar{B}(\mathcal{A}))) \cong \mathcal{A}$.

For general (not necessarily Koszul) chiral algebras, the classical
bar-cobar adjunction is acyclic in the \emph{connected} graded
setting (\cite{LV12}, Theorem~2.2.5), but the chiral
analog requires the Koszul or pro-nilpotent hypothesis: at
admissible affine levels and minimal-model loci, the counit
$\Omega(\bar{B}(\mathcal{A})) \to \mathcal{A}$ is not a
quasi-isomorphism (Remark~\ref{rem:sl2-admissible}).
\end{proof}

\begin{remark}
The bar-cobar construction is inversion:
\begin{align*}
\mathcal{A} &\xrightarrow{\bar{B}} \text{bar coalgebra}\\
&\xrightarrow{\Omega} \text{cobar algebra}\\
&\xrightarrow{\epsilon} \mathcal{A}.
\end{align*}

The Hochschild complex applies
$\mathrm{Hom}_{\mathcal A^e}(-,\mathcal A)$ to the diagonal
two-sided bar resolution.  Verdier duality forms
$K_X(\mathcal A)=\mathbb D_{\Ran}\bar B(\mathcal A)$, and the
Koszul-pair map~$\nu_\mathcal A$ supplies~$\mathcal A^!$.
\end{remark}

\subsection{Chiral Hochschild cohomology formula}

\begin{theorem}[Fulton--MacPherson model for chiral Hochschild cochains;
\ClaimStatusProvedHere]\label{thm:HH-config-space-formula}
Under the hypotheses of Theorem~\textup{\ref{thm:hochschild-bar-cobar}},
the chiral derived centre
\[
Z^{\mathrm{der}}_{\mathrm{ch}}(\mathcal A)
   =R\operatorname{Hom}_{\mathcal A^e}(\mathcal A,\mathcal A)
   =\ChirHoch^\bullet(\mathcal A,\mathcal A)
\]
is computed by the total complex
\[
C^n_{\mathrm{FM}}(\mathcal A)
=
\Gamma\!\left(
\overline{C}_{n+2}(X),
\mathcal{H}om_{\mathcal{D}_{X^{n+2}}}
\!\left(j_*j^*\mathcal A^{\boxtimes(n+2)},\Delta_*\mathcal A\right)
\otimes \Omega^n_{\overline{C}_{n+2}(X)}(\log D)
\right),
\]
with total differential
\begin{equation}\label{eq:d-Hoch-decomposition}
d_{\mathrm{Hoch}}
=d_{\mathrm{int}}+d_{\mathrm{fact}}+d_{\mathrm{form}} .
\end{equation}
Here $d_{\mathrm{int}}$ is the internal differential of the
$\mathcal D_X$-module $\mathcal A$, $d_{\mathrm{fact}}$ is the
alternating sum of chiral products along collision divisors, and
$d_{\mathrm{form}}$ is the logarithmic de~Rham differential on
$\overline C_{n+2}(X)$. This computes the Hochschild cochain object.
The separate Verdier algebra and its Koszul partner are
\[
K_X(\mathcal A)=\mathbb D_{\Ran}\bar B^{\mathrm{ch}}(\mathcal A),
\qquad
\nu_\mathcal A\colon K_X(\mathcal A)\longrightarrow\mathcal A^!.
\]
\end{theorem}

\begin{proof}
Theorem~\ref{thm:hochschild-bar-cobar} identifies
$R\operatorname{Hom}_{\mathcal A^e}(\mathcal A,\mathcal A)$ with
\[
\operatorname{Hom}_{\mathcal A^e}
\bigl(\operatorname{Bar}_{\mathcal A^e}(\mathcal A),\mathcal A\bigr),
\qquad
\operatorname{Bar}_{\mathcal A^e}(\mathcal A)
=\mathcal A\otimes\bar B^{\mathrm{ch}}(\mathcal A)\otimes\mathcal A .
\]
The Fulton--MacPherson compactification records the same two-sided bar
configuration geometrically: the two endpoint copies of $\mathcal A$
and the $n$ internal bar inputs give $n+2$ marked points, the
diagonal target is $\Delta_*\mathcal A$, and logarithmic forms encode
boundary collisions. Applying
$\mathcal H om_{\mathcal D_{X^{n+2}}}(-,\Delta_*\mathcal A)$ and
taking global sections gives $C^\bullet_{\mathrm{FM}}(\mathcal A)$.

The three summands in \eqref{eq:d-Hoch-decomposition} are the three
summands of the two-sided chiral bar differential after applying
$\operatorname{Hom}_{\mathcal A^e}(-,\mathcal A)$: internal
$\mathcal D$-module differential, collision residue/chiral product,
and logarithmic de~Rham differential. The equality
$d_{\mathrm{Hoch}}^2=0$ is Theorem~\ref{thm:chiral-hochschild-differential}.
Thus the FM complex computes the derived chiral centre and introduces
no map from the Hochschild cochain object to $\mathcal A^!$.
\end{proof}

\begin{remark}[Decomposition of the differential]
The decomposition \eqref{eq:d-Hoch-decomposition} mirrors the bar differential:
\begin{center}
\begin{tabular}{c|c|c}
\textbf{Bar} & \textbf{chiral Hochschild} & \textbf{Meaning} \\
\hline
$d_{\text{internal}}$ & $d_{\text{internal}}$ & Differential on $\mathcal{A}$ \\
$d_{\text{residue}}$ & $d_{\text{factor}}$ & Extract boundary data \\
$d_{\text{form}}$ & $d_{\text{form}}$ & Configuration space geometry \\
\end{tabular}
\end{center}
\end{remark}

\subsection{Explicit computation: free boson (Heisenberg algebra)}

\begin{example}[Bounded cohomology of the rank-one even superboson;
\ClaimStatusProvedElsewhere]
\label{ex:HH-heisenberg-complete}
Let \(B_{\mathfrak h}\) be the superboson vertex algebra associated to
a one-dimensional even space \(\mathfrak h\).  Bakalov--De Sole--Kac
compute
\[
 H^n_{\mathrm{ch},b}(B_{\mathfrak h},B_{\mathfrak h})
 \cong
 \bigl(S^n(\Pi\mathfrak h)\bigr)^*
 \oplus
 \bigl(S^{n+1}(\Pi\mathfrak h)\bigr)^*
 \qquad\text{\cite[Theorem~7.4]{BDSK21}}.
\]
Since \(\Pi\mathfrak h\) is a one-dimensional odd space, its symmetric
powers occur in degrees zero and one.  Hence
\[
 \dim H^n_{\mathrm{ch},b}(B_{\mathfrak h},B_{\mathfrak h})
 =
 \begin{cases}
  2,&n=0,\\
  1,&n=1,\\
  0,&n\geq2.
 \end{cases}
\]
For the associated chiral algebra \(\cA_{B_{\mathfrak h}}\), a
bounded-to-chart quasi-isomorphism
\[
 \chi_{B_{\mathfrak h}}^{\mathrm{bd}}\colon
 C^\bullet_{\mathrm{ch},b}(B_{\mathfrak h},B_{\mathfrak h})
 \xrightarrow{\ \sim\ }
 C^\bullet_{\mathrm{ch}}
 (\cA_{B_{\mathfrak h}},\cA_{B_{\mathfrak h}})
\]
transports the calculation to
\[
 P_{\cA_{B_{\mathfrak h}}}(t)=2+t.
\]
For a fixed nonzero Heisenberg level, rescaling the generator after a
choice of square root identifies the corresponding fibres.  A level
direction belongs to a chosen parameter family and its relative
deformation complex.
\end{example}

\subsection{Explicit computation: free fermion}

\begin{example}[Free-fermion chart datum; \ClaimStatusConditional]
\label{ex:HH-fermion-complete}
Let \(\mathcal F\) be the free-fermion vertex algebra with
\[
 \psi(z)\psi^*(w)\sim\frac{1}{z-w}.
\]
Its curve-level Hochschild object is the chart complex
\[
 C^\bullet_{\mathrm{ch}}(\cA_{\mathcal F},\cA_{\mathcal F}).
\]
A family calculation consists of an incidence-compatible retract
\(H_H(\cA_{\mathcal F};S_{\mathcal F})\); Theorem~H then gives
\[
 \ChirHoch^n(\cA_{\mathcal F})
 \cong H^n(K_{\cA_{\mathcal F},S_{\mathcal F}}).
\]
Equivalently, a bounded vertex-cohomology computation transports to
the chart after the map
\[
 \chi_{\mathcal F}^{\mathrm{bd}}\colon
 C^\bullet_{\mathrm{ch},b}(\mathcal F,\mathcal F)
 \longrightarrow
 C^\bullet_{\mathrm{ch}}(\cA_{\mathcal F},\cA_{\mathcal F})
\]
has been proved to be a quasi-isomorphism.  The simple-pole residue
enters the stratum differential of this complex.  The spin parameter
\((h_\psi,h_{\psi^*})=(\lambda,1-\lambda)\) varies the chosen conformal
vector and belongs to the enriched conformal atlas.
\end{example}

\subsection{Koszul duality and HH* pairing}

The curve-level support theorem is realized by a family support datum
$H_H(\cA;S)$.  Its Koszul-dual reflection is realized by the additional
perfect chain pairing of Theorem~\ref{thm:main-koszul-hoch}.

\begin{example}[$bc$/$\beta\gamma$ chart pairing;
\ClaimStatusConditional]
\label{ver:boson-fermion-HH}
Let \(\mathcal F_{bc}\) and \(\mathcal{BG}\) be the quadratic
\(bc\)/\(\beta\gamma\) Koszul pair on the curve.  Suppose their
quadratic comparison maps \(q_{\mathcal F_{bc}}\) and
\(q_{\mathcal{BG}}\) are quasi-isomorphisms, both chart complexes carry
family support data, and the residue evaluation defines a perfect
degree-\(d\) chain pairing.  Then Theorem~H gives
\[
 \ChirHoch^n(\mathcal F_{bc})
 \cong
 \ChirHoch^{d-n}(\mathcal{BG})^\vee\otimes\omega_X.
\]
The support sets and the integer \(d\) are determined by the explicit
collision retracts and the total degree of the residue pairing.  The
spin parameter varies the chosen conformal vector and belongs to the
enriched free-field atlas.
\end{example}

\subsection{Comparison with classical Hochschild cohomology}

\begin{remark}[Chiral vs.\ classical]\label{rem:chiral-vs-classical-hochschild}
Chiral $\ChirHoch^*$ differs from classical algebraic Hochschild
$HH^n_{\text{classical}}(A)=\Ext^n_{A^e}(A,A)$ by locality,
configuration-space integrals, differential-form coefficients, and the
spatial direction from~$X$. Only in the flat constant-coefficient
model, after discarding all collision-residue and $\mathcal D_X$
data, does one recover the classical algebraic Hochschild complex of
the coefficient algebra. This specialization is a comparison of
models, not a definition of chiral Hochschild cohomology.
\end{remark}

\subsection{The chiral Gerstenhaber bracket from configuration spaces}

\begin{theorem}[Chiral Gerstenhaber structure on $\ChirHoch^*$;
\ClaimStatusConditional]\label{thm:gerstenhaber-structure}
\index{Gerstenhaber bracket!chiral|textbf}
The single-insertion residue sum of
Construction~\textup{\ref{const:gerstenhaber-bracket}}
induces a degree-$(-1)$ pre-Lie product
\[
\circ \colon
\ChirHoch^p(\mathcal{A}) \otimes \ChirHoch^q(\mathcal{A})
\longrightarrow
\ChirHoch^{p+q-1}(\mathcal{A}),
\]
and its antisymmetrization
\[
[f,g]_{\mathrm G}
\;:=\;
f \circ g
\;-\;
(-1)^{(|f|-1)(|g|-1)}\, g \circ f
\]
defines a degree-$(-1)$ bracket
\[
[-,-]_{\mathrm G} \colon
\ChirHoch^p(\mathcal{A}) \otimes \ChirHoch^q(\mathcal{A})
\longrightarrow
\ChirHoch^{p+q-1}(\mathcal{A}).
\]
For homogeneous classes $f,g,h$, this bracket is graded
skew-symmetric, satisfies the graded Jacobi identity, and is a
derivation of the cup product:
\[
[f, g \smile h]_{\mathrm G}
=
[f,g]_{\mathrm G}\smile h
+
(-1)^{(|f|-1)|g|}\, g \smile [f,h]_{\mathrm G}.
\]
Hence $\ChirHoch^*(\mathcal{A})$ is a Gerstenhaber algebra.
This is the chiral analogue of Gerstenhaber's classical
bracket~\cite{Ger63} on associative Hochschild cochains, realized
here by OPE residues on configuration spaces.
\end{theorem}

\begin{proof}
Choose homogeneous cocycle representatives, still denoted
$f,g,h$, in the geometric chiral Hochschild complex. By
Proposition~\ref{prop:geometric-algebraic-hochschild}, under its
logarithmic/completed finite-piece hypotheses, this geometric complex
is identified with the algebraic cochain complex
$C^\bullet_{\mathrm{ch}}(A,A)$ of
Definition~\ref{def:chiral-hochschild-cochain-brace}. In that
model the single brace $f\{g\}$ is the one-fold insertion of
Example~\ref{ex:single-brace-chiral}, and
Definition~\ref{def:gerstenhaber-bracket-chiral} defines
\[
[f,g]_{\mathrm G}
=
f\{g\}
-(-1)^{(|f|-1)(|g|-1)}g\{f\}.
\]
Construction~\ref{const:gerstenhaber-bracket} gives the same
operation geometrically:
\[
f \circ g
\;:=\;
\sum_{i=1}^{p+1} (-1)^{(q-1)(i-1)}\, f \circ_i g,
\]
where $f \circ_i g$ is the residue along the boundary divisor where
the output point of the $g$-configuration collides with the
$i$-th marked point of the $f$-configuration. Under the
comparison of Proposition~\ref{prop:geometric-algebraic-hochschild}, each
residue term $f \circ_i g$ is the $i$-th chiral insertion, so
$f \circ g = f\{g\}$. Therefore the residue bracket of
Construction~\ref{const:gerstenhaber-bracket} is exactly the
antisymmetrization of the single brace.

Proposition~\ref{prop:pre-lie-chiral} gives the pre-Lie identity
for the single brace. Rewriting it in the residue notation yields
\[
(f \circ g)\circ h - f \circ (g \circ h)
=
(-1)^{(|g|-1)(|h|-1)}
\bigl((f \circ h)\circ g - f \circ (h \circ g)\bigr).
\]
Antisymmetrizing a pre-Lie product gives a graded Lie bracket, so
$[-,-]_{\mathrm G}$ is graded skew-symmetric and satisfies the
graded Jacobi identity on cochains.

The chiral Hochschild differential is $\delta=[m,-]$ on
$C^\bullet_{\mathrm{ch}}(A,A)$. By
Theorem~\ref{thm:brace-dg-algebra}, $\delta$ is a derivation of the
brace operations. Since $f \circ g = f\{g\}$, one obtains
\[
\delta [f,g]_{\mathrm G}
=
[\delta f,g]_{\mathrm G}
+
(-1)^{|f|-1}[f,\delta g]_{\mathrm G}.
\]
Hence the bracket descends to cohomology.

For the Leibniz rule, use the same brace dg algebra together with
the cup product of Remark~\ref{rem:cup-product}. The proof of
Corollary~\ref{cor:gerstenhaber-cohomology} shows that the brace
identity~\textup{(}B1\textup{)} gives
\[
[f, g \smile h]_{\mathrm G}
=
[f,g]_{\mathrm G}\smile h
+
(-1)^{(|f|-1)|g|}\, g \smile [f,h]_{\mathrm G}
\]
on cochains, hence on cohomology. Equivalently,
Theorem~\ref{thm:chiral-deligne-tamarkin} places the pair
$(C^\bullet_{\mathrm{ch}}(A,A),A)$, whose cohomology is
$Z^{\mathrm{der}}_{\mathrm{ch}}(\mathcal{A})$, in the local chiral
Swiss-cheese operad, so the cup product and the single brace are
the closed and mixed operations on the derived-center pair.  Their
bar realization is carried by the corresponding comparison map.

For the Virasoro computation, regard the stress tensor $T$ as the unary
cochain given by insertion of the Virasoro field. Its OPE is
\[
T(z)T(w)
\sim
\frac{c/2}{(z-w)^4}
+
\frac{2T(w)}{(z-w)^2}
+
\frac{\partial T(w)}{z-w}.
\]
The residue defining the single insertion therefore extracts only
the simple pole:
\[
T \circ T
=
\operatorname{Res}_{z \to w} T(z)T(w)
=
\partial T.
\]
The opposite insertion gives the same unary cochain, so
$[T,T]_{\mathrm G}=0$.  This proves the bracket identity for the
displayed unary class.  A degree-$2$ Virasoro deformation class has
self-bracket in degree~$3$.  Bakalov--De Sole--Kac compute a
one-dimensional bounded group in that degree
\cite[Theorem~7.2]{BDSK21}; its bracket value is therefore a separate
cochain calculation, transported to the chart by
$\chi_{\operatorname{Vir}_c}^{\mathrm{bd}}$.

Therefore the residue construction defines a degree-$(-1)$
Gerstenhaber bracket on
$\ChirHoch^\bullet(\mathcal{A}) = Z^{\mathrm{der}}_{\mathrm{ch}}
(\mathcal{A})$.
\end{proof}

\begin{construction}[Geometric realization of bracket]\label{const:gerstenhaber-bracket}
The Gerstenhaber bracket has a natural geometric interpretation via configuration spaces.

For $f \in \ChirHoch^p$ and $g \in \ChirHoch^q$, represented as:
\begin{align*}
f &\in \Gamma(\overline{C}_{p+2}(X), \ldots)\\
g &\in \Gamma(\overline{C}_{q+2}(X), \ldots)
\end{align*}

For each insertion place $1 \leq i \leq p+1$, let
$f \circ_i g$ be the residue obtained by colliding the
output point of the $g$-configuration with the $i$-th marked
point of the $f$-configuration. The total single-insertion
sum
\[
f \circ g
\;:=\;
\sum_{i=1}^{p+1} (-1)^{(q-1)(i-1)}\, f \circ_i g.
\]
is the geometric pre-Lie product on chiral Hochschild cochains.
The chiral Gerstenhaber bracket is its antisymmetrization:
\[
[f,g]_{\mathrm G}
\;:=\;
f \circ g
\;-\;
(-1)^{(p-1)(q-1)}\, g \circ f.
\]

Equivalently, the construction proceeds in three stages:
\begin{enumerate}
\item \emph{Diagonal insertion}: insert one configuration into a marked point of the other
\item \emph{Summation}: Sum over all possible insertion points
\item \emph{Residue}: Extract the coefficient of singular terms
\end{enumerate}

Explicitly, the first insertion sum has the form
\[
f \circ g
=
\sum_{i=1}^{p+1} (-1)^{\epsilon_i}
\operatorname{Res}_{z_0 \to z_i}
\left[
f(\ldots, z_i, \ldots)\cdot
g(z_0, \ldots, z_q)
\right],
\]
and $g \circ f$ is defined by the same formula with $f$ and $g$
interchanged. The residue extracts the collision behavior along the
relevant boundary divisor in the Fulton--MacPherson compactification,
so $f \circ g$ is the geometric realization of the single brace
operation and $[f,g]_{\mathrm G}$ is the corresponding
Gerstenhaber bracket.
\end{construction}

\begin{example}[Support constraint on the Gerstenhaber bracket]
\label{ex:gerstenhaber-heisenberg}
Let \(\cA\) carry \(H_H(\cA;S)\), and let
\(x\in\ChirHoch^p(\cA)\), \(y\in\ChirHoch^q(\cA)\).  The bracket has
degree \(p+q-1\):
\[
 [x,y]_{\mathrm G}\in\ChirHoch^{p+q-1}(\cA).
\]
Consequently
\[
 p+q-1\notin S
 \quad\Longrightarrow\quad
 [x,y]_{\mathrm G}=0.
\]
For the rank-one even superboson, the bounded support is
\(S=\{0,1\}\); the same conclusion holds for the curve chart whenever
\(\chi_{B_{\mathfrak h}}^{\mathrm{bd}}\) is a quasi-isomorphism.
\end{example}

\subsection{\texorpdfstring{Higher structure: $L_\infty$ operations}{Higher structure: L-infinity operations}}

\begin{remark}[Beyond Gerstenhaber: completed \texorpdfstring{$L_\infty$}{L-infinity}]\label{rem:L-infty-structure}
Configuration space geometry encodes a completed $L_\infty$ structure
$\ell_n\colon \ChirHoch^*(\cA)^{\otimes n} \to \ChirHoch^{*+2-n}(\cA)$ via
\[
\ell_n(f_1, \ldots, f_n) = \int_{C_{n,1}(X)} f_1 \wedge \cdots \wedge f_n
\wedge K_n,
\]
where $K_n \in \Omega^{n-2}(\overline{C}_{n+1}(X), \mathrm{End}(\cA))$
is the propagator kernel ($K_2 = dz/(z_1-z_2)$ recovers the Gerstenhaber
bracket). Formal deformations are MC elements
$\gamma \in \ChirHoch^2(\cA)[[\hbar]]$ satisfying
$\sum_{n \geq 1} \frac{1}{n!}\ell_n(\gamma, \ldots, \gamma) = 0$.
\end{remark}

\subsection{Diagonal bar computation}

\begin{remark}[Computing \texorpdfstring{$\ChirHoch^*$}{CH^*} in practice]\label{rem:compute-HH}
For a chiral algebra given by generators and OPE data, first construct
the two-sided bar complex
$\mathcal A\otimes\bar B(\mathcal A)\otimes\mathcal A$ on the FM
configuration spaces. Then apply
$\operatorname{Hom}_{\mathcal A^e}(-,\mathcal A)$ and compute
$\ker(d)/\operatorname{im}(d)$ in each degree.  On the Koszul locus,
the comparison
$q_\mathcal A\colon\mathcal A^i\to\bar B(\mathcal A)$ is a
quasi-isomorphism.  Verdier duality gives the Koszul dual algebra
$\mathcal A^!=\mathbb D_X(\mathcal A^i)$. This separation is what
reduces the computation to the curve-level Ext complexes of Theorem~H.
\end{remark}

\begin{remark}[Noncommutative Hodge degeneration]\label{rem:nc-hodge-degeneration}
\index{noncommutative Hodge theory}
The pole-order filtration on the chiral Hochschild complex degenerates
at $E_1$ for Koszul algebras
(Theorem~\ref{thm:spectral-sequence-collapse}), giving a mixed
structure analogous to Kaledin's noncommutative Hodge
theory~\cite{Kal08}. Theorem~H supplies the curve-level Verdier
duality for chiral Hochschild cochains, exchanging the Hodge
filtrations of $\cA$ and~$\cA^!$ on the PBW-completed Koszul locus.
Theorem~C is the separate shifted-symplectic ambient complementarity
theorem: it uses this Hochschild/Verdier comparison as one input but
lives on cyclic deformation complexes and kernel formal moduli, not on
the Hochschild cochain complex alone.
\end{remark}

\subsection{Diagonal bar calculus}

The diagonal bar resolution gives explicit configuration-space
formulas for $\ChirHoch^*$, the Gerstenhaber bracket, the
$L_\infty$ operations, and the family support retract.  A perfect
degree-$d$ chain pairing supplies the separate Koszul reflection
$\ChirHoch^n(\cA)\cong
\ChirHoch^{d-n}(\cA^!)^\vee\otimes\omega_X$.  The support set~$S$
locates the Kuranishi targets in the shifted deformation complex;
higher transferred brackets are governed by the brace-compatible
formality input of Proposition~\ref{prop:e2-formality-hochschild}.

\medskip

The deformation-obstruction theory developed above applies to the
PBW Koszul chiral algebras in the stated ambient. For every
major family of vertex algebras (Heisenberg, free fermion, $\beta\gamma$-$bc$,
affine Kac--Moody, Virasoro, $W$-algebras, Yangians) we compute the bar
complex, identify the Koszul dual, determine the obstruction coefficient
$\kappa(\cA)$, and prove the complementarity and genus universality theorems
explicitly. The shared discriminant phenomenon
(Theorem~\ref{thm:ds-bar-gf-discriminant}) and the $\hat{A}$-genus appearance
in genus expansions emerge as consequences of the modular characteristic
package, not as isolated coincidences.

%================================================================
% COSIMPLICIAL STRUCTURE AND HOMOTOPY CHIRAL ALGEBRAS
%================================================================

\section{Cosimplicial descent and the chiral Hochschild--\v{C}ech comparison}
\label{sec:cosimplicial-hochschild}
\index{cosimplicial!Hochschild complex}
\index{homotopy chiral algebra!Hochschild perspective}

The chiral Hochschild complex $C^\bullet_{\mathrm{chiral}}(\cA)$
is built from configuration spaces of points on~$X$: the complex
in cochain degree~$n$ involves $(n+2)$ marked points
(\S\ref{sec:hochschild-via-bar-cobar-complete}). When the points are
restricted to lie in a fixed affine open, the complex degenerates
to local data. Globalizing via an open cover introduces a second
direction of complexity, the \v{C}ech direction, and the
interplay between the chiral Hochschild and \v{C}ech differentials is
controlled by the homotopy chiral algebra structure of
Malikov--Schechtman~\cite{MS24}.

\begin{remark}[Eilenberg--Zilber operad as genus-$0$ convolution control]
\label{rem:ez-operad-convolution}
\index{Eilenberg--Zilber operad!convolution interpretation}
The Eilenberg--Zilber operad
$\mathcal{Y}(n) = \mathcal{Z}(n) \otimes_k \mathrm{Lie}(n)$
of \cite{MS24} is the genus-$0$ projection of the modular
convolution algebra
$\mathfrak{g}^{\mathrm{mod}}_\cA$
(Definition~\ref{def:modular-convolution-dg-lie}). The key
identifications are:
\begin{enumerate}[label=\textup{(\roman*)}]
\item The Eilenberg--Zilber chain complex $\mathcal{Z}(n) =
 \operatorname{Tot}\operatorname{Hom}_\Delta(Z, Z^{\otimes n})$
 encodes the simplicial combinatorics of $n$-fold intersections
 in a cover. It is acyclic: $H^i(\mathcal{Z}(n)) = 0$ for
 $i \neq 0$, $H^0(\mathcal{Z}(n)) = k$
 (\cite[\S2.5]{MS24}, originating in Dold~\cite{Dold61}).
\item At genus~$0$, the role of $\mathcal{Z}(n)$ is equivalent
 to that of the cellular chains on the associahedron
 $C_*(K_n)$, which in turn model
 $C_*(\overline{\cM}_{0,n+1})$. The augmentation
 $\varepsilon\colon \mathcal{Z}(n) \to k$ and the cellular
 augmentation $C_*(K_n) \to k$ are quasi-isomorphisms controlling
 the same homotopy-coherent data.
\item At genus $g \geq 1$, the role of $\mathcal{Z}(n)$ is
 replaced by $C_*(\overline{\cM}_{g,n+1})$, and the HCA
 structure extends to the modular Feynman transform. The
 genus-$0$ Eilenberg--Zilber operations $F_n$ are the tree-level
 truncation of the modular $L_\infty$ brackets $\ell_n^{(0)}$
 (Construction~\ref{constr:explicit-convolution-linfty}).
\item The Malikov--Schechtman operadic model
 ($\tau_{\leq 0}\mathcal{Y}$-algebras, \cite[\S2.6]{MS24}) and
 the Loday--Vallette coalgebraic model (codifferentials on cofree
 coalgebras, \cite[\S10.3]{LV12}) are equivalent over
 characteristic-$0$ fields: both encode $L_\infty$ structures
 with the same MC moduli
 (Remark~\ref{rem:conv-functoriality-infrastructure}).
\end{enumerate}
\end{remark}

\subsection{The bicomplex}
\label{subsec:hochschild-cech-bicomplex}

Fix an affine cover $\mathcal{U} = \{U_i\}$ of~$X$. Define
the bicomplex
\[
E_0^{p,q}
= \bigoplus_{i_0 < \cdots < i_p}
 C^q_{\mathrm{chiral}}(\cA|_{U_{i_0\cdots i_p}})
\]
with horizontal differential $d_{\check{C}}$ (\v{C}ech) and
vertical differential $d_{\mathrm{Hoch}}$ (chiral Hochschild,
Theorem~\ref{thm:chiral-hochschild-differential}). Both
differentials square to zero individually, but they do not
strictly anticommute: the failure is measured by the \v{C}ech
Jacobiator (Definition~\ref{def:jacobiator}).

\begin{proposition}[chiral Hochschild--\v{C}ech spectral sequence;
\ClaimStatusProvedHere]
\label{prop:hochschild-cech-ss}
The bicomplex $E_0^{p,q}$ gives rise to two spectral sequences:
\begin{enumerate}[label=\textup{(\roman*)}]
\item \emph{\v{C}ech-first:} $E_2^{p,q} =
 \check{H}^p(X, \mathcal{H}^q_{\mathrm{ch}}(\cA))
 \Rightarrow \ChirHoch^{p+q}(\cA)$,
 where $\mathcal{H}^q_{\mathrm{ch}}(\cA)$ is the sheaf of
 local chiral Hochschild cohomology.
\item \emph{chiral-Hochschild-first:} $E_2^{p,q} =
 H^q_{\mathrm{ch}}(\check{C}^p(\mathcal{U};\cA))
 \Rightarrow \ChirHoch^{p+q}(\cA)$.
\end{enumerate}
Both converge for quasi-coherent $\cA$ on a separated scheme~$X$
(by the standard bounded-below complete convergence theorem).
\end{proposition}

\begin{proof}
The convergence follows from the vanishing of \v{C}ech cohomology
of quasi-coherent sheaves on affines: for any affine
$U_{i_0\cdots i_p}$, the local chiral Hochschild complex computes
$H^q_{\mathrm{ch}}(\cA|_U)$ and the \v{C}ech differentials are
well-defined on cohomology sheaves. Both spectral sequences are
bounded in the horizontal direction (the cover is finite) and
bounded below in the vertical direction.
\end{proof}

\subsection{HCA enhancement of the bicomplex}
\label{subsec:hca-hochschild}

The \v{C}ech-first spectral sequence sees only the \v{C}ech
cohomology of the cohomology sheaves; it does not see the
\emph{algebraic structure} on the resolution.
Malikov--Schechtman's theorem
(Theorem~\ref{thm:cech-hca}) enriches the \v{C}ech direction
with HCA operations $\{F_n\}$, which produce a stronger
convergence statement.

\begin{conjecture}[HCA structure on local chiral Hochschild cochains;
\ClaimStatusConjectured]
\label{conj:hca-hochschild}
The \v{C}ech complex of the presheaf
$U \mapsto C^\bullet_{\mathrm{chiral}}(\cA|_U)$
carries an HCA MC element
$\Phi_{\ChirHoch} \in
\operatorname{MC}(\mathfrak{g}^{\check{C}}_{\ChirHoch})$
in the \v{C}ech convolution algebra
\textup{(}Definition~\textup{\ref{def:cech-convolution})},
with components:
\begin{enumerate}[label=\textup{(\roman*)}]
\item $F_2$ is the brace operation
 (Theorem~\ref{thm:chiral-gerstenhaber-kps}$\,$), encoding the
 Gerstenhaber bracket;
\item $F_3$ is the Jacobiator homotopy for the brace bracket
 (nonzero only when $\check{C}^2 \neq 0$, i.e., for covers
 with three or more overlapping opens);
\item the full MC element $\Phi_{\ChirHoch}$
 satisfies $D\Phi + \tfrac{1}{2}[\Phi,\Phi] = 0$, whose
 projections to each degree give the $L_\infty$ relations
 on $\ChirHoch^\bullet(\cA)$.
\end{enumerate}
In particular, the Gerstenhaber algebra structure on
$\ChirHoch^\bullet(\cA)$ is the cohomological shadow of the
MC element $\Phi_{\ChirHoch}$: the bracket is the degree-$2$
projection, and Jacobi holds on cohomology because $F_3$ is
exact.
\end{conjecture}

\begin{remark}[Comparison with the modular deformation complex]
\label{rem:hca-vs-modular-deformation}
The HCA structure on the chiral Hochschild--\v{C}ech bicomplex is the
genus-$0$ shadow of the modular cyclic deformation complex
$\mathrm{Def}_{\mathrm{cyc}}^{\mathrm{mod}}(\cA)$
(Definition~\ref{def:modular-cyclic-deformation-complex}).
At genus~$0$, the operadic control is via the Eilenberg--Zilber
operad $\mathcal{Y}(n)$ acting on the \v{C}ech direction.
At genus~$g \geq 1$, the modular operad replaces
$\mathcal{Y}(n)$ with the Feynman transform $\mathrm{FCom}$,
and the cosimplicial identities that make $d_{\check{C}}^2 = 0$
are replaced by the boundary relations on
$\overline{\mathcal{M}}_{g,n}$ that make
$d_{\mathrm{fib}}^2 = \kappa \cdot \omega_g$.
At each finite order, the shadow obstruction tower
(Definition~\ref{def:shadow-postnikov-tower}) unifies
both the genus-$0$ HCA relations and the higher-genus
Maurer--Cartan structure; the all-degree limit
$D\Theta + \tfrac{1}{2}[\Theta,\Theta] = 0$
on $\mathrm{Def}_{\mathrm{cyc}}^{\mathrm{mod}}(\cA)$
is the full master equation, proved by the bar-intrinsic
construction (Theorem~\ref{thm:mc2-bar-intrinsic}).
\end{remark}

%================================================================
% AMBIENT COMPLEMENTARITY
%================================================================

\section{Ambient complementarity: the shifted-symplectic theorem}
\label{sec:ambient-complementarity}
\index{complementarity!ambient|textbf}
\index{shifted symplectic!ambient complementarity|textbf}
\index{Lagrangian!complementary|textbf}

On the finite-window PBW/perfect locus of Theorem~C, the genus-$g$
realization contains complementary summands
$Q_g(\cA) \oplus Q_g(\cA^!)
\subset R\Gamma(\overline{\mathcal{M}}_g, Z(\cA))$.
For non-generic class-\(\mathsf M\) simple quotients, critical affine
level, and unrestricted product completions, this displayed inclusion
is not the theorem; the ambient object is a completed or coderived
replacement equipped with the Mittag--Leffler and perfectness
hypotheses used in Theorem~H. The Lagrangian refinement has two
layers. The genus-$g$ Verdier pairing gives the proved complementary
Lagrangian polarization of the two summands on that finite-window
locus. The ambient cyclic deformation problem gives a conditional
$(-1)$-shifted symplectic formal moduli problem with Lagrangian maps
from the two dual packages, under the perfectness, nondegeneracy, and
uniform-weight BV hypotheses stated below.

This section starts from cyclic deformation complexes and the tangent
fibre of the universal kernel. Theorem~H identifies and dualizes
$Z^{\mathrm{der}}_{\mathrm{ch}}(\cA)=\ChirHoch^\bullet(\cA,\cA)$ on
the PBW-completed curve-level chart. The cochain-amplitude theorem is
the curve-level input. Theorem~C adds the cyclic pairing, the adjoint
kernel maps, and the second-variation form needed to place the two
deformation packages as complementary Lagrangians in the common
ambient problem.

\begin{definition}[Kernel deformation complex]
\label{def:kernel-deformation-complex}
\index{twisting morphism!kernel deformation}
Let $(\cA, \cA^!)$ be a cyclic modular Koszul chiral pair
on~$X$, and let
$\tau_\cA \in \mathrm{MC}(\cA^! \,\widehat{\otimes}\, \cA)$
be the universal twisting kernel
(Convention~\ref{conv:bar-coalgebra-identity}). Define
\[
K_\cA := (\cA^! \,\widehat{\otimes}\, \cA)[1],
\]
with differential twisted by~$\tau_\cA$. The linearizations
of~$\tau_\cA$ in the $\cA$- and $\cA^!$-directions give maps
\[
\nabla_\cA \colon \mathrm{Def}_{\mathrm{cyc}}(\cA)
\to K_\cA,
\qquad
\nabla_{\cA^!} \colon \mathrm{Def}_{\mathrm{cyc}}(\cA^!)
\to K_\cA.
\]
\end{definition}

\begin{definition}[Ambient complementarity tangent complex; \ClaimStatusDefinitional]
\label{def:ambient-tangent-complex}
Define
\begin{equation}\label{eq:ambient-tangent}
T_{\mathrm{comp}}(\cA)
:= \mathrm{fib}\bigl(
 \nabla_\cA - \nabla_{\cA^!} \colon
 \mathrm{Def}_{\mathrm{cyc}}(\cA) \oplus
 \mathrm{Def}_{\mathrm{cyc}}(\cA^!)
 \longrightarrow K_\cA
\bigr).
\end{equation}
A tangent vector is a simultaneous first-order deformation
of~$\cA$, of~$\cA^!$, and of~$\tau_\cA$, constrained by the
linearized Maurer--Cartan equation. The formal moduli problem
of triples $(\cA', \alpha, B')$ with $\cA'$ a deformation
of~$\cA$, $B'$ a deformation of~$\cA^!$, and
$\alpha \in \mathrm{MC}(B' \widehat{\otimes} \cA')$ a
deformation of~$\tau_\cA$, has tangent
complex~$T_{\mathrm{comp}}(\cA)$.
\end{definition}

\begin{proposition}[Self-duality of the kernel fibre; \ClaimStatusProvedHere]
\label{prop:ambient-self-duality}
Let
$D_\cA=\mathrm{Def}_{\mathrm{cyc}}(\cA)$ and
$D_{\cA^!}=\mathrm{Def}_{\mathrm{cyc}}(\cA^!)$ be perfect
filtered cyclic complexes. Assume:
\begin{enumerate}[label=\textup{(\roman*)}]
\item the cyclic/Koszul residue pairing identifies
 $D_{\cA^!}$ with $D_\cA^\vee[-1]$;
\item $K_\cA$ carries a perfect residue pairing compatible with
 the involution exchanging the two tensor factors
 $\cA^!$ and~$\cA$;
\item the linearized kernel maps
 $\nabla_\cA$ and $\nabla_{\cA^!}$ are adjoint for these
 pairings, with the Koszul sign dictated by the shift in
 $K_\cA=(\cA^!\widehat{\otimes}\cA)[1]$.
\end{enumerate}
Then the homotopy fibre
$T_{\mathrm{comp}}(\cA)$ in
\eqref{eq:ambient-tangent} carries a canonical
nondegenerate pairing of degree~$-1$.
\end{proposition}

\begin{proof}
Write the fibre as the cone model
$T_{\mathrm{comp}}(\cA)=
(D_\cA\oplus D_{\cA^!})\oplus K_\cA[-1]$ with differential
twisted by $\nabla_\cA-\nabla_{\cA^!}$. Pair the two outer
summands by the cyclic/Koszul pairing and pair the shifted
kernel summands by the residue pairing on~$K_\cA$, with the
standard cone sign. The adjointness hypothesis is precisely the
identity saying that the twisted cone differential is
skew-adjoint. The pairing therefore descends to a closed
degree-$-1$ pairing on the fibre. Perfectness of the two outer
complexes and of~$K_\cA$ makes the induced map
$T_{\mathrm{comp}}(\cA)\to T_{\mathrm{comp}}(\cA)^\vee[-1]$
a quasi-isomorphism.
\end{proof}

\begin{proposition}[One-sided isotropy criterion; \ClaimStatusProvedHere]
\label{prop:one-sided-isotropy}
Assume the hypotheses of
Proposition~\textup{\ref{prop:ambient-self-duality}}. Suppose in
addition that the universal twisting kernel satisfies the
second-variation identities
\[
\big\langle \nabla_\cA u,\nabla_\cA v\big\rangle_{K_\cA}
\simeq 0,
\qquad
\big\langle \nabla_{\cA^!} u',\nabla_{\cA^!} v'
\big\rangle_{K_\cA}
\simeq 0
\]
for tangent vectors
$u,v\in D_\cA$ and
$u',v'\in D_{\cA^!}$. Then the one-sided deformation maps
\[
i_\cA \colon \mathcal{M}_\cA \to \mathcal{M}_{\mathrm{comp}}(\cA),
\qquad
i_{\cA^!} \colon \mathcal{M}_{\cA^!}
 \to \mathcal{M}_{\mathrm{comp}}(\cA),
\]
obtained by varying only~$\cA$ or only~$\cA^!$ together with the
induced first-order variation of~$\tau_\cA$, are isotropic.
\end{proposition}

\begin{proof}
The one-sided deformation maps
have tangent complexes represented by the graphs
$u\mapsto (u,0,\nabla_\cA u)$ and
$u'\mapsto (0,u',-\nabla_{\cA^!}u')$ in the cone model for
$T_{\mathrm{comp}}(\cA)$. On the first graph the outer
cyclic/Koszul summand pairs $D_\cA$ only with
$D_{\cA^!}$, so it contributes zero. The remaining term is
$\langle\nabla_\cA u,\nabla_\cA v\rangle_{K_\cA}$, which is
null-homotopic by the first second-variation identity. The
second graph is identical with $\cA$ and~$\cA^!$ exchanged.
\end{proof}

\begin{theorem}[Ambient complementarity criterion; \ClaimStatusConditional]
\label{thm:ambient-complementarity}
\index{ambient complementarity theorem|textbf}
Assume the perfect cyclic/Koszul pairings, kernel residue
pairing, adjointness, and second-variation identities of
Propositions~\textup{\ref{prop:ambient-self-duality}} and
\textup{\ref{prop:one-sided-isotropy}}. Assume also:
\begin{enumerate}[label=\textup{(\alph*)}]
\item the quotient of $T_{\mathrm{comp}}(\cA)$ by the
 one-sided tangent of~$\mathcal M_\cA$ is identified by
 bar--cobar inversion and Verdier duality with
 $T_{\mathcal M_{\cA^!}}^\vee[-1]$, and symmetrically;
\item the filtered cyclic $L_\infty$ enhancement integrates
 $T_{\mathrm{comp}}(\cA)$ to a pointed formal moduli problem
 carrying the closed $(-1)$-shifted $2$-form represented by
 the pairing above;
\item the genus-$g$ realization functor sends the cyclic
 kernel pairing to the BV/Verdier pairing used in
 Theorem~\textup{\ref{thm:quantum-complementarity-main}}.
\end{enumerate}
Then:
\begin{enumerate}[label=\textup{(\roman*)}]
\item $T_{\mathrm{comp}}(\cA)$ is canonically self-dual
 of degree~$-1$
 \textup{(}Proposition~\textup{\ref{prop:ambient-self-duality}}\textup{)}.
\item The one-sided deformation maps define complementary
 Lagrangian subcomplexes inside~$T_{\mathrm{comp}}(\cA)$
 \textup{(}Proposition~\textup{\ref{prop:one-sided-isotropy}}
 gives isotropy; assumption \textup{(a)} gives maximality\textup{)}.
\item The filtered cyclic $L_\infty$ enhancement defines a
 pointed $(-1)$-shifted symplectic formal moduli problem
 $\mathcal{M}_{\mathrm{comp}}(\cA)$, and the induced maps
 $\mathcal{M}_\cA \to \mathcal{M}_{\mathrm{comp}}(\cA)
 \leftarrow \mathcal{M}_{\cA^!}$ are Lagrangian.
\item The normal complex to~$\mathcal{M}_\cA$ in
 $\mathcal{M}_{\mathrm{comp}}(\cA)$ is quasi-isomorphic to
 $T_{\mathcal{M}_{\cA^!}}^\vee[-1]$, and symmetrically
 with $\cA$ and~$\cA^!$ exchanged.
\item After applying the genus-$g$ realization functor, one
 recovers the proved Lagrangian polarization
 $Q_g(\cA),\, Q_g(\cA^!)
 \subset R\Gamma(\overline{\mathcal{M}}_g, Z(\cA))$
 of Theorem~\textup{\ref{thm:quantum-complementarity-main}}.
\end{enumerate}
\end{theorem}

\begin{proof}
The ambient deformation functor is the homotopy pullback of the
deformation functors of~$\cA$, $\cA^!$, and the
Maurer--Cartan functor of~$\tau_\cA$; differentiating this
pullback gives the fibre complex
$T_{\mathrm{comp}}(\cA)$. Part~\textup{(i)} is
Proposition~\ref{prop:ambient-self-duality}. Part~\textup{(ii)}
combines Proposition~\ref{prop:one-sided-isotropy} with
assumption~\textup{(a)}, which identifies each normal complex
with the shifted dual tangent complex of the opposite side.
Assumption~\textup{(b)} is exactly the integration datum required
by the formal moduli/shifted symplectic correspondence, giving
part~\textup{(iii)}. Part~\textup{(iv)} restates the
normal-complex identification in assumption~\textup{(a)}.
Applying the genus-$g$ realization functor and then
assumption~\textup{(c)} gives the BV/Verdier Lagrangian
polarization of Theorem~\ref{thm:quantum-complementarity-main},
which is part~\textup{(v)}.
\end{proof}

\begin{definition}[Complementarity potential: cotangent normal form]
\label{def:complementarity-potential-cotangent}
\index{complementarity potential|textbf}
The shifted cotangent normal form
(formal Darboux lemma for $(-1)$-shifted symplectic stacks)
identifies $\mathcal{M}_{\mathrm{comp}}(\cA)$ with
$T^*[-1]\mathcal{M}_\cA$ near the base point. Under this
identification, $\mathcal{M}_{\cA^!}$ is the graph of a
closed $1$-form $\alpha_\cA = dS_\cA$ for a formal function
\[
S_\cA \in \widehat{\mathcal{O}}(\mathcal{M}_\cA)^0,
\]
the \emph{complementarity potential} of the Koszul pair
$(\cA, \cA^!)$.
\end{definition}

\begin{proposition}[Derived critical locus in a shifted cotangent chart;
\ClaimStatusConditional]
\label{prop:derived-critical-locus-chk}
Assume the formal moduli problem of
Theorem~\textup{\ref{thm:ambient-complementarity}} admits the
shifted cotangent normal form of
Definition~\textup{\ref{def:complementarity-potential-cotangent}},
and that $\mathcal{M}_{\cA^!}$ is the graph of the exact closed
$1$-form $dS_\cA$ in this chart. Then there is a canonical
equivalence
\[
\mathcal{M}_\cA \times^h_{\mathcal{M}_{\mathrm{comp}}(\cA)}
\mathcal{M}_{\cA^!}
\;\simeq\;
\mathrm{Crit}(S_\cA).
\]
The simultaneous deformations of~$\cA$ and~$\cA^!$ that
preserve complementarity are governed by the derived
critical locus of the complementarity potential.
\end{proposition}

\begin{proof}
Under the shifted cotangent normal form,
$\mathcal{M}_{\mathrm{comp}}(\cA) \simeq T^*[-1]\mathcal{M}_\cA$
near the base point, $\mathcal{M}_\cA$ is the zero section, and
$\mathcal{M}_{\cA^!}$ is the graph of~$dS_\cA$. For a smooth
formal moduli problem~$Y$ and a formal function~$S$, the derived
intersection of the zero section $Z$ with the graph
$\Gamma_{dS}$ inside $T^*[-1]Y$ is
\[
Z \times^h_{T^*[-1]Y} \Gamma_{dS}
\;\simeq\;
\operatorname{Spec}\bigl(
\operatorname{Sym}(\Omega_Y[1]),\; dS\bigr)
\;=\;
\mathrm{Crit}(S),
\]
the Koszul complex of~$dS$
\textup{(}cf.~\cite[Proposition~2.1]{PTVV13}\textup{)}.
Taking $Y = \mathcal{M}_\cA$ and $S = S_\cA$ gives the stated
equivalence.
\end{proof}

\begin{proposition}[Formal Legendre duality; \ClaimStatusConditional]
\label{prop:formal-legendre}
Assume the two one-sided Lagrangians of
Theorem~\textup{\ref{thm:ambient-complementarity}} admit exact
shifted cotangent charts
$T^*[-1]\mathcal{M}_\cA$ and
$T^*[-1]\mathcal{M}_{\cA^!}$ near the base point, and that the
transition between the charts is the induced exact
symplectomorphism. Let $S_\cA$ and $S_{\cA^!}$ be the two
generating functions. Then $S_\cA$ and~$S_{\cA^!}$ are formal
Legendre transforms of one another.
\end{proposition}

\begin{proof}
Both generating functions describe the same Lagrangian subspace
inside $\mathcal{M}_{\mathrm{comp}}(\cA)$, viewed from opposite
shifted cotangent charts. In exact shifted symplectic linear
algebra, changing the zero section from one complementary
Lagrangian to the other replaces the generating function by the
formal Legendre transform. The asserted relation is precisely this
chart transition.
\end{proof}

\begin{remark}[Fake complementarity criterion]
\label{rem:fake-complementarity}
The complementarity potential $S_\cA$ detects genuine
complementarity: $S_\cA$ is quadratic if and only if the
dual Lagrangian $\mathcal{M}_{\cA^!}$ is linear in
formal Darboux coordinates, which holds if and only if the
derived critical locus $\mathrm{Crit}(S_\cA)$ is governed
entirely by its Hessian complex. The cubic and higher Taylor
coefficients of~$S_\cA$ are the first obstructions to fake
complementarity, the regime where an eigenspace splitting
mimics Lagrangian geometry without being one. The
Hessian $\mathrm{Hess}_0(S_\cA)$ is exactly the linear
complementarity isomorphism
$T_0\mathcal{M}_\cA \xrightarrow{\sim}
T_0\mathcal{M}_{\cA^!}^\vee[-1]$.
\end{remark}

\begin{remark}[Modular extension and the true invariant]
\label{rem:complementarity-potential-modular}
Once the universal Maurer--Cartan class~$\Theta_\cA$ is
constructed and shown to satisfy clutching and Verdier
compatibilities, the closed $1$-forms~$dS_\cA$ inherit
the same compatibilities. After choosing pointed
primitives, the complementarity potentials globalize over
the modular operad; their scalar trace recovers the
established tower
$\kappa(\cA) \cdot \lambda_g$ (uniform-weight). The true invariant is not
the additive decomposition
$Q_g(\cA) \oplus Q_g(\cA^!)$ but the formal generating
function~$S_\cA$. Its Hessian is the linear
complementarity (already proved). Its cubic term is the
first place where genuine complementarity departs from a
fake linear splitting. A natural target: compute the
cubic term of~$S_\cA$ for Heisenberg,
$\widehat{\mathfrak{sl}}_2$, and $\beta\gamma$.
\end{remark}

\begin{remark}[Nonlinear shadow calculus]
\label{rem:nonlinear-shadow-pointer}
The first nonlinear jets of~$S_\cA$ are developed systematically
in Chapter~\ref{app:nonlinear-modular-shadows}, where the quartic
shadow cocycle
$\Theta_\cA^{\le 4} = H_\cA + \mathfrak C_\cA + \mathfrak Q_\cA$
is constructed, its master and clutching equations are proved,
and the three frame families are shown to realize the Gaussian,
Lie/tree, and contact/quartic archetypes respectively.
The modular quartic resonance class~$\mathfrak R_{4,g,n}^{\mathrm{mod}}(\cA)$
constructed there is the first Picard-valued nonlinear characteristic of the theory.
\end{remark}

\begin{remark}[Branch-line reductions]
\label{rem:branch-line-pointer}
Appendix~\ref{app:branch-line-reductions} extracts exact
finite-dimensional quotients of the full modular deformation problem: the
scalar Maurer--Cartan skeleton, the spectral cumulant hierarchy, and
rank-two branch-line reductions. On every one-line branch quotient
with scalar square, genus-$2$ is transparent and the first primitive
traceless coefficient begins in genus~$3$ with the universal
constant~$7$ on the shared
$\widehat{\mathfrak{sl}}_2$/Virasoro/$\beta\gamma$ spectral sheet.
\end{remark}

\begin{remark}[The categorical logarithm and its Taylor jets]
\label{rem:categorical-logarithm-jets}
The ambient complementarity section, together with the two preceding
pointer remarks, articulates a single hierarchy rooted in the
categorical logarithm paradigm
(Remark~\ref{rem:categorical-logarithm-paradigm}). The
complementarity potential $S_\cA$ is the connected generating
function of the categorical logarithm applied to the Koszul pair
$(\cA,\cA^!)$. Its Taylor jets have precise identities:
\begin{itemize}
\item the Hessian $H_\cA$ is the quadratic term, the proved linear
 complementarity of Theorem~C;
\item the cubic coefficient $\mathfrak C_\cA$ is the first connected
 cumulant, the first obstruction to fake complementarity, whose
 cohomology class $[\mathfrak C_\cA]$ is a canonical modular
 invariant;
\item the quartic resonance class
 $\mathfrak R_{4,g,n}^{\mathrm{mod}}(\cA)$ is the first
 Picard-valued nonlinear characteristic, with clutching correction
 governed by the cubic tree divisor.
\end{itemize}
The branch-line reductions
(Appendix~\ref{app:branch-line-reductions}) are the spectral
cumulants of the logarithm restricted to one-dimensional
sub-problems: the universal genus-$3$ constant~$7$ is the first
primitive traceless cumulant on the shared
$\widehat{\mathfrak{sl}}_2$/Virasoro/$\beta\gamma$ spectral sheet.
The canonical homotopy degree of the subregular hook family
is the canonical
Taylor degree of the logarithm in the non-principal
$\mathcal W$-algebra sector: it measures how many Taylor jets are
needed before the duality structure stabilizes.
\end{remark}


%% ===========================================================
%% FACTORIZATION ENVELOPE SHADOW FUNCTOR
%% ===========================================================

\section{The factorization envelope shadow functor}
\label{sec:envelope-shadow-functor}

\index{factorization envelope!shadow functor|textbf}
\index{shadow tower!factorization envelope|textbf}
\index{Lie conformal algebra!shadow data|textbf}

A Lie conformal algebra~$R$ over a smooth curve~$X$ determines,
via the factorization envelope construction
\textup{(}\cite[\S3.7]{BD04}, \cite{CG17}\textup{)}, a vertex algebra
$V(R)$, the enveloping vertex algebra of~$R$.
The Maurer--Cartan shadow obstruction tower of~$V(R)$
(Definition~\ref{def:shadow-algebra}) is an invariant of the
vertex algebra; the content of the following proposition is that
this invariant is \emph{already determined at the level of the
Lie conformal input~$R$}, without passing through the full
vertex algebra.

\begin{proposition}[Factorization envelope shadow functor; \ClaimStatusConditional]
\label{prop:envelope-shadow}
\index{factorization envelope!shadow determination|textbf}
\index{modular characteristic!factorization envelope}
\index{shadow depth!from Lie conformal data}
Let $R$ be a cyclically admissible Lie conformal algebra
\textup{(}Definition~\textup{\ref{def:v1-cyclically-admissible})},
and let $V(R)$ denote its factorization envelope.
\begin{enumerate}[label=\textup{(\roman*)}]
\item \textup{(Modular characteristic from $\lambda$-bracket.)}
 The modular characteristic $\kappa(V(R))$ is determined by
 the degree-$2$ part of the bar complex of $V(R)$, which depends
 only on the self-pairing data of~$R$:
 \[
 \kappa(V(R))
 \;=\;
 \kappa_R,
 \]
 where $\kappa_R$ is determined by the $\lambda$-bracket
 $[{-}_\lambda\, {-}]$ and the invariant bilinear form of~$R$.
 % A universal closed-form
 % supertrace formula for kappa(V(R)) from the lambda-bracket alone is
 % not available; kappa depends on the full Sugawara denominator (k+h^v),
 % not just the lambda-bracket coefficients. The explicit family-by-family
 % values in the table below are computed by the compute layer
 % and by complementarity). A universal formula requires specifying
 % the Sugawara construction for each family, which is the content of
 % Theorem D.
 Explicitly:
 \begin{center}
 \renewcommand{\arraystretch}{1.25}
 \begin{tabular}{lll}
 Lie conformal $R$ & Envelope $V(R)$ & $\kappa_R$ \\
 \hline
 Free $\operatorname{Lie}\{x_h\}$ at level~$k$, $h = 1$
 & Heisenberg $\mathfrak{H}_k$ & $k$ \\
 $\operatorname{Cur}(\mathfrak{g})$ at level~$k$
 & $V_k(\mathfrak{g})$
 & $\dim(\mathfrak{g})\,(k + h^\vee)/(2h^\vee)$ \\
 Virasoro LCA at $c$
 & $\mathrm{Vir}_c$
 & $c/2$
 \end{tabular}
 \end{center}

\item \textup{(Shadow depth from Lie conformal structure.)}
 The shadow depth class
 \textup{(}Definition~\textup{\ref{def:shadow-depth-classification})} of
 $V(R)$ is determined by the structural complexity of the
 $\lambda$-bracket of~$R$:
 \begin{itemize}
 \item Class~G \textup{(}$r_{\max} = 2$\textup{)}: $R$ abelian
 \textup{(}$[a_\lambda\, b] = k\,(a, b)\, \lambda$ for some
 pairing\textup{)}. Envelope = free field.
 \item Class~L \textup{(}$r_{\max} = 3$\textup{)}: $R = \operatorname{Cur}(\mathfrak{g})$
 for a simple Lie algebra~$\mathfrak{g}$
 \textup{(}the Lie bracket of~$\mathfrak{g}$ generates cubic
 shadows via tree-level graphs\textup{)}.
 \item Class~C \textup{(}$r_{\max} = 4$\textup{)}: $R$ has a
 generator of weight~$\geq 2$ with nonzero quartic contact
 invariant but no higher forced shadows
 \textup{(}$\beta\gamma$ system\textup{)}.
 \item Class~M \textup{(}$r_{\max} = \infty$\textup{)}:
 $V(R)$ has a higher-pole OPE \textup{(}pole order $\geq 4$
 in the self-OPE of a generator\textup{)}, forcing an infinite
 shadow obstruction tower. Virasoro and $\cW_N$ for $N \geq 2$.
 \end{itemize}

\item \textup{(Functoriality.)}
 The assignment $R \mapsto (\kappa_R,\, \text{shadow class})$ is
 \emph{functorial}: a morphism $f \colon R_1 \to R_2$ of Lie
 conformal algebras induces a compatible map on shadow data.
 In particular:
 \begin{enumerate}[label=\textup{(\alph*)}]
 \item For a subalgebra inclusion
 $\operatorname{Cur}(\mathfrak{h}) \hookrightarrow
 \operatorname{Cur}(\mathfrak{g})$ at the same level~$k$,
 the shadow map sends $\kappa_{\mathfrak{h}}$ to
 $\kappa_{\mathfrak{g}}$.
 \item For the Drinfeld--Sokolov reduction
 $\operatorname{Cur}(\mathfrak{g}) \twoheadrightarrow R_f$
 \textup{(}with $R_f$ the $\cW$-algebra Lie conformal
 data\textup{)}, the shadow map sends
$\kappa(\widehat{\mathfrak{g}}_k)$ to
$\kappa(\cW_k(\mathfrak{g}, f)) = \rho(\mathfrak{g})\, c_{\cW}$,
where $\rho(\mathfrak{g}) = H_N - 1$ is the anomaly ratio
for the principal type~$A_{N-1}$ algebra $\mathfrak{sl}_N$.
 \end{enumerate}
\end{enumerate}
\end{proposition}

\begin{proof}
\textbf{Part~(i).}
The modular characteristic $\kappa(V(R))$ is defined as the
degree-$2$ scalar projection of the Maurer--Cartan element
$\Theta_{V(R)}$ in the modular convolution algebra
$\gAmod$
(Definition~\ref{def:modular-convolution-dg-lie}).
At degree~$2$, the bar complex $B(V(R))_2$ involves exactly the
binary OPE of~$V(R)$, which is determined by the
$\lambda$-bracket of~$R$ together with the symmetric invariant
form. The factorization envelope construction
\cite[\S3.7]{BD04} builds $V(R)$ as the universal
vertex algebra generated by~$R$ with OPE given by the
$\lambda$-bracket. No higher-degree vertex algebra data enters
the degree-$2$ bar component.

The explicit values follow from the modular characteristic
computations proved in the body of the monograph.
For the current algebra $\operatorname{Cur}(\mathfrak{g})$
at level~$k$, the weight-$1$ space is $\mathfrak{g}$ with the
Sugawara correction, giving
$\kappa = \dim(\mathfrak{g})(k + h^\vee)/(2h^\vee)$
(Theorem~\ref{thm:modular-characteristic}).
For the free Lie conformal algebra on one weight-$1$ generator
at level~$k$ (Heisenberg at level~$k$): $\kappa = k$
(Chapter~\ref{chap:free-fields}).
For Virasoro: the weight-$1$ space is empty, and the
anomaly $\kappa = c/2$ comes from the weight-$2$ generator via
the $d\log$ pole-order reduction
(the bar propagator has weight~$1$;
Remark~\ref{rem:propagator-weight-universality}).

\textbf{Part~(ii).}
The shadow depth of $V(R)$ is determined by the vanishing
pattern $\{S_r(V(R))\}_{r \geq 2}$
(Definition~\ref{def:shadow-depth-classification}).
For an abelian $R$ (class~G): the $\lambda$-bracket has no Lie
component, so the OPE of $V(R)$ is purely bilinear.
The bar complex is formal ($A_\infty$ operations $m_r = 0$ for
$r \geq 3$ on bar cohomology), and all shadow invariants $S_r$
for $r \geq 3$ vanish.

For $R = \operatorname{Cur}(\mathfrak{g})$ (class~L): the Lie
bracket of~$\mathfrak{g}$ generates a nonzero cubic shadow
$S_3$ via tree-level degree-$3$ Feynman graphs. The Jacobi
identity of~$\mathfrak{g}$ forces the quartic shadow $S_4 = 0$
(the one-loop quartic graph sum vanishes by the Jacobi identity
applied to the four-valent vertex), so the tower terminates at
degree~$3$.

For class~C ($\beta\gamma$ at weight $\geq 2$): the quartic
contact invariant $Q^{\mathrm{contact}} \neq 0$ (from the
higher-pole structure), but a rank-one rigidity argument
(Theorem~\ref{thm:betagamma-rank-one-rigidity}) shows that the
self-bracket of the quartic class exits the complex, terminating
the tower at degree~$4$.

For class~M (Virasoro, $\cW_N$): the self-OPE has poles of order
$\geq 4$ (for Virasoro: $T_{(3)} T = c/2$), producing nonzero
shadow invariants at every degree. The discriminant
$\Delta = 8\kappa S_4 \neq 0$
(Theorem~\ref{thm:single-line-dichotomy}),
which by the Riccati algebraicity theorem
(Theorem~\ref{thm:riccati-algebraicity}) forces the shadow obstruction tower
to be infinite.

\textbf{Part~(iii).}
Functoriality at the level of $\kappa$ follows from the
degree-$2$ naturality of the bar construction: a morphism
$f \colon R_1 \to R_2$ induces a morphism
$V(f) \colon V(R_1) \to V(R_2)$ of vertex algebras, hence a
morphism $B(V(f)) \colon B(V(R_1)) \to B(V(R_2))$ of bar
complexes. At degree~$2$, this is the map on self-pairings
induced by~$f$, which sends $\kappa_{R_1}$ to $\kappa_{R_2}$.

For part~(a): the subalgebra inclusion
$\mathfrak{h} \hookrightarrow \mathfrak{g}$ induces
$\operatorname{Cur}(\mathfrak{h}) \hookrightarrow
\operatorname{Cur}(\mathfrak{g})$ at the same level, and the
induced map on envelopes is the inclusion
$V_k(\mathfrak{h}) \hookrightarrow V_k(\mathfrak{g})$.
The shadow map sends $\kappa(\mathfrak{h}, k)$ to
$\kappa(\mathfrak{g}, k)$ by restricting the supertrace from
$\mathfrak{g}$ to $\mathfrak{h}$ (the Casimir correction
restricts compatibly).

For part~(b): the DS reduction at the level of Lie conformal
algebras is the BRST quotient
$R_f = H^0_{\mathrm{BRST}}(\operatorname{Cur}(\mathfrak{g}),
\chi_f)$. The shadow map sends $\kappa(\widehat{\mathfrak{g}}_k)
= \dim(\mathfrak{g})(k + h^\vee)/(2h^\vee)$ to the $\cW$-algebra
modular characteristic
$\kappa(\cW_k) = \rho \cdot c_{\cW}$, where
$\rho = H_N - 1 = \sum_{i=1}^{N-1} 1/(i+1)$ is the anomaly
ratio and $c_{\cW}$ is the $\cW$-algebra central charge.
The compatibility
$\kappa(\cW_k) = \rho \cdot c_{\cW}$ follows by explicit
computation for $\mathfrak{sl}_2$ ($\rho = 1/2$,
$c_{\cW} = c_{\mathrm{Vir}}$, $\kappa = c/2$) and
$\mathfrak{sl}_3$ ($\rho = 5/6$, computed in the compute module
\texttt{factorization\_envelope\_shadow.py}).
\end{proof}


%% ===========================================================
%% KONTSEVICH GRAPH COMPLEX AND SHADOW TOWER DEFORMATIONS
%% ===========================================================

\section{The graph complex and deformation theory of the shadow obstruction tower}
\label{sec:graph-complex-shadow-deformation}

\index{graph complex!Kontsevich|textbf}
\index{Grothendieck--Teichm\"uller!and shadow tower|textbf}
\index{shadow tower!deformation theory|textbf}

The shadow obstruction tower of a modular Koszul algebra encodes the
finite-order projections of the universal Maurer--Cartan element
$\Theta_\cA$. A graph-complex deformation theory for this tower
requires a map from the modular convolution dGLA to
$\operatorname{GC}_2$ compatible with the shadow projections. This is
not a statement about algebraic Hochschild chains, categorical traces,
or $\mathrm{THH}$: the source is the modular chiral convolution
algebra, and the input needed to build the map is an $\Etwo$ formality
structure on the bar complex. Once such a map exists, Willwacher's
theorem identifies the target cohomology with
$\mathfrak{grt}_1$.
\begin{conjecture}[Graph complex controls shadow obstruction tower deformations; \ClaimStatusConjectured]
\label{conj:graph-complex-shadow}
\index{graph complex!shadow tower deformation control|textbf}
\index{Grothendieck--Teichm\"uller!isomorphism with $H^0(\operatorname{GC}_2)$|textbf}
\index{wheel cocycle|textbf}
Let $\operatorname{GC}_n$ denote the Kontsevich graph complex, whose
chain groups are spanned by isomorphism classes of connected
graphs with all vertex valences $\geq 3$ and whose differential
is edge contraction. The cohomological degree of a graph
$\Gamma$ is $|\Gamma| = |E(\Gamma)| - 2|V(\Gamma)|$.
\begin{enumerate}[label=\textup{(\roman*)}]
\item \textup{(Shadow-to-graph-complex map.)}
 There exists a Lie algebra map
 \[
 \Phi \colon
 \gAmod \;\longrightarrow\;
 \operatorname{GC}_2
 \]
 from the modular convolution algebra to the graph complex,
 induced by the Kontsevich formality quasi-isomorphism
 $\mathcal{U} \colon T_{\mathrm{poly}} \xrightarrow{\sim}
 D_{\mathrm{poly}}$ applied to the $\Etwo$ structure on the bar
 complex. Under~$\Phi$, the odd-degree shadow
 invariants $S_{2k+1}(\cA)$ for $k \geq 1$ map to the
 wheel cocycles $\sigma_{2k+1} \in H^*(\operatorname{GC}_2)$:
 \[
 \Phi(S_{2k+1}(\cA))
 \;=\;
 a_{2k+1}(\cA) \cdot [\sigma_{2k+1}],
 \]
 where $a_{2k+1}(\cA) \in \bC$ is a scalar coefficient
 depending on the algebra~$\cA$, and $\sigma_{2k+1}$ is
 represented by the wheel graph $W_{2k+1}$
 \textup{(}$2k+1$ spokes, $2k+2$ vertices, $2(2k+1)$ edges,
 loop order $2k+1$\textup{)}.
 The even-degree shadows $S_{2k}(\cA)$ map to exact forms
 in $\operatorname{GC}_2$ \textup{(}boundaries, not
 cocycles\textup{)}.

\item \textup{(Willwacher's theorem and $\mathfrak{grt}_1$.)}
 By Willwacher~\textup{\cite{Willwacher15}}, there is an
 isomorphism of graded Lie algebras
 \[
 H^0(\operatorname{GC}_2) \;\cong\; \mathfrak{grt}_1,
 \]
 where $\mathfrak{grt}_1$ is the
 Grothendieck--Teichm\"uller Lie algebra.
 The generators of $\mathfrak{grt}_1$ are the
 Deligne--Drinfeld elements $\sigma_{2k+1}$ for $k \geq 1$,
 each corresponding to the odd zeta value $\zeta(2k+1)$.
 The known dimensions of $\mathfrak{grt}_1$ at each weight
 \textup{(}Ihara, Schneps, Brown\textup{)} are:
 \begin{center}
 \renewcommand{\arraystretch}{1.15}
 \begin{tabular}{c|cccccccccccc}
 weight & $3$ & $4$ & $5$ & $6$ & $7$ & $8$ & $9$ & $10$ & $11$ & $12$ & $13$ & $14$ \\
 \hline
 $\dim$ & $1$ & $0$ & $1$ & $0$ & $1$ & $1$ & $1$ & $1$ & $2$ & $2$ & $3$ & $3$
 \end{tabular}
 \end{center}
 At even weights, all elements arise from Lie brackets:
 $[\sigma_3, \sigma_5]$ at weight~$8$,
 $[\sigma_3, \sigma_7]$ at weight~$10$.
 At weight~$12$, the two brackets
 $[\sigma_3, \sigma_9]$ and $[\sigma_5, \sigma_7]$ satisfy
 an Ihara relation, leaving one independent bracket; the
 cusp form $\Delta_{12} \in S_{12}(\mathrm{SL}_2(\bZ))$
 contributes a second generator at depth~$2$, giving
 $\dim \mathfrak{grt}_1(12) = 2$.
 The first odd weight with depth-$3$ contributions is
 weight~$11$: $\dim \mathfrak{grt}_1(11) = 2$
 \textup{(}$\sigma_{11}$ at depth~$1$ and
 $[\sigma_3, [\sigma_3, \sigma_5]]$ at depth~$3$\textup{)}.

\item \textup{(Depth classification via truncation.)}
 The shadow depth classification
 \textup{(}Definition~\textup{\ref{def:shadow-depth-classification})}
 corresponds to the truncation level in
 $\mathfrak{grt}_1$:
 \begin{center}
 \renewcommand{\arraystretch}{1.15}
 \begin{tabular}{llll}
 Class & $r_{\max}$ & $\mathfrak{grt}_1$ components & Families \\
 \hline
 G & $2$ & none
 & Heisenberg, lattice \\
 L & $3$ & $\sigma_3$ only
 & affine Kac--Moody \\
 C & $4$ & $\sigma_3$ only
 & $\beta\gamma$ \\
 M & $\infty$ & all $\sigma_{2k+1}$
 & Virasoro, $\cW_N$
 \end{tabular}
 \end{center}
 The $\operatorname{GC}_2$ components of an algebra of depth~$d$
 are supported on wheel cocycles $\sigma_{2k+1}$
 arising from odd-degree shadow invariants $S_{2k+1}$ with
 $2k + 1 \leq d$. For class~G ($d = 2$), no wheel cocycle
 component is activated. For class~L ($d = 3$) and class~C
 ($d = 4$), only~$\sigma_3$ is activated (the even-degree
 shadow $S_4$ maps to a boundary, not a cocycle).
 For class~M, all wheel cocycles $\sigma_{2k+1}$ are nonzero.

\item \textup{(Deformation control.)}
 The space of infinitesimal deformations of the shadow obstruction tower
 of~$\cA$ \textup{(}varying the algebra~$\cA$ within the
 moduli of modular Koszul algebras\textup{)} is controlled by
 $\mathfrak{grt}_1$ via the map~$\Phi$:
 a deformation of $\cA$ that changes $S_{2k+1}(\cA)$
 activates the corresponding $\sigma_{2k+1}$ component
 in $\mathfrak{grt}_1$. The $\mathfrak{grt}_1$ relations
 \textup{(}e.g., the Ihara relation at weight~$12$\textup{)} impose
 constraints on which shadow obstruction towers can arise from actual
 algebras: not every sequence $(a_3, a_5, a_7, \ldots)$ is
 realizable.
\end{enumerate}
\end{conjecture}

\begin{remark}[Graph-complex construction obligations]
The external theorem in
Conjecture~\ref{conj:graph-complex-shadow} is Willwacher's
identification
\[
H^0(\operatorname{GC}_2)\cong\mathfrak{grt}_1
\]
\cite{Willwacher15}. The two internal obligations are independent
of Theorem~H:
\begin{enumerate}[label=\textup{(\roman*)}]
\item Construct a filtered dg Lie map
\[
\Phi\colon \gAmod\longrightarrow\operatorname{GC}_2
\]
from the $\Etwo$-formality structure on the bar complex
$B(\cA)$, compatible with the modular convolution differential and
with Maurer--Cartan gauge equivalence.
\item Prove the shadow-evaluation formula
\[
\Phi(S_{2k+1}(\cA))
=a_{2k+1}(\cA)[\sigma_{2k+1}]
\]
in graph cohomology, with even shadows mapping to boundaries. This
requires fixing the graph-complex signs, the normalization of the
wheel representatives, and the comparison between the scalar
shadow coefficient and the Kontsevich weight integral.
\end{enumerate}
After these two obligations are discharged, the depth table follows
formally from the vanishing pattern of the shadow tower: class~G
activates no wheel component, classes~L and~C activate only the
$\sigma_3$ component, and class~M activates all odd wheel components
for which the corresponding $S_{2k+1}$ is nonzero. For Virasoro the
input coefficients are the recursively computed rational functions
with seeds
$S_2=c/2$, $S_3=2$, and $S_4=10/(c(5c+22))$
\textup{(}Theorem~\ref{thm:riccati-algebraicity}\textup{)}; their
generic odd-degree nonvanishing is the computation performed by
\texttt{kontsevich\_graph\_complex.py}. A deformation
$\cA_t$ then maps to
$\dot a_{2k+1}=(d/dt)S_{2k+1}(\cA_t)$ in the
$\sigma_{2k+1}$ direction, and the Ihara relations in
$\mathfrak{grt}_1$ become the stated constraints on realizable
shadow sequences.
\end{remark}

\begin{remark}[Normalized shadow zeta values]
\label{rem:normalized-shadow-zeta}
\index{shadow zeta values!normalized}
The \emph{normalized shadow zeta values} of an algebra~$\cA$
of class~M are the dimensionless ratios
\[
\zeta^{\mathrm{sh}}_r(\cA)
\;=\;
\frac{S_r(\cA)}{|\kappa(\cA)|^{r/2}}.
\]
For the Virasoro algebra at central charge~$c$, these are
rational functions of~$c$ (computed exactly in
\texttt{kontsevich\_graph\_complex.py},
function~\texttt{shadow\_zeta\_exact}).
The terminology ``shadow zeta values'' reflects the
correspondence $\sigma_{2k+1} \leftrightarrow \zeta(2k+1)$
in $\mathfrak{grt}_1$: the shadow coefficient
$S_{2k+1}(\cA)$ maps to the generator $\sigma_{2k+1}$ that
corresponds to the odd zeta value $\zeta(2k+1)$. The shadow
zeta values are \emph{rational} (they are algebraic invariants
of the vertex algebra), not transcendental; the connection
to the actual zeta values $\zeta(2k+1) \in \bR$ is through
the graph weights in the Kontsevich integral, which are periods
(integrals of algebraic forms over semi-algebraic sets) that
evaluate to multiple zeta values.
\end{remark}

\section{K3 categorical input and curve-level chiral Hochschild scope}
\label{sec:chk-supplement}

\begin{remark}[K3 categorical Hochschild bi-grading and the curve
specialisation]
\label{rem:chk-K3-cy2-shift}
At $\cC = D^b\mathrm{Coh}(K3)$ the \emph{categorical}
Hochschild cohomology is HKR bi-graded
$\mathrm{HH}^\bullet(D^b\mathrm{Coh}(K3)) = \bigoplus_{p+q=\bullet} H^p(K3, \Lambda^q T_{K3})$,
and the categorical Serre shift equals the CY dimension of the input:
$[2]$.  The Vol~I curve chart is a separate Hochschild functor.  After
a Vol~III two-stage specialization to a reference curve, a specified
categorical-to-chart comparison and a family support datum
$H_H(\cA;S)$ identify its curve-level support with~$S$.
Under $\widetilde{M}_{24}$-equivariance the Serre functor acquires an
anomaly cocycle of order $6$ in
$H^3(M_{24}, \mathrm{U}(1)) \supseteq \mathbb{Z}/12$
(generalised-Mathieu-moonshine 't~Hooft anomaly,
Gaberdiel--Persson--Ronellenfitsch--Volpato 2013;
Bridgeland--Maciocia $2001$ on K3 autoequivalences). The Vol~III
K3 chiral Hall--Drinfeld bridge supplies a bar-cohomology coalgebra:
\[
(\mathbf{H}_{\Delta_5})^i
:=H^\bullet B^{\mathrm{ord}}(\mathbf{H}_{\Delta_5})
\;\simeq\;
V(\mathfrak{g}_{\Delta_5})^{\mathrm{coalg}}[2]
\otimes \eta_{\widetilde{M}_{24}}.
\]
A reference-curve specialization followed by finite-type continuous
Verdier duality forms the Koszul dual algebra
$(\mathbf{H}_{\Delta_5})^!$.  This is an additional Vol~III bridge
input to the curve-level Theorem~H datum.
\end{remark}

\begin{remark}[Mukai doubling $K^{\kappa_{\mathrm{ch}}} = 8$]
\label{rem:chk-Kkappa-8}
On the Mukai-enhanced K3 Hochschild-Koszul datum, the conductor is
$K^{\kappa_{\mathrm{ch}}}(\mathcal{H}_{\mathrm{Muk}}(K3)) = 2\,c_+(\mathrm{Mukai}(K3)) = 8$.
Universal identity $\hbar^2 \cdot K^{\kappa_{\mathrm{ch}}} = -1$ pins the Lusztig
$u_\zeta$-specialisation at $\ell = 8$, matching the order of the monodromy of
$\mathcal{L}^{\Delta_5}$ around Humbert $H_1 \subset \overline{\mathcal{A}_2}$
(Bruinier $2002$ Prop~$5.1$). Vol~I Theorem~C bucket enlarges on the
$\mathcal{B}$-family to $\kappa + \kappa^! \in \{0, 8, 13, 250/3, 25/3\}$.
\end{remark}

\section{K3 Hochschild boundary data}
\label{sec:chk-deep}

\begin{remark}[Curve-level support for $\mathbf H_{\Delta_5}$]
\label{rem:chk-deep-1}
Regard \(\mathbf H_{\Delta_5}\) as a chiral algebra on a fixed smooth
projective curve.  An explicit family support datum
\(H_H(\mathbf H_{\Delta_5};S_{\Delta_5})\) gives
\[
 \operatorname{Supp}
 \mathrm{ChirHoch}^\bullet(\mathbf H_{\Delta_5})
 \subseteq S_{\Delta_5}.
\]
The curve-level calculation then has three further concrete tasks:
\begin{itemize}
\item degree zero is the chiral centre
\(Z(\mathbf H_{\Delta_5})\); the Hall--Drinfeld model determines its
Casimir and Igusa-form generators;
\item degree one is the quotient of chiral derivations by inner chiral
derivations, computed in the same curve-level model;
\item an identification of a degree-two chart class with
\(H^2(\mathfrak g_{\Delta_5})^{\mathbb Z/2,K(1)}\) is carried by an
Etingof--Kazhdan comparison map for the Vol~III Hall--Drinfeld
construction.
\end{itemize}
The CY-\(3\) volume class belongs to the product-ambient complex after
the Vol~III \(\Phi_3^{(\Sigma_2,C)}\) specialization.  Its relation to
the curve chart is governed by the corresponding product-to-chart
comparison map.
\end{remark}

\begin{remark}[Four-piece Hochschild differential]
\label{rem:chk-deep-2}
In the Vol~III product-ambient Hall--Drinfeld model, the Hochschild
differential decomposes as
\[
d_{\mathrm{Hoch}} \;=\; d_{\mathrm{Hall,Hoch}} \;+\; d_{\mathrm{Sieg,Hoch}} \;+\; \hbar\, d_{\mathrm{assoc,Hoch}} \;+\; \hbar^2\, d_{\Phi_{10}/\eta^{24},\mathrm{Hoch}},
\]
the Hochschild dual of the four-piece bar differential (Vol.~II,
\S\ref{V2-sec:bar-cobar-review-supplement}). The four pieces live in
distinct $\mathrm{Sp}_4(\mathbb Z)$-weights: $d_{\mathrm{Hall,Hoch}}$
in weight $0$, $d_{\mathrm{Sieg,Hoch}}$ in weight $5$,
$d_{\mathrm{assoc,Hoch}}$ in weight $10$, and
$d_{\Phi_{10}/\eta^{24},\mathrm{Hoch}}$ in weight $10 - 24 = -14$
(the anti-holomorphic Dedekind contribution). This is not the
curve-only differential used in Theorem~H unless the product-ambient
pieces have first been restricted to the reference curve and the
PBW chiral Koszul hypotheses have been proved there.
\end{remark}

\begin{remark}[Etingof-Kazhdan Part V and K3 comparison]
\label{rem:chk-deep-3}
The Etingof-Kazhdan 2007 Part V super-category extension classifies
deformations of the BKM Lie bialgebra $\mathfrak g_{\Delta_5}$ into
its universal enveloping chiral algebra. Applied to the Vol~III
Hall--Drinfeld construction, it supplies the comparison input
\[
\mathrm{ChirHoch}^2(\mathbf H_{\Delta_5})^{\mathbb Z/2,K(1)}
\;\cong\; \mathbb C\cdot\Delta_5.
\]
The isomorphism identifies the deformation direction only after the
curve-level chiral Hochschild complex has been compared with the BKM
super-category complex. It does not by itself prove uniqueness of
$\mathbf H_{\Delta_5}$ as a chiral bialgebra; uniqueness requires a
separate rigidity theorem for the Hall--Drinfeld presentation. The
non-projective global lift to $\bar{\mathcal A}_2$ requires pinning
$\hbar$ to the
Lusztig value, which is why convergence occurs at $\hbar^2 = -1/8$.
Primary-lit: Etingof-Kazhdan 2007 Part~V, Morava 1985 K-theory.
\end{remark}
